% !TeX root = ../Book2.tex
\BookPart{Fourier 分析与逼近论}{b2:part:01}
\BookPartGuide

Fourier 分析把函数的变化分解成不同频率的振动。周期问题首先引出 Fourier 级数，整个欧氏空间上的问题则引出 Fourier 变换。分解之后，还要说明怎样收敛、怎样恢复函数，以及截取有限项时会留下多大误差。经典逼近论进一步比较不同的近似函数，研究最佳误差与光滑性的关系。

\ProseParagraphBreak

本编预设 Lebesgue 积分、基本 Hilbert 空间理论及\chapref{b2:ch:01}的必要微积分工具。\chapref{b2:ch:02}从正交展开研究周期函数；\chapref{b2:ch:approximation}承接周期卷积与正核，组织一致最佳逼近和定量逼近。\chapref{b2:ch:03}在可积函数上建立 Fourier 变换，\chapref{b2:ch:04}用平方积分与插值扩大函数空间，\chapref{b2:ch:05}通过快速下降测试函数与温和分布容纳更多对象。\chapref{b2:ch:06}再把离散频率和连续频率联系起来。主要阅读依赖为
\[
\begin{tikzcd}[column sep=2.5em,row sep=1.8em]
\text{第\ref{b2:ch:02}章}\arrow[r]\arrow[rrd]
 &\text{第\ref{b2:ch:approximation}章}&\\
 &&\text{第\ref{b2:ch:06}章}\\
\text{第\ref{b2:ch:03}章}\arrow[r]
 &\text{第\ref{b2:ch:04}章}\arrow[r]
 &\text{第\ref{b2:ch:05}章}\arrow[u]
\end{tikzcd}
\]
两条起始支线共同使用第一卷的积分与函数空间工具。逼近论中的 Jackson 构造沿用周期核；它不反过来承担 Fourier 系完备性的证明。Poisson 求和使用周期 Fourier 系数及连续 Fourier 变换，其所需的后者已经沿第二条路线建立。逼近论章仍有待撰写部分，现有 Fourier 理论的证明不依赖这些占位内容。

\ProseParagraphBreak

以 Fourier 基础为目标，可从级数转入变换，随后补读经典逼近论；若希望系统理解有限近似与误差速度，则按本编章序阅读。以经典 PDE 为目标，可在掌握基本级数与变换、并补齐第一编的 ODE 工具后进入第三编。\chapref{b2:ch:06}的带限采样与频率混叠另列为\subsecref{b2:subsec:060105}的选读专题，其中使用\chapref{b2:ch:04}的 Plancherel 理论；求和公式及后续热核论证均不以它为前置。点态收敛的深层结果适合第二遍阅读。

\ProseParagraphBreak

Elias M. Stein 与 Rami Shakarchi 的《Fourier Analysis: An Introduction》\cite{Stein2003Fourier}从级数讲到变换，并通过具体应用解释频率分解。积分换序与函数空间的基础可回查第一卷相应章节。

% !TeX root = ../../Book2.tex
\chapter{Fourier 级数}
\label{b2:ch:02}
\BookChapterNumber{b2:ch:02}
% Axler2020Measure：实读11A的系数、正核与光滑衰减，11B的完备性、卷积、唯一性；
% p.352仅用来定位Carleson与Kolmogorov概览，不冒称原书给出深层证明。
\begin{chapterabstract}
周期函数可以用不同频率的复指数描述，但计算出系数只是展开的开始。本章从 Hilbert 空间中的正交投影出发，将部分和写成 Dirichlet 核的卷积，分别研究局部有界变差函数的点态收敛和平方可积函数的范数收敛。Fejér 平均通过正核给出更稳定的逼近。随后把光滑性与系数衰减联系起来，说明积分、经典求导和弱求导各自的条件。选读部分以 Gibbs 过冲和算子范数增长解释普通部分和的局限。
\end{chapterabstract}

\begin{learninggoals}
\begin{itemize}
\item\label{b2:item:02-goal-projection}用有限误差恒等式理解 Bessel、Parseval 与最佳平方逼近，区分正交性和完备性。
\item\label{b2:item:02-goal-kernel}计算 Dirichlet 核、Fejér 核与周期 Poisson 核，并判断振荡、非负性和质量集中各起什么作用。
\item\label{b2:item:02-goal-convergence}按函数的局部正则性、可积指数和系数衰减选择合适的收敛结论。
\item\label{b2:item:02-goal-operation}掌握逐项运算的条件，用带权平方和刻画周期弱导数。
\item\label{b2:item:02-goal-boundary}（选读）解释 Gibbs 过冲、几乎处处收敛的例外集及对称部分和在 \(L^1\) 端点的范数障碍，避免将存在性反例误读为所有函数都发散。
\end{itemize}
\end{learninggoals}

一根弦的初始形状通常不是单独的正弦曲线，而每个正弦振动都有一个易于求解的时间演化。把初始形状分解成这些振动之后，便可以分别求解再叠加。这个想法马上产生两个问题：怎样选取每个振动的大小，有限叠加趋向无穷叠加时又保留了什么？在最小平方误差的意义下，大小由正交投影确定；若关心某一点的位移，尤其是有折角或跳跃的位置，平方误差公式还不能回答全部问题。

\ProseParagraphBreak

这里还要区分固定端点与周期边界。长度为 \(L>0\)、两端固定的弦使用空间模态 \(\sin(n\pi x/L)\)，其中 \(n\geq1\) 为整数；对满足固定端点条件的初始形状，先从 \((0,L)\) 作奇延拓，再按 \(2L\) 周期延拓，就把正弦展开放进了周期 Fourier 展开的框架。重新缩放坐标后，便可使用本章的单位周期约定。我们先研究一般周期函数；边界条件决定选用哪些模态，不能在换成周期模型时将它们忽略。

\ProseParagraphBreak

本章预设 Lebesgue 积分、\(L^p\) 完备性及 Hilbert 空间内积的基本性质。所需正交展开在\secref{b2:sec:0201}重新建立；选读部分关于发散的存在性证明用到第一卷定理 8.3.3（一致有界原理）。周期归一化固定为 \(1\)，内积仍对第一变量线性。

\ProseParagraphBreak

\secref{b2:sec:0202}和\secref{b2:sec:0203}分别讨论点态收敛与平方范数收敛，\secref{b2:sec:0204}利用正核研究连续函数与可积函数的逼近。\secref{b2:sec:0205}与\secref{b2:sec:0206}准备后续方程所需的收敛与求导工具。初次阅读可暂缓\secref{b2:sec:0207}，但应注意，平方范数收敛并不直接推出整列点态收敛。

\ProseParagraphBreak

Sheldon Axler 的《Measure, Integration \& Real Analysis》\cite[Sections 11A--11B]{Axler2020Measure}讨论系数唯一性、正核和 Hilbert 空间方法。Elias M. Stein 与 Rami Shakarchi 的《Fourier Analysis: An Introduction》\cite{Stein2003Fourier}把 Fourier 级数用于周期振动与热传导；研究更深入的三角级数问题时，可查阅 Antoni Zygmund 的《Trigonometric Series》\cite{Zygmund2002Trigonometric}。

% !TeX root = ../../Book2.tex
\section{Hilbert 空间中的正交系统}
\label{b2:sec:0201}

把一个向量沿互相垂直的方向分解，是有限维线性代数中熟悉的计算。对函数进行同样的分解时，有限个坐标变成一列积分，而重建函数要经过极限。本节先把这两步分开：正交性给出误差公式；只有正交系足够丰富，误差才必然趋于零。

\subsection{Fourier 系数}
\label{b2:subsec:020101}

本章固定圆周 \(\mathbb T\coloneqq\mathbb R/\mathbb Z\)，把它上的函数看作周期为 \(1\) 的函数。积分取任一长度为 \(1\) 的区间上的 Lebesgue 积分，通常写为 \(\int_{\mathbb T}\)，因此 \(\int_{\mathbb T}1\,\mathrm{d}x=1\)。函数空间 \(L^p(\mathbb T)\) 中的元素仍按几乎处处相等取商；需要点值时另行指定代表。内积固定为
\[
 \langle f,g\rangle\coloneqq\int_{\mathbb T}f(x)\cdot\overline{g(x)}\,\mathrm{d}x,
\]
对第一变量线性。周期为 \(L>0\) 的问题可令 \(x=Lt\) 化到本章约定，系数中的积分同时变为 \(L^{-1}\int_0^L\)。
\symidx{\(\mathbb T\)：单位周期圆周}
\symidx{\(e_k\)：圆周上的指数函数}

令 \(e_k(x)\coloneqq e^{2\pi\mathrm{i}kx}\)。直接积分给出 \(\langle e_k,e_l\rangle=\delta_{kl}\)。若实或复 Hilbert 空间 \(H\) 中一族向量 \((e_k)_{k\in J}\) 满足 \(\langle e_k,e_l\rangle=\delta_{kl}\)，称其为\term[biaozhun-zhengjiao-xi]{标准正交系}[orthonormal system]；不要求每个向量具有单位长度时，只称正交系。此处的“标准”正好消去投影系数中的长度分母。

\begin{definition}\label{b2:def:02-fourier-coefficient}
取 \(\mathbb T=\mathbb R/\mathbb Z\) 的归一化 Lebesgue 测度，\(f\in L^1(\mathbb T;\mathbb C)\)，并令 \(e_k(x)=e^{2\pi\mathrm i kx}\)。对每个 \(k\in\mathbb Z\)，积分
\[
 \widehat f(k)\coloneqq\int_{\mathbb T}f(x)\cdot e^{-2\pi\mathrm{i}kx}\,\mathrm{d}x
\]
称为 \(f\) 的第 \(k\) 个\term[fourier-xishu]{Fourier 系数}[Fourier coefficient]。形式级数 \(\sum_{k\in\mathbb Z}\widehat f(k)e_k\) 称为它的\term[fourier-jishu]{Fourier 级数}[Fourier series]。在确定收敛意义以前，不把这个形式级数直接写成等于 \(f\)。
\end{definition}
\symidx{\(\widehat f(k)\)：周期函数的 Fourier 系数}

估计 \(|\widehat f(k)|\leq\|f\|_1\) 说明每个系数都定义在 \(L^1\) 等价类上。若 \(f\in L^2\)，系数又等于 \(\langle f,e_k\rangle\)。例如常函数只占据频率 \(0\)；\(\cos(2\pi mx)\) 在 \(m\) 与 \(-m\) 两个频率上的系数均为 \(1/2\)。复指数的正负频率不是两个独立的实振动：实值函数满足 \(\widehat f(-k)=\overline{\widehat f(k)}\)。

\subsection{Bessel 不等式}
\label{b2:subsec:020102}

下列误差恒等式同时给出\term[bessel-budengshi]{Bessel 不等式}[Bessel's inequality]及有限维逼近的几何意义。

\begin{proposition}[Bessel]\label{b2:prop:02-bessel}
设 \((u_j)_{j\geq1}\) 是 实或复 Hilbert 空间 \(H\)（复数时内积对第一变量线性） 中的标准正交系。对 \(v\in H\)、有限集 \(F\subset\mathbb N\) 及任意标量 \((a_j)_{j\in F}\)，有
\[
 \left\|v-\sum_{j\in F}a_ju_j\right\|^2
 =\|v\|^2-\sum_{j\in F}|\langle v,u_j\rangle|^2
   +\sum_{j\in F}|a_j-\langle v,u_j\rangle|^2.
\]
特别地，\(\sum_{j\geq1}|\langle v,u_j\rangle|^2\leq\|v\|^2\)。
\end{proposition}
\begin{proof}
记 \(P_Fv=\sum_{j\in F}\langle v,u_j\rangle u_j\)。对每个 \(l\in F\)，内积第一变量的线性性给出 \(\langle v-P_Fv,u_l\rangle=0\)。于是 \(v-P_Fv\) 与 \(\sum_{j\in F}(\langle v,u_j\rangle-a_j)u_j\) 正交。对这两项使用勾股恒等式，再利用标准正交性计算后一项的范数，就得到所述等式。取 \(a_j=\langle v,u_j\rangle\)，右侧剩余的误差范数非负；最后让有限集递增穷尽 \(\mathbb N\)。
\end{proof}

只知道一组方向互相垂直，并不能断言它们能恢复所有向量。例如只取偶数频率 \((e_{2k})_{k\in\mathbb Z}\)，函数 \(e_1\) 的全部相应系数都为零。Bessel 不等式没有失效，而是这组方向遗漏了整个正交补。

\subsection{抽象正交展开与 Parseval 定理}
\label{b2:subsec:020103}

\begin{definition}\label{b2:def:02-complete-system}
设 \(H\) 为实或复 Hilbert 空间，\((u_j)_{j\geq1}\) 是其中的标准正交系。若其有限线性组合在 \(H\) 中稠密，则称它是\term[wanbei-zhengjiao-xi]{完备正交系}[complete orthonormal system]，也称\term[biaozhun-zhengjiao-ji]{标准正交基}[orthonormal basis]。
\end{definition}

\begin{theorem}[正交展开与 Parseval]\label{b2:thm:02-abstract-parseval}
设 \(H\) 为实或复 Hilbert 空间（复数时内积对第一变量线性），\((u_j)_{j\geq1}\) 是其中的标准正交系，\(M\) 是其线性张成空间的闭包。对任意 \(v\in H\)，级数 \(\sum_j\langle v,u_j\rangle u_j\) 在 \(H\) 中收敛到唯一满足 \(w\in M\)、\(v-w\perp M\) 的向量 \(w\)。该系完备当且仅当每个 \(v\in H\) 都满足
\[
 \|v\|^2=\sum_{j\geq1}|\langle v,u_j\rangle|^2.
\]
这一等式称为\term[parseval-hengdengshi]{Parseval 恒等式}[Parseval's identity]。
\end{theorem}
\begin{proof}
Bessel 不等式使系数的平方和收敛。两段部分和之差的范数平方恰为相应尾部平方和，故部分和是 Cauchy 列，由 \(H\) 完备得到极限 \(w\in M\)。对每个固定 \(u_l\) 取内积极限可得 \(\langle v-w,u_l\rangle=0\)，再用连续性推广到 \(M\)。若另有 \(w'\) 满足同样条件，则 \(w-w'\) 同时属于 \(M\) 和 \(M^\perp\)，只能为零。有限误差公式取极限后成为
\[
 \|v-w\|^2=\|v\|^2-\sum_j|\langle v,u_j\rangle|^2.
\]
若 \(M=H\)，则 \(w=v\)，等式成立。反之，若等式对每个 \(v\) 成立，误差为零，所以每个 \(v\) 都属于 \(M\)。
\end{proof}

这一证明还说明，只要有限频率集最终包含每个频率，逼近次序就不影响 Hilbert 空间中的极限。对称部分和是研究点态行为时的自然选择，却不是平方范数收敛所必需的排列约定。

\subsection{正交投影}
\label{b2:subsec:020104}

前一定理用可数标准正交系构造了到其闭线性张成空间的投影。一般闭子空间的存在性则由第一卷定理 8.4.6 保证：取标量域 \(\mathbb K=\mathbb R\) 或 \(\mathbb C\)，设 \(H\) 为 \(\mathbb K\) 上的 Hilbert 空间、\(M\subseteq H\) 为闭线性子空间，复情形内积仍对第一变量线性。每个 \(v\in H\) 都有唯一分解
\[
 v=P_Mv+(v-P_Mv),\quad P_Mv\in M,\quad v-P_Mv\in M^\perp.
\]
映射 \(P_M:H\to M\) 称为到 \(M\) 的\term[zhengjiao-touying]{正交投影}[orthogonal projection]。所引定理先以距离最小化和平行四边形恒等式建立最近点，再证明误差正交，因此不要求 \(M\) 已经由某个可数正交系生成。
% 跨卷引用：第一卷定理8.4.6，b1:thm:hilbert-projection，Book1/Part02/Chapter08/Section0804.tex。

\ProseParagraphBreak

线性性可以直接从分解的唯一性读出。对任意 \(v_1,v_2\in H\) 与 \(a,b\in\mathbb K\)，向量 \(aP_Mv_1+bP_Mv_2\) 属于 \(M\)，而
\[
 (av_1+bv_2)-(aP_Mv_1+bP_Mv_2)
 =a(v_1-P_Mv_1)+b(v_2-P_Mv_2)\in M^\perp.
\]
最后的包含关系来自内积对第一变量的线性性：右端与任意 \(m\in M\) 的内积都为零。由唯一性，\(P_M(av_1+bv_2)=aP_Mv_1+bP_Mv_2\)。勾股恒等式再给出
\[
 \|v\|^2=\|P_Mv\|^2+\|v-P_Mv\|^2,
\]
故 \(\|P_Mv\|\leq\|v\|\)。对任意 \(z\in M\)，误差 \(v-P_Mv\) 与 \(P_Mv-z\) 正交，因而
\[
 \|v-z\|^2=\|v-P_Mv\|^2+\|P_Mv-z\|^2.
\]
因此 \(P_Mv\) 是 \(M\) 中唯一的最佳逼近。这里的距离来自 Hilbert 范数；换成一致范数，最佳逼近一般不再由 Fourier 系数直接给出。

\ProseParagraphBreak

对圆周上的函数，记
\[
 S_Nf\coloneqq\sum_{|k|\leq N}\widehat f(k)e_k.
\]
当 \(f\in L^2(\mathbb T)\) 时，它是到 \(\operatorname{span}\{e_{-N},\ldots,e_N\}\) 的正交投影，且 \(\|S_Nf\|_2\leq\|f\|_2\)。对一般 \(L^1\) 函数，这个有限和仍有意义，但正交投影解释已超出其定义范围。下一节先研究这个有限和的积分核，再回到整个三角系统是否完备的问题。
\symidx{\(S_Nf\)：Fourier 级数的对称部分和}

% !TeX root = ../../Book2.tex
\section{Dirichlet 核与经典收敛}
\label{b2:sec:0202}

一个部分和含有许多 Fourier 系数。把各系数的积分合并，可以把“同时截取许多频率”改写成对原函数的一次加权积分。权重的振荡既产生抵消，也决定跳跃附近的行为。

\subsection{Dirichlet 核}
\label{b2:subsec:020201}

\begin{definition}\label{b2:def:02-dirichlet-kernel}
对非负整数 \(N\)，函数
\[
 D_N(t)\coloneqq\sum_{|k|\leq N}e^{2\pi\mathrm{i}kt}
\]
称为\term[dirichlet-he]{Dirichlet 核}[Dirichlet kernel]。
\end{definition}
\symidx{\(D_N\)：Dirichlet 核}

有限几何级数给出
\[
 D_N(t)=\dfrac{\sin((2N+1)\pi t)}{\sin(\pi t)}\quad(t\notin\mathbb Z),
 \quad D_N(0)=2N+1.
\]
它是实值偶函数，积分为 \(1\)，但并不处处非负。峰值随 \(N\) 增高，并不能单独保证卷积趋向原函数；还要控制旁边不断增加的正负振荡。

\subsection{部分和算子}
\label{b2:subsec:020202}

取 \(\mathbb T=\mathbb R/\mathbb Z\) 的归一化 Lebesgue 测度，\(f,g\in L^1(\mathbb T;\mathbb C)\)。圆周上的\term[zhouqi-juanji]{周期卷积}[periodic convolution]记为
\[
 (f*g)(x)\coloneqq\int_{\mathbb T}f(x-t)\cdot g(t)\,\mathrm{d}t,
\]
在积分绝对收敛的点定义。有限换序直接给出
\begin{equation}\label{b2:eq:02-dirichlet-convolution}
 S_Nf(x)=(f*D_N)(x)
 =\int_{-1/2}^{1/2}f(x-t)\cdot D_N(t)\,\mathrm{d}t.
\end{equation}
由于 \(D_N\) 有界，这个等式对 \(f\in L^1\) 的每个 \(x\) 都有意义，右端的值不随代表的零测修改改变。

\begin{lemma}[振荡积分趋零]\label{b2:lem:02-oscillation}
若 \(h\in L^1(a,b)\)，其中 \(a<b\) 有限，则
\[
 \int_a^b h(t)\cdot e^{\mathrm{i}\lambda t}\,\mathrm{d}t\longrightarrow0
 \quad(|\lambda|\to\infty).
\]
\end{lemma}
\begin{proof}
对区间示性函数，积分等于两个单位复数之差除以 \(\mathrm{i}\lambda\)，模不超过 \(2/|\lambda|\)。有限阶梯函数也满足结论。任意 \(L^1\) 函数可先截断幅值，再用简单函数逼近；有限测度可测集由有限个区间的并在测度上逼近，故阶梯函数在 \(L^1(a,b)\) 中稠密。选阶梯函数 \(q\) 使 \(\|h-q\|_1<\varepsilon\)，两积分之差对所有 \(\lambda\) 都不超过 \(\varepsilon\)，再令 \(|\lambda|\to\infty\) 和 \(\varepsilon\downarrow0\)。
\end{proof}

上述区间逼近只用 Lebesgue 测度的正则性：先取开集逼近可测集，再保留开集区间分解的有限项。将引理用于 \(\lambda=2\pi k\) 可知 \(\widehat f(k)\to0\)。这称为周期情形的\term[riemann-lebesgue-yinli]{Riemann--Lebesgue 引理}[Riemann--Lebesgue lemma]。它只说单个系数变小，不说系数绝对可和。

\subsection{Dirichlet--Jordan 定理}
\label{b2:subsec:020203}

先说明处理有界变差所需的抵消估计。有限闭区间 \([a,b]\)、\(a<b\) 上实值或复值函数的变差是所有有限分割下相邻函数值差的绝对值之和的上确界；上确界有限时称为\term[youjie-biancha-hanshu]{有界变差函数}[function of bounded variation]。开区间上的变差也取全部落在区间内部的有限分割，不要求端点属于函数的定义域。这里使用的是一个指定代表的点值变差，不是在等价类中对代表取下确界。

所需估计的核心可以先概括如下：Dirichlet 核虽然正负振荡，但它在任意子区间上的积分有同一个常数界。若 \(g\) 是阶梯函数，分部求和就把 \(\int gD_N\) 改写成一个端点值与各次跳变量乘以这些核积分；跳变量的绝对值之和受总变差控制。对于一般有界变差函数，再用阶梯函数逼近即可。最后缩小原点附近的区间：\(g(0+)=0\) 使端点值趋零，有限总变差则使小区间中的变差趋零。下面依次实现这几个步骤。

先说明下文出现的单侧极限确实存在。有界变差的实值函数在一个区间中可固定内点 \(c\)，令 \(V(t)=\operatorname{Var}_{[c,t]}g\) 于 \(t>c\)，\(V(c)=0\)，\(V(t)=-\operatorname{Var}_{[t,c]}g\) 于 \(t<c\)。相邻区间的变差可加，故 \((V+g)/2\) 与 \((V-g)/2\) 都是有界非减函数，其差为 \(g\)。这证明可测性，也保证内部点及区间端点的单侧极限存在。复值函数的实部、虚部各自有界变差，因而有相同性质；不需要预先把函数改成右连续代表。

\begin{lemma}\label{b2:lem:02-bv-oscillation}
设 \(0<\delta_0<1/2\)，实值或复值函数 \(g:(0,\delta_0)\to\mathbb C\) 满足 \(\operatorname{Var}_{(0,\delta_0)}g<\infty\) 及 \(g(0+)=0\)。存在与 \(g,\delta_0,N,\delta\) 无关的常数 \(C\)，使每个整数 \(N\geq0\) 及 \(0<\delta\leq\delta_0\) 都满足
\[
 \left|\int_0^\delta g(t)\cdot D_N(t)\,\mathrm{d}t\right|
 \leq C\bigl(|g(\delta-)|+\operatorname{Var}_{(0,\delta)}g\bigr).
\]
其中 \(g(\delta-)\) 是从区间内部取得的左极限。对同一个函数 \(g\) 缩小 \(\delta\) 时，右侧趋于零，因而左侧对全部 \(N\) 一致趋零。
\end{lemma}
\begin{proof}
先证明核积分的一致界。记 \(\lambda=(2N+1)\pi\)。在 \(0<t<1/\lambda\) 上，\(|\sin(\lambda t)|\leq\lambda t\)、\(\sin(\pi t)\geq2t\)，所以任何落在这段内的子区间，其 \(|D_N|\) 积分至多 \(1/2\)。若 \(1/\lambda\leq u<v\leq1/2\)，令
\[
 w(t)=\frac1{\sin(\pi t)},\qquad
 A(t)=\int_u^t\sin(\lambda s)\,\mathrm ds.
\]
此时 \(w\) 为正且单调不增，\(|A(t)|\leq2/\lambda\)。普通分部积分给出
\[
 \begin{aligned}
 \left|\int_u^vD_N(t)\,\mathrm dt\right|
 &=\left|w(v)A(v)-\int_u^vw'(t)A(t)\,\mathrm dt\right|\\
 &\leq\frac2\lambda\left(w(v)+\int_u^v|w'(t)|\,\mathrm dt\right)\\
 &=\frac{2w(u)}\lambda\leq\frac1{\lambda u}\leq1.
 \end{aligned}
\]
把一个任意区间在 \(1/\lambda\) 处分开，得到
\[
 \left|\int_a^bD_N(t)\,\mathrm dt\right|\leq C_0
 \qquad(0\leq a<b\leq1/2),
\]
其中可取绝对常数 \(C_0=2\)，与 \(N,a,b\) 均无关。

接着固定 \(0<a<b<\delta\)，只用这个闭区间上的点值，不改换 \(g\) 的代表。把 \([a,b]\) 等分为 \(a=t_0<\cdots<t_m=b\)，令阶梯函数 \(q_m(t)=g(t_{j-1})\) 于 \(t_{j-1}<t<t_j\)，有限个分点处的值不影响积分。由变差定义，
\[
 \|q_m-g\|_{L^1(a,b)}
 \leq\frac{b-a}{m}\sum_{j=1}^m
       \operatorname{Var}_{[t_{j-1},t_j]}g
 =\frac{b-a}{m}\operatorname{Var}_{[a,b]}g\longrightarrow0.
\]
这里用到变差在相邻闭区间上可加：把分割在公共端点处分开给出一个方向，把两个分割接起来给出另一个方向。这个论证对复数值同样成立，因为只使用模的三角不等式。

令 \(P_N(t)=\int_a^tD_N(s)\,\mathrm ds\)，则 \(P_N(a)=0\)、\(|P_N|\leq C_0\)。有限求和的相消给出精确恒等式
\[
 \begin{aligned}
 \int_a^b q_m(t)D_N(t)\,\mathrm dt
 &=\sum_{j=1}^m g(t_{j-1})\bigl(P_N(t_j)-P_N(t_{j-1})\bigr)\\
 &=g(t_{m-1})P_N(b)
   -\sum_{j=1}^{m-1}\bigl(g(t_j)-g(t_{j-1})\bigr)P_N(t_j).
 \end{aligned}
\]
把 \(g(t_{m-1})\) 写成 \(g(b)+(g(t_{m-1})-g(b))\)，最后一个增量与求和中的增量一起受总变差控制，故上式的模不超过
\(C_0(|g(b)|+\operatorname{Var}_{[a,b]}g)\)。对固定 \(N\)，\(D_N\) 有界，上述 \(L^1\) 逼近允许令 \(m\to\infty\)，从而
\[
 \left|\int_a^b g(t)D_N(t)\,\mathrm dt\right|
 \leq C_0\bigl(|g(b)|+\operatorname{Var}_{[a,b]}g\bigr).
\]
再令 \(a\downarrow0\)、\(b\uparrow\delta\)。\(g\) 在 \((0,\delta_0)\) 上有界，因而对固定 \(N\) 积分收敛；\(g(b)\to g(\delta-)\)，且 \(\operatorname{Var}_{[a,b]}g\leq\operatorname{Var}_{(0,\delta)}g\)。这就得到引理所列的、对全部 \(N\) 统一的估计。端点项由 \(g(0+)\) 和 \(g(\delta-)\) 的实际单侧极限确定。

最后核对小区间变差确实趋零。取固定 \(c\in(0,\delta_0)\)，令
\(V(t)=\operatorname{Var}_{[t,c]}g\)、\(0<t<c\)。随 \(t\downarrow0\)，\(V(t)\) 单调增加且有上界，故有有限极限 \(L\)。当 \(0<a<b<\delta<c\) 时，相邻区间的可加性给出
\[
 \operatorname{Var}_{[a,b]}g=V(a)-V(b)
 \leq L-V(\delta).
\]
对全部内部有限分割取上确界，得到
\(\operatorname{Var}_{(0,\delta)}g\leq L-V(\delta)\to0\)。这也控制原代表在孤立点的跳动，不必先删除这些点。另一方面，\(g(0+)=0\) 使 \(g(\delta-)\to0\)，于是右侧随 \(\delta\downarrow0\) 趋零。
\end{proof}

\begin{theorem}[Dirichlet--Jordan]\label{b2:thm:02-dirichlet-jordan}
设 \(f\in L^1(\mathbb T;\mathbb C)\)，\(x\in\mathbb T\)，且某一周期代表在点 \(x\) 两侧的一个邻域内有界变差，则
\[
 S_Nf(x)\longrightarrow\dfrac{f(x+)+f(x-)}2.
\]
\end{theorem}
\begin{proof}
令 \(A=(f(x+)+f(x-))/2\)，并在 \(0<t<1/2\) 记
\[
 g(t)=f(x+t)+f(x-t)-2A.
\]
利用核的偶性及积分为 \(1\)，\eqref{b2:eq:02-dirichlet-convolution}给出 \(S_Nf(x)-A=\int_0^{1/2}g(t)\cdot D_N(t)\,\mathrm{d}t\)。取固定的 \(0<\delta_0<1/2\)，使所选代表在 \((x-\delta_0,x)\) 与 \((x,x+\delta_0)\) 上的变差有限；圆周坐标在这里按周期延伸理解。则 \(g\) 在 \((0,\delta_0)\) 上的变差不超过这两个变差之和，且 \(g(0+)=f(x+)+f(x-)-2A=0\)。因此前一引理使 \((0,\delta)\) 上的积分在 \(\delta\downarrow0\) 时对全部 \(N\) 一致趋小。在 \([\delta,1/2]\) 上，\(g(t)/\sin(\pi t)\) 可积，\cref{b2:lem:02-oscillation}使其与 \(\sin((2N+1)\pi t)\) 的积分趋零。先固定 \(\delta\) 令 \(N\to\infty\)，再令 \(\delta\downarrow0\)。
\end{proof}

定理名称\term[dirichlet-jordan]{Dirichlet--Jordan 定理}[Dirichlet--Jordan theorem]对应的是局部有界变差条件。远处只需可积性，点 \(x\) 本身被赋予什么值也不影响结论；跳跃点恢复的是两侧值的平均。

\subsection{分段光滑函数}
\label{b2:subsec:020204}

若一个周期内有有限分割，函数在每段内部为 \(C^1\)，且函数与一阶导数均可连续延拓到该段端点，则它在每段上的变差不超过导数绝对值的积分，加上有限个跳跃仍为有界变差。于是 Fourier 级数在连续点收敛到函数值，在跳跃点收敛到中点值。这里端点是圆周上的同一个点，检查周期拼接处时必须比较 \(f(0+)\) 与 \(f(1-)\)。

\ProseParagraphBreak

例如取锯齿函数 \(f(x)=x-1/2\) 于 \(0<x<1\)。分部积分给出
\[
 \widehat f(0)=0,\quad \widehat f(k)=-\dfrac1{2\pi\mathrm{i}k}\quad(k\ne0).
\]
因此
\[
 x-\dfrac12=-\dfrac1\pi\sum_{k=1}^\infty\dfrac{\sin(2\pi kx)}k
 \quad(0<x<1).
\]
在整数处，每个对称部分和都为零，正好等于跳跃两侧的平均。连续的部分和不可能在整个圆周上一致趋向这个不连续代表；点态收敛结论与一致收敛之间仍有实质距离。

% !TeX root = ../../Book2.tex
\section{\texorpdfstring{\(L^2\)}{L²} 收敛}
\label{b2:sec:0203}

上一节在局部正则性假设下控制单点误差。平方积分理论不要求函数有单侧极限，它把全部误差汇总成一个非负积分。此时困难集中在一个问题：三角多项式是否足以逼近所有平方可积函数。

\subsection{三角系统的完备性}
\label{b2:subsec:020301}

稠密性的关键是构造一个非负、质量集中在原点附近的三角多项式核，用它平均连续函数。记 \(\mathrm{d}_{\mathbb T}(x,0)=\min_{k\in\mathbb Z}|x-k|\)。

\begin{lemma}\label{b2:lem:02-trigonometric-density}
取圆周 \(\mathbb T=\mathbb R/\mathbb Z\) 上总质量为一的 Lebesgue 测度。复三角多项式在 \(C(\mathbb T;\mathbb C)\) 的一致范数中稠密，并在每个 \(L^p(\mathbb T;\mathbb C)\)、\(1\leq p<\infty\) 中稠密。实值输入可以用实值三角多项式逼近。
\end{lemma}
\begin{proof}
取 \(B_m(t)=c_m\cdot\cos^{2m}(\pi t)\)，其中 \(c_m>0\) 使圆周积分为 \(1\)。偶数幂保证它是周期为 \(1\) 的非负三角多项式。对固定 \(0<\delta<1/2\)，比较外侧 \(|t|\geq\delta\) 与内侧 \(|t|\leq\delta/2\)，得到
\[
 \int_{\delta\leq|t|\leq1/2}B_m(t)\,\mathrm{d}t
 \leq\dfrac1\delta
 \left(\dfrac{\cos(\pi\delta)}{\cos(\pi\delta/2)}\right)^{2m}
 \longrightarrow0.
\]
若 \(f\in C(\mathbb T)\)，由一致连续性选 \(\delta\) 使 \(|f(x-t)-f(x)|<\varepsilon\) 对全部 \(x\) 及 \(|t|<\delta\) 成立。于是
\[
 |(f*B_m)(x)-f(x)|
 \leq\varepsilon+2\|f\|_\infty\int_{|t|\geq\delta}B_m(t)\,\mathrm{d}t.
\]
右端与 \(x\) 无关，而 \(f*B_m\) 是有限三角多项式，故一致稠密性成立。

对有限 \(p\)，简单函数首先在 \(L^p\) 中稠密。圆周可测集的示性函数可用有限区间并逼近，再在有限个端点附近用连续的线性过渡代替跳跃；过渡只发生在总长度可任意小的区域，函数误差不超过 \(1\)，故其 \(p\) 次积分也可任意小。因而连续周期函数在 \(L^p\) 中稠密。

现在固定 \(1\leq p<\infty\)，给定任意 \(f\in L^p(\mathbb T;\mathbb C)\) 与 \(\varepsilon>0\)。先取 \(g\in C(\mathbb T;\mathbb C)\)，使 \(\|f-g\|_p<\varepsilon/2\)；再由已经证明的一致稠密性取三角多项式 \(q\)，使 \(\|g-q\|_\infty<\varepsilon/2\)。由于圆周测度为一，
\[
 \|f-q\|_p
 \leq\|f-g\|_p+\|g-q\|_p
 \leq\|f-g\|_p+\|g-q\|_\infty
 <\varepsilon.
\]
连续函数 \(g\) 将一般 \(L^p\) 输入与一致逼近连接起来，两段误差分别在各自适用的范数中控制。若 \(f\) 实值，最后把 \(q\) 换成 \(\operatorname{Re}q\)，仍是三角多项式，并且逐点有 \(|f-\operatorname{Re}q|\leq|f-q|\)；一致逼近的实值情形也作同样处理。
\end{proof}

这个构造没有要求预先知道 \(f\) 的 Fourier 级数收敛。它只利用非负质量向原点集中，再通过有限多项式的卷积产生近似函数，因此可以用来证明三角系的完备性。

\subsection{\texorpdfstring{\(L^2\)}{L²} 收敛}
\label{b2:subsec:020302}

\begin{theorem}\label{b2:thm:02-l2-convergence}
指数函数系 \((e_k)_{k\in\mathbb Z}\) 是 \(L^2(\mathbb T)\) 的标准正交基。每个 \(f\in L^2(\mathbb T)\) 满足
\[
 \lim_{N\to\infty}\|S_Nf-f\|_2=0.
\]
\end{theorem}
\begin{proof}
标准正交性已由直接积分验证，完备性来自\cref{b2:lem:02-trigonometric-density}。也可直接检查误差：给定 \(\varepsilon>0\)，取次数为 \(M\) 的三角多项式 \(q\)，使 \(\|f-q\|_2<\varepsilon\)。对 \(N\geq M\)，\(S_Nq=q\)，且 \(S_N\) 是范数至多为 \(1\) 的正交投影，所以
\[
 \|f-S_Nf\|_2\leq\|f-q\|_2+\|S_N(q-f)\|_2<2\varepsilon.
\]
\end{proof}

范数收敛保证存在几乎处处收敛的子列：可取 \(N_j\) 使 \(\|S_{N_j}f-f\|_2^2\leq2^{-3j}\)，再估计误差大于 \(2^{-j}\) 的集合测度，所得上界 \(2^{-j}\) 可和。删去这些集合的上极限，即得逐点收敛。整列部分和的几乎处处收敛是更强的问题，单靠范数收敛无法取得。

\subsection{Parseval 恒等式}
\label{b2:subsec:020303}

\begin{corollary}[Parseval]\label{b2:cor:02-parseval}
对 \(f,g\in L^2(\mathbb T)\)，有
\[
 \|f\|_2^2=\sum_{k\in\mathbb Z}|\widehat f(k)|^2,
 \quad
 \langle f,g\rangle=\sum_{k\in\mathbb Z}\widehat f(k)\cdot\overline{\widehat g(k)}.
\]
映射 \(f\mapsto(\widehat f(k))_{k\in\mathbb Z}\) 是 \(L^2(\mathbb T)\) 到 \(\ell^2(\mathbb Z)\) 的满射线性等距映射。
\end{corollary}
\begin{proof}
第一式由完备性和\cref{b2:thm:02-abstract-parseval}得到。第二式先对有限部分和成立，再在内积中取极限；右侧的绝对收敛由序列 Cauchy--Schwarz 不等式保证。若 \((a_k)\in\ell^2\)，则有限和 \(\sum_{|k|\leq N}a_ke_k\) 在 \(L^2\) 中为 Cauchy 列。其极限 \(f\) 的每个系数为 \(a_k\)，因为取固定系数是连续线性泛函。这证明满射性。
\end{proof}

对上一节的锯齿函数，左侧为 \(\int_0^1(x-1/2)^2\,\mathrm{d}x=1/12\)，右侧为 \(2\sum_{k\geq1}(4\pi^2k^2)^{-1}\)。于是得到 \(\sum_{k\geq1}k^{-2}=\pi^2/6\)。这个数项级数来自锯齿函数的平方积分与 Fourier 系数平方和之间的等式。

\subsection{最佳逼近}
\label{b2:subsec:020304}

\begin{proposition}\label{b2:prop:02-best-approximation}
设 \(N\geq0\) 为整数。若 \(f\in L^2(\mathbb T)\)，则在所有次数不超过 \(N\) 的三角多项式中，\(S_Nf\) 唯一地使平方误差最小，且
\[
 \|f-S_Nf\|_2^2=\sum_{|k|>N}|\widehat f(k)|^2.
\]
若又有 \(\sum_k|k|^{2s}|\widehat f(k)|^2<\infty\)，其中 \(s>0\)，则
\[
 \|f-S_Nf\|_2
 \leq (N+1)^{-s}\left(\sum_k|k|^{2s}|\widehat f(k)|^2\right)^{1/2}.
\]
\end{proposition}
\begin{proof}
最小性和唯一性来自\cref{b2:prop:02-bessel}的有限误差公式。再由 Parseval 将总能量减去已保留的频率能量，得到尾和。对每个 \(|k|>N\)，有 \(|k|^{-2s}\leq(N+1)^{-2s}\)，代入尾和后取平方根。
\end{proof}

因此高频衰减直接控制有限频率近似的误差。一个函数即使有跳跃，平方误差仍可趋零；但若希望得到更快的逼近率，便需要额外的频率尾部控制。后面将把这个控制与导数联系起来。

% !TeX root = ../../Book2.tex
\section{Fejér 定理}
\label{b2:sec:0204}

对称部分和对所有频率一刀截断，所得 Dirichlet 核有正有负。若把前面各次部分和平均，高频会逐渐减轻权重，积分核也会变成非负函数。这个简单改变使连续函数的一致逼近和可积函数的范数逼近都获得直接证明。

\subsection{Fejér 核}
\label{b2:subsec:020401}

\begin{definition}\label{b2:def:02-fejer-kernel}
对非负整数 \(N\)，定义\term[fejer-he]{Fejér 核}[Fejér kernel]为
\[
 F_N(t)\coloneqq\dfrac1{N+1}\sum_{j=0}^N D_j(t)
 =\sum_{|k|\leq N}\left(1-\dfrac{|k|}{N+1}\right)e^{2\pi\mathrm{i}kt}.
\]
\end{definition}
\symidx{\(F_N\)：Fejér 核}

对平方模展开，频率 \(k\) 恰好出现 \(N+1-|k|\) 次，故
\[
 F_N(t)=\dfrac1{N+1}\left|\sum_{j=0}^N e^{2\pi\mathrm{i}jt}\right|^2
 =\dfrac1{N+1}\left(\dfrac{\sin((N+1)\pi t)}{\sin(\pi t)}\right)^2.
\]
整数处按连续性补为 \(N+1\)。因而 \(F_N\geq0\)、\(\int_{\mathbb T}F_N=1\)，且对固定 \(0<\delta<1/2\)，
\[
 \sup_{\delta\leq|t|\leq1/2}F_N(t)
 \leq\dfrac1{(N+1)\sin^2(\pi\delta)}\longrightarrow0.
\]
这里不只是质量集中，还得到离开原点后逐点统一变小的估计。

\ProseParagraphBreak

\cref{b2:fig:02-dirichlet-fejer-kernels}把同一截断次数的两个核放在相同坐标尺度下比较。Dirichlet 核依靠正负抵消使积分保持为一；Fejér 核处处非负，它在原点附近的峰与远处的小值共同分配这一个单位的质量。随 \(N\) 增长，后者的质量集中由上面的尾部估计保证。

% 自拟函数图：N=4，直接绘制有限余弦和，避免整数点处商式的0/0。
\begin{figure}[htbp]
\centering
\begin{tikzpicture}[x=5.3cm,y=0.27cm,>=stealth,font=\small]
 \begin{scope}
  \draw[->,AnalysisGray] (-0.53,0)--(0.56,0) node[right] {\(t\)};
  \draw[->,AnalysisGray] (0,-2.5)--(0,10);
  \foreach \x/\lab in {-0.5/{-1/2},-0.25/{-1/4},0.25/{1/4},0.5/{1/2}} {
   \draw[AnalysisGray] (\x,-0.15)--(\x,0.15);
   \node[AnalysisGray,anchor=north] at (\x,-2.9) {\(\lab\)};
  }
  \foreach \y in {-2,5,9}
   \draw[AnalysisGray] (-0.015,\y)--(0.015,\y) node[left=2pt] {\(\y\)};
  \draw[thick,AnalysisBlue] plot[domain=-0.5:0.5,samples=201]
   (\x,{1+2*cos(360*\x)+2*cos(720*\x)+2*cos(1080*\x)+2*cos(1440*\x)});
  \node[anchor=south west] at (-0.5,9.8) {\(D_4(t)\)};
 \end{scope}
 \begin{scope}[xshift=7cm]
  \draw[->,AnalysisGray] (-0.53,0)--(0.56,0) node[right] {\(t\)};
  \draw[->,AnalysisGray] (0,-2.5)--(0,10);
  \foreach \x/\lab in {-0.5/{-1/2},-0.25/{-1/4},0.25/{1/4},0.5/{1/2}} {
   \draw[AnalysisGray] (\x,-0.15)--(\x,0.15);
   \node[AnalysisGray,anchor=north] at (\x,-2.9) {\(\lab\)};
  }
  \foreach \y in {-2,5,9}
   \draw[AnalysisGray] (-0.015,\y)--(0.015,\y) node[left=2pt] {\(\y\)};
  \draw[thick,AnalysisBlue] plot[domain=-0.5:0.5,samples=201]
   (\x,{1+1.6*cos(360*\x)+1.2*cos(720*\x)+0.8*cos(1080*\x)+0.4*cos(1440*\x)});
  \node[anchor=south west] at (-0.5,9.8) {\(F_4(t)\)};
 \end{scope}
\end{tikzpicture}
\caption[Dirichlet 核与 Fejér 核的比较]{Dirichlet 核与 Fejér 核的比较。两图均取 \(N=4\)，横轴为一个周期 \([-1/2,1/2]\)，纵轴尺度也相同；两核的积分都为 \(1\)。左图出现负瓣，右图始终非负。图由各核的有限余弦和计算，原点处分别取 \(D_4(0)=9\)、\(F_4(0)=5\)。}
\label{b2:fig:02-dirichlet-fejer-kernels}
\end{figure}

\subsection{Cesàro 求和}
\label{b2:subsec:020402}

\begin{definition}\label{b2:def:02-cesaro}
对 \(f\in L^1(\mathbb T)\) 和非负整数 \(N\)，记
\[
 \sigma_Nf\coloneqq\dfrac1{N+1}\sum_{j=0}^N S_jf=f*F_N,
\]
称为 Fourier 级数的\term[cesaro-pingjun]{Cesàro 平均}[Cesàro mean]，也称 Fejér 平均。
\end{definition}
\symidx{\(\sigma_Nf\)：Fourier 级数的 Cesàro 平均}

若数列 \(a_j\to a\)，则其算术平均也趋于 \(a\)：先使尾部每项误差小于 \(\varepsilon\)，再将有限头部除以 \(N+1\)。因此普通收敛蕴含 Cesàro 收敛，反向一般不成立。对 Fourier 级数，这个平均不是随意重新排列项；每个 \(\widehat f(k)\) 被乘以清楚给定的三角形权重。

\subsection{正逼近恒等}
\label{b2:subsec:020403}

\begin{definition}\label{b2:def:02-positive-identity}
设 \(K_j\in L^1(\mathbb T)\)，且 \(K_j\geq0\)、\(\int K_j=1\)。若对每个 \(\delta>0\) 都有
\[
 \int_{\mathrm{d}_{\mathbb T}(t,0)\geq\delta}K_j(t)\,\mathrm{d}t\longrightarrow0,
\]
则称 \((K_j)\) 为一族\term[zheng-bijin-hengdeng]{正逼近恒等}[positive approximate identity]。
\end{definition}

\begin{proposition}\label{b2:prop:02-positive-approximation}
设 \((K_j)_{j\geq1}\) 是 \(\mathbb T\) 上的正逼近恒等，即 \(K_j\in L^1(\mathbb T)\)、\(K_j\geq0\)、\(\int_{\mathbb T}K_j=1\)，且每个 \(\delta>0\) 都有 \(\int_{\mathrm d_{\mathbb T}(t,0)\geq\delta}K_j(t)\,\mathrm dt\to0\)。则\(K_j*f\to f\) 在 \(C(\mathbb T)\) 中一致成立，也在每个 \(L^p(\mathbb T)\)、\(1\leq p<\infty\) 中按范数成立。
\end{proposition}
\begin{proof}
连续情形将
\[
 (K_j*f)(x)-f(x)=\int_{\mathbb T}K_j(t)\cdot(f(x-t)-f(x))\,\mathrm{d}t
\]
分成 \(|t|<\delta\) 和其余部分。前一部分由一致连续性小于 \(\varepsilon\)，后一部分不超过 \(2\|f\|_\infty\) 乘核的尾部质量。

对 \(1\leq p<\infty\)，Jensen 不等式与 Tonelli 给出
\[
 \|K_j*f\|_p^p
 \leq\int_{\mathbb T}\int_{\mathbb T}K_j(t)\cdot|f(x-t)|^p\,\mathrm{d}t\,\mathrm{d}x
 =\|f\|_p^p.
\]
其中先对绝对值使用正核，便保证卷积在几乎处处有定义。取连续 \(g\) 满足 \(\|f-g\|_p<\varepsilon\)，则
\[
 \|K_j*f-f\|_p\leq2\varepsilon+\|K_j*g-g\|_p.
\]
最后一项由连续情形趋零，故结论成立。
\end{proof}

\subsection{Fejér 收敛与 Fourier 系数唯一性}
\label{b2:subsec:020404}

\begin{theorem}[Fejér]\label{b2:thm:02-fejer}
设 \(\mathbb T=\mathbb R/\mathbb Z\) 取总质量为 \(1\) 的 Lebesgue 测度，函数可取实值或复值，\(\sigma_Nf=f*F_N\) 是非负整数 \(N\) 对应的 Fejér 平均。若 \(f\in C(\mathbb T)\)，则 \(\sigma_Nf\to f\) 一致成立。若 \(f\in L^p(\mathbb T)\)、\(1\leq p<\infty\)，则 \(\|\sigma_Nf-f\|_p\to0\)。此外，若 \(f\in L^1(\mathbb T)\)，并选定一个在 \(x\in\mathbb T\) 连续的代表，仍记为 \(f\)，则 \(\sigma_Nf(x)\to f(x)\)。
\end{theorem}
\begin{proof}
Fejér 核满足前一命题的全部条件，故前两项成立。单点结论须另查远处可能无界的函数。给定 \(\varepsilon>0\)，取 \(\delta\) 使 \(|f(x-t)-f(x)|<\varepsilon\) 对 \(|t|<\delta\) 成立。近处误差至多 \(\varepsilon\)；远处误差不超过
\[
 \sup_{|t|\geq\delta}F_N(t)\cdot(\|f\|_1+|f(x)|),
\]
由本节的显式估计趋零。
\end{proof}

卷积平均在每一点的值只取决于 \(L^1\) 等价类，因为改变零集上的函数值不改变积分。定理的点态结论还指定了连续代表：若只把 \(f(x)\) 单独改成另一个数，平均并不会随之改变，但新代表通常已不在该点连续。

\ProseParagraphBreak

Fejér 核的显式表达还可以量化逼近速度。设 \(n=N+1\)，在 \(0<|t|\leq1/2\) 上同时有
\[
 F_N(t)\leq n,\qquad F_N(t)\leq\frac1{4nt^2}.
\]
第一式由有限和的平方模得到，第二式使用 \(\sin(\pi t)\geq2|t|\)。若 \(|f(x-t)-f(x)|\leq L|t|^\alpha\)，其中 \(0<\alpha\leq1\)，则卷积误差不超过 \(L\int|t|^\alpha F_N(t)\,\mathrm dt\)。当 \(n\geq2\) 时在 \(t=1/n\) 分段，有
\[
 \int_{-1/2}^{1/2}|t|^\alpha F_N(t)\,\mathrm dt
 \leq 2n\int_0^{1/n}t^\alpha\,\mathrm dt
       +\frac1{2n}\int_{1/n}^{1/2}t^{\alpha-2}\,\mathrm dt.
\]
因此在 \(0<\alpha<1\) 时误差至多为 \(C_\alpha L n^{-\alpha}\)，而在 \(\alpha=1\) 时至多为 \(CL(1+\log n)/n\)。端点的对数来自核尾部的积分，不能把非端点公式中依赖 \(\alpha\) 的常数直接代入 \(\alpha=1\)。具体函数的误差可以小于这个统一上界。

\ProseParagraphBreak

例如对 \(f(x)=\cos(2\pi x)\)，平均可以直接算出 \(\sigma_Nf=(1-1/n)f\)，误差正好为 \(1/n\)。这个例子也说明，一般正核逼近的速度估计要同时看核和被逼近函数，统一上界与单个函数的实际误差不同。

\ProseParagraphBreak

这就是\term[fejer-dingli]{Fejér 定理}[Fejér's theorem]。它还给出一个不要求平方可积性的唯一性结论。

\begin{corollary}[Fourier 系数的唯一性]\label{b2:cor:02-uniqueness}
若 \(f,g\in L^1(\mathbb T)\) 的全部 Fourier 系数相同，则 \(f=g\) 几乎处处。
\end{corollary}
\begin{proof}
函数 \(h=f-g\) 的所有 Fejér 平均都为零，而这些平均在 \(L^1\) 中趋于 \(h\)，故 \(\|h\|_1=0\)。
\end{proof}

对于跳跃函数，\(L^\infty\) 范数逼近仍会失败。例如 \(f=\mathbf1_{(0,1/2)}\)，任何实值或复值连续周期函数 \(g\) 与它的本性上确界距离至少为 \(1/2\)。若 \(\|g-f\|_\infty=\varepsilon<1/2\)，则在跳跃两侧分别有 \(\operatorname{Re}g\geq1-\varepsilon\) 和 \(\operatorname{Re}g\leq\varepsilon\)。这些界先几乎处处成立，由连续性又在两侧开弧上处处成立；趋向同一跳跃点时，\(1-\varepsilon>\varepsilon\) 使两侧界不相容。正核消除了部分和的不稳定，却没有把不连续对象变成连续函数的一致极限。

% !TeX root = ../../Book2.tex
\section{点态、一致与绝对收敛}
\label{b2:sec:0205}

Fourier 系数唯一决定可积函数，并不意味着对称部分和在每种意义下都恢复它。本节逐项说明点态、一致和绝对收敛分别需要什么控制，尤其把控制函数值与控制导数两种方法联系起来。

\subsection{点态收敛}
\label{b2:subsec:020501}

\begin{theorem}[Dini 判据]\label{b2:thm:02-dini}
设 \(f\in L^1(\mathbb T)\)，\(x\in\mathbb T\)、\(A\in\mathbb C\)。若存在 \(0<\delta<1/2\) 使
\[
 \int_0^\delta\dfrac{|f(x+t)+f(x-t)-2A|}{t}\,\mathrm{d}t<\infty,
\]
则 \(S_Nf(x)\to A\)。
\end{theorem}
\begin{proof}
令 \(g(t)=f(x+t)+f(x-t)-2A\)。核的偶性使误差等于 \(\int_0^{1/2}g(t)\cdot D_N(t)\,\mathrm{d}t\)。由 \(\sin(\pi t)\geq2t\)，假设保证 \(g/\sin(\pi t)\) 在原点附近可积；远离原点时，可积性由 \(f\in L^1\) 给出。对整个区间应用\cref{b2:lem:02-oscillation}。
\end{proof}

上述条件称为\term[dini-panju]{Dini 判据}[Dini's criterion]。若在 \(x\) 附近有 \(|f(x+t)-f(x)|\leq C|t|^\alpha\)，其中 \(\alpha>0\)，则可取 \(A=f(x)\)，因为 \(t^{\alpha-1}\) 可积。另一方面，仅知道模连续度趋零并不足以直接套用这个判据：例如 \(1/|\log t|\) 虽趋零，除以 \(t\) 后的积分仍发散。判据不满足并不等于级数必发散。

\subsection{一致收敛}
\label{b2:subsec:020502}

\begin{proposition}\label{b2:prop:02-uniform-dini}
设 \(f\in C(\mathbb T)\)，模连续度
\[
 \omega_f(t)\coloneqq\sup_{\mathrm{d}_{\mathbb T}(x,y)\leq t}|f(x)-f(y)|
\]
满足 \(\int_0^{1/2}\omega_f(t)t^{-1}\,\mathrm{d}t<\infty\)，则 \(S_Nf\to f\) 一致成立。
\end{proposition}
\begin{proof}
误差仍为对称差与 \(D_N\) 的积分。在 \((0,\delta)\) 上，绝对误差至多 \(\int_0^\delta\omega_f(t)t^{-1}\,\mathrm{d}t\)，因而对 \(x,N\) 一致趋小。对固定 \(\delta\)，函数族
\[
 h_x(t)=\dfrac{f(x+t)+f(x-t)-2f(x)}{\sin(\pi t)}\quad(\delta\leq t\leq1/2)
\]
在 \(L^1\) 中是紧集：映射 \(x\mapsto h_x\) 由 \(f\) 的一致连续性而连续，参数圆周紧。给这个紧集取有限 \(\varepsilon\)-网，对网中每个函数应用振荡积分引理；任意其他函数的积分误差至多为其 \(L^1\) 距离。这使远处的积分也对全部 \(x\) 一致趋零。
\end{proof}

例如任何正指数 Hölder 连续周期函数都满足此条件。证明中的有限网补足了从“每个 \(x\) 分别收敛”到“对全部 \(x\) 一致收敛”的量词变化，不能只把单点判据逐点使用后省去这一步。

\subsection{绝对收敛}
\label{b2:subsec:020503}

\begin{proposition}\label{b2:prop:02-absolute}
若 \(f\in L^1(\mathbb T)\) 且 \(\sum_k|\widehat f(k)|<\infty\)，则其 Fourier 级数绝对一致收敛到某个连续周期函数 \(g\)，且 \(f=g\) 几乎处处。若 \(f\) 本身连续，则处处相等。
\end{proposition}
\begin{proof}
每个指数因子的模为 \(1\)，故绝对可和的系数直接控制一致尾部。极限 \(g\) 连续，逐项积分表明它与 \(f\) 的全部系数一致。由\cref{b2:cor:02-uniqueness}，二者几乎处处相等。若两者连续而在一点不同，连续性会使差在某个开弧上与零保持正距离，与几乎处处相等矛盾。
\end{proof}

例如给定 \(0<r<1\)，几何级数相加得到
\[
 \sum_{k\in\mathbb Z}r^{|k|}e^{2\pi\mathrm{i}kx}
 =\dfrac{1-r^2}{1-2r\cos(2\pi x)+r^2}.
\]
这个核记为 \(P_r(x)\)，称为圆周上的\term[poisson-he]{Poisson 核}[Poisson kernel]。对每个 \(0<r<1\)，它非负、积分为 \(1\)。这里的逼近方向是 \(r\uparrow1\)：对任意固定的 \(0<\delta<1/2\)，分母可写成 \((1-r)^2+2r(1-\cos(2\pi x))\)，所以
\[
 \sup_{\mathrm d_{\mathbb T}(x,0)\geq\delta}P_r(x)
 \leq\frac{1-r^2}{2r(1-\cos(2\pi\delta))}
 \longrightarrow0\quad(r\uparrow1).
\]
圆周的总测度为一，此估计也控制原点以外的尾部质量。因此任取 \(r_j\in(0,1)\)、\(r_j\to1\)，特别是任取 \(r_j\uparrow1\)，序列 \((P_{r_j})\) 都满足\cref{b2:def:02-positive-identity}。由\cref{b2:prop:02-positive-approximation}，\(P_r*f\to f\) 对连续周期函数一致成立，对 \(f\in L^p(\mathbb T)\)、\(1\leq p<\infty\) 按 \(L^p\) 范数成立。因为任意趋于一的参数序列都得到同一极限，这些结论是完整的 \(r\uparrow1\) 极限，不只针对某一个预先选定的序列。
\symidx{\(P_r\)：圆周 Poisson 核}

\ProseParagraphBreak

固定 \(0<r<1\) 时，Poisson 平均还提高了光滑性。对实值或复值 \(f\in L^1(\mathbb T)\)，核的绝对一致收敛级数允许交换求和与积分，给出
\[
 (P_r*f)(x)=\sum_{k\in\mathbb Z}r^{|k|}\widehat f(k)e^{2\pi\mathrm i kx}.
\]
系数满足 \(|\widehat f(k)|\leq\|f\|_1\)，而对每个整数 \(m\geq0\)，都有 \(\sum_k(1+|k|)^m r^{|k|}<\infty\)：正负频率的尾部都由比值极限为 \(r<1\) 的级数控制。因此各阶形式导数级数绝对一致收敛。先把有限部分和的导数积分回原函数，再取一致极限，就识别出这些极限是实际导数。于是 \(P_r*f\in C^\infty(\mathbb T)\)。固定 \(r<1\) 的光滑化与 \(r\uparrow1\) 时恢复原函数，是两个不同的结论；其导数界可以随 \(r\) 趋于一而增大。

\subsection{光滑性与 Fourier 系数衰减}
\label{b2:subsec:020504}

\begin{proposition}\label{b2:prop:02-smooth-decay}
设 \(m\geq1\) 为整数。若 \(f\) 是周期为一的 \(C^m\) 实值或复值函数，则
\[
 \widehat {f^{(m)}}(k)=(2\pi\mathrm{i}k)^m\widehat f(k),
 \quad |\widehat f(k)|\leq\dfrac{\|f^{(m)}\|_1}{(2\pi|k|)^m}\quad(k\ne0).
\]
若 \(f\) 绝对连续、\(f(0)=f(1)\) 且 \(f'\in L^2(0,1)\)，则 \(\sum_k|\widehat f(k)|<\infty\)，从而其 Fourier 级数一致收敛到 \(f\)。
\end{proposition}
\begin{proof}
逐次分部积分，周期端点的边界项相消，给出第一式。取绝对值便得到系数估计。绝对连续情形同样可作一次分部积分；再用 Cauchy--Schwarz 与 Bessel 不等式，
\[
 \sum_{k\ne0}|\widehat f(k)|
 \leq\left(\sum_{k\ne0}\dfrac1{4\pi^2k^2}\right)^{1/2}
       \left(\sum_{k\ne0}|\widehat {f'}(k)|^2\right)^{1/2}
 \leq\dfrac1{\sqrt{12}}\|f'\|_2.
\]
绝对收敛结论随后由前一命题给出。
\end{proof}

因而 \(C^1\) 已足以保证绝对一致收敛，尽管对每个系数分别估计只给出 \(O(1/|k|)\)。改善来自把全部导数系数一起放进平方和，而不是逐个使用相同上界。若在周期拼接处发生跳跃，分部积分会留下边界项，通常正是 \(1/k\) 衰减的来源。

\ProseParagraphBreak

把一个可跟算的折角例子与上一节的锯齿比较。令 \(v(x)=|x|\) 于 \([-1/2,1/2]\)，再作周期延拓。它连续，原点和半整数处有折角，但没有跳跃。偶性使正弦系数为零，且 \(\widehat v(0)=1/4\)。对非零整数 \(k\)，两次一元积分给出
\[
 \begin{aligned}
 \widehat v(k)&=2\int_0^{1/2}x\cos(2\pi kx)\,\mathrm dx\\
 &=\left[\frac{x\sin(2\pi kx)}{\pi k}
          +\frac{\cos(2\pi kx)}{2\pi^2k^2}\right]_0^{1/2}
 =\frac{(-1)^k-1}{2\pi^2k^2}.
 \end{aligned}
\]
这里积分边界的正弦项消失，而斜率在端点的变化留下 \(k^{-2}\) 项。由于系数绝对可和，\cref{b2:prop:02-absolute}给出每个点都成立的展开
\[
 v(x)=\frac14-\frac2{\pi^2}
       \sum_{\substack{k\geq1\\k\text{ 奇}}}\frac{\cos(2\pi kx)}{k^2}.
\]
在原点代入得 \(\sum_{k\geq1,\,k\text{ 奇}}k^{-2}=\pi^2/8\)。连续折角对应的平方衰减比锯齿跳跃的 \(1/k\) 快，但形式导数的系数仍只有 \(1/k\)；在折角处不能据原级数的一致收敛断言可微。判断衰减时，函数值的跳跃与导数的跳跃需要分别计入。

\ProseParagraphBreak

\cref{b2:tab:02-summation-convergence}集中比较三种求和的条件与结论。表中均使用总质量为一的圆周 Lebesgue 测度，函数可取实值或复值；点态行固定 \(x\in\mathbb T\)，局部值和单侧极限应对所选代表理解。条件的强弱不能只按名字排列：局部变差控制普通部分和的抵消，正核则能在不要求局部变差的情况下给出有限 \(p\) 的范数逼近。

\begin{table}[htbp]
\centering
\caption[三种 Fourier 求和的条件与收敛方式]{三种 Fourier 求和的条件与收敛方式。每行只列已经建立的充分条件；点态、范数和一致收敛分别判断。}
\label{b2:tab:02-summation-convergence}
\begin{tabularx}{\textwidth}{@{}p{0.13\textwidth}p{0.29\textwidth}p{0.28\textwidth}>{\raggedright\arraybackslash}X@{}}
\toprule
求和方法 & 条件 & 收敛方式与极限 & 适用范围及端点限制 \\
\midrule
\(S_Nf\) & \(f\in L^2(\mathbb T)\) & \(\|S_Nf-f\|_2\to0\) & 平方范数收敛本身不指定逐点极限，也不提供 \(L^1\) 或 \(L^\infty\) 范数结论。 \\
\addlinespace
\(S_Nf(x)\) & \(f\in L^1(\mathbb T)\)，所选代表在 \(x\) 两侧局部有界变差；或\cref{b2:thm:02-dini}的积分条件对某个 \(A\in\mathbb C\) 成立 & 前一条件给出 \((f(x+)+f(x-))/2\)；后一条件给出 \(A\) & 改动单点不改变部分和；跳跃点的极限不由任意指定的 \(f(x)\) 决定。 \\
\addlinespace
\(\sigma_Nf\) & \(f\in C(\mathbb T)\)；或 \(f\in L^p(\mathbb T)\)、\(1\leq p<\infty\) & 分别一致趋于 \(f\)、按 \(L^p\) 范数趋于 \(f\)；\(L^1\) 代表在 \(x\) 连续时还在该点收敛 & 不能把有限 \(p\) 的结论扩到全部 \(L^\infty\)；跳跃函数与连续平均的本性上确界距离可保持正值。 \\
\addlinespace
\(P_r*f\) & \(f\in C(\mathbb T)\)；或 \(f\in L^p(\mathbb T)\)、\(1\leq p<\infty\)，参数 \(0<r<1\)、\(r\uparrow1\) & 分别一致趋于 \(f\)、按 \(L^p\) 范数趋于 \(f\)；固定 \(0<r<1\) 时为光滑函数 & 这里取核参数极限；对跳跃函数，\(L^\infty\) 范数逼近同样会失败。 \\
\bottomrule
\end{tabularx}
\end{table}

表中最后两行的 \(L^\infty\) 限制可用同一个例子核验：若 \(f=\mathbf1_{(0,1/2)}\)，任一连续周期函数与它的本性上确界距离至少为 \(1/2\)。Fejér 平均和固定 \(0<r<1\) 的 Poisson 平均都连续，因此即使它们在每个有限 \(p\) 的范数中逼近 \(f\)，也不能在 \(L^\infty\) 范数中逼近。对普通部分和，若还要求一致或绝对收敛，应另查本节的模连续度积分条件或系数绝对可和条件，不能从平方范数收敛直接推出。

% !TeX root = ../../Book2.tex
\section{函数项级数与逐项运算}
\label{b2:sec:0206}

一旦 Fourier 级数被用于求解方程，就会遇到逐项积分与逐项微分。积分通常改善系数的衰减，微分则放大高频。两种运算不能用同一收敛假设来保证。

\subsection{逐项积分}
\label{b2:subsec:020601}

\begin{proposition}\label{b2:prop:02-integrated-series}
设 \(f\in L^2(\mathbb T)\)，且 \(\widehat f(0)=0\)。函数
\[
 A(x)=\int_0^x f(t)\,\mathrm{d}t
\]
可周期延拓，且
\[
 A(x)=\sum_{k\ne0}\widehat f(k)\cdot\dfrac{e^{2\pi\mathrm{i}kx}-1}{2\pi\mathrm{i}k}
 \quad(0\leq x\leq1),
\]
右侧绝对一致收敛。
\end{proposition}
\begin{proof}
零平均保证 \(A(1)=A(0)\)。对部分和积分，有限运算给出右侧的截断和。另一方面，
\[
 \sup_{0\leq x\leq1}\left|\int_0^x(f-S_Nf)(t)\,\mathrm{d}t\right|
 \leq\|f-S_Nf\|_1\leq\|f-S_Nf\|_2\longrightarrow0.
\]
最后，\(\sum_{k\ne0}|\widehat f(k)|/|k|<\infty\) 由序列 Cauchy--Schwarz 与 Parseval 保证，故所写级数还绝对一致收敛。
\end{proof}

若平均不为零，应先减去 \(\widehat f(0)\)，原函数再加上线性项 \(x\widehat f(0)\)。这个线性项不能周期化；它准确记录了周期原函数存在的相容条件。

\subsection{逐项求导}
\label{b2:subsec:020602}

\begin{proposition}\label{b2:prop:02-termwise-differentiation}
设 \((a_k)_{k\in\mathbb Z}\) 为复数列，\(m\geq1\) 为整数。若 \(\sum_{k\in\mathbb Z}|k|^m|a_k|<\infty\)，则级数 \(\sum_ka_ke^{2\pi\mathrm{i}kx}\) 定义一个 \(C^m\) 周期函数，其不超过 \(m\) 阶的导数可逐项计算。
\end{proposition}
\begin{proof}
除 \(k=0\) 的一项外，假设控制每个不超过 \(m\) 阶的导数级数，所以这些级数一致收敛。为识别极限，先看一阶：部分和满足 \(S_N(x)-S_N(0)=\int_0^xS_N'(t)\,\mathrm{d}t\)。在两侧取一致极限，得到极限函数等于其原点值加导数级数极限的积分，故可微且导数就是该极限。对导数级数重复这个论证，直到 \(m\) 阶。
\end{proof}

例如 \(\sum_{k\geq1}k^{-2}\cos(2\pi kx)\) 绝对一致收敛，但形式导数的系数仅为 \(O(1/k)\)，不能再用绝对一致收敛作保证。原级数收敛良好并不自动许可任意次逐项微分。

\subsection{弱意义下的求导}
\label{b2:subsec:020603}

这里沿用\subsecref{b2:subsec:010405}的弱导数定义，并将测试函数取为光滑周期函数：若 \(f,g\in L^2(\mathbb T)\) 且对全部 \(\varphi\in C^\infty(\mathbb T)\) 有 \(\int f\varphi'=-\int g\varphi\)，则记 \(f'=g\)。积分配对不加复共轭。

\begin{theorem}\label{b2:thm:02-weak-derivative}
设 \(f\in L^2(\mathbb T)\)。存在平方可积的周期弱导数，当且仅当
\[
 \sum_{k\in\mathbb Z}k^2|\widehat f(k)|^2<\infty.
\]
此时 \(\widehat {f'}(k)=2\pi\mathrm{i}k\widehat f(k)\)，且 \((S_Nf)'\to f'\) 于 \(L^2\)。
\end{theorem}
\begin{proof}
若弱导数 \(g\) 存在，将 \(e^{-2\pi\mathrm{i}kx}\) 代入测试恒等式，便得系数关系，Parseval 随即给出必要条件。

反之，由平方可和性和\cref{b2:cor:02-parseval}的满射性，存在 \(g\in L^2\)，系数为 \(2\pi\mathrm{i}k\widehat f(k)\)。有限部分和满足 \((S_Nf)'=S_Ng\)。对光滑周期 \(\varphi\) 分部积分，再利用 \(S_Nf\to f\)、\(S_Ng\to g\) 的平方范数收敛，得到 \(\int f\varphi'=-\int g\varphi\)。导数部分和的收敛也同时成立。
\end{proof}

这个结论还可以给出函数的具体代表，而不仅是一个带权序列。设 \(g=f'\in L^2(\mathbb T)\)。以常函数测试，得到 \(\int_{\mathbb T}g=0\)。因此 \(A(x)=\int_0^x g(t)\,\mathrm dt\) 的端点值相同，可以连续地周期延拓。积分的绝对连续性和分部积分表明它的周期弱导数为 \(g\)。对 \(f-A\) 的非零 Fourier 系数使用前一定理，全部得到零；再由系数唯一性，存在一个常数 \(c\)，使 \(f=A+c\) 几乎处处。

\ProseParagraphBreak

所得代表在每一段不跨越周期接缝的区间上绝对连续。沿圆周上较短的弧从 \(x\) 积分到 \(y\)，Cauchy--Schwarz 不等式给出
\[
 |f(y)-f(x)|\leq\|g\|_2\cdot\mathrm d_{\mathbb T}(x,y)^{1/2}.
\]
若该弧经过接缝，就把积分拆成两段；两端值相同，仍得到同一式。因此平方可积的一阶弱导数确实给出连续代表，而含跳跃的锯齿不可能属于这个空间。折角函数 \(v=\mathrm d_{\mathbb T}(\cdot,0)\) 却属于它：斜率为分段常数，属于 \(L^2\)，斜率本身的不连续不妨碍一阶弱导数存在。

\ProseParagraphBreak

从此可把周期一阶 Sobolev 空间写成系数带权平方可和的空间。函数值与弱导数不是逐点对每一项取极限得来的，而是通过测试恒等式和范数收敛共同确定。

\subsection{正则化与逼近}
\label{b2:subsec:020604}

对实值或复值 \(f\in L^p(\mathbb T)\)、\(1\leq p<\infty\)，Fejér 平均是三角多项式，并在 \(N\to\infty\) 时按 \(L^p\) 范数趋于 \(f\)。Poisson 平均 \(P_r*f\) 对每个 \(0<r<1\) 为光滑函数，也在 \(r\uparrow1\) 时按 \(L^p\) 范数趋于 \(f\)。前者由\cref{b2:thm:02-fejer}给出，后者由\subsecref{b2:subsec:020503}中已核验的正逼近恒等性质给出。

\ProseParagraphBreak

如果还需要逼近导数，就必须使用导数本身的可积性。具有平方可积周期弱导数的实值或复值 \(L^2(\mathbb T)\) 函数组成周期空间 \(H^1(\mathbb T)\)，规定
\[
 \|f\|_{H^1(\mathbb T)}
 \coloneqq\left(\|f\|_2^2+\|f'\|_2^2\right)^{1/2}.
\]
这里的弱导数以全部光滑周期函数为测试函数，采用\subsecref{b2:subsec:020603}的不加复共轭配对。
\symidx{\(H^1(\mathbb T)\)：具有平方可积周期弱导数的空间}

\ProseParagraphBreak

设 \(f\in H^1(\mathbb T)\)、\(g=f'\in L^2(\mathbb T)\)。由\cref{b2:thm:02-weak-derivative}，\(\widehat g(k)=2\pi\mathrm i k\widehat f(k)\)。对 Fejér 平均的有限和求导，得到
\[
 \begin{aligned}
 (\sigma_Nf)'(x)
 &=\sum_{|k|\leq N}\left(1-\frac{|k|}{N+1}\right)
      2\pi\mathrm i k\widehat f(k)e^{2\pi\mathrm i kx}\\
 &=\sigma_Ng(x).
 \end{aligned}
\]
对每个固定 \(0<r<1\)，\(f,g\in L^1(\mathbb T)\)，而 Poisson 权重的指数衰减使下面两个级数绝对一致收敛。按\subsecref{b2:subsec:020503}已证明的逐项求导合法性，有
\[
 (P_r*f)'(x)
 =\sum_{k\in\mathbb Z}r^{|k|}\cdot2\pi\mathrm i k\widehat f(k)e^{2\pi\mathrm i kx}
 =(P_r*g)(x).
\]
这些等式表明正则化与求导相交换；其依据是两种正则化都只给第 \(k\) 个系数乘一个标量权重，而弱求导乘以 \(2\pi\mathrm i k\)。把前面的 \(L^2\) 逼近分别用于 \(f\) 和 \(g\)，便得到
\[
 \|\sigma_Nf-f\|_2+\|(\sigma_Nf)'-f'\|_2
 \longrightarrow0\quad(N\to\infty),
\]
以及
\[
 \|P_r*f-f\|_2+\|(P_r*f)'-f'\|_2
 \longrightarrow0\quad(r\uparrow1).
\]
因此两种平均都在 \(H^1(\mathbb T)\) 范数中逼近 \(f\)，导数的逼近并不是由函数的 \(L^2\) 收敛单独推出的。

\ProseParagraphBreak

对一般实值或复值 \(f\in L^1(\mathbb T)\)，对称部分和即使不按 \(L^1\) 范数收敛，也仍满足
\[
 \int_{\mathbb T}S_Nf\cdot\varphi\,\mathrm{d}x
 =\int_{\mathbb T}f\cdot S_N\varphi\,\mathrm{d}x
 \longrightarrow\int_{\mathbb T}f\cdot\varphi\,\mathrm{d}x
 \quad(\varphi\in C^\infty(\mathbb T)).
\]
等式来自有限展开与核的偶性；极限来自光滑 \(\varphi\) 的 Fourier 级数一致收敛。任给整数 \(m\geq0\)，对光滑三角多项式 \(S_Nf\) 分部积分，周期端点项相消，给出
\[
 \begin{aligned}
 \int_{\mathbb T}(S_Nf)^{(m)}\varphi\,\mathrm dx
 &=(-1)^m\int_{\mathbb T}S_Nf\cdot\varphi^{(m)}\,\mathrm dx\\
 &=(-1)^m\int_{\mathbb T}f\cdot S_N(\varphi^{(m)})\,\mathrm dx\\
 &\longrightarrow(-1)^m\int_{\mathbb T}f\cdot\varphi^{(m)}\,\mathrm dx.
 \end{aligned}
\]
最后一步的误差至多为
\(\|f\|_1\cdot\|S_N(\varphi^{(m)})-\varphi^{(m)}\|_\infty\)，它趋于零，因为 \(\varphi^{(m)}\) 仍是光滑周期函数，具有绝对一致收敛的 Fourier 级数。这就给出形式微分级数在每个光滑周期测试函数上的极限，即 \((S_Nf)^{(m)}\) 趋于 \(f\) 的 \(m\) 阶分布导数。这种收敛按测试函数配对理解；分布导数未必由局部可积函数表示，微分后的部分和也未必逐点收敛。测试函数的光滑性吸收了频率增长，这一机制将在\chapref{b2:ch:05}系统化。

% !TeX root = ../../Book2.tex
\OptionalSection{Fourier 级数的深层收敛现象}{b2:sec:0207}

跳跃附近的 Fourier 部分和会发生局部过冲，Dirichlet 核绝对质量的增长会使算子范数增大；这些现象揭示了逐点逼近与整体逼近的差别。另一方面，几乎处处收敛还可以在更广的函数类中成立。比较这些结论，能看清局部正则性、积分范数与收敛方式各自的作用。

\subsection{Gibbs 现象}
\label{b2:subsec:020701}

仍取 \(f(x)=x-1/2\) 于 \((0,1)\)。固定 \(c>0\)，在距跳跃约 \(1/N\) 的点 \(x_N=c/(2\pi N)\) 处，
\[
 S_Nf(x_N)=-\dfrac1\pi\sum_{k=1}^N\dfrac{\sin(kc/N)}k
 \longrightarrow-\dfrac1\pi\int_0^c\dfrac{\sin u}{u}\,\mathrm{d}u.
\]
这是连续函数 \(\sin u/u\)（原点补为 \(1\)）的 Riemann 和。另一方面，\(f(x_N)\to-1/2\)。当 \(c=\pi\) 时，右侧约为 \(-0.58949\)，超出了右侧函数值的极限。这种过冲称为\term[gibbs-xianxiang]{Gibbs 现象}[Gibbs phenomenon]。

\ProseParagraphBreak

为看清过冲确实为正，先注意 \(\int_0^\infty\sin u/u\,\mathrm{d}u=\pi/2\)。可在积分中加入 \(e^{-\varepsilon u}\)，对频率参数求导后积分得到 \(\arctan(1/\varepsilon)\)，再由分部积分的一致尾部界令 \(\varepsilon\downarrow0\)。从 \(\pi\) 开始，正负半波的绝对积分逐段减小，第一段为负，故剩余尾积分严格为负。因此 \(\pi^{-1}\int_0^\pi\sin u/u\,\mathrm{d}u-1/2>0\)。过冲的高度不消失，发生过冲的区域却不断收缩；这与平方误差趋零并不矛盾。

\ProseParagraphBreak

这一计算对应于\cref{b2:fig:02-gibbs}。图中的两种截断使用同一个坐标尺度，可同时比较过冲的高度与区域宽度。
% 本书原创：按第三章锯齿函数部分和精确公式生成。
\begin{figure}[htbp]
\centering
\begin{tikzpicture}[x=1cm,y=1cm,>=stealth,every node/.style={font=\small},line width=.65pt]
\begin{scope}[xshift=0.0cm]
\draw[->,gray] (-.1,2.275)--(5.6,2.275) node[right]{\(x\)};
\draw[->,gray] (2.64,-.05)--(2.64,4.7);
\node at (2.64,5.0) {\(N=8\)};
\draw[gray] (2.57,0.5250)--(2.71,0.5250);\node[left] at (-0.12,0.5250) {\(-1/2\)};
\draw[gray] (2.57,2.2750)--(2.71,2.2750);\node[left] at (-0.12,2.2750) {\(0\)};
\draw[gray] (2.57,4.0250)--(2.71,4.0250);\node[left] at (-0.12,4.0250) {\(1/2\)};
\draw[gray,dashed,thick] plot coordinates {(0.00000,3.60500) (0.05387,3.61357) (0.10775,3.62214) (0.16162,3.63071) (0.21549,3.63928) (0.26937,3.64785) (0.32324,3.65642) (0.37711,3.66500) (0.43098,3.67357) (0.48486,3.68214) (0.53873,3.69071) (0.59260,3.69928) (0.64648,3.70785) (0.70035,3.71642) (0.75422,3.72499) (0.80810,3.73356) (0.86197,3.74213) (0.91584,3.75070) (0.96972,3.75927) (1.02359,3.76784) (1.07746,3.77641) (1.13133,3.78498) (1.18521,3.79356) (1.23908,3.80213) (1.29295,3.81070) (1.34683,3.81927) (1.40070,3.82784) (1.45457,3.83641) (1.50845,3.84498) (1.56232,3.85355) (1.61619,3.86212) (1.67006,3.87069) (1.72394,3.87926) (1.77781,3.88783) (1.83168,3.89640) (1.88556,3.90498) (1.93943,3.91355) (1.99330,3.92212) (2.04718,3.93069) (2.10105,3.93926) (2.15492,3.94783) (2.20880,3.95640) (2.26267,3.96497) (2.31654,3.97354) (2.37041,3.98211) (2.42429,3.99068) (2.47816,3.99925) (2.53203,4.00782) (2.58591,4.01639) (2.63978,4.02496)};
\draw[gray,dashed,thick] plot coordinates {(2.64022,0.52504) (2.69409,0.53361) (2.74797,0.54218) (2.80184,0.55075) (2.85571,0.55932) (2.90959,0.56789) (2.96346,0.57646) (3.01733,0.58503) (3.07120,0.59360) (3.12508,0.60217) (3.17895,0.61074) (3.23282,0.61931) (3.28670,0.62788) (3.34057,0.63645) (3.39444,0.64502) (3.44832,0.65360) (3.50219,0.66217) (3.55606,0.67074) (3.60994,0.67931) (3.66381,0.68788) (3.71768,0.69645) (3.77155,0.70502) (3.82543,0.71359) (3.87930,0.72216) (3.93317,0.73073) (3.98705,0.73930) (4.04092,0.74787) (4.09479,0.75644) (4.14867,0.76502) (4.20254,0.77359) (4.25641,0.78216) (4.31028,0.79073) (4.36416,0.79930) (4.41803,0.80787) (4.47190,0.81644) (4.52578,0.82501) (4.57965,0.83358) (4.63352,0.84215) (4.68740,0.85072) (4.74127,0.85929) (4.79514,0.86786) (4.84902,0.87643) (4.90289,0.88501) (4.95676,0.89358) (5.01063,0.90215) (5.06451,0.91072) (5.11838,0.91929) (5.17225,0.92786) (5.22613,0.93643) (5.28000,0.94500)};
\draw[AnalysisBlue,thick] plot coordinates {(0.00000,3.43226) (0.02039,3.43461) (0.04077,3.43740) (0.06116,3.44063) (0.08154,3.44430) (0.10193,3.44843) (0.12232,3.45301) (0.14270,3.45806) (0.16309,3.46357) (0.18347,3.46954) (0.20386,3.47598) (0.22425,3.48289) (0.24463,3.49027) (0.26502,3.49811) (0.28541,3.50640) (0.30579,3.51515) (0.32618,3.52435) (0.34656,3.53399) (0.36695,3.54407) (0.38734,3.55456) (0.40772,3.56548) (0.42811,3.57679) (0.44849,3.58849) (0.46888,3.60057) (0.48927,3.61300) (0.50965,3.62579) (0.53004,3.63890) (0.55042,3.65231) (0.57081,3.66602) (0.59120,3.68000) (0.61158,3.69423) (0.63197,3.70868) (0.65236,3.72334) (0.67274,3.73818) (0.69313,3.75317) (0.71351,3.76829) (0.73390,3.78351) (0.75429,3.79881) (0.77467,3.81415) (0.79506,3.82951) (0.81544,3.84486) (0.83583,3.86016) (0.85622,3.87539) (0.87660,3.89052) (0.89699,3.90552) (0.91737,3.92034) (0.93776,3.93497) (0.95815,3.94936) (0.97853,3.96348) (0.99892,3.97731) (1.01931,3.99080) (1.03969,4.00393) (1.06008,4.01666) (1.08046,4.02895) (1.10085,4.04077) (1.12124,4.05209) (1.14162,4.06288) (1.16201,4.07310) (1.18239,4.08272) (1.20278,4.09171) (1.22317,4.10004) (1.24355,4.10766) (1.26394,4.11456) (1.28432,4.12071) (1.30471,4.12606) (1.32510,4.13060) (1.34548,4.13429) (1.36587,4.13711) (1.38625,4.13903) (1.40664,4.14002) (1.42703,4.14007) (1.44741,4.13914) (1.46780,4.13721) (1.48819,4.13426) (1.50857,4.13027) (1.52896,4.12523) (1.54934,4.11910) (1.56973,4.11188) (1.59012,4.10355) (1.61050,4.09410) (1.63089,4.08351) (1.65127,4.07176) (1.67166,4.05886) (1.69205,4.04479) (1.71243,4.02954) (1.73282,4.01312) (1.75320,3.99550) (1.77359,3.97670) (1.79398,3.95670) (1.81436,3.93552) (1.83475,3.91315) (1.85514,3.88959) (1.87552,3.86485) (1.89591,3.83895) (1.91629,3.81187) (1.93668,3.78364) (1.95707,3.75427) (1.97745,3.72377) (1.99784,3.69215) (2.01822,3.65943) (2.03861,3.62563) (2.05900,3.59077) (2.07938,3.55486) (2.09977,3.51793) (2.12015,3.48000) (2.14054,3.44110) (2.16093,3.40125) (2.18131,3.36049) (2.20170,3.31883) (2.22208,3.27632) (2.24247,3.23297) (2.26286,3.18883) (2.28324,3.14392) (2.30363,3.09829) (2.32402,3.05197) (2.34440,3.00499) (2.36479,2.95739) (2.38517,2.90921) (2.40556,2.86050) (2.42595,2.81128) (2.44633,2.76161) (2.46672,2.71151) (2.48710,2.66105) (2.50749,2.61025) (2.52788,2.55916) (2.54826,2.50783) (2.56865,2.45630) (2.58903,2.40461) (2.60942,2.35281) (2.62981,2.30095) (2.65019,2.24905) (2.67058,2.19719) (2.69097,2.14539) (2.71135,2.09370) (2.73174,2.04217) (2.75212,1.99084) (2.77251,1.93975) (2.79290,1.88895) (2.81328,1.83849) (2.83367,1.78839) (2.85405,1.73872) (2.87444,1.68950) (2.89483,1.64079) (2.91521,1.59261) (2.93560,1.54501) (2.95598,1.49803) (2.97637,1.45171) (2.99676,1.40608) (3.01714,1.36117) (3.03753,1.31703) (3.05792,1.27368) (3.07830,1.23117) (3.09869,1.18951) (3.11907,1.14875) (3.13946,1.10890) (3.15985,1.07000) (3.18023,1.03207) (3.20062,0.99514) (3.22100,0.95923) (3.24139,0.92437) (3.26178,0.89057) (3.28216,0.85785) (3.30255,0.82623) (3.32293,0.79573) (3.34332,0.76636) (3.36371,0.73813) (3.38409,0.71105) (3.40448,0.68515) (3.42486,0.66041) (3.44525,0.63685) (3.46564,0.61448) (3.48602,0.59330) (3.50641,0.57330) (3.52680,0.55450) (3.54718,0.53688) (3.56757,0.52046) (3.58795,0.50521) (3.60834,0.49114) (3.62873,0.47824) (3.64911,0.46649) (3.66950,0.45590) (3.68988,0.44645) (3.71027,0.43812) (3.73066,0.43090) (3.75104,0.42477) (3.77143,0.41973) (3.79181,0.41574) (3.81220,0.41279) (3.83259,0.41086) (3.85297,0.40993) (3.87336,0.40998) (3.89375,0.41097) (3.91413,0.41289) (3.93452,0.41571) (3.95490,0.41940) (3.97529,0.42394) (3.99568,0.42929) (4.01606,0.43544) (4.03645,0.44234) (4.05683,0.44996) (4.07722,0.45829) (4.09761,0.46728) (4.11799,0.47690) (4.13838,0.48712) (4.15876,0.49791) (4.17915,0.50923) (4.19954,0.52105) (4.21992,0.53334) (4.24031,0.54607) (4.26069,0.55920) (4.28108,0.57269) (4.30147,0.58652) (4.32185,0.60064) (4.34224,0.61503) (4.36263,0.62966) (4.38301,0.64448) (4.40340,0.65948) (4.42378,0.67461) (4.44417,0.68984) (4.46456,0.70514) (4.48494,0.72049) (4.50533,0.73585) (4.52571,0.75119) (4.54610,0.76649) (4.56649,0.78171) (4.58687,0.79683) (4.60726,0.81182) (4.62764,0.82666) (4.64803,0.84132) (4.66842,0.85577) (4.68880,0.87000) (4.70919,0.88398) (4.72958,0.89769) (4.74996,0.91110) (4.77035,0.92421) (4.79073,0.93700) (4.81112,0.94943) (4.83151,0.96151) (4.85189,0.97321) (4.87228,0.98452) (4.89266,0.99544) (4.91305,1.00593) (4.93344,1.01601) (4.95382,1.02565) (4.97421,1.03485) (4.99459,1.04360) (5.01498,1.05189) (5.03537,1.05973) (5.05575,1.06711) (5.07614,1.07402) (5.09653,1.08046) (5.11691,1.08643) (5.13730,1.09194) (5.15768,1.09699) (5.17807,1.10157) (5.19846,1.10570) (5.21884,1.10937) (5.23923,1.11260) (5.25961,1.11539) (5.28000,1.11774)};
\end{scope}
\begin{scope}[xshift=6.6cm]
\draw[->,gray] (-.1,2.275)--(5.6,2.275) node[right]{\(x\)};
\draw[->,gray] (2.64,-.05)--(2.64,4.7);
\node at (2.64,5.0) {\(N=32\)};
\draw[gray] (2.57,0.5250)--(2.71,0.5250);\node[left] at (-0.12,0.5250) {\(-1/2\)};
\draw[gray] (2.57,2.2750)--(2.71,2.2750);\node[left] at (-0.12,2.2750) {\(0\)};
\draw[gray] (2.57,4.0250)--(2.71,4.0250);\node[left] at (-0.12,4.0250) {\(1/2\)};
\draw[gray,dashed,thick] plot coordinates {(0.00000,3.60500) (0.05387,3.61357) (0.10775,3.62214) (0.16162,3.63071) (0.21549,3.63928) (0.26937,3.64785) (0.32324,3.65642) (0.37711,3.66500) (0.43098,3.67357) (0.48486,3.68214) (0.53873,3.69071) (0.59260,3.69928) (0.64648,3.70785) (0.70035,3.71642) (0.75422,3.72499) (0.80810,3.73356) (0.86197,3.74213) (0.91584,3.75070) (0.96972,3.75927) (1.02359,3.76784) (1.07746,3.77641) (1.13133,3.78498) (1.18521,3.79356) (1.23908,3.80213) (1.29295,3.81070) (1.34683,3.81927) (1.40070,3.82784) (1.45457,3.83641) (1.50845,3.84498) (1.56232,3.85355) (1.61619,3.86212) (1.67006,3.87069) (1.72394,3.87926) (1.77781,3.88783) (1.83168,3.89640) (1.88556,3.90498) (1.93943,3.91355) (1.99330,3.92212) (2.04718,3.93069) (2.10105,3.93926) (2.15492,3.94783) (2.20880,3.95640) (2.26267,3.96497) (2.31654,3.97354) (2.37041,3.98211) (2.42429,3.99068) (2.47816,3.99925) (2.53203,4.00782) (2.58591,4.01639) (2.63978,4.02496)};
\draw[gray,dashed,thick] plot coordinates {(2.64022,0.52504) (2.69409,0.53361) (2.74797,0.54218) (2.80184,0.55075) (2.85571,0.55932) (2.90959,0.56789) (2.96346,0.57646) (3.01733,0.58503) (3.07120,0.59360) (3.12508,0.60217) (3.17895,0.61074) (3.23282,0.61931) (3.28670,0.62788) (3.34057,0.63645) (3.39444,0.64502) (3.44832,0.65360) (3.50219,0.66217) (3.55606,0.67074) (3.60994,0.67931) (3.66381,0.68788) (3.71768,0.69645) (3.77155,0.70502) (3.82543,0.71359) (3.87930,0.72216) (3.93317,0.73073) (3.98705,0.73930) (4.04092,0.74787) (4.09479,0.75644) (4.14867,0.76502) (4.20254,0.77359) (4.25641,0.78216) (4.31028,0.79073) (4.36416,0.79930) (4.41803,0.80787) (4.47190,0.81644) (4.52578,0.82501) (4.57965,0.83358) (4.63352,0.84215) (4.68740,0.85072) (4.74127,0.85929) (4.79514,0.86786) (4.84902,0.87643) (4.90289,0.88501) (4.95676,0.89358) (5.01063,0.90215) (5.06451,0.91072) (5.11838,0.91929) (5.17225,0.92786) (5.22613,0.93643) (5.28000,0.94500)};
\draw[AnalysisBlue,thick] plot coordinates {(0.00000,3.56851) (0.02039,3.57759) (0.04077,3.58787) (0.06116,3.59911) (0.08154,3.61103) (0.10193,3.62333) (0.12232,3.63570) (0.14270,3.64780) (0.16309,3.65932) (0.18347,3.66994) (0.20386,3.67939) (0.22425,3.68743) (0.24463,3.69385) (0.26502,3.69852) (0.28541,3.70136) (0.30579,3.70235) (0.32618,3.70155) (0.34656,3.69905) (0.36695,3.69506) (0.38734,3.68979) (0.40772,3.68353) (0.42811,3.67661) (0.44849,3.66939) (0.46888,3.66222) (0.48927,3.65549) (0.50965,3.64958) (0.53004,3.64481) (0.55042,3.64150) (0.57081,3.63992) (0.59120,3.64027) (0.61158,3.64269) (0.63197,3.64726) (0.65236,3.65397) (0.67274,3.66273) (0.69313,3.67340) (0.71351,3.68573) (0.73390,3.69945) (0.75429,3.71419) (0.77467,3.72957) (0.79506,3.74516) (0.81544,3.76052) (0.83583,3.77521) (0.85622,3.78880) (0.87660,3.80089) (0.89699,3.81113) (0.91737,3.81923) (0.93776,3.82496) (0.95815,3.82816) (0.97853,3.82880) (0.99892,3.82688) (1.01931,3.82255) (1.03969,3.81601) (1.06008,3.80757) (1.08046,3.79758) (1.10085,3.78650) (1.12124,3.77482) (1.14162,3.76305) (1.16201,3.75174) (1.18239,3.74142) (1.20278,3.73262) (1.22317,3.72582) (1.24355,3.72142) (1.26394,3.71979) (1.28432,3.72118) (1.30471,3.72574) (1.32510,3.73353) (1.34548,3.74447) (1.36587,3.75840) (1.38625,3.77501) (1.40664,3.79392) (1.42703,3.81462) (1.44741,3.83654) (1.46780,3.85905) (1.48819,3.88147) (1.50857,3.90310) (1.52896,3.92323) (1.54934,3.94122) (1.56973,3.95642) (1.59012,3.96832) (1.61050,3.97645) (1.63089,3.98050) (1.65127,3.98027) (1.67166,3.97570) (1.69205,3.96691) (1.71243,3.95414) (1.73282,3.93781) (1.75320,3.91847) (1.77359,3.89680) (1.79398,3.87360) (1.81436,3.84978) (1.83475,3.82627) (1.85514,3.80409) (1.87552,3.78423) (1.89591,3.76766) (1.91629,3.75529) (1.93668,3.74794) (1.95707,3.74629) (1.97745,3.75086) (1.99784,3.76199) (2.01822,3.77981) (2.03861,3.80424) (2.05900,3.83493) (2.07938,3.87133) (2.09977,3.91263) (2.12015,3.95779) (2.14054,4.00557) (2.16093,4.05453) (2.18131,4.10306) (2.20170,4.14945) (2.22208,4.19186) (2.24247,4.22844) (2.26286,4.25731) (2.28324,4.27667) (2.30363,4.28477) (2.32402,4.28003) (2.34440,4.26104) (2.36479,4.22661) (2.38517,4.17581) (2.40556,4.10801) (2.42595,4.02288) (2.44633,3.92042) (2.46672,3.80097) (2.48710,3.66521) (2.50749,3.51414) (2.52788,3.34909) (2.54826,3.17167) (2.56865,2.98375) (2.58903,2.78741) (2.60942,2.58494) (2.62981,2.37873) (2.65019,2.17127) (2.67058,1.96506) (2.69097,1.76259) (2.71135,1.56625) (2.73174,1.37833) (2.75212,1.20091) (2.77251,1.03586) (2.79290,0.88479) (2.81328,0.74903) (2.83367,0.62958) (2.85405,0.52712) (2.87444,0.44199) (2.89483,0.37419) (2.91521,0.32339) (2.93560,0.28896) (2.95598,0.26997) (2.97637,0.26523) (2.99676,0.27333) (3.01714,0.29269) (3.03753,0.32156) (3.05792,0.35814) (3.07830,0.40055) (3.09869,0.44694) (3.11907,0.49547) (3.13946,0.54443) (3.15985,0.59221) (3.18023,0.63737) (3.20062,0.67867) (3.22100,0.71507) (3.24139,0.74576) (3.26178,0.77019) (3.28216,0.78801) (3.30255,0.79914) (3.32293,0.80371) (3.34332,0.80206) (3.36371,0.79471) (3.38409,0.78234) (3.40448,0.76577) (3.42486,0.74591) (3.44525,0.72373) (3.46564,0.70022) (3.48602,0.67640) (3.50641,0.65320) (3.52680,0.63153) (3.54718,0.61219) (3.56757,0.59586) (3.58795,0.58309) (3.60834,0.57430) (3.62873,0.56973) (3.64911,0.56950) (3.66950,0.57355) (3.68988,0.58168) (3.71027,0.59358) (3.73066,0.60878) (3.75104,0.62677) (3.77143,0.64690) (3.79181,0.66853) (3.81220,0.69095) (3.83259,0.71346) (3.85297,0.73538) (3.87336,0.75608) (3.89375,0.77499) (3.91413,0.79160) (3.93452,0.80553) (3.95490,0.81647) (3.97529,0.82426) (3.99568,0.82882) (4.01606,0.83021) (4.03645,0.82858) (4.05683,0.82418) (4.07722,0.81738) (4.09761,0.80858) (4.11799,0.79826) (4.13838,0.78695) (4.15876,0.77518) (4.17915,0.76350) (4.19954,0.75242) (4.21992,0.74243) (4.24031,0.73399) (4.26069,0.72745) (4.28108,0.72312) (4.30147,0.72120) (4.32185,0.72184) (4.34224,0.72504) (4.36263,0.73077) (4.38301,0.73887) (4.40340,0.74911) (4.42378,0.76120) (4.44417,0.77479) (4.46456,0.78948) (4.48494,0.80484) (4.50533,0.82043) (4.52571,0.83581) (4.54610,0.85055) (4.56649,0.86427) (4.58687,0.87660) (4.60726,0.88727) (4.62764,0.89603) (4.64803,0.90274) (4.66842,0.90731) (4.68880,0.90973) (4.70919,0.91008) (4.72958,0.90850) (4.74996,0.90519) (4.77035,0.90042) (4.79073,0.89451) (4.81112,0.88778) (4.83151,0.88061) (4.85189,0.87339) (4.87228,0.86647) (4.89266,0.86021) (4.91305,0.85494) (4.93344,0.85095) (4.95382,0.84845) (4.97421,0.84765) (4.99459,0.84864) (5.01498,0.85148) (5.03537,0.85615) (5.05575,0.86257) (5.07614,0.87061) (5.09653,0.88006) (5.11691,0.89068) (5.13730,0.90220) (5.15768,0.91430) (5.17807,0.92667) (5.19846,0.93897) (5.21884,0.95089) (5.23923,0.96213) (5.25961,0.97241) (5.28000,0.98149)};
\end{scope}
\end{tikzpicture}
\caption[Gibbs 过冲与尺度收缩]{Gibbs 过冲与尺度收缩。灰色虚线为锯齿函数在跳跃两侧的值，蓝线为对称部分和；两图的坐标尺度相同。截断阶数增大后，过冲集中于更窄的区域，其相对跳跃高度仍不趋零。}
\label{b2:fig:02-gibbs}
\end{figure}


\subsection{对称部分和的\texorpdfstring{\(L^1\)}{L¹}收敛障碍}
\label{b2:subsec:020702}

\begin{proposition}\label{b2:prop:02-lebesgue-constants}
存在绝对常数 \(c,C>0\)，使 \(N\geq2\) 时
\[
 c\cdot\log(N+1)\leq\|D_N\|_1\leq C\cdot\log(N+1).
\]
此外，\(\|S_N\|_{L^1\to L^1}=\|D_N\|_1\)。
\end{proposition}
\begin{proof}
在 \(0<t\leq1/2\) 上，\(|D_N(t)|\leq\min\{2N+1,(2t)^{-1}\}\)，分成 \((0,1/(2N+1))\) 和其余部分积分，得到上界。下界利用 \(\sin(\pi t)\leq\pi t\)；令 \(u=(2N+1)\pi t\)，可得
\[
 \|D_N\|_1\geq\dfrac2\pi\int_0^{(2N+1)\pi/2}\dfrac{|\sin u|}{u}\,\mathrm{d}u.
\]
在每个完整区间 \([j\pi,(j+1)\pi]\) 上，以 \(((j+1)\pi)^{-1}\) 下估 \(1/u\)，而 \(\int_{j\pi}^{(j+1)\pi}|\sin u|\,\mathrm{d}u=2\)。求和即得调和级数的对数下界。

卷积和 Tonelli 给出算子范数不超过 \(\|D_N\|_1\)。反向取非负单位积分的窄区间核 \(q_\varepsilon=\varepsilon^{-1}\mathbf1_{(0,\varepsilon)}\)。固定 \(N\) 时，\(q_\varepsilon*D_N\to D_N\) 一致成立，因为 \(D_N\) 连续；且 \(\|q_\varepsilon\|_1=1\)。故算子范数至少为 \(\|D_N\|_1\)。
\end{proof}

\begin{corollary}\label{b2:cor:02-l1-failure}
存在 \(f\in L^1(\mathbb T)\)，使 \(\sup_N\|S_Nf\|_1=\infty\)。因而其 Fourier 部分和不在 \(L^1\) 中收敛。
\end{corollary}
\begin{proof}
若每个 \(f\) 的部分和范数都有界，由第一卷定理 8.3.3（一致有界原理），算子族 \((S_N)\) 的范数也应一致有界，这与前一命题矛盾。
\end{proof}

同一个函数的 Fejér 平均仍按 \(L^1\) 范数收敛。失败属于所选求和方法，不是 Fourier 系数失去了唯一确定函数的能力。

\subsection{发散例子}
\label{b2:subsec:020703}

\begin{proposition}\label{b2:prop:02-continuous-divergence}
对任意预先指定的 \(x_0\in\mathbb T\)，存在连续周期函数 \(f\)，使 \(\sup_N|S_Nf(x_0)|=\infty\)。
\end{proposition}
\begin{proof}
将 \(T_Nf=S_Nf(x_0)\) 看作 \(C(\mathbb T)\) 上的连续线性泛函。积分公式给出 \(\|T_N\|\leq\|D_N\|_1\)。为证明反向不等式，取连续函数
\[
 f_\varepsilon(x_0-t)=\dfrac{D_N(t)}{|D_N(t)|+\varepsilon}.
\]
它的上确界范数至多为 \(1\)，且 \(T_Nf_\varepsilon=\int D_N^2/(|D_N|+\varepsilon)\to\|D_N\|_1\)，其中控制函数为 \(|D_N|\)。故 \(\|T_N\|=\|D_N\|_1\to\infty\)。再用第一卷定理 8.3.3，存在固定的 \(f\) 使 \((T_Nf)\) 无界。
\end{proof}

这个存在性证明没有给出一串闭式系数，但确实产生了一个对所有 \(N\) 固定的连续函数。只为每个 \(N\) 另选一个坏函数，尚不能证明存在单个发散例子；一致有界原理在这里承担的正是这一量词转换。

\subsection{Carleson 定理概览}
\label{b2:subsec:020704}

\begin{claim}[Carleson，深层收敛定理]\label{b2:claim:02-carleson}
取 \(\mathbb T=\mathbb R/\mathbb Z\) 的归一化 Lebesgue 测度。若 \(f\in L^2(\mathbb T;\mathbb C)\)，令 \(\widehat f(k)=\int_{\mathbb T}f(t)e^{-2\pi\mathrm i kt}\,\mathrm dt\)、\(S_Nf(x)=\sum_{|k|\le N}\widehat f(k)e^{2\pi\mathrm i kx}\)，\(N\in\mathbb N_0\)。则对称 Fourier 部分和 \(S_Nf(x)\) 在几乎处处收敛到 \(f(x)\)。
\end{claim}

这一\term[carleson-dingli]{Carleson 定理}[Carleson's theorem]的证明超出了本节的范围。Sheldon Axler 的《Measure, Integration \& Real Analysis》\cite[p. 352, Section 11A]{Axler2020Measure}介绍了该结论与连续函数发散例子的关系。其困难在于同时控制全部截断频率；固定一个 \(N\) 的平方范数估计，以及本章抽子列的办法，都不能代替这种控制。

\ProseParagraphBreak

“几乎处处”允许存在例外零测集，所以定理与在指定点发散的连续函数不冲突。对于一般的 \(L^1(\mathbb T)\) 函数，不能保证对称 Fourier 部分和几乎处处收敛。Kolmogorov 构造了一个 \(f\in L^1(\mathbb T)\)，使对称部分和序列 \((S_Nf(x))_{N\geq0}\) 在几乎处处发散；上述资料同页介绍了这一结果。这是存在性反例，并非每个可积函数都发散：例如三角多项式的部分和在充分大时就等于它自身。

\ProseParagraphBreak

Kolmogorov 的具体构造在此略去。这个反例涉及几乎处处发散，与前面讨论的 \(L^1\) 范数收敛失败是两个不同结论。读者首先应掌握普通部分和、正核平均、范数收敛与点态收敛之间的区别，再进入深层收敛理论。


\begin{chaptersummary}
Fourier 系数把周期函数投向离散频率。正交性先给出有限维误差公式，三角系的稠密性再把它提升为完整的平方范数展开；系数尾部就是舍弃频率后的能量误差。点态收敛还需另一类论证，局部有界变差和 Dini 条件分别用变差控制或可积控制约束振荡积分。

\ProseParagraphBreak

正核平均避免了 Dirichlet 核绝对质量增长带来的不稳定。Fejér 平均对连续函数一致逼近，对所有有限指数的可积函数按范数逼近，并由此证明 Fourier 系数的唯一性。微分放大高频，积分减弱高频；经典逐项运算、弱导数与分布配对应按各自的收敛条件使用。这些区别在欧氏空间上的 Fourier 变换中仍会出现。
\end{chaptersummary}

\BookExercises

前两题练习频率与对称性、相关积分的关系；随后处理由系数恢复函数、正则性边界及带最优常数的估计。\Cref{b2:ex:02-fejer-operator}与\cref{b2:ex:02-periodic-heat}把频率截断和指数衰减看作算子；\cref{b2:ex:02-step-convergence}与\cref{b2:ex:02-dini-holder-boundary}则直接比较跳跃处的点态极限、范数收敛和不同的局部正则性条件。

% 自拟：将平移对称性与系数唯一性组合，区分偶频和奇频。
\begin{exercise}\label{b2:ex:02-antiperiodic}
设 \(f\in L^1(\mathbb T)\)。证明 \(f(x+1/2)=-f(x)\) 几乎处处，当且仅当全部偶数频率的 Fourier 系数为零。再证明 \(f\) 几乎处处为实值，当且仅当 \(\widehat f(-k)=\overline{\widehat f(k)}\) 对全部 \(k\) 成立。
\end{exercise}
\begin{solution}
换元给出平移函数的第 \(k\) 个系数为 \((-1)^k\widehat f(k)\)。因而 \(f(\cdot+1/2)+f\) 的系数等于 \(((-1)^k+1)\widehat f(k)\)；它们全为零等价于偶频系数全为零，再用唯一性。另有 \(\widehat{\overline f}(k)=\overline{\widehat f(-k)}\)。所给共轭对称条件等价于 \(f\) 与 \(\overline f\) 系数一致，故等价于两者几乎处处相等。
\end{solution}

% 自拟：由区间示性函数转为平移相关，给出含参数的平方和。
\begin{exercise}\label{b2:ex:02-interval-correlation}
设 \(0<a<1\)，\(f=\mathbf1_{(0,a)}\)。计算其 Fourier 系数，证明
\[
 \sum_{k=1}^\infty\dfrac{\sin^2(\pi ka)}{k^2}=\dfrac{\pi^2}2a(1-a).
\]
再将 \(C(h)=\int_{\mathbb T}f(x+h)\cdot f(x)\,\mathrm{d}x\) 写成绝对一致收敛的 Fourier 级数，并解释它的几何意义。
\end{exercise}
\begin{solution}
直接积分得到 \(\widehat f(0)=a\)，且对 \(k\ne0\)，
\[
 \widehat f(k)=e^{-\pi\mathrm{i}ka}\cdot\dfrac{\sin(\pi ka)}{\pi k}.
\]
Parseval 给出 \(a=a^2+2\sum_{k\geq1}\sin^2(\pi ka)/(\pi^2k^2)\)，整理即得所求。用内积形式的 Parseval，
\[
 C(h)=\sum_{k\in\mathbb Z}|\widehat f(k)|^2e^{2\pi\mathrm{i}kh}.
\]
系数总和为 \(a\)，故绝对一致收敛。积分等于圆周上弧段 \((0,a)\) 与其平移的重叠长度，故 \(0\leq C(h)\leq a\)、\(C(0)=a\)。为覆盖全部 \(0<a<1\)，记
\[
 d=\mathrm d_{\mathbb T}(h,0)\in[0,1/2],\qquad t_+=\max\{t,0\}.
\]
相关函数满足 \(C(-h)=C(h)\)：在积分中作平移换元即可验证。因此可把平移量取为 \(d\)。在 \((0,a)\) 上检查 \(x+d\) 是否仍落在同一弧段：未越过周期接缝的部分为 \(0<x<a-d\)，长度是 \((a-d)_+\)；越过 \(1\) 后回到弧段的部分为 \(1-d<x<a\)，长度是 \((a+d-1)_+\)。两段不交，故
\[
 C(h)=(a-d)_++(a+d-1)_+.
\]
第二项正是圆周接缝带来的额外重叠。当 \(a\leq1/2\) 时，它恒为零，才可简化为 \(C(h)=(a-d)_+\)。例如 \(a=3/4\)、\(h=1/2\) 时，两项各为 \(1/4\)，总重叠长度为 \(1/2\)。
\end{solution}

% 改编：Axler2020Measure, Exercises11B,7（印刷361页）；增加范数和重建说明。
\begin{exercise}\label{b2:ex:02-l1-l2-criterion}
设 \(f\in L^1(\mathbb T)\)。证明 \(f\) 有平方可积代表，当且仅当 \(\sum_k|\widehat f(k)|^2<\infty\)。在充分性证明中，不得直接对尚未知道属于 \(L^2\) 的 \(f\) 使用 Parseval 恒等式。
\end{exercise}
\begin{solution}
必要性由 Parseval 得到。反之，用平方可和系数构造 \(L^2\) 收敛的三角多项式列，其极限记为 \(g\)。连续的系数泛函给出 \(\widehat g(k)=\widehat f(k)\)。由于圆周总测度为 \(1\)，\(g\in L^1\)，故可对 \(f-g\) 使用系数唯一性，得到 \(f=g\) 几乎处处。此时才可写 \(\|f\|_2^2=\sum_k|\widehat f(k)|^2\)。
\end{solution}

% 自拟：绝对收敛不保证可微，利用稀疏频率的非负对称差。
\begin{exercise}\label{b2:ex:02-lacunary-nondifferentiable}
证明
\[
 f(x)=\sum_{j=1}^\infty2^{-j}\cos(2\pi2^jx)
\]
定义一个连续周期函数，Fourier 级数绝对一致收敛，但 \(f\) 在 \(0\) 不可微。
\end{exercise}
\begin{solution}
各项绝对值不超过 \(2^{-j}\)，故一致收敛且极限连续。逐项积分可辨认这些就是所写函数的 Fourier 系数。函数为偶函数，若原点可微，导数只能为零。取 \(h_m=2^{-m}\)、\(m\geq2\)。由于每项 \(1-\cos(2\pi2^jh_m)\) 非负，保留 \(j=m-1\) 一项便得
\[
 \dfrac{f(0)-f(h_m)}{h_m}
 \geq 2^m\cdot2^{-(m-1)}\bigl(1-\cos\pi\bigr)=4.
\]
这排除了导数为零的可能，所以导数不存在。
\end{solution}

% 自拟：周期弱导数与最佳常数；取等分类提供正文之外的增量。
\begin{exercise}\label{b2:ex:02-sharp-poincare}
设 \(f\) 有平方可积周期弱导数。证明
\[
 \|f-\widehat f(0)\|_2\leq\dfrac1{2\pi}\|f'\|_2,
\]
常数最佳，并求全部取等函数。
\end{exercise}
\begin{solution}
Parseval 和弱导数的系数关系给出
\[
 \|f-\widehat f(0)\|_2^2=\sum_{k\ne0}|\widehat f(k)|^2
 \leq\sum_{k\ne0}k^2|\widehat f(k)|^2=(2\pi)^{-2}\|f'\|_2^2.
\]
取等当且仅当全部 \(|k|\geq2\) 的系数为零，故函数为 \(c+ae^{2\pi\mathrm{i}x}+be^{-2\pi\mathrm{i}x}\) 几乎处处。取 \(a\ne0,b=0\) 即证常数不能缩小。
\end{solution}

% 自拟：频率方法求周期强制方程，检查零频与解的意义。
\begin{exercise}\label{b2:ex:02-periodic-resolvent}
设 \(a>0\)、\(f\in L^2(\mathbb T)\)。用 Fourier 系数构造唯一满足 \(-u''+a^2u=f\) 的周期函数 \(u\)，其中 \(u,u',u''\in L^2\)，导数均取弱意义。证明 \(\|u\|_2\leq a^{-2}\|f\|_2\)，并说明 \(a=0\) 时存在性与唯一性如何改变。
\end{exercise}
\begin{solution}
设 \(\widehat u(k)=\widehat f(k)/(4\pi^2k^2+a^2)\)。系数及其乘以 \(k\)、\(k^2\) 后均平方可和，所以连续两次使用弱导数刻画可得所需导数。逐频率检验方程，再用系数唯一性，便知方程在 \(L^2\) 中成立。任何两个解之差的每个系数都乘以非零数 \(4\pi^2k^2+a^2\) 后等于零，故解唯一。分母至少为 \(a^2\)，由 Parseval 得到范数界。若 \(a=0\)，零频方程要求 \(\widehat f(0)=0\)；满足时按非零频率同样构造，而 \(\widehat u(0)\) 可任意选取，故只在相差常数的意义下唯一。
\end{solution}

% 自拟：强收敛、算子范数与有限秩的区分。
\begin{exercise}\label{b2:ex:02-fejer-operator}
把 \(\sigma_N\) 看作 \(L^2(\mathbb T)\) 上的算子。证明它自伴、非负、有限秩，且 \(\|\sigma_N\|=1\)、\(\|I-\sigma_N\|=1\)。这为何不与对每个 \(f\) 有 \(\sigma_Nf\to f\) 矛盾？
\end{exercise}
\begin{solution}
在指数基上，算子乘以实权重 \(w_N(k)=(1-|k|/(N+1))_+\)，该权重介于 \(0\) 和 \(1\) 之间，且仅有限多个非零。这分别给出自伴、非负与有限秩。常数函数对应权重 \(1\)，故范数为 \(1\)；对 \(|k|>N\)，\((I-\sigma_N)e_k=e_k\)，加上所有权重差不超过 \(1\)，得到第二个范数等式。对固定 \(f\)，平方误差为 \(\sum_k|1-w_N(k)|^2|\widehat f(k)|^2\)，由控制收敛趋零。算子范数则允许随 \(N\) 选择不同的高频测试函数，量词不同。
\end{solution}

% 自拟：指数频率阻尼、半群与热方程；不作为后文基本核证明的唯一输入。
\begin{optionalexercise}\label{b2:ex:02-periodic-heat}
设 \(f\in L^2(\mathbb T)\)，\(t>0\)。证明
\[
 u(t,x)=\sum_{k\in\mathbb Z}e^{-4\pi^2k^2t}\widehat f(k)e^{2\pi\mathrm{i}kx}
\]
在每个 \(t\geq\varepsilon>0\) 的范围内可任意次逐项求导，满足 \(u_t=u_{xx}\)，且 \(u(t,\cdot)\to f\) 于 \(L^2\)。证明相应算子满足 \(H_tH_s=H_{t+s}\)，并求 \(t\to\infty\) 的极限及平均为零时的衰减率。
\end{optionalexercise}
\begin{solution}
任意混合导数只给系数增加一个固定次幂 \(k^m\)。由 Cauchy--Schwarz，
\[
 \sum_k|k|^m e^{-4\pi^2k^2\varepsilon}\cdot|\widehat f(k)|<\infty,
\]
所以导数级数一致收敛，求导合法。逐项计算得热方程。初值收敛由
\[
 \|u(t,\cdot)-f\|_2^2
 =\sum_k|e^{-4\pi^2k^2t}-1|^2|\widehat f(k)|^2\longrightarrow0
\]
得到。两次作用使指数因子相乘，故有半群等式。零频保持不变，非零频衰减；因此 \(u(t,\cdot)\to\widehat f(0)\) 于 \(L^2\)，且
\[
 \|u(t,\cdot)-\widehat f(0)\|_2
 \leq e^{-4\pi^2t}\|f-\widehat f(0)\|_2.
\]
只含频率 \(\pm1\) 时取等，因而这一一致衰减指数不能增大。
\end{solution}

% 自拟：同一阶跃函数用两种点态判据及正核处理，另检查L∞范数中的障碍。
\Needspace{6\baselineskip}
\begin{exercise}\label{b2:ex:02-step-convergence}
设 \(0<a<1\)，\(b_0,b_a\in\mathbb R\)。在单位圆周上选定代表
\[
 f(x)=\begin{cases}
 1,&0<x<a,\\
 -1,&a<x<1,\\
 b_0,&x=0,\\
 b_a,&x=a.
 \end{cases}
\]
\begin{enumerate}
 \item\label{b2:item:02-step-pointwise} 分别用 Dirichlet--Jordan 定理与 Dini 判据求每个点的 \(\lim_{N\to\infty}S_Nf(x)\)。说明改变 \(b_0,b_a\) 是否改变部分和及其极限。
 \item\label{b2:item:02-step-fejer} 求每个点的 \(\lim_{N\to\infty}\sigma_Nf(x)\)。在跳跃点直接使用 Fejér 核的偶性与质量集中，不把仅对连续点陈述的定理用于跳跃点。
 \item\label{b2:item:02-step-linfty} 证明任意 \(q\in C(\mathbb T;\mathbb C)\) 都有 \(\|q-f\|_\infty\geq1\)，这里取 \(L^\infty\) 的本性上确界范数。由此说明两种部分和都不能按 \(L^\infty\) 范数趋于 \(f\)，并解释这为何不与本章的有限指数范数收敛结论冲突。
\end{enumerate}
\end{exercise}
\begin{solution}
函数只有两处跳跃，在每个点的邻域内都有界变差。在两个跳跃点，左右极限依次为 \(-1,1\) 或 \(1,-1\)，平均值均为零。因此 Dirichlet--Jordan 定理给出
\[
 \lim_{N\to\infty}S_Nf(x)=\ell(x),\qquad
 \ell(x)=\begin{cases}1,&0<x<a,\\-1,&a<x<1,\\0,&x=0\text{ 或 }x=a.\end{cases}
\]
也可直接检验 Dini 条件。在不含跳跃的充分小邻域内，取 \(A=f(x)\)，对称差 \(f(x+t)+f(x-t)-2A\) 为零。在跳跃点取 \(A=0\)，只要 \(0<t<\delta<\min\{a,1-a\}\)，两侧的函数值也是一正一负，对称差仍为零，所以 Dini 积分有限。\(b_0,b_a\) 只影响两个零测点，全部 Fourier 系数、普通部分和与 Fejér 平均均不变；极限由两侧行为决定，不必等于这两个预先指定的点值。

在连续点，Fejér 定理直接给出 \(\sigma_Nf(x)\to f(x)\)。在跳跃点 \(x_0\in\{0,a\}\)，取上述 \(\delta\)，利用核的偶性写成
\[
 \sigma_Nf(x_0)=\int_0^{1/2}
 \bigl(f(x_0-t)+f(x_0+t)\bigr)F_N(t)\,\mathrm dt.
\]
\((0,\delta)\) 上的和为零；其余部分的绝对值至多
\[
 2\int_\delta^{1/2}F_N(t)\,\mathrm dt\longrightarrow0.
\]
所以 Fejér 平均也在每点趋于 \(\ell\)。这里利用了具体的左右常值结构，而非把连续点结论直接套在跳跃点上。

最后令 \(z=q(0)\)。沿 \(0\) 两侧的开弧趋近零点，连续性给出
\[
 \|q-f\|_\infty\geq\max\{|z-1|,|z+1|\}\geq1.
\]
第一个不等式对本性上确界仍成立：若某侧的连续误差在一点超过该上界，它就在一个正测度开弧上超过上界；再取趋近零点的极限即可。第二个不等式来自 \(2\leq|z-1|+|z+1|\)。普通部分和与 Fejér 平均均为连续三角多项式，所以同受这个下界约束。不过 \(f\in L^2\)，仍有 \(S_Nf\to f\) 于 \(L^2\)；它又属于每个有限指数的 \(L^p\)，仍有 \(\sigma_Nf\to f\) 于这些 \(L^p\)。这些范数允许误差集中在越来越窄的区域，\(L^\infty\) 范数则继续测量那里保留的高度。
\end{solution}

% 自拟：Hölder给出Dini充分条件；连续BV的对数尖点表明Dini不是必要条件。
\Needspace{6\baselineskip}
\begin{exercise}\label{b2:ex:02-dini-holder-boundary}
\begin{enumerate}
 \item\label{b2:item:02-holder-implies-dini} 设 \(f\in L^1(\mathbb T;\mathbb C)\)，选定其代表及一点 \(x_0\)。若存在 \(\alpha>0\)、\(C\geq0\)、\(0<\delta<1/2\)，使 \(|f(x_0+t)-f(x_0)|\leq C|t|^\alpha\) 对 \(|t|<\delta\) 成立，证明 Dini 判据适用，并指出由此得到的是哪一种收敛。
 \item\label{b2:item:02-log-cusp} 记 \(d(x)=\mathrm d_{\mathbb T}(x,0)\)，令
 \[
 v(x)=\begin{cases}\dfrac1{\log(e/d(x))},&d(x)>0,\\0,&d(x)=0.\end{cases}
 \]
 证明 \(v\) 是连续的有界变差周期函数，但在原点不满足任何正指数的局部 Hölder 条件，也不满足取 \(A=v(0)\) 时的 Dini 积分条件。
 \item\label{b2:item:02-log-cusp-convergence} 尽管上一问的 Dini 条件失败，证明仍有 \(S_Nv(0)\to0\) 与 \(\sigma_Nv(0)\to0\)。说明这对判据的必要性意味着什么。
\end{enumerate}
\end{exercise}
\begin{solution}
第一问中，对称差的绝对值不超过 \(2Ct^\alpha\)，所以
\[
 \int_0^\delta\frac{|f(x_0+t)+f(x_0-t)-2f(x_0)|}{t}\,\mathrm dt
 \leq2C\int_0^\delta t^{\alpha-1}\,\mathrm dt
 =\frac{2C\delta^\alpha}{\alpha}<\infty.
\]
Dini 判据给出 \(S_Nf(x_0)\to f(x_0)\)。这是所指定一点的普通部分和收敛；题设没有对全部 \(x_0\) 给出一致常数，不能据此声称一致收敛。

令 \(\psi(0)=0\)、\(\psi(t)=1/\log(e/t)\)，其中 \(0<t\leq1/2\)。它连续且严格递增，在零点的右极限为零。函数 \(v\) 在 \([-1/2,0]\) 上递减，在 \([0,1/2]\) 上递增，两端值相等。因此它可连续地周期延拓，且该周期上的总变差为
\[
 2\psi(1/2)=\frac2{\log(2e)}<\infty.
\]
对任意 \(\alpha>0\)，取 \(t_n=e^{-n}\)，则
\[
 \frac{|v(t_n)-v(0)|}{t_n^\alpha}
 =\frac{e^{\alpha n}}{1+n}\longrightarrow\infty,
\]
所以不存在所需的局部 Hölder 常数。由于 \(v(t)=v(-t)\)，对任何 \(0<\delta<1/2\) 都有
\[
 \int_0^\delta\frac{|v(t)+v(-t)-2v(0)|}{t}\,\mathrm dt
 =2\int_0^\delta\frac{\mathrm dt}{t\log(e/t)}=\infty;
\]
最后的发散可用 \(s=\log(e/t)\) 换元化为 \(2\int_{\log(e/\delta)}^\infty s^{-1}\,\mathrm ds\)。

另一方面，\(v\) 在原点两侧都有界变差，且左右极限都为零，故 Dirichlet--Jordan 定理给出 \(S_Nv(0)\to0\)。连续性与 Fejér 定理又给出 \(\sigma_Nv(0)\to0\)。因此 Dini 条件是足以保证点态收敛的一种条件，未满足它并不能推出发散；有界变差则提供了这里仍可使用的另一条证明路线。
\end{solution}

% !TeX root = ../../Book2.tex
\chapter{经典逼近论}
\label{b2:ch:approximation}
\BookChapterNumber{b2:ch:approximation}

\begin{chapterabstract}
用多项式逼近连续函数，既要判断误差能否趋于零，也要研究给定次数下能达到多小的误差。本章先以 Bernstein 多项式给出可计算的一致逼近，再用 Chebyshev 交错定理刻画实区间上的最佳多项式。周期卷积的正核构造导出 Jackson 直接定理；三角多项式的导数估计则反过来从逼近速度辨认函数的光滑性。最后讨论高阶差分、整数光滑阶的限制和区间端点的特殊尺度，说明几种看似相近的逼近结论为何不能混用。
\end{chapterabstract}

\begin{learninggoals}
\begin{itemize}
\item\label{b2:item:approximation-goal-modulus}用连续模统一控制紧集上连续函数的局部变化。
\item\label{b2:item:approximation-goal-bernstein}掌握一元与多元 Bernstein 多项式的构造，以非负权重和二阶矩估计证明一致逼近。
\item\label{b2:item:approximation-goal-planned}证明实区间最佳多项式的存在唯一性，用交错极值检验最佳性，并区别一致范数与平方范数的优化目标。
\item\label{b2:item:approximation-goal-jackson}从正核的次数、归一化和一阶矩推出 Jackson 估计，解释构造性逼近与最佳逼近速度的区别。
\item\label{b2:item:approximation-goal-inverse}用 Bernstein 导数估计和二进尺度分解证明周期 Hölder 刻画，辨认整数阶与区间端点的限制。
\end{itemize}
\end{learninggoals}

用有限个数描述一个连续函数，最直接的办法是在一些点上取样。怎样把这些取样值组合成一个多项式，使它在取样点之间也接近原函数？Bernstein 构造把附近的取样值赋予较大权重，提供了一种直观的做法。得到近似之后，还可以继续追问：同样的次数，是否能把最大误差降得更低？这一问题把注意力引向误差最大的那些点，最终产生 Chebyshev 的交错判据。

\ProseParagraphBreak

增加多项式的次数，通常会使近似更精确；但误差究竟下降多快，还取决于原函数怎样变化。平滑的曲线是否比带尖点的曲线更容易逼近？如果已经知道误差下降很快，又能从中读出函数的哪些性质？Jackson 定理与逆定理分别回答这两个方向的问题。章末再考察高阶变化，以及闭区间端点附近与内部不同的逼近尺度。

\ProseParagraphBreak

Elias M. Stein 与 Rami Shakarchi 的《Fourier Analysis: An Introduction》\cite{Stein2003Fourier}可用于回顾前章的周期分析背景；Antoni Zygmund 的《Trigonometric Series》\cite{Zygmund2002Trigonometric}可作为继续研究三角级数与函数光滑性的进阶读物。若想了解区间之外的连续函数逼近，可另读第一卷第29.2节的一般 Stone--Weierstrass 定理。

% !TeX root = ../../Book2.tex
\section{逼近空间、误差与连续模}
\label{b2:sec:approximation-spaces-errors}

\subsection{逼近空间}
\label{b2:subsec:approximation-spaces}

待撰写。

\subsection{连续模与逼近误差}
\label{b2:subsec:020901}

Bernstein 构造在有限个点上取样，再以非负权重求平均。估计取样误差时，需要同时控制所有相距较近的点上的函数值之差。

\begin{definition}
\label{b2:def:approximation-modulus-continuity}
设 \(d\) 为正整数，\(K\subset\mathbb R^d\) 为非空紧集，\(\mathbb F=\mathbb R\) 或 \(\mathbb C\)，且 \(f\in C(K;\mathbb F)\)。对 \(\delta\geq0\)，定义
\[
 \omega_f(\delta)
 \coloneqq\sup_{\substack{x,y\in K\\|x-y|\leq\delta}}
             |f(x)-f(y)|.
\]
这里 \(|x-y|\) 为欧氏距离。函数 \(\omega_f\) 称为 \(f\) 的\term[lianxu-mo]{连续模}[modulus of continuity]。
\end{definition}
\symidx{\(\omega_f\)：函数的连续模}

连续模单调不减，且 \(0\leq\omega_f(\delta)\leq2\lVert f\rVert_\infty\)。紧集上的一致连续性给出 \(\omega_f(\delta)\to0\)（\(\delta\downarrow0\)）；这一步把逐对点的连续性变为统一的误差控制。

\subsection{最佳逼近误差}
\label{b2:subsec:approximation-best-error}

待撰写。

% !TeX root = ../../Book2.tex
\section{Bernstein 构造性逼近}
\label{b2:sec:approximation-bernstein}

要实际计算一个逼近多项式，可以先从等距网格取得函数值
\(f(0),f(1/n),\ldots,f(1)\)，把它们组合成一个明确的多项式。
构造的关键是使各个权重非负、总和为一，而且对固定的 \(x\)，主要权重落在
\(k/n\) 接近 \(x\) 的取样点上。这样，连续性便能把取样点的接近转化为函数值的接近。

\ProseParagraphBreak
这一办法不要求多项式经过所有取样点。插值规定每个节点上的值都完全正确，
这里则允许在节点上有误差，以换取容易控制的加权平均。

\subsection{一元 Bernstein 多项式}
\label{b2:subsec:approximation-bernstein-univariate}

下面的\term[bernstein-bijin-dingli]{Bernstein 逼近定理}[Bernstein approximation theorem]用二项式权重构造连续函数的多项式一致逼近。

\begin{theorem}[Bernstein 逼近：一元情形]
\label{b2:thm:approximation-bernstein-univariate}
设 \(\mathbb F=\mathbb R\) 或 \(\mathbb C\)，且 \(f\in C([0,1];\mathbb F)\)。对每个正整数 \(n\)，定义
\[
 b_{n,k}(x)\coloneqq\binom nk x^k(1-x)^{n-k},
 \quad
 B_nf(x)\coloneqq\sum_{k=0}^nf(k/n)\cdot b_{n,k}(x),
 \quad 0\leq x\leq1.
\]
这里零次幂按 \(1\) 计算。\(B_nf\) 是次数不超过 \(n\) 的\term[bernstein-duoxiangshi]{Bernstein 多项式}[Bernstein polynomial]。以
\[
 \omega_f(\delta)\coloneqq
 \sup\{|f(s)-f(t)|:s,t\in[0,1],\ |s-t|\leq\delta\}
 \quad(\delta>0)
\]
表示 \(f\) 关于通常距离的连续模，则对 \(0<\delta\leq1\)，
\[
 \lVert B_nf-f\rVert_\infty
 \leq\omega_f(\delta)
      +\frac{\lVert f\rVert_\infty}{2n\delta^2}.
\]
因此 \(B_nf\to f\) 一致成立。此外，若 \(a<b\) 为实数，\(h\in C([a,b];\mathbb F)\)，则存在次数不超过 \(n\) 的多项式 \(q_n\)，使 \(q_n\to h\) 在 \([a,b]\) 上一致成立。
\end{theorem}
\begin{proof}
这些权重非负。对固定的 \(x\in[0,1]\)，引入辅助变量 \(z\)，由二项式公式，
\[
 \sum_{k=0}^n b_{n,k}(x)z^k=(1-x+xz)^n.
\]
分别令 \(z=1\)，以及对 \(z\) 求一、二阶导数后令 \(z=1\)，得到
\[
\begin{aligned}
 &\sum_{k=0}^n b_{n,k}(x)=1,\quad
   \sum_{k=0}^n k\cdot b_{n,k}(x)=nx,\\
 &\sum_{k=0}^n k(k-1)\cdot b_{n,k}(x)=n(n-1)x^2.
\end{aligned}
\]
当 \(n=1\) 时，第三个等式两边均为零。这些都是关于 \(x\) 的多项式恒等式，
所以端点 \(x=0,1\) 同样适用。利用 \(k^2=k(k-1)+k\) 展开平方，可得
\begin{equation}
 \label{b2:eq:approximation-bernstein-variance}
 \sum_{k=0}^n(k/n-x)^2\cdot b_{n,k}(x)
 =\frac{x(1-x)}n\leq\frac1{4n}.
\end{equation}
这说明按这些权重平均，取样点与 \(x\) 的平方距离很小。
固定 \(0<\delta\leq1\)，将求和分为 \(|k/n-x|\leq\delta\) 与其补集。邻近部分的函数值之差至多为 \(\omega_f(\delta)\)，而远处权重满足
\[
 \sum_{|k/n-x|>\delta}b_{n,k}(x)
 \leq\frac1{\delta^2}
       \sum_{k=0}^n(k/n-x)^2\cdot b_{n,k}(x)
 \leq\frac1{4n\delta^2}.
\]
由于权重之和为一，
\[
\begin{aligned}
 |B_nf(x)-f(x)|
 &\leq\sum_{k=0}^n|f(k/n)-f(x)|\cdot b_{n,k}(x)\\
 &\leq\omega_f(\delta)
       +2\lVert f\rVert_\infty\cdot\frac1{4n\delta^2}.
\end{aligned}
\]
对 \(x\) 取上确界即得误差估计。取 \(\delta=n^{-1/4}\)，右端成为
\[
 \omega_f(n^{-1/4})+\frac{\lVert f\rVert_\infty}{2\sqrt n},
\]
由一致连续性趋于零。

对一般区间上的 \(h\)，令 \(H(x)=h(a+(b-a)x)\)，其中 \(0\leq x\leq1\)，再设
\[
 q_n(t)\coloneqq B_nH\!\left(\frac{t-a}{b-a}\right),
 \quad a\leq t\leq b.
\]
仿射换元保持次数不超过 \(n\)，而且
\(\lVert q_n-h\rVert_{C[a,b]}=\lVert B_nH-H\rVert_{C[0,1]}\to0\)，得到所需逼近。
\end{proof}

三个权重恒等式还给出可以直接核验的例子。若 \(\iota(x)=x\)，则
\[
 B_n1=1,\quad B_n\iota=\iota,\quad
 B_n(\iota^2)(x)=x^2+\frac{x(1-x)}n.
\]
因此常数和仿射函数被完全保持；对 \(x^2\)，误差在两端为零，在中点达到
\(1/(4n)\)。这也说明 \(B_n\) 一般不是投影到次数不超过 \(n\) 的多项式空间：
即使 \(n\geq2\)，它仍会改变已在该空间中的 \(x^2\)。对所有连续 \(f\)，
端点处只有一个权重非零，所以 \(B_nf(0)=f(0)\)、\(B_nf(1)=f(1)\)；
内部取样点通常没有这种精确相等。

\subsection{连续模与有效取样尺度}
\label{b2:subsec:approximation-bernstein-rate}

前面的远近拆分便于证明收敛，但取 \(\delta=n^{-1/4}\) 并没有反映构造本身的尺度。
由\eqref{b2:eq:approximation-bernstein-variance}，平方距离的平均是 \(O(n^{-1})\)，
所以距离的典型大小应为 \(n^{-1/2}\)。要把这个信息用于误差，
可以将较远的两点之间的变化拆成若干段短距离上的变化。

\begin{lemma}[区间连续模的分段估计]
\label{b2:lem:approximation-modulus-subdivision}
设 \(a<b\) 为实数，\(\mathbb F=\mathbb R\) 或 \(\mathbb C\)，
\(f\in C([a,b];\mathbb F)\)，并令
\(\omega_f(\delta)=\sup_{s,t\in[a,b],\,|s-t|\leq\delta}|f(s)-f(t)|\)。
则对每个 \(\delta>0\) 和 \(s,t\in[a,b]\)，
\[
 |f(s)-f(t)|\leq\left(1+\frac{|s-t|}{\delta}\right)\omega_f(\delta).
\]
\end{lemma}
\begin{proof}
当 \(s=t\) 时左端为零。否则取
\(m=\lceil|s-t|/\delta\rceil\)，把从 \(s\) 到 \(t\) 的线段分成 \(m\) 个等长小段。
每段长不超过 \(\delta\)，所有分点仍在 \([a,b]\) 中。
沿分点逐次使用三角不等式，得到
\(|f(s)-f(t)|\leq m\omega_f(\delta)\)。最后使用
\(m\leq1+|s-t|/\delta\) 即可。
\end{proof}

\begin{proposition}[Bernstein 多项式的连续模估计]
\label{b2:prop:approximation-bernstein-modulus}
设 \(\mathbb F=\mathbb R\) 或 \(\mathbb C\)，\(f\in C([0,1];\mathbb F)\)，
\(n\geq1\) 为整数，且
\[
 B_nf(x)=\sum_{k=0}^nf(k/n)\binom nk x^k(1-x)^{n-k},
\]
端点的零次幂按 \(1\) 计算。以 \(\omega_f\) 表示关于通常距离的连续模。
对每个 \(\delta>0\) 和 \(x\in[0,1]\)，有
\[
 |B_nf(x)-f(x)|
 \leq\left(1+\frac{\sqrt{x(1-x)}}{\sqrt n\,\delta}\right)\omega_f(\delta).
\]
特别地，
\[
 \lVert B_nf-f\rVert_\infty\leq\frac32\omega_f(n^{-1/2}).
\]
若还存在 \(0<\alpha\leq1\) 和 \(L\geq0\)，使
\(|f(s)-f(t)|\leq L|s-t|^\alpha\) 对所有 \(s,t\in[0,1]\) 成立，则
\[
 |B_nf(x)-f(x)|\leq L\left(\frac{x(1-x)}n\right)^{\alpha/2},\quad
 \lVert B_nf-f\rVert_\infty\leq\frac{L}{2^\alpha n^{\alpha/2}}.
\]
\end{proposition}
\begin{proof}
先将\cref{b2:lem:approximation-modulus-subdivision}应用于 \(s=k/n,t=x\)，再按权重求和：
\[
 |B_nf(x)-f(x)|
 \leq\omega_f(\delta)\left(1+\frac1\delta
       \sum_{k=0}^n|k/n-x|\cdot b_{n,k}(x)\right).
\]
由 Cauchy--Schwarz 不等式、权重和为一及\eqref{b2:eq:approximation-bernstein-variance}，
\[
 \begin{aligned}
 \sum_{k=0}^n|k/n-x|\cdot b_{n,k}(x)
 &\leq\left(\sum_{k=0}^n(k/n-x)^2\cdot b_{n,k}(x)\right)^{1/2}
       \left(\sum_{k=0}^nb_{n,k}(x)\right)^{1/2}\\
 &=\sqrt{\frac{x(1-x)}n}.
 \end{aligned}
\]
代回即得第一式。令 \(\delta=n^{-1/2}\)，并使用 \(x(1-x)\leq1/4\)，得到一致范数的估计。

在 Hölder 条件下，直接由有限和的 Hölder 不等式，以共轭指数
\(2/\alpha\) 与 \(2/(2-\alpha)\) 估计，得到
\[
 \begin{aligned}
 |B_nf(x)-f(x)|
 &\leq L\sum_{k=0}^n|k/n-x|^\alpha\cdot b_{n,k}(x)\\
 &\leq L\left(\sum_{k=0}^n(k/n-x)^2\cdot b_{n,k}(x)\right)^{\alpha/2}
       \left(\sum_{k=0}^nb_{n,k}(x)\right)^{1-\alpha/2}.
 \end{aligned}
\]
再用权重恒等式即可。
\end{proof}

这个估计把近似的精度与函数的连续性联系起来：Hölder 指数为 \(\alpha\) 时，
它给出 \(n^{-\alpha/2}\) 的误差上界。它只是这一具体构造的保证，
不是最佳逼近误差的下界，也不能据此断言所有函数的实际误差都恰好有这个量级。
前面的 \(x^2\) 就有更快的 \(n^{-1}\) 误差。\secref{b2:sec:approximation-jackson}将用更集中的核，
把一般 Hölder 函数的逼近上界改进到 \(n^{-\alpha}\)。

\subsection{正性、稳定性与形状保持}
\label{b2:subsec:approximation-bernstein-shape}

非负权重除了简化误差估计，还使构造具有直接的稳定性。
如果每个取样值都有至多 \(\varepsilon\) 的误差，加权平均后的误差仍不超过
\(\varepsilon\)，不会因节点数增加而放大。对于实值函数，
这些权重还保留单调性与凸性。

\begin{proposition}[Bernstein 算子的稳定性与保形性]
\label{b2:prop:approximation-bernstein-shape}
设 \(\mathbb F=\mathbb R\) 或 \(\mathbb C\)，\(n\geq1\) 为整数，并对
\(f\in C([0,1];\mathbb F)\) 定义
\[
 B_nf(x)=\sum_{k=0}^nf(k/n)\binom nk x^k(1-x)^{n-k},\quad 0\leq x\leq1,
\]
端点的零次幂按 \(1\) 计算。则 \(B_n\) 是线性映射，且对所有
\(f,g\in C([0,1];\mathbb F)\)，
\[
 \lVert B_nf-B_ng\rVert_\infty\leq\lVert f-g\rVert_\infty.
\]
当 \(\mathbb F=\mathbb R\) 时，\(f\geq0\) 蕴含 \(B_nf\geq0\)；
若 \(f\) 单调不减，则 \(B_nf\) 单调不减；若 \(f\) 为凸函数，则 \(B_nf\) 为凸函数，
并且 \(B_nf(x)\geq f(x)\) 对全部 \(x\in[0,1]\) 成立。
\end{proposition}
\begin{proof}
线性由有限求和直接给出。由非负权重之和为一，
\[
 |B_nf(x)-B_ng(x)|
 \leq\sum_{k=0}^n|f(k/n)-g(k/n)|\cdot b_{n,k}(x)
 \leq\lVert f-g\rVert_\infty.
\]
对 \(x\) 取上确界即得稳定性；正性同样由各项非负得到。

为考察形状，暂将越界的 \(b_{m,k}\) 定义为零，并令 \(b_{0,0}=1\)。
对权重求导并利用二项式系数的关系，有
\[
 b'_{n,k}=n(b_{n-1,k-1}-b_{n-1,k}).
\]
两部分分别换指标，端点的越界项为零，便得到
\[
 (B_nf)'(x)=n\sum_{k=0}^{n-1}
       \bigl(f((k+1)/n)-f(k/n)\bigr)b_{n-1,k}(x).
\]
当 \(f\) 单调不减时，每个差均非负，所以导数非负。
当 \(n\geq2\) 时，记
\(d_k=f((k+2)/n)-2f((k+1)/n)+f(k/n)\)。再次求导，得到
\[
 (B_nf)''(x)=n(n-1)\sum_{k=0}^{n-2}d_k\cdot b_{n-2,k}(x).
\]
对凸函数，中点处的值不超过两端值的平均，所以二阶差 \(d_k\) 非负，
从而二阶导数非负。\(n=1\) 时 \(B_1f\) 为仿射函数，也为凸函数。
最后，由有限凸组合的凸性不等式及 \(\sum_k(k/n)b_{n,k}(x)=x\)，有
\[
 f(x)=f\!\left(\sum_{k=0}^n(k/n)\cdot b_{n,k}(x)\right)
 \leq\sum_{k=0}^nf(k/n)\cdot b_{n,k}(x)=B_nf(x).
\]
有限凸组合的这一不等式由凸函数定义中的两点不等式逐次合并得到，
零权重项直接略去即可。
\end{proof}

取 \(f(x)=|x-1/2|\)，原函数在中点不可微，而 \(B_nf\) 是处处光滑的凸多项式，
并始终位于 \(f\) 的上方。这里不需要对 \(f\) 求导：多项式的导数来自离散取样值的差。
因此连续函数的某些几何性质，可以在它尚无导数时就传给光滑近似。

\subsection{多元 Bernstein 多项式}
\label{b2:subsec:approximation-bernstein-multivariate}

在立方体上，可沿每个坐标分别使用同样的二项式权重，再取乘积。
所得网格共有 \((n+1)^d\) 个取样点，每个坐标的集中程度仍由一元计算控制。
这里要区分每个变量的次数与总次数：乘积可能同时含有各坐标的 \(n\) 次幂，
所以总次数上界是 \(dn\)，而不是 \(n\)。

\begin{theorem}[Bernstein 逼近：多元情形]
\label{b2:thm:approximation-bernstein-multivariate}
设 \(d\) 为正整数，\(\mathbb F=\mathbb R\) 或 \(\mathbb C\)，且 \(f\in C([0,1]^d;\mathbb F)\)。对正整数 \(n\)、\(\mathbf k=(k_1,\ldots,k_d)\in\{0,\ldots,n\}^d\) 与 \(x\in[0,1]^d\)，定义
\[
\begin{aligned}
 w_{n,\mathbf k}(x)
 &\coloneqq\prod_{j=1}^d\binom n{k_j}x_j^{k_j}(1-x_j)^{n-k_j},\\
 B_nf(x)
 &\coloneqq\sum_{\mathbf k\in\{0,\ldots,n\}^d}
       f(\mathbf k/n)\cdot w_{n,\mathbf k}(x).
\end{aligned}
\]
零次幂按 \(1\) 计算。\(B_nf\) 在每个变量上的次数均不超过 \(n\)，总次数不超过 \(dn\)。以
\[
 \omega_f(\delta)\coloneqq
 \sup\{|f(s)-f(t)|:s,t\in[0,1]^d,\ |s-t|\leq\delta\}
\]
表示 \(f\) 关于欧氏距离的连续模，则对每个 \(\delta>0\)，
\[
 \lVert B_nf-f\rVert_\infty
 \leq\omega_f(\delta)
       +\frac{d\lVert f\rVert_\infty}{2n\delta^2}.
\]
特别地，\(B_nf\to f\) 在 \([0,1]^d\) 上一致成立；还有较细的估计
\[
 \lVert B_nf-f\rVert_\infty\leq\frac32\omega_f(\sqrt{d/n}).
\]
\end{theorem}
\begin{proof}
乘积权重非负。对每个坐标依次求和，并使用一元二项式恒等式，得到
\[
 \sum_{\mathbf k}w_{n,\mathbf k}(x)=1,
 \quad
 \sum_{\mathbf k}\frac{k_j}{n}\cdot w_{n,\mathbf k}(x)=x_j,
 \quad 1\leq j\leq d.
\]
同样计算每个坐标的二阶矩，
\[
 \sum_{\mathbf k}(k_j/n-x_j)^2\cdot w_{n,\mathbf k}(x)
 =\frac{x_j(1-x_j)}n.
\]
把这些等式相加，便有
\[
 \sum_{\mathbf k}|\mathbf k/n-x|^2\cdot w_{n,\mathbf k}(x)
 =\sum_{j=1}^d\frac{x_j(1-x_j)}n
 \leq\frac d{4n}.
\]
所以对任意 \(\delta>0\)，远处取样点的总权重满足
\[
 \sum_{|\mathbf k/n-x|>\delta}w_{n,\mathbf k}(x)
 \leq\frac d{4n\delta^2}.
\]
在 \(|\mathbf k/n-x|\leq\delta\) 的部分，用连续模控制函数值之差；在补集上，用 \(2\lVert f\rVert_\infty\) 控制。由此得到
\[
\begin{aligned}
 |B_nf(x)-f(x)|
 &\leq\sum_{\mathbf k}|f(\mathbf k/n)-f(x)|\cdot w_{n,\mathbf k}(x)\\
 &\leq\omega_f(\delta)
       +2\lVert f\rVert_\infty\cdot\frac d{4n\delta^2}.
\end{aligned}
\]
取上确界即得所述估计。先由一致连续性取足够小的 \(\delta\)，再令 \(n\) 足够大，就得到一致收敛。

较细的估计使用与一元情形相同的分段办法。立方体是凸集，
连接其中两点的线段仍留在立方体内，故按长度至多 \(\delta\) 分段可以得到
\[
 |f(s)-f(t)|\leq\left(1+\frac{|s-t|}{\delta}\right)\omega_f(\delta).
\]
代入加权和，再用 Cauchy--Schwarz 不等式及上述平方距离估计，得到
\[
 \begin{aligned}
 |B_nf(x)-f(x)|
 &\leq\left(1+\frac1\delta
 \sqrt{\sum_{\mathbf k}|\mathbf k/n-x|^2\cdot w_{n,\mathbf k}(x)}\right)
 \omega_f(\delta)\\
 &\leq\left(1+\frac{\sqrt d}{2\sqrt n\,\delta}\right)\omega_f(\delta).
 \end{aligned}
\]
取 \(\delta=\sqrt{d/n}\) 并对 \(x\) 取上确界，便得第二个估计。
\end{proof}

对于任意非退化闭矩形 \(Q=\prod_{j=1}^d[a_j,b_j]\)，逐坐标作仿射换元
\(x_j=(t_j-a_j)/(b_j-a_j)\)，便把任意 \(C(Q)\) 函数转成单位立方体上的连续函数。
把所得 Bernstein 多项式换回原变量，仍是多项式，且一致误差完全相同。
所以多项式在 \(C(Q)\) 中一致稠密。这个推论正是后文证明某些正交展开完备性时所需的具体工具。

\ProseParagraphBreak
若函数最初只定义在球面 \(S^{d-1}\) 上，则还不能直接对它取立方体网格上的值。
对 \(d\geq2\) 和 \(f\in C(S^{d-1};\mathbb F)\)，其中 \(\mathbb F=\mathbb R\) 或 \(\mathbb C\)，
可取连续函数 \(\chi:[0,\infty)\to[0,1]\)，使
\(\chi(1)=1\)、\(\chi(r)=0\) 于 \(r\leq1/2\)，并令
\[
 F(x)=\chi(|x|)\cdot f(x/|x|)\quad(x\ne0),\quad F(0)=0.
\]
这个函数在原点附近恒为零，在其余位置由连续函数复合而成，因而在包住球面的立方体上连续，
且限制到球面恰好等于 \(f\)。对 \(F\) 作多项式一致逼近，再限制到球面，误差只会减小。
\secref{b2:sec:B203}的球谐展开将实际使用这条路径：先作连续延拓，
再作多项式逼近，最后处理多项式的调和分解。这样既说明多元定理的用途，
也交代了从矩形到球面所需的额外一步。

% !TeX root = ../../Book2.tex
\section{一致最佳逼近与 Chebyshev 理论}
\label{b2:sec:approximation-chebyshev}
\label{b2:sec:0208}

待撰写。

\subsection{一致范数下的最佳逼近}
\label{b2:subsec:020801}

待撰写。

\subsection{Chebyshev 多项式与极值性质}
\label{b2:subsec:020802}

待撰写。

\subsection{交错定理与唯一性}
\label{b2:subsec:020803}

待撰写。

\subsection{与正交最佳逼近的比较}
\label{b2:subsec:020804}

待撰写。

% !TeX root = ../../Book2.tex
\section{Jackson 定理与直接逼近}
\label{b2:sec:approximation-jackson}
\label{b2:sec:0209}

Bernstein 多项式给出了逐点取样的显式算法。一致收敛以后，还有一个更细的问题：若已知函数在距离 \(\delta\) 内的振幅不超过 \(\omega_f(\delta)\)，次数为 \(n\) 的多项式能否把误差降到 \(\omega_f(1/n)\) 的量级？本节先在圆周上用卷积构造这种逼近，再通过余弦变换处理闭区间。

\ProseParagraphBreak

这里的\term[zhijie-bijin-dingli]{直接逼近定理}[direct approximation theorem]是指由函数的连续性或光滑性推出逼近误差的上界。证明要同时保住两个条件：卷积核的质量足够集中，所得函数的次数也不能超过预算。只让核越来越窄，而不计算它的次数，并不能回答给定 \(n\) 时的问题。

\subsection{Jackson 核与直接定理}
\label{b2:subsec:020902}

回顾\subsecref{b2:subsec:020401}的约定：圆周为 \(\mathbb T=\mathbb R/\mathbb Z\)，积分采用总质量为 \(1\) 的 Lebesgue 测度，而 Fejér 核 \(F_N\) 的次数为 \(N\)。因此，若用正整数 \(m\) 表示集中尺度，应该写
\[
 F_{m-1}(t)=\frac1m
  \left(\frac{\sin(m\pi t)}{\sin(\pi t)}\right)^2,
 \quad F_{m-1}(0)=m.
\]
Fejér 核已经非负且积分为 \(1\)，但它离开原点以后只按 \(|t|^{-2}\) 衰减。要估计连续模，需要控制核的第一矩；这个衰减速度会带来一个对数。把核平方，再恢复总质量为 \(1\)，可以改善尾部而不失去非负性。

\begin{definition}
\label{b2:def:approximation-jackson-kernel}
设 \(m\geq1\) 为整数，\(F_{m-1}\) 是单位周期圆周 \(\mathbb T=\mathbb R/\mathbb Z\) 上的 Fejér 核。记
\[
 A_m\coloneqq\int_{\mathbb T}F_{m-1}(t)^2\,\mathrm dt,
 \quad
 J_m(t)\coloneqq\frac{F_{m-1}(t)^2}{A_m}.
\]
函数 \(J_m\) 称为本节使用的\term[jackson-he]{Jackson 核}[Jackson kernel]。
\end{definition}
\symidx{\(J_m\)：Jackson 核}

\begin{lemma}
\label{b2:lem:approximation-jackson-kernel}
设 \(m\geq1\) 为整数，\(F_{m-1}\) 是 \(\mathbb T=\mathbb R/\mathbb Z\) 上的 Fejér 核，积分取归一化 Lebesgue 测度，\(A_m=\int_{\mathbb T}F_{m-1}^2\)、\(J_m=F_{m-1}^2/A_m\)。则 \(J_m\) 是实值、偶、非负的三角多项式，次数不超过 \(2(m-1)\)，且
\[
 \int_{\mathbb T}J_m(t)\,\mathrm dt=1,
 \quad A_m=\frac{2m^2+1}{3m}.
\]
存在与 \(m\) 无关的常数 \(C>0\)，使
\[
 \begin{aligned}
 J_m(t)&\leq\frac{Cm}{(1+m|t|)^4}
 \quad(-1/2\leq t\leq1/2),\\
 \int_{-1/2}^{1/2}|t|J_m(t)\,\mathrm dt&\leq\frac Cm.
 \end{aligned}
\]
这里 \(|t|\) 就是所选代表到圆周原点的最短弧距。
\end{lemma}
\begin{proof}
Fejér 核的有限展开为
\[
 F_{m-1}(t)=\sum_{|k|<m}\left(1-\frac{|k|}{m}\right)e^{2\pi\mathrm i kt}.
\]
两个次数不超过 \(m-1\) 的三角多项式相乘，频率相加，所以平方的次数不超过 \(2(m-1)\)。实值、偶性和非负性直接保留。利用有限指数函数的正交性，有
\[
 A_m=1+2\sum_{k=1}^{m-1}\left(1-\frac{k}{m}\right)^2
 =1+\frac{2}{m^2}\cdot\frac{(m-1)m(2m-1)}6
 =\frac{2m^2+1}{3m}.
\]
因而 \(2m/3\leq A_m\leq m\)，特别地，归一化分母严格为正。

从有限指数和的平方模表达得到 \(F_{m-1}\leq m\)。另一方面，在 \(0<|t|\leq1/2\) 上，\(|\sin(\pi t)|\geq2|t|\) 且 \(|\sin(m\pi t)|\leq1\)，所以
\[
 0\leq F_{m-1}(t)\leq
 \min\left\{m,\frac1{4mt^2}\right\}.
\]
平方后除以 \(A_m\geq2m/3\)，便有
\[
 J_m(t)\leq C m\min\{1,(m|t|)^{-4}\}
 \leq\frac{C' m}{(1+m|t|)^4}.
\]
最后令 \(u=m|t|\)，把正半轴上的积分上界延长到无穷，得到
\[
 \int_{-1/2}^{1/2}|t|J_m(t)\,\mathrm dt
 \leq\frac{2C'}m\int_0^\infty\frac{u}{(1+u)^4}\,\mathrm du
 \leq\frac{C''}m.
\]
积分有限，因为被积函数在无穷远按 \(u^{-3}\) 衰减。质量为 \(1\) 则来自 \(A_m\) 的定义。
\end{proof}

核的平方把次数增加一倍，也把远处衰减从二次提高到四次；归一化以后，它的高度仍为 \(m\) 的量级，主要质量仍落在宽度约 \(1/m\) 的区间内。第一矩估计量化了这个宽度，下面的误差证明只需要质量与第一矩这两项信息。

\begin{theorem}[Jackson：周期情形]
\label{b2:thm:approximation-jackson-periodic}
设 \(\mathbb F=\mathbb R\) 或 \(\mathbb C\)，\(f\in C(\mathbb T;\mathbb F)\)，其中 \(\mathbb T=\mathbb R/\mathbb Z\) 取归一化 Lebesgue 测度。以 \(\omega_f\) 表示关于最短弧距的连续模，以 \(\mathcal E_n(f)\) 表示次数不超过 \(n\) 的三角多项式的一致最佳逼近误差。存在绝对常数 \(C>0\)，使对每个整数 \(n\geq1\)，都有
\[
 \mathcal E_n(f)\leq C\omega_f(1/n).
\]
具体地，取 \(m=\lfloor n/2\rfloor+1\)，则
\[
 Q_nf\coloneqq f*J_m
\]
是次数不超过 \(n\) 的三角多项式，并满足
\[
 \lVert f-Q_nf\rVert_\infty\leq C\omega_f(1/n),
 \quad \lVert Q_nf\rVert_\infty\leq\lVert f\rVert_\infty.
\]
算子 \(Q_n\) 保持常数；对实值函数还保持非负性。
\end{theorem}
\begin{proof}
把 \(J_m\) 展为有限 Fourier 和，逐项积分即可写出
\[
 (f*J_m)(x)=\sum_{|k|\leq2(m-1)}
     \widehat f(k)\widehat J_m(k)e^{2\pi\mathrm i kx}.
\]
这里不需要 \(f\) 本身的 Fourier 级数收敛，因为被展开的是一个有限多项式。由 \(2(m-1)\leq n\)，所得近似函数符合次数限制。

利用 \(J_m\geq0\) 和 \(\int J_m=1\)，有
\[
 |f(x)-(f*J_m)(x)|
 \leq\int_{-1/2}^{1/2}|f(x)-f(x-t)|J_m(t)\,\mathrm dt.
\]
沿圆周最短弧把长度 \(|t|\) 分成至多 \(1+m|t|\) 段，每段长度不超过 \(1/m\)，由三角不等式得到
\[
 |f(x)-f(x-t)|\leq(1+m|t|)\omega_f(1/m).
\]
代回后，\cref{b2:lem:approximation-jackson-kernel}给出
\[
 \lVert f-f*J_m\rVert_\infty
 \leq\omega_f(1/m)\left(1+m\int_{-1/2}^{1/2}|t|J_m(t)\,\mathrm dt\right)
 \leq C\omega_f(1/m).
\]
由于 \(m\geq n/2\)，再把长度不超过 \(2/n\) 的短弧分为两段，便有 \(\omega_f(1/m)\leq2\omega_f(1/n)\)。这些估计也包括 \(n=1\)：此时 \(m=1\)、\(J_1=1\)，输出就是 \(f\) 的平均值。收缩性、保持常数和保持非负性都由非负核的单位质量直接得到。
\end{proof}

这一结果称为\term[jackson-dingli]{Jackson 定理}[Jackson's theorem]。构造出的 \(Q_nf\) 给出了最佳误差的一个上界；具体函数的最佳误差还可能更小。下面将它与前章的 Fejér 平均比较，看看核的选择怎样影响收敛速度。

\ProseParagraphBreak

与 Fejér 平均的差别可以在一个具体函数上看清。取 \(f(t)=\mathrm d_{\mathbb T}(t,0)\)，它是 Lipschitz 常数为 \(1\) 的连续周期函数。在原点，Fejér 逼近的误差恰为
\[
 (f*F_{m-1})(0)-f(0)
 =\int_{-1/2}^{1/2}|t|F_{m-1}(t)\,\mathrm dt.
\]
将积分分为 \(|t|\leq1/m\) 与其余部分，上述核估计给出上界 \(C\log(m+1)/m\)。这个对数不是估计方法造成的假象：对 \(m\geq4\)，在
\[
 I_k=\left[\frac{k+1/4}{m},\frac{k+3/4}{m}\right],
 \quad 1\leq k\leq\lfloor m/2\rfloor-1,
\]
有 \(\sin^2(m\pi t)\geq1/2\)，而 \(\sin^2(\pi t)\leq\pi^2t^2\)。因而每个 \(I_k\) 对第一矩的贡献至少为 \(c/(m(k+1))\)；对 \(k\) 求和得到 \(c'\log(m+1)/m\) 的下界。相比之下，Jackson 构造对同一函数的误差至多为 \(C/n\)。Fejér 平均很好地解决了收敛问题，若还要得到一般 Lipschitz 函数的最佳次数量级，就需要更集中的核。

\subsection{代数多项式的直接逼近}
\label{b2:subsec:approximation-jackson-algebraic}

闭区间上的函数一般不满足端点值相等，直接作周期延拓可能产生跳跃。余弦变换提供另一条路线：它让圆周先沿区间走一遍，再沿原路返回，因而两端都自动接上。

\begin{theorem}[Jackson：闭区间情形]
\label{b2:thm:approximation-jackson-algebraic}
设 \(a<b\) 为实数，\(\mathbb F=\mathbb R\) 或 \(\mathbb C\)，且 \(f\in C([a,b];\mathbb F)\)。令 \(E_n(f)\) 为次数不超过 \(n\) 的代数多项式在 \([a,b]\) 上的一致最佳逼近误差，\(\omega_f\) 为通常距离下的连续模。存在绝对常数 \(C>0\)，使对所有整数 \(n\geq1\)，
\[
 E_n(f)\leq C\omega_f\left(\frac{b-a}{n}\right).
\]
特别地，若 \(0<\alpha\leq1\)、\(H\geq0\)，且 \(|f(x)-f(y)|\leq H|x-y|^\alpha\) 对所有 \(x,y\in[a,b]\) 成立，则
\[
 E_n(f)\leq CH\left(\frac{b-a}{n}\right)^\alpha.
\]
\end{theorem}
\begin{proof}
设 \(c=(a+b)/2\)、\(r=(b-a)/2\)，并在单位圆周上定义
\[
 g(t)=f(c+r\cos(2\pi t)).
\]
这是连续、偶的周期函数。余弦的导数绝对值不超过 \(1\)，所以沿最短圆弧取代表后，有
\[
 |(c+r\cos(2\pi t))-(c+r\cos(2\pi s))|
 \leq2\pi r\mathrm d_{\mathbb T}(t,s).
\]
从而 \(\omega_g(\delta)\leq\omega_f(\pi(b-a)\delta)\)。这里连续模的参数允许超过区间长度；超过以后，它仅仅等于函数在整个区间上的最大振幅。

对 \(g\) 使用周期 Jackson 构造，得到 \(q_n=g*J_m\)，其中 \(m=\lfloor n/2\rfloor+1\)。核与 \(g\) 都是偶函数，换元 \(t\mapsto-t\) 便知 \(q_n\) 也偶，因此它的有限三角展开只有余弦项：
\[
 q_n(t)=a_0+\sum_{k=1}^na_k\cos(2\pi kt).
\]
余弦递推式
\[
 \cos((k+1)\theta)=2\cos\theta\cos(k\theta)-\cos((k-1)\theta)
\]
表明 \(\cos(k\theta)=T_k(\cos\theta)\)，其中 \(T_0=1\)、\(T_1(u)=u\)、\(T_{k+1}(u)=2uT_k(u)-T_{k-1}(u)\)。因此
\[
 p_n(x)=a_0+\sum_{k=1}^na_kT_k\left(\frac{x-c}{r}\right)
\]
是次数不超过 \(n\) 的代数多项式，而且 \(p_n(c+r\cos(2\pi t))=q_n(t)\)。由于这个余弦参数化覆盖整个 \([a,b]\)，两种一致误差恰好相同。于是
\[
 \lVert f-p_n\rVert_{\infty,[a,b]}
 =\lVert g-q_n\rVert_{\infty,\mathbb T}
 \leq C\omega_g(1/n)
 \leq C\omega_f\left(\frac{\pi(b-a)}n\right).
\]
在区间中把任意一段分为四段，连续模的次可加性给出
\(\omega_f(\pi\delta)\leq4\omega_f(\delta)\)。把 \(4\) 吸收入绝对常数，就得到第一式；代入 Hölder 条件即得第二式。
\end{proof}

例如 \(f(x)=|x|\) 在 \([-1,1]\) 上的连续模满足 \(\omega_f(\delta)\leq\delta\)，所以存在次数为 \(n\) 的代数多项式使一致误差不超过 \(C/n\)，尽管 \(f\) 在原点不可微。这里的构造先把尖点变为 \(|\cos(2\pi t)|\) 的周期尖点，用 Jackson 核平滑，再把偶三角多项式还原为代数多项式。

\ProseParagraphBreak

Bernstein 取样构造与 Jackson 构造的用途有所不同：前者直接使用网格上的函数值，并有清楚的非负权重；后者通过卷积得到更好的通用次数估计。下一节把问题反过来：如果最佳误差本来就衰减得快，函数本身是否必须有相应的连续模？回答这个问题需要先限制有限次数多项式的变化速度。

% !TeX root = ../../Book2.tex
\section{Bernstein 估计与逆定理}
\label{b2:sec:approximation-inverse}

待撰写。

\subsection{三角多项式的 Bernstein 估计}
\label{b2:subsec:020903}

待撰写。

\subsection{逼近速度与光滑性的逆向关系}
\label{b2:subsec:approximation-inverse-smoothness}

待撰写。

% !TeX root = ../../Book2.tex
\OptionalSection{高阶逼近与端点问题}{b2:sec:approximation-higher-order}

一阶连续模只测量函数值的变化。即使函数无限次可微，只要不是常数，它的一阶连续模通常也只有一阶小量；因而仅用它不能分辨二阶、三阶光滑性。上节在 \(0<\alpha<1\) 时把逼近速度与 Hölder 连续性联系起来；光滑阶达到整数时，二者的关系还会出现新的现象。本节先说明高阶差分怎样消去低阶变化，再分别处理光滑指数的整数端点和空间区间的几何端点。这两种“端点”是不同的问题。

\subsection{高阶连续模}
\label{b2:subsec:approximation-higher-moduli}

从二阶开始看最直接。若 \(f\) 在一小段上近似为仿射函数，则 \(f(x+2h)-2f(x+h)+f(x)\) 比 \(f(x+h)-f(x)\) 更小：常数项和一次项恰好抵消。高阶差分把这一消去过程重复进行，而不要求预先知道导数存在。

\begin{definition}\label{b2:def:approximation-higher-modulus}
设 \(\mathbb F=\mathbb R\) 或 \(\mathbb C\)，\(f\in C(\mathbb T;\mathbb F)\)，其中 \(\mathbb T=\mathbb R/\mathbb Z\)，且 \(r\geq1\) 为整数。对 \(h\in\mathbb R\)，定义 \(r\) 阶差分及相应的\term[gaojie-lianxu-mo]{高阶连续模}[higher order modulus of smoothness]为
\[
 \Delta_h^r f(x)=\sum_{j=0}^r(-1)^{r-j}\binom rj f(x-jh),
 \quad
 \omega_r(f,\delta)=\sup_{|h|\leq\delta}\|\Delta_h^r f\|_\infty
 \quad(\delta\geq0).
\]
自变量均按模 \(1\) 理解。特别地，\(\omega_1(f,\delta)=\omega_f(\delta)\) 是圆周上的一阶连续模。
\end{definition}
\symidx{\(\omega_r(f,\delta)\)：周期函数的高阶连续模}

这里允许负步长，且 \(\delta\) 可以大于半个周期；这使伸缩估计无需另外分段。\(\Delta_h^r=(\tau_h-I)^r\)，其中 \(\tau_hf(x)=f(x-h)\)，所以差分是平移的有限组合。连续性与紧性保证上式范数有限，且 \(\omega_r(f,\delta)\to0\) 当 \(\delta\downarrow0\)。

\begin{proposition}\label{b2:prop:approximation-difference-estimates}
设 \(r\geq1\) 为整数，\(\mathbb F=\mathbb R\) 或 \(\mathbb C\)，且 \(f\in C(\mathbb T;\mathbb F)\)。对 \(\delta>0\)、\(\lambda>0\)，有
\[
 \omega_r(f,\lambda\delta)\leq(1+\lambda)^r\omega_r(f,\delta),
 \quad
 \omega_r(f,\delta)\leq2^{r-1}\omega_1(f,\delta).
\]
若进一步 \(f\in C^r(\mathbb T;\mathbb F)\)，则
\[
 \omega_r(f,\delta)\leq\delta^r\|f^{(r)}\|_\infty.
\]
\end{proposition}
\begin{proof}
对任意整数 \(m\geq1\)，平移算子相互交换，并且
\[
 \tau_{mh}-I=(I+\tau_h+\cdots+\tau_{(m-1)h})(\tau_h-I).
\]
取 \(r\) 次幂。每个平移在一致范数下都是等距算子，括号中有 \(m\) 项，故 \(\|\Delta_{mh}^r f\|_\infty\leq m^r\|\Delta_h^r f\|_\infty\)。对 \(|u|\leq\lambda\delta\)，取 \(m=\max\{1,\lceil\lambda\rceil\}\)、\(h=u/m\)，得到第一式。第二式来自 \(\Delta_h^r=(\tau_h-I)^{r-1}\Delta_h\) 和 \(\|\tau_h-I\|\leq2\)。

对 \(C^r\) 函数，反复使用微积分基本定理得到
\[
 \Delta_h^r f(x)=\int_0^h\cdots\int_0^h
 (-1)^rf^{(r)}(x-t_1-\cdots-t_r)\,\mathrm dt_1\cdots\mathrm dt_r.
\]
负的 \(h\) 按有向积分解释。取绝对值后，积分区域的体积是 \(|h|^r\)，于是得到最后的估计。
\end{proof}

高阶差分消去了低阶变化，也改变了误差的估计方式。为了利用这种抵消，需要让核的尾部衰减得更快，再用它平均高阶差分。将 Fejér 核取更高的整数次幂，就能得到所需的衰减。

\begin{theorem}[Jackson：高阶周期形式]\label{b2:thm:approximation-jackson-higher}
对每个整数 \(r\geq1\)，存在只依赖 \(r\) 的常数 \(C_r>0\)，使任意 \(\mathbb F=\mathbb R\) 或 \(\mathbb C\)、\(f\in C(\mathbb T;\mathbb F)\) 及整数 \(n\geq1\) 均满足
\[
 \mathcal E_n(f)\leq C_r\omega_r(f,1/n).
\]
这里 \(\mathcal E_n(f)\) 是用次数不超过 \(n\) 的三角多项式作一致逼近的最佳误差。因此，若 \(f\in C^r(\mathbb T;\mathbb F)\)，则 \(\mathcal E_n(f)\leq C_r n^{-r}\|f^{(r)}\|_\infty\)。
\end{theorem}
\begin{proof}
固定整数 \(s\geq2\)，使 \(2s>r+1\)。沿用\cref{b2:def:02-fejer-kernel}的编号，\(F_{m-1}\) 的次数为 \(m-1\)。对 \(m\geq1\)，定义
\[
 a_{m,s}=\int_{-1/2}^{1/2}F_{m-1}(t)^s\,\mathrm dt,
 \quad K_{m,s}(t)=\frac{F_{m-1}(t)^s}{a_{m,s}}.
\]
这是非负、积分为 \(1\)、次数为 \(s(m-1)\) 的三角多项式。

先核对所需估计。由正弦商公式，在 \(|t|\leq1/2\) 上
\[
 F_{m-1}(t)\leq\min\left\{m,\frac1{4mt^2}\right\}
 \leq\frac{C m}{(1+m|t|)^2}.
\]
在 \(|t|\leq1/(2m)\) 上，用 \(|\sin(m\pi t)|\geq2m|t|\)、\(|\sin(\pi t)|\leq\pi|t|\)，可得 \(F_{m-1}(t)\geq4m/\pi^2\)，零点处按连续性理解。因而下界通过在这段长 \(1/m\) 的区间积分取得，上界通过代换 \(u=mt\) 取得：
\[
 c_s m^{s-1}\leq a_{m,s}\leq C_s m^{s-1},
 \quad K_{m,s}(t)\leq\frac{C_s m}{(1+m|t|)^{2s}}.
\]
这些估计也包含 \(m=1\)。由 \(2s>r+1\)，进一步有
\[
 \int_{-1/2}^{1/2}(1+m|t|)^rK_{m,s}(t)\,\mathrm dt\leq C_r.
\]
这里及以下选定的 \(s\) 只依赖 \(r\)，常数可统一记为 \(C_r\)。

现在取 \(m=\lfloor n/s\rfloor+1\)，使 \(s(m-1)\leq n\) 且 \(n/m\leq s\)，并令
\[
 Q_n(x)=f(x)-(-1)^r\int_{-1/2}^{1/2}\Delta_t^r f(x)K_{m,s}(t)\,\mathrm dt.
\]
为什么这个表达式是三角多项式，而不是仍含未知频率的 \(f\)？差分中 \(j=0\) 的项是 \((-1)^rf(x)\)，积分后恰好与前面的 \(f(x)\) 抵消。因此 \(Q_n\) 是下列 \(r\) 个函数的有限线性组合：
\[
 G_j(x)=\int_{\mathbb T}f(x-jt)K_{m,s}(t)\,\mathrm dt,
 \quad 1\leq j\leq r.
\]
连续函数在紧圆周上有界，故 Fubini 可用于计算其 Fourier 系数。换元 \(y=x-jt\) 给出
\[
 \widehat G_j(k)=\widehat f(k)\widehat K_{m,s}(jk).
\]
当 \(|jk|>s(m-1)\) 时右端为零。取由其有限个非零系数组成的三角多项式，它与 \(G_j\) 的全部 Fourier 系数一致；\cref{b2:cor:02-uniqueness}给出几乎处处相等，而二者连续，故处处相等。这样 \(G_j\) 的次数不超过 \(\lfloor s(m-1)/j\rfloor\)，\(Q_n\) 的次数不超过 \(n\)。实值输入也给出实值 \(Q_n\)。

最后直接从构造估计误差。由连续模的伸缩估计和核的非负性，
\[
 \begin{aligned}
 \|f-Q_n\|_\infty
 &\leq\int_{-1/2}^{1/2}\omega_r(f,|t|)K_{m,s}(t)\,\mathrm dt\\
 &\leq\omega_r(f,1/n)
       \int_{-1/2}^{1/2}(1+n|t|)^rK_{m,s}(t)\,\mathrm dt
 \leq C_r\omega_r(f,1/n).
 \end{aligned}
\]
最后一步使用 \(n/m\leq s\) 及刚建立的矩估计。取最佳误差即得结论，再用\cref{b2:prop:approximation-difference-estimates}得到 \(C^r\) 情形。
\end{proof}

这里 \(Q_n\) 在 \(r\geq2\) 时一般不是正平均；正的是用来估计差分的核。正性负责误差控制，差分的系数负责消去低阶变化，二者承担不同任务。差分估计还适用于没有经典导数的连续函数；下一小节的例子将显示这种描述的用处。

\subsection{逆定理与高阶逼近的边界}
\label{b2:subsec:020904}

先回顾\secref{b2:sec:approximation-inverse}中二进尺度求和的关键位置。误差 \(2^{-j\alpha}\) 乘上导数所需的次数因子 \(2^j\)，得到 \(2^{j(1-\alpha)}\)。当 \(0<\alpha<1\) 时，这个几何级数由最后一项控制；在 \(\alpha=1\) 时，每一项同阶，累积项数变成对数。这说明旧证明只给出 \(h\log(1/h)\) 量级。下面的函数表明这项损失确实可能出现。

\begin{proposition}\label{b2:prop:approximation-integer-counterexample}
函数
\[
 f(x)=\sum_{j=1}^{\infty}2^{-j}\sin(2\pi 2^j x),\quad x\in\mathbb T,
\]
是实值连续周期函数，满足 \(\mathcal E_n(f)\leq2/n\)（整数 \(n\geq1\)），但不是 Lipschitz 连续的。它的二阶连续模却满足 \(\omega_2(f,\delta)\leq C\delta\)（\(0<\delta\leq1/2\)），其中 \(C\) 为绝对常数。
\end{proposition}
\begin{proof}
各项一致范数之和为 \(\sum_{j\geq1}2^{-j}=1\)，故级数一致收敛，和函数连续。若 \(2^N\leq n<2^{N+1}\)，取前 \(N\) 项作为次数不超过 \(n\) 的逼近，尾项范数至多 \(2^{-N}\leq2/n\)；\(n=1\) 时取零多项式。

为检查 Lipschitz 性，取 \(h_N=2^{-N-2}\)。当 \(1\leq j\leq N\) 时，\(2\pi2^jh_N\in[0,\pi/2]\)，由 \(\sin u\geq2u/\pi\) 得
\[
 2^{-j}\sin(2\pi2^jh_N)\geq4h_N.
\]
当 \(j=N+1\) 时正弦自变量为 \(\pi\)，随后为 \(\pi\) 的整数倍，所有尾项恰好为零。于是 \(f(0)=0\) 且
\[
 \frac{|f(h_N)-f(0)|}{h_N}\geq4N\longrightarrow\infty.
\]
不存在统一的 Lipschitz 常数。

另一方面，指数函数的二阶差分乘子是 \((e^{-2\pi\mathrm i 2^jh}-1)^2\)，所以每项的二阶差分范数至多
\[
 2^{-j}\min\{4,(2\pi2^j h)^2\}.
\]
对 \(0<|h|\leq1/2\)，选 \(N\geq1\) 使 \(2^{-N-1}<|h|\leq2^{-N}\)。低频之和至多 \(C h^2\sum_{j\leq N}2^j\leq C|h|\)，高频之和至多 \(4\sum_{j>N}2^{-j}\leq C|h|\)。对 \(|h|\leq\delta\) 取上确界即得结论。
\end{proof}

\(\omega_2(f,\delta)=O(\delta)\) 常用于描述整数光滑阶附近比 Lipschitz 更宽的函数类。这个例子已经足够说明：把上节 Hölder 等价中的 \(\alpha<1\) 改成 \(\alpha\leq1\)，并保持右侧为一阶 Hölder 条件，会得到错误命题。若研究更高整数阶，也必须先选合适的高阶差分尺度，而不是继续使用一阶连续模。

\ProseParagraphBreak

还有一个更基础的限制：若区间上的函数满足 \(|f(x)-f(y)|\leq C|x-y|^\alpha\) 且 \(\alpha>1\)，把任意线段分成 \(m\) 等份可得
\[
 |f(x)-f(y)|\leq C|x-y|^\alpha m^{1-\alpha}\longrightarrow0.
\]
它只能是常数。光滑函数的一阶变化通常仍是线性的，所以高阶光滑性不可能用“一阶差小于距离的高次幂”来定义。高阶差分和导数的 Hölder 条件正是为避免这一矛盾而引入的。

\subsection{端点情形}
\label{b2:subsec:approximation-endpoints}

区间上的代数多项式可以写成 \(p(\cos\theta)\)，再使用三角多项式工具；但 \(x=\cos\theta\) 不是处处保持距离尺度的变换。区间内部 \(\sin\theta\) 不为零，角度变化与位置变化成比例；在 \(x=1\) 附近则有 \(1-\cos\theta\sim\theta^2/2\)。因此同一个角度小量，在端点变成二次的位置小量。

\begin{proposition}\label{b2:prop:approximation-algebraic-derivative}
设 \(p\) 为实或复系数代数多项式，次数不超过整数 \(n\geq1\)，范数取在 \([-1,1]\) 上。则
\[
 \sqrt{1-x^2}\,|p'(x)|\leq n\|p\|_\infty\quad(-1<x<1),
 \quad \|p'\|_\infty\leq\frac\pi2 n^2\|p\|_\infty.
\]
此外，Chebyshev 多项式满足 \(\|T_n\|_\infty=1\) 和 \(T_n'(1)=n^2\)，故全区间导数估计中的 \(n^2\) 不能统一换成 \(n\)。第二个不等式在此只说明次数的正确尺度，不主张常数 \(\pi/2\) 最优。
\end{proposition}
\begin{proof}
令 \(v(\theta)=p(\cos\theta)\)。这是角度周期 \(2\pi\)、次数不超过 \(n\) 的三角多项式。把\cref{b2:thm:approximation-bernstein-trigonometric}中的自变量按 \(\theta=2\pi t\) 换元，得到 \(\|v'\|_\infty\leq n\|p\|_\infty\) 及 \(\|v''\|_\infty\leq n^2\|p\|_\infty\)。链式法则给出
\[
 v'(\theta)=-\sin\theta\,p'(\cos\theta),
\]
因而得到第一个估计。

由于 \(v'(0)=v'(\pi)=0\)，微积分基本定理同时给出
\[
 |v'(\theta)|\leq n^2\|p\|_\infty\min\{\theta,\pi-\theta\}
 \quad(0\leq\theta\leq\pi).
\]
使用 \(\sin\theta\geq(2/\pi)\min\{\theta,\pi-\theta\}\)，相除后得到第二个估计，端点由 \(p'\) 连续性补上。最后，\(T_n(\cos\theta)=\cos(n\theta)\) 给出
\[
 T_n'(1)=\lim_{\theta\to0}\frac{\cos(n\theta)-1}{\cos\theta-1}=n^2,
\]
而 \(|\cos(n\theta)|\leq1\) 并在端点取等号，所以其一致范数为 \(1\)。
\end{proof}

这也影响逆定理的表述。令 \(f(x)=\sqrt{1-x^2}\)（\(-1\leq x\leq1\)）。对应的周期函数 \(g(t)=f(\cos(2\pi t))=|\sin(2\pi t)|\) 是 Lipschitz 函数，常数可取 \(2\pi\)。对 \(g\) 使用周期 Jackson 逼近，再取偶部并把余弦换成 Chebyshev 多项式，就像\secref{b2:sec:approximation-jackson}中的代数转移证明一样，得到 \(E_n(f)\leq C/n\)。然而
\[
 \frac{|f(1-h)-f(1)|}{h^\alpha}
 =h^{1/2-\alpha}\sqrt{2-h}\longrightarrow\infty
 \quad(\alpha>1/2).
\]
因此即使固定 \(1/2<\alpha<1\)，代数最佳误差 \(E_n(f)=O(n^{-\alpha})\) 也不能推出闭区间通常距离下的 \(\alpha\)-Hölder 连续性。周期逆定理没有错；错误发生在忽略了余弦变换的端点退化。

\ProseParagraphBreak

在闭子区间 \([-1+\varepsilon,1-\varepsilon]\) 上，\(\sqrt{1-x^2}\) 有正下界，可以恢复与角度相当的局部尺度。若希望在整个闭区间获得精确的高阶正逆理论，则需要与端点相适应的加权连续模：差分步长也要随位置调整，才能同时反映内部与端点附近的变化。


\begin{chaptersummary}
Bernstein 多项式通过离散正权重把连续函数变为可计算的多项式，二阶矩控制一致误差；Chebyshev 交错定理则通过误差的等幅换号判定实区间上的最佳性，并给出唯一性。两者分别解决“怎样构造”和“怎样判断最优”，不能把一个具体算子的误差自动视为最佳误差。

周期问题中，平方 Fejér 核同时保留有限次数、非负性与单位质量，又改善尾部衰减，给出 Jackson 估计。反向推导用 Bernstein 导数估计控制每个频率尺度的变化，得到 \(0<\alpha<1\) 时 \(\mathcal E_n(f)=O(n^{-\alpha})\) 与 \(\alpha\)-Hölder 连续性的等价。高阶差分继续消去低阶变化，但整数光滑阶须另作判断；余弦换元又在区间端点改变距离尺度。这些限制都是结论的一部分，应当与定理一起记住。
\end{chaptersummary}

\BookExercises

% 来源：自拟。目标：以子空间反例区别有限维最佳存在与多项式空间的唯一性。
\begin{exercise}[子空间中的非唯一性]\label{b2:ex:approximation-nonunique}
在 \(C([-1,1];\mathbb R)\) 中取 \(V=\{ax^2:a\in\mathbb R\}\)。求常函数 \(1\) 到 \(V\) 的一致最佳误差及全部最佳逼近。解释这为什么不与 Chebyshev 交错定理及其唯一性推论矛盾。
\end{exercise}
\begin{solution}
所有 \(ax^2\) 都在 \(x=0\) 取零，所以误差至少为 \(1\)。令 \(t=x^2\in[0,1]\)，仿射函数 \(1-at\) 的绝对值在端点达到最大，因而
\[
 \|1-ax^2\|_\infty=\max\{1,|1-a|\}.
\]
最小值为 \(1\)，且恰好 \(a\in[0,2]\) 时取到。\(V\) 虽然有限维，却不是全部次数不超过某个整数的多项式空间。这里的非零基函数 \(x^2\) 在区间中有零点，不能用“一维空间中非零函数无零点”的性质完成交错唯一性论证。有限维紧性给出最优者存在，唯一性还需要空间的额外结构。
\end{solution}

% 来源：自拟。目标：同一目标函数的二次一致逼近与正交逼近比较。
\begin{exercise}[绝对值函数的两种二次逼近]\label{b2:ex:approximation-absolute-value}
求 \(f(x)=|x|\) 在 \([-1,1]\) 上的最佳一致二次逼近，以及对 Lebesgue 测度的最佳 \(L^2\) 二次逼近。分别计算两者的一致误差。
\end{exercise}
\begin{solution}
先取 \(p(x)=x^2+1/8\)。在 \(0\leq x\leq1\) 上，误差为 \(x-x^2-1/8\)，最小值 \(-1/8\) 在 \(0,1\) 取到，最大值 \(1/8\) 在 \(1/2\) 取到。误差是偶函数，于是它在
\[
 -1,-1/2,0,1/2,1
\]
依次取 \(-1/8,1/8,-1/8,1/8,-1/8\)。任取其中连续的四点即满足二次交错判据。因此 \(p\) 是唯一的最佳一致逼近，误差为 \(1/8\)。

对 \(L^2\) 问题，奇偶正交性使最佳逼近的一次项系数为零。写成 \(q(x)=a x^2+c\)，误差与 \(1,x^2\) 正交给出
\[
 2c+\frac23a=1,\quad \frac23c+\frac25a=\frac12.
\]
解得 \(a=15/16,c=3/16\)。它的 \(L^2\) 最佳性由正交投影定理保证。\(0\leq x\leq1\) 上误差 \(x-15x^2/16-3/16\) 的内部最大值在 \(x=8/15\) 取到，为 \(19/240\)；两个端点的值分别为 \(-3/16,-1/8\)。故 \(\|f-q\|_\infty=3/16>1/8\)。减小平方平均误差会允许某些点的误差更大。
\end{solution}

% 来源：自拟。目标：固定次数反复应用正算子，与次数趋于无穷的逼近相区分。
\begin{exercise}[反复应用 Bernstein 算子]\label{b2:ex:approximation-bernstein-iterates}
固定整数 \(n\geq2\)，以 \(B_n^k\) 表示将 \([0,1]\) 上的 Bernstein 算子复合 \(k\) 次。对任意 \(f\in C([0,1];\mathbb R)\)，证明
\[
 B_n^k f\longrightarrow (1-x)f(0)+xf(1)
\]
一致成立。提示：利用 \(u(x)=x(1-x)\) 及 \(B_nu=(1-1/n)u\)。
\end{exercise}
\begin{solution}
令 \(\ell(x)=(1-x)f(0)+xf(1)\)、\(g=f-\ell\)，则 \(g(0)=g(1)=0\)，且 \(B_n\ell=\ell\)。给定 \(\varepsilon>0\)，由端点连续性，在两端各一小段上 \(|g|\leq\varepsilon\)；剩下的紧区间上 \(u\) 有正下界。因而存在有限的 \(C_\varepsilon\geq0\)，使整个区间上
\[
 |g(x)|\leq\varepsilon+C_\varepsilon u(x).
\]
二阶矩计算给出 \(B_n(x^2)=x^2+x(1-x)/n\)，所以提示中的特征关系成立。正性、保常数及逐次应用产生
\[
 |B_n^k g|\leq B_n^k|g|
 \leq\varepsilon+C_\varepsilon(1-1/n)^ku.
\]
先令 \(k\to\infty\)，再令 \(\varepsilon\downarrow0\)，便得一致收敛。固定采样次数反复平均会逐渐抹去内部形状，只留下端点仿射插值；这与 \(n\to\infty\) 时 \(B_nf\to f\) 是两个不同的极限。
\end{solution}

% 来源：自拟。目标：Fourier系数作逼近下界，构造任意幂次的实际最佳误差。
\begin{exercise}[可精确判断量级的逼近率]\label{b2:ex:approximation-lacunary-rate}
证明：对 \(f\in C(\mathbb T;\mathbb C)\) 及 \(|k|>n\)，有 \(\mathcal E_n(f)\geq|\widehat f(k)|\)。再对任意 \(\alpha>0\)，定义
\[
 f_\alpha(x)=\sum_{j=1}^\infty2^{-j\alpha}\cos(2\pi2^jx).
\]
证明它连续，且存在只依赖 \(\alpha\) 的 \(c_\alpha,C_\alpha>0\)，使 \(c_\alpha n^{-\alpha}\leq\mathcal E_n(f_\alpha)\leq C_\alpha n^{-\alpha}\) 对整数 \(n\geq1\) 成立。
\end{exercise}
\begin{solution}
对任意 \(p\in\mathcal T_n\)，高于次数的系数为零，所以
\[
 |\widehat f(k)|=|\widehat{f-p}(k)|
 \leq\int_{\mathbb T}|f-p|\leq\|f-p\|_\infty.
\]
取下确界得到下界工具。定义 \(f_\alpha\) 的级数由收敛的正项几何级数控制，因而一致收敛，且可以逐项计算 Fourier 系数。

取 \(N\geq0\) 使 \(2^N\leq n<2^{N+1}\)。前 \(N\) 项是可容许的逼近，尾和至多
\[
 \frac{2^{-(N+1)\alpha}}{1-2^{-\alpha}}\leq C_\alpha n^{-\alpha}.
\]
另一方面，在频率 \(2^{N+1}>n\) 处，系数恰好为 \(2^{-1-(N+1)\alpha}\)。由于 \(n\geq2^N\)，它至少为 \(2^{-1-\alpha}n^{-\alpha}\)，给出下界。\(N=0\) 时前零项理解为零函数，论证也成立。这里下界检查了任何次数不超过 \(n\) 的候选者，而不只检查某一种截断算法。
\end{solution}

% 来源：自拟。目标：从Lagrange插值构造最佳误差的有限节点证书。
\begin{exercise}[节点值给出的误差下界]\label{b2:ex:approximation-dual-certificate}
设 \(a\leq x_0<\cdots<x_{n+1}\leq b\)、\(a<b\)，其中 \(n\geq0\) 为整数。令
\[
 w_j=\frac1{\prod_{\ell\ne j}|x_j-x_\ell|},\quad
 \lambda_j=\frac{w_j}{\sum_{\ell=0}^{n+1}w_\ell},\quad
 L(f)=\sum_{j=0}^{n+1}(-1)^{n+1-j}\lambda_j f(x_j).
\]
证明对任意 \(f\in C([a,b];\mathbb R)\)，有 \(E_n(f)\geq|L(f)|\)。说明当这些点是某个最佳误差的交错极值点时，为什么下界取等。
\end{exercise}
\begin{solution}
对 \(p\in\mathcal P_n\)，用 \(n+2\) 个节点写 Lagrange 插值：
\[
 p(x)=\sum_{j=0}^{n+1}p(x_j)
 \frac{\prod_{\ell\ne j}(x-x_\ell)}{\prod_{\ell\ne j}(x_j-x_\ell)}.
\]
右侧与左侧之差次数至多为 \(n+1\)，又在全部 \(n+2\) 个节点为零，故恒为零。比较 \(x^{n+1}\) 的系数得到
\[
 \sum_{j=0}^{n+1}\frac{p(x_j)}{\prod_{\ell\ne j}(x_j-x_\ell)}=0.
\]
分母的符号为 \((-1)^{n+1-j}\)，所以 \(L(p)=0\)。又因 \(\lambda_j>0\)、\(\sum_j\lambda_j=1\)，
\[
 |L(f)|=|L(f-p)|\leq\sum_j\lambda_j|f(x_j)-p(x_j)|
 \leq\|f-p\|_\infty.
\]
取下确界即可。若 \(f(x_j)-p_*(x_j)=\varepsilon(-1)^jE_n(f)\)，则所有带符号的求和项都具有同一符号，得 \(L(f)=\varepsilon(-1)^{n+1}E_n(f)\)。因此有限个取样值便能证明任何候选都不可能优于该误差。
\end{solution}

% 来源：自拟。目标：最佳一致误差经Parseval与分块Cauchy--Schwarz控制Fourier绝对收敛。
\begin{exercise}[逼近率与 Fourier 绝对收敛]\label{b2:ex:approximation-fourier-absolute}
设 \(f\in C(\mathbb T;\mathbb C)\)，且
\[
 \sum_{j=0}^\infty 2^{j/2}\mathcal E_{2^j}(f)<\infty.
\]
证明 \(\sum_{k\in\mathbb Z}|\widehat f(k)|<\infty\)，并证明 Fourier 级数一致收敛到 \(f\)。由此推出：若 \(f\) 是 \(\alpha\)-Hölder 连续函数，\(1/2<\alpha<1\)，则其 Fourier 级数绝对且一致收敛。
\end{exercise}
\begin{solution}
选 \(p_j\in\mathcal T_{2^j}\) 为一致最佳逼近。对块 \(A_j=\{k:2^j<|k|\leq2^{j+1}\}\)，\(p_j\) 的相应系数为零。Bessel 不等式与圆周测度为一给出
\[
 \sum_{k\in A_j}|\widehat f(k)|^2
 \leq\|f-p_j\|_2^2\leq\mathcal E_{2^j}(f)^2.
\]
该块有 \(2^{j+1}\) 个整数，因此 Cauchy--Schwarz 给出
\[
 \sum_{k\in A_j}|\widehat f(k)|\leq2^{(j+1)/2}\mathcal E_{2^j}(f).
\]
对块求和，再加上三个低频 \(k=0,\pm1\)，即得绝对收敛。所得级数一致收敛到某个连续函数 \(g\)，逐项积分说明其 Fourier 系数与 \(f\) 相同。\cref{b2:cor:02-uniqueness}给出几乎处处相等，再由连续性得到处处相等。

若 \(f\) 是所述 Hölder 函数，周期 Jackson 定理给出 \(\mathcal E_{2^j}(f)\leq C2^{-j\alpha}\)，所需级数由 \(\sum_j2^{-j(\alpha-1/2)}\) 控制。该论证在 \(\alpha=1/2\) 时不再给出收敛，因而不能把本题的严格不等号直接删去。
\end{solution}

% 来源：自拟。目标：可求和的导数增量恢复经典可微性，不依赖形式逐项微分。
\begin{optionalexercise}[由更快逼近恢复导数]\label{b2:ex:approximation-derivative-recovery}
设 \(f\in C(\mathbb T;\mathbb C)\)，且 \(\sum_{j=0}^\infty2^j\mathcal E_{2^j}(f)<\infty\)。令 \(p_j\) 为次数不超过 \(2^j\) 的一致最佳逼近。证明 \(f\in C^1(\mathbb T;\mathbb C)\)，且 \(p_j'\to f'\) 一致成立。若进一步 \(\mathcal E_n(f)\leq A n^{-\alpha}\)、\(1<\alpha<2\)，证明 \(f'\) 为 \((\alpha-1)\)-Hölder 连续。
\end{optionalexercise}
\begin{solution}
最佳误差随次数下降，故
\[
 \|p_{j+1}-p_j\|_\infty\leq2\mathcal E_{2^j}(f).
\]
Bernstein 导数估计给出 \(\|p_{j+1}'-p_j'\|_\infty\leq C2^j\mathcal E_{2^j}(f)\)。假设使这些导数增量可求和，所以 \(p_j'\) 一致趋于一个连续周期函数 \(g\)。另一方面，Fejér 一致逼近说明 \(\mathcal E_n(f)\to0\)，从而 \(p_j\to f\) 一致。对 \(x\in[0,1]\)，在
\[
 p_j(x)-p_j(0)=\int_0^xp_j'(t)\,\mathrm dt
\]
中取极限，得到 \(f(x)-f(0)=\int_0^xg(t)\,\mathrm dt\)。因此 \(f'=g\)，端点也按周期延拓相容。这个积分论证才使极限可以求导。

在幂次假设下，对导数增量尾和估计得
\[
 \|f'-p_J'\|_\infty\leq C_\alpha A2^{-J(\alpha-1)}.
\]
若 \(2^J\leq n<2^{J+1}\)，则 \(p_J'\in\mathcal T_n\)，故 \(\mathcal E_n(f')\leq C_\alpha A n^{-(\alpha-1)}\)。由于 \(0<\alpha-1<1\)，对连续函数 \(f'\) 使用主线的周期逆定理，得到所需 Hölder 性。这里先建立导数存在，随后才把逆定理应用于导数。
\end{solution}

% 来源：自拟。目标：把端点导数n²与可见的极值点距离和缩放极限联系起来。
\begin{optionalexercise}[端点附近的振荡尺度]\label{b2:ex:approximation-endpoint-scale}
对 Chebyshev 多项式 \(T_n\)，证明从 \(x=1\) 到相邻极值点的距离为 \(1-\cos(\pi/n)\sim\pi^2/(2n^2)\)。再证明：对任意固定 \(S>0\)，当 \(n\) 足够大时，
\[
 T_n(1-s/n^2)\longrightarrow\cos\sqrt{2s}
\]
在 \(0\leq s\leq S\) 上一致成立，并解释为何端点的自然尺度是 \(n^{-2}\)。
\end{optionalexercise}
\begin{solution}
恒等式 \(T_n(\cos\theta)=\cos(n\theta)\) 给出相邻极值点 \(1\) 和 \(\cos(\pi/n)\)，两点函数值从 \(1\) 变到 \(-1\)。利用 \(1-\cos u=2\sin^2(u/2)\) 即得距离渐近式。

在 \(0\leq s\leq S\) 且 \(n\) 足够大时，令 \(u=\sqrt{s/2}/n\leq1/2\)。由半角公式，
\[
 n\arccos(1-s/n^2)=2n\arcsin u.
\]
由 \(\arcsin u-u=\int_0^u((1-t^2)^{-1/2}-1)\,\mathrm dt\)，且括号在 \([0,1/2]\) 上至多为 \(Ct^2\)，可得
\[
 \bigl|2n\arcsin u-\sqrt{2s}\bigr|\leq C S^{3/2}/n^2.
\]
余弦函数是 Lipschitz 的，于是得到一致收敛。次数虽只有 \(n\)，靠近端点时却能在长为 \(n^{-2}\) 的区间内产生固定大小的变化；这与 \(T_n'(1)=n^2\) 一致，也解释了为何不能把内部的 \(n^{-1}\) 尺度直接移到端点。
\end{solution}

% !TeX root = ../../Book2.tex
\chapter{\texorpdfstring{\(L^1\)}{L¹} Fourier 变换与卷积}
\label{b2:ch:03}
\BookChapterNumber{b2:ch:03}
% Axler2020Measure Section11C：本轮实读印刷363--367、374--377页。
% 本章改用d维Gauss正则化组织反演；先证变换公式，不预设反演或Plancherel。
\begin{chapterabstract}
从周期函数转到欧氏空间上的函数，频率也由离散变为连续。可积性首先保证变换有定义，Riemann--Lebesgue 引理给出连续性与无穷远衰减。卷积将空间中的平移平均转化为频率乘法，Young 不等式给出相应的范数估计。通过逼近恒等与 Gauss 核，反演公式从合法的绝对收敛积分中逐步建立；最后区分普通反演、Lebesgue 点处的正则化反演和仅依赖零变换的唯一性。
\end{chapterabstract}
\begin{learninggoals}
\begin{itemize}
\item\label{b2:item:03-goal-normalization}按统一的 \(2\pi\) 约定计算平移、调制、线性变换及微分公式，说明每个公式的可积条件。
\item\label{b2:item:03-goal-decay}说明空间矩控制频率正则性、空间导数控制高频衰减，区分固定函数的衰减与整个函数族的统一衰减。
\item\label{b2:item:03-goal-convolution}证明卷积的存在、等价类良定义及指数估计，核对换序所需的绝对可积性。
\item\label{b2:item:03-goal-inversion}用 Gauss 正则化建立反演和唯一性，说明函数代表的选择以及连续点、Lebesgue 点条件在反演中的作用。
\end{itemize}
\end{learninggoals}

周期越长，相邻振动频率之间的距离越小。把一个紧支集函数限制在长度为 \(L\) 的区间内，再作周期延拓，它的第 \(k\) 个系数是原函数在频率 \(k/L\) 处的积分除以 \(L\)。当 \(L\) 增大时，频率网格逐渐变密，离散叠加提示了连续频率积分。另一方面，若一个过程把不同位置的输入按同一种规则平均，所得输出是卷积；将卷积变成频率上的乘法，就能把这个过程逐频率分析。本章要证明这些对应关系确实保留原函数，而非只作形式替换。

\ProseParagraphBreak

本章在 \(\mathbb R^d\) 上使用 Lebesgue 测度。第一卷的连续紧支集稠密性、有限指数的平移连续性及 Lebesgue 点理论，在使用时给出相应编号；多元换元和积分号下求导沿用\chapref{b2:ch:01}。卷积的基本估计在本章按所需的一般指数完整推导。

\ProseParagraphBreak

先读\secref{b2:sec:0301}的归一化和运算规则，再把\secref{b2:sec:0302}的衰减与\secref{b2:sec:0303}的乘法对应连起来。\secref{b2:sec:0304}所构造的核既是逼近工具，也是\secref{b2:sec:0305}反演证明中的收敛控制。Gauss 变换的证明仅用常微分方程与实积分，因此不会借用尚未建立的复积分公式。

\ProseParagraphBreak

Sheldon Axler 的《Measure, Integration \& Real Analysis》\cite[Section 11C]{Axler2020Measure}采用一维 Poisson 核证明反演公式，可与本章的 Gauss 核方法比较。Elias M. Stein 与 Rami Shakarchi 的《Fourier Analysis: An Introduction》\cite{Stein2003Fourier}从周期级数讲到欧氏空间上的变换，便于比较离散与连续两种频率表示。

% !TeX root = ../../Book2.tex
\section{Fourier 变换}
\label{b2:sec:0301}

周期函数的频率由整数标记。若函数定义在整个欧氏空间上而不再周期重复，频率也变为连续变量。离散系数的求和于是让位于积分，但可积性、归一化与极限交换仍须逐项检查。

\subsection{定义与规范}
\label{b2:subsec:030101}

\begin{definition}\label{b2:def:03-fourier-transform}
设 \(d\) 为正整数，\(f\in L^1(\mathbb R^d;\mathbb C)\)。函数
\begin{equation}\label{b2:eq:03-fourier-definition}
 \widehat f(\xi)\coloneqq\mathcal Ff(\xi)
 =\int_{\mathbb R^d}f(x)\cdot e^{-2\pi\mathrm{i}x\cdot\xi}\,\mathrm{d}x
\end{equation}
称为 \(f\) 的\term[fourier-bianhuan]{Fourier 变换}[Fourier transform]。对 \(g\in L^1(\mathbb R^d)\)，另记
\[
 \check g(x)\coloneqq\int_{\mathbb R^d}g(\xi)\cdot e^{2\pi\mathrm{i}x\cdot\xi}\,\mathrm{d}\xi.
\]
此时只定义了正号积分；它何时反演 \(\mathcal F\)，将在\secref{b2:sec:0305}证明。
\end{definition}
\symidx{\(\mathcal Ff=\widehat f\)：欧氏空间 Fourier 变换}
\symidx{\(\check g\)：正号 Fourier 积分}

积分对每个 \(\xi\) 绝对收敛，且 \(|\widehat f(\xi)|\leq\|f\|_1\)，所以变换不依赖 \(f\) 的代表。这里 \(x\cdot\xi\) 是欧氏内积，虚数单位写作 \(\mathrm{i}\)。采用这个 \(2\pi\) 约定后，后面的反演公式没有额外常数，平方范数也将保持不变。若使用 \(e^{-\mathrm{i}x\cdot\xi}\) 的其他教材，需要同时转换频率尺度，不能只删去指数中的 \(2\pi\)。

\ProseParagraphBreak

在一维，取 \(a>0\)，区间示性函数给出一个可直接计算的模型：
\[
 \widehat{\mathbf1_{[-a,a]}}(\xi)=\dfrac{\sin(2\pi a\xi)}{\pi\xi}\quad(\xi\ne0),
 \quad \widehat{\mathbf1_{[-a,a]}}(0)=2a.
\]
原函数紧支集，并不使其变换紧支集；跳跃还使变换仅按 \(1/|\xi|\) 的尺度衰减。相反，对 \(a>0\)，分成正、负半轴积分可得
\[
 \mathcal F(e^{-a|x|})(\xi)=\dfrac{2a}{a^2+4\pi^2\xi^2}.
\]
这些公式同时提供了后面检验常数和反演条件的例子。

上面的连续频率可以用有限区间的级数作具体比较。在一维，设 \(f\) 连续且支撑在 \([-A,A]\)，并取 \(L>2A\)。将 \(f\) 在 \([-L/2,L/2]\) 上周期延拓，长度 \(L\) 的归一化 Fourier 系数为
\[
 c_L(k)=\frac1L\int_{-L/2}^{L/2}f(x)\cdot e^{-2\pi\mathrm i kx/L}\,\mathrm dx
       =\frac1L\widehat f(k/L).
\]
因此保留 \(|k|\leq RL\) 的有限频率和是
\[
 \frac1L\sum_{|k|\leq\lfloor RL\rfloor}
 \widehat f(k/L)\cdot e^{2\pi\mathrm i xk/L}.
\]
对固定 \(R>0\) 与 \(x\)，变换定义及控制收敛给出被求和函数在 \([-R,R]\) 连续；这一个和在 \(L\to\infty\) 时是频率步长 \(1/L\) 的 Riemann 和，趋于 \(\int_{-R}^R\widehat f(\xi)e^{2\pi\mathrm i x\xi}\,\mathrm d\xi\)。有限频带的极限只解释了积分为何自然出现；随后再令 \(R\to\infty\) 需要单独的收敛依据，不能从这个计算直接推出反演。归一化中的 \(1/L\) 正好成为频率积分元，所以采用 \(2\pi\) 后不再出现额外常数。

\subsection{平移}
\label{b2:subsec:030102}

对 \(f\in L^1(\mathbb R^d;\mathbb C)\)、\(h\in\mathbb R^d\)，沿用第一卷的平移约定 \(\tau_hf(x)\coloneqq f(x-h)\)。变量替换 \(y=x-h\) 给出
\[
 \widehat{\tau_hf}(\xi)=e^{-2\pi\mathrm{i}h\cdot\xi}\cdot\widehat f(\xi).
\]
平移保持 \(L^1\) 范数，故两侧都有定义。空间中的位置变化只改变频率中的相位，不改变 \(|\widehat f|\)。因而仅保留变换的模，通常不能恢复函数所在的位置。
\symidx{\(\tau_hf\)：平移，\(\tau_hf(x)=f(x-h)\)}

\subsection{调制}
\label{b2:subsec:030103}

对 \(\eta\in\mathbb R^d\)、\(f\in L^1(\mathbb R^d;\mathbb C)\)，算子
\[
 M_\eta f(x)\coloneqq e^{2\pi\mathrm{i}\eta\cdot x}\cdot f(x)
\]
称为\term[tiaozhi]{调制}[modulation]。合并指数便得
\[
 \widehat{M_\eta f}(\xi)=\widehat f(\xi-\eta).
\]
因此空间中的调制将频率整体平移。两种操作的次序要保留：\(\tau_hM_\eta=e^{-2\pi\mathrm{i}\eta\cdot h}M_\eta\tau_h\)，一般并不交换。
\symidx{\(M_\eta\)：调制算子}

\subsection{伸缩}
\label{b2:subsec:030104}

\begin{proposition}\label{b2:prop:03-linear-change}
设 \(d\) 为正整数，\(f\in L^1(\mathbb R^d;\mathbb C)\)，\(A\) 是可逆实 \(d\times d\) 矩阵，\(g(x)=f(Ax)\)，则
\[
 \widehat g(\xi)=|\det A|^{-1}\cdot\widehat f(A^{-\mathsf T}\xi).
\]
特别地，对 \(a>0\)，若 \(f_a(x)\coloneqq a^{-d}\cdot f(x/a)\)，则 \(\widehat{f_a}(\xi)=\widehat f(a\xi)\)。
\end{proposition}
\begin{proof}
在定义积分中令 \(y=Ax\)。体积元为 \(|\det A|^{-1}\mathrm{d}y\)，而 \(x\cdot\xi=y\cdot A^{-\mathsf T}\xi\)。绝对可积性由同一换元公式保证；取 \(A=a^{-1}I\) 再乘以归一化因子，即得后一式。
\end{proof}

一个在空间中压缩到宽度 \(a\) 的单位质量函数，其变换扩展到约 \(a^{-1}\) 的频率范围。这一换元公式精确记录了空间伸缩与频率伸缩的相反方向。


\begin{proposition}\label{b2:prop:03-symmetries}
设 \(d\) 为正整数，\(f\in L^1(\mathbb R^d;\mathbb C)\)，则 \(\widehat{\overline f}(\xi)=\overline{\widehat f(-\xi)}\)。若 \(f\) 实值且 \(f(-x)=f(x)\) 几乎处处，则其变换为实值偶函数。若对每个实正交矩阵 \(Q\) 都有 \(f(Qx)=f(x)\) 几乎处处，则 \(\widehat f(Q\xi)=\widehat f(\xi)\) 对每个 \(Q\) 及所有 \(\xi\in\mathbb R^d\) 成立。
\end{proposition}
\begin{proof}
第一式由对定义积分取共轭得到。偶性由反射换元得到，再与第一式组合即知实偶函数的变换为实值。对正交矩阵 \(Q\)，线性换元公式中的绝对行列式为 \(1\)、逆转置为 \(Q\)。若 \(f(Qx)=f(x)\)，便有 \(\widehat f(Q\xi)=\widehat f(\xi)\)。
\end{proof}

例如矩形 \(Q=\prod_j[-a_j,a_j]\) 的变换为各坐标区间变换的乘积，频率中的零点形成相应的超平面。球的示性函数则有旋转不变的变换；但旋转不变性并不提供其径向表达式，还须计算角积分。这里先由对称性确定变换的结构，具体特殊函数留待附录处理。

\ProseParagraphBreak

反射与共轭也是校验计算的工具。若一个实值偶函数被算出了非实的变换，或径向函数被算出了依赖坐标方向的变换，应先回查相位和换元。另一方面，实值函数不一定有实值变换：只有共轭对称是普遍结论，偶性是额外条件。

\subsection{微分与乘法}
\label{b2:subsec:030105}

\begin{proposition}\label{b2:prop:03-moments}
设 \(d\) 为正整数，\(m\) 为非负整数。若 \(f\in L^1(\mathbb R^d;\mathbb C)\) 且 \(|x|^m f\in L^1(\mathbb R^d)\)，则 \(\widehat f\in C^m(\mathbb R^d;\mathbb C)\)，并对 \(|\alpha|\leq m\) 有
\[
 \partial_\xi^\alpha\widehat f(\xi)
 =(-2\pi\mathrm{i})^{|\alpha|}\mathcal F(x^\alpha f)(\xi).
\]
单独考虑一个坐标 \(1\leq j\leq d\) 时，只要 \(f,x_jf\in L^1(\mathbb R^d;\mathbb C)\)，就有连续的偏导数 \(\partial_{\xi_j}\widehat f=-2\pi\mathrm{i}\mathcal F(x_jf)\)，无须假设其他坐标的一阶矩可积。
\end{proposition}
\begin{proof}
一阶差商中指数之差除以步长的模不超过 \(2\pi|x_j|\)。因此 \(|x_jf|\) 提供共同的可积控制，可以把差商极限移入积分。导数的连续性由相同控制下的控制收敛得到，这就证明了仅需一个坐标矩的结论。若具有 \(m\) 阶绝对矩，迭代时每次增加一个坐标因子；所有不超过 \(m\) 阶的单项式都由常数乘 \(1+|x|^m\) 控制，故每一步都合法。
\end{proof}

\begin{proposition}\label{b2:prop:03-derivative}
设 \(d\) 为正整数，\(1\leq j\leq d\)。若 \(f\in C^1(\mathbb R^d;\mathbb C)\)，且 \(f,\partial_jf\in L^1(\mathbb R^d)\)，则
\[
 \mathcal F(\partial_jf)(\xi)=2\pi\mathrm{i}\xi_j\widehat f(\xi).
\]
若 \(f\in L^1(\mathbb R^d;\mathbb C)\) 且其第 \(j\) 个弱导数属于 \(L^1(\mathbb R^d;\mathbb C)\)，公式也成立。
\end{proposition}
\begin{proof}
取光滑紧支集函数 \(\chi\)，在单位球上等于 \(1\)，并令 \(\chi_R(x)=\chi(x/R)\)。先对紧支集函数 \(\chi_Rf\) 分部积分，得到
\[
 \mathcal F(\chi_R\partial_jf)(\xi)
 +\mathcal F(f\partial_j\chi_R)(\xi)
 =2\pi\mathrm{i}\xi_j\mathcal F(\chi_Rf)(\xi).
\]
第一项及右侧用控制收敛取极限；中间误差的绝对值至多为 \(R^{-1}\|\partial_j\chi\|_\infty\cdot\|f\|_1\)，故趋零。弱导数情形直接把 \(\chi_Re^{-2\pi\mathrm{i}x\cdot\xi}\) 代入弱导数的测试恒等式，得到同一公式和估计。
\end{proof}

于是空间中的可积导数产生频率衰减，空间中的可积矩产生频率光滑性。两类假设各控制一个不同方向，不能把“变换连续”误认为原函数已经可微。

% !TeX root = ../../Book2.tex
\section{Riemann--Lebesgue 引理}
\label{b2:sec:0302}

周期情形中 Fourier 系数趋零；在欧氏空间中，对应的问题是频率沿任意方向走向无穷时，变换是否趋零。单纯观察某个积分公式不足以说明所有方向，必须用可积逼近将有限计算推广。

\subsection{基本衰减}
\label{b2:subsec:030201}

\begin{theorem}[Riemann--Lebesgue]\label{b2:thm:03-riemann-lebesgue}
设 \(d\) 为正整数。Fourier 变换将复值空间 \(L^1(\mathbb R^d;\mathbb C)\) 连续地映入 \(C_0(\mathbb R^d;\mathbb C)\)，且对每个 \(f\in L^1(\mathbb R^d;\mathbb C)\) 有
\[
 \|\widehat f\|_\infty\leq\|f\|_1.
\]
这里 \(C_0\) 表示在无穷远处趋零的连续函数空间。
\end{theorem}
\begin{proof}
对有界矩形 \(Q=\prod_{j=1}^d[a_j,b_j]\)，Fubini 把 \(\widehat{\mathbf1_Q}\) 分解成一维区间变换的乘积。每个因子连续且有界；当 \(|\xi|\to\infty\) 时，至少一个坐标的绝对值不小于 \(|\xi|/\sqrt d\)，相应因子趋零，所以整个乘积趋零。

有限个矩形示性函数的线性组合在 \(L^1\) 中稠密。具体可先用第一卷命题 10.2.4 的连续紧支集稠密性逼近 \(f\)，再把包含其支集的一个大矩形细分成网格；一致连续性使每个网格内用常数代替函数的总积分误差趋零。记所得阶梯函数为 \(q\)，则
\[
 \|\widehat f-\widehat q\|_\infty\leq\|f-q\|_1.
\]
因此 \(\widehat f\) 是 \(C_0\) 函数的一致极限，仍属于 \(C_0\)。范数界直接来自积分三角不等式。
\end{proof}

该结论沿用\subsecref{b2:subsec:020202}的 Riemann--Lebesgue 名称。对 \(f_n\to f\) 于 \(L^1\)，变换在全部频率上一致收敛，而不只是每个固定频率分别收敛。

\subsection{一致连续性与频率正则性}
\label{b2:subsec:030202}

前一小节已经证明了变换的连续性。现在进一步考察：频率改变一个小增量时，变换的变化能否在所有频率中心处同时变小？设 \(f\in L^1(\mathbb R^d;\mathbb C)\)。对任意 \(\xi,h\in\mathbb R^d\)，定义式给出
\[
 |\widehat f(\xi+h)-\widehat f(\xi)|
 \leq\int_{\mathbb R^d}|f(x)|\cdot|e^{-2\pi\mathrm{i}x\cdot h}-1|\,\mathrm{d}x.
\]
右侧不依赖 \(\xi\)。其中被积函数随 \(h\to0\) 逐点趋零，并由可积函数 \(2|f|\) 控制，故整个积分趋零。这就得到 \(\widehat f\) 的一致连续性，而不仅是在每个固定频率处的连续性。

若还满足一阶绝对矩条件 \(\int_{\mathbb R^d}|x|\cdot|f(x)|\,\mathrm{d}x<\infty\)，利用 \(|e^{\mathrm i t}-1|\leq|t|\) 及 \(|x\cdot h|\leq|x||h|\)，可将同一估计写成
\[
 |\widehat f(\xi+h)-\widehat f(\xi)|
 \leq 2\pi|h|\int_{\mathbb R^d}|x|\cdot|f(x)|\,\mathrm{d}x.
\]
这给出了明确的 Lipschitz 常数。更一般地，\cref{b2:prop:03-moments}说明：若 \(m\) 是非负整数，且 \(f,|x|^m f\in L^1(\mathbb R^d)\)，则 \(\widehat f\in C^m\)，对每个 \(|\alpha|\leq m\) 有
\[
 \|\partial_\xi^\alpha\widehat f\|_\infty
 \leq(2\pi)^{|\alpha|}\|x^\alpha f\|_1.
\]
空间矩因而控制变换在频率变量上的光滑程度及导数大小。这些估计描述的是频率改变时函数值如何变化；是否能进一步控制整个函数族在高频处的大小，还须另作判断。

\ProseParagraphBreak

注意 \(L^1\) 假设并未强迫原函数有界或连续。例如 \(f(x)=|x|^{-a}\mathbf1_{\{|x|<1\}}\) 在 \(0<a<d\) 时可积，其变换仍连续且趋零。反过来，紧支集函数的变换也未必可积；区间示性函数的 \(1/|\xi|\) 尾部正好落在一维不可积的边界。

\subsection{无穷远行为}
\label{b2:subsec:030203}

\begin{proposition}\label{b2:prop:03-no-uniform-decay}
设 \(d\) 为正整数。在 \(L^1(\mathbb R^d;\mathbb C)\) 的单位球上，不存在只依赖 \(|\xi|\)、并在无穷远处趋零的共同上界 \(|\widehat f(\xi)|\leq r(|\xi|)\)。即使另外要求函数族的每一阶空间绝对矩各有一个一致的上界，也不能保证统一的高频衰减。
\end{proposition}
\begin{proof}
固定一个非负函数 \(g\in C_c^\infty(\mathbb R^d)\)，使 \(\int_{\mathbb R^d}g=1\)。对每个 \(\eta\in\mathbb R^d\)，令
\[
 f_\eta(x)\coloneqq e^{2\pi\mathrm i\eta\cdot x}g(x)=M_\eta g(x).
\]
因为复指数的模为 \(1\)，\(\|f_\eta\|_1=1\)；而且对任意非负整数 \(m\)，都有
\[
 \int_{\mathbb R^d}|x|^m|f_\eta(x)|\,\mathrm dx
 =\int_{\mathbb R^d}|x|^m g(x)\,\mathrm dx<\infty.
\]
右侧可依赖于 \(m\)，但不依赖于 \(\eta\)。因此这个族同时满足每一固定阶绝对矩的一致有界性。

另一方面，\subsecref{b2:subsec:030103}的调制公式给出
\[
 \widehat{f_\eta}(\xi)=\widehat g(\xi-\eta),
 \qquad \widehat{f_\eta}(\eta)=\widehat g(0)=1.
\]
令 \(|\eta|\to\infty\)，便与任何共同趋零上界矛盾。更明确地，对每个 \(A>0\)，先取 \(|\eta|\geq A\)，再在 \(\xi=\eta\) 处取值，结合范数界可得
\[
 \sup_{\eta\in\mathbb R^d}\ \sup_{|\xi|\geq A}|\widehat{f_\eta}(\xi)|=1.
\]
\end{proof}

每个固定的 \(f_\eta\) 仍然满足 Riemann--Lebesgue 引理。这里失败的是对全族统一选取高频阈值：调制不改变空间中的绝对值，却把频率中的峰移到了 \(\eta\) 附近。因此，空间矩的一致控制可以给出一致的频率正则性，却不能单独保证统一的高频衰减。

\ProseParagraphBreak

高频衰减的定量估计来自另一类假设，即空间导数的 \(L^1\) 控制。一阶公式已在\cref{b2:prop:03-derivative}证明。高阶情形若只知道最高阶纯弱导数可积，不能直接反复套用一阶公式，因为尚未假设中间各阶弱导数也可积。下面直接把全部导数移到一个光滑紧支集测试函数上，保留这个较弱的假设。

\begin{proposition}\label{b2:prop:03-higher-derivative-decay}
设 \(d,m\) 为正整数，\(f\in L^1(\mathbb R^d;\mathbb C)\)。假设对每个 \(j=1,\ldots,d\)，\(f\) 的分布意义下的纯坐标弱导数 \(g_j\coloneqq\partial_j^m f\) 属于 \(L^1(\mathbb R^d)\)。则对每个 \(\xi\in\mathbb R^d\) 有
\[
 \widehat{g_j}(\xi)=(2\pi\mathrm i\xi_j)^m\widehat f(\xi),
 \qquad 1\leq j\leq d,
\]
并且
\[
 |\widehat f(\xi)|
 \leq\frac{d^{m/2}}{(2\pi)^m}|\xi|^{-m}
       \sum_{j=1}^d\|\partial_j^m f\|_1
 \quad(\xi\ne0).
\]
\end{proposition}
\begin{proof}
最高阶弱导数的假设是指：对每个 \(\varphi\in C_c^\infty(\mathbb R^d)\)，有
\[
 \int_{\mathbb R^d}g_j(x)\varphi(x)\,\mathrm dx
 =(-1)^m\int_{\mathbb R^d}f(x)\partial_j^m\varphi(x)\,\mathrm dx.
\]
这里采用不加共轭的分布配对；复值测试函数由实部、虚部分别测试即可。指数函数本身不紧支集，不能直接代入这个恒等式，所以先作空间截断。

取 \(\chi\in C_c^\infty(\mathbb R^d)\)，满足 \(0\leq\chi\leq1\)、在单位球上等于 \(1\)，并令 \(\chi_R(x)\coloneqq\chi(x/R)\)，其中 \(R\geq1\)。固定 \(\xi\in\mathbb R^d\) 及一个坐标 \(j\)，把
\[
 \varphi_R(x)\coloneqq\chi_R(x)e^{-2\pi\mathrm i x\cdot\xi}
\]
代入弱导数恒等式。由 Leibniz 公式，
\[
 \partial_j^m\varphi_R(x)
 =\sum_{k=0}^m\binom{m}{k}R^{-k}
   (\partial_j^k\chi)(x/R)
   (-2\pi\mathrm i\xi_j)^{m-k}e^{-2\pi\mathrm i x\cdot\xi}.
\]
其中 \(k=0\) 项给出所需乘子；把其他项分别记为
\[
 \begin{aligned}
 E_{R,k}\coloneqq{}&(-1)^m\binom{m}{k}
 (-2\pi\mathrm i\xi_j)^{m-k}R^{-k}\\
 &{}\times\int_{\mathbb R^d}f(x)(\partial_j^k\chi)(x/R)
                   e^{-2\pi\mathrm i x\cdot\xi}\,\mathrm dx,
 \qquad 1\leq k\leq m.
 \end{aligned}
\]
于是
\[
 \begin{aligned}
 &\int_{\mathbb R^d}g_j(x)\chi_R(x)e^{-2\pi\mathrm i x\cdot\xi}\,\mathrm dx\\
 &\qquad=(2\pi\mathrm i\xi_j)^m
   \int_{\mathbb R^d}f(x)\chi_R(x)e^{-2\pi\mathrm i x\cdot\xi}\,\mathrm dx
   +\sum_{k=1}^m E_{R,k}.
 \end{aligned}
\]
每个截断误差都有单独的估计
\[
 |E_{R,k}|
 \leq\binom{m}{k}(2\pi|\xi_j|)^{m-k}R^{-k}
       \|\partial_j^k\chi\|_\infty\|f\|_1
 \longrightarrow0\qquad(R\to\infty).
\]
当 \(k=m\) 时，频率因子为 \(1\)。这里先固定 \(\xi\)，故有限个误差项都趋零；所有截断误差都由 \(\|f\|_1\) 控制。又因为 \(\chi_R\to1\) 逐点成立、\(|\chi_R|\leq1\)，左右两个主项分别由 \(|g_j|\)、\(|f|\) 控制。控制收敛便给出
\[
 \widehat{g_j}(\xi)=(2\pi\mathrm i\xi_j)^m\widehat f(\xi).
\]

最后，对 \(\xi\ne0\)，至少有一个坐标 \(j\) 满足 \(|\xi_j|\geq|\xi|/\sqrt d\)。在这个坐标上使用刚证的公式及 \(|\widehat{g_j}(\xi)|\leq\|g_j\|_1\)，得到
\[
 |\widehat f(\xi)|
 \leq(2\pi|\xi_j|)^{-m}\|g_j\|_1
 \leq\frac{d^{m/2}}{(2\pi)^m}|\xi|^{-m}\|g_j\|_1.
\]
再以所有 \(\|g_j\|_1\) 的和控制这一项，即得所述估计。
\end{proof}

因此，当一个函数族的这些最高阶纯弱导数范数之和有共同上界时，就得到全族的 \(|\xi|^{-m}\) 高频衰减估计。常数由维数、阶数和指定的导数范数控制；空间矩则继续承担上一小节所说明的频率正则性作用。

% !TeX root = ../../Book2.tex
\section{卷积}
\label{b2:sec:0303}

频率上的乘法对应空间中的卷积。要使用这个对应关系，先须证明卷积本身存在，并明确它落在哪个函数空间。相应的可积性估计还将保证 Fourier 计算中的积分换序。

\subsection{卷积代数}
\label{b2:subsec:030301}

\begin{definition}\label{b2:def:03-convolution}
设 \(d\) 为正整数。对可测函数 \(f,g:\mathbb R^d\to\mathbb C\) 和 \(x\in\mathbb R^d\)，若积分绝对收敛，定义它们在 \(x\) 处的\term[juanji]{卷积}[convolution]为
\[
 (f*g)(x)\coloneqq\int_{\mathbb R^d}f(x-y)\cdot g(y)\,\mathrm{d}y.
\]
\end{definition}
\symidx{\(f*g\)：欧氏空间卷积}

\begin{proposition}\label{b2:prop:03-l1-algebra}
设 \(d\) 为正整数。若 \(f,g\in L^1(\mathbb R^d;\mathbb C)\)，则卷积几乎处处绝对收敛，属于 \(L^1(\mathbb R^d;\mathbb C)\)，且
\[
 \|f*g\|_1\leq\|f\|_1\cdot\|g\|_1.
\]
卷积在 \(L^1\) 等价类上良定义，满足交换律、结合律和双线性。
\end{proposition}
\begin{proof}
可测函数可选取 Borel 代表；函数 \((x,y)\mapsto f(x-y)g(y)\) 因而联合可测。对绝对值用 Tonelli 和平移不变性，
\[
 \int\!\int|f(x-y)|\cdot|g(y)|\,\mathrm{d}y\,\mathrm{d}x
 =\|f\|_1\cdot\|g\|_1<\infty.
\]
这既保证几乎处处绝对收敛，又给出范数界。改变任一代表所产生的差，按同一估计为零，故等价类良定义。令 \(z=x-y\) 得交换律。三个函数的绝对值重积分等于三个范数之积，所以可用 Fubini 重排；两种卷积括号都等于
\[
 \int\!\int f(x-y-z)\cdot g(y)h(z)\,\mathrm{d}y\,\mathrm{d}z
\]
几乎处处。线性来自积分的线性。
\end{proof}

结合 \(L^1\) 的完备性，这使 \(L^1\) 成为卷积运算下的 Banach 代数。它没有单位元：若 \(e*f=f\) 对所有 \(f\in L^1\) 成立，后面将由变换证明 \(\widehat e\) 必须恒等于 \(1\)，与 Riemann--Lebesgue 引理矛盾；在分布空间中才会出现相应的 Dirac 单位。

\subsection{Young 不等式}
\label{b2:subsec:030302}

\begin{theorem}[Young，卷积形式]\label{b2:thm:03-young}
设 \(d\) 为正整数，\(1\leq p,q,r\leq\infty\) 且 \(1/p+1/q=1+1/r\)，其中 \(1/\infty=0\)。若 \(f\in L^p(\mathbb R^d;\mathbb C)\)、\(g\in L^q(\mathbb R^d;\mathbb C)\)，则卷积几乎处处绝对收敛，并且
\[
 \|f*g\|_r\leq\|f\|_p\cdot\|g\|_q.
\]
\end{theorem}
\begin{proof}
固定 \(f,g\) 的有限值 Borel 可测代表，使被积函数关于 \((x,y)\) 联合可测。对实部和虚部的正负部分分别使用 Tonelli，可知参数积分在绝对收敛处可测；在其余点赋零仍为可测函数。对每个固定的 \(x\)，改变代表只会在 \(y\) 的零测集上改变被积函数，因此不影响存在的卷积积分。若 \(r=\infty\)，指数 \(p,q\) 共轭，对每个 \(x\) 使用 Hölder 即得积分的绝对收敛与一致上界。以下设 \(r<\infty\)；若某个范数为零，卷积处处为零，故只须处理两个范数均为正的情形。

指数关系给出
\[
 \frac1p-\frac1r=1-\frac1q\geq0,
 \qquad
 \frac1q-\frac1r=1-\frac1p\geq0.
\]
因此 \(p,q\leq r<\infty\)。所要使用的三个 Hölder 指数的倒数是
\[
 \frac1r,\qquad \frac1p-\frac1r,\qquad \frac1q-\frac1r;
 \qquad
 \frac1r+\left(\frac1p-\frac1r\right)
             +\left(\frac1q-\frac1r\right)=1.
\]
它们非负，因而确实给出合法的 Hölder 指数。下面先确定哪些 \(x\) 可以使用这一步。令
\[
 I(x)=\int_{\mathbb R^d}|f(x-y)|^p\cdot|g(y)|^q\,\mathrm dy.
\]
被积函数联合可测且非负，Tonelli 定理允许先对 \(x\) 积分；平移不变性给出
\[
 \int_{\mathbb R^d}I(x)\,\mathrm dx
 =\int_{\mathbb R^d}|g(y)|^q
    \left(\int_{\mathbb R^d}|f(x-y)|^p\,\mathrm dx\right)\mathrm dy
 =\|f\|_p^p\cdot\|g\|_q^q<\infty.
\]
所以 \(I(x)<\infty\) 对几乎处处的 \(x\) 成立。对这样的 \(x\)，把绝对值乘积拆成
\[
 |f(x-y)|\cdot|g(y)|
 =\bigl(|f(x-y)|^p\cdot|g(y)|^q\bigr)^{1/r}
    |f(x-y)|^{1-p/r}\cdot|g(y)|^{1-q/r}.
\]
若第二或第三个倒指数为零，则相应的 \(1-p/r\) 或 \(1-q/r\) 也为零；此时直接省去该幂因子，等价于在 Hölder 中使用常数 \(1\) 这个 \(L^\infty\) 因子，其范数为 \(1\)。当 \(p<r\) 时，第二因子的 Hölder 指数为 \(pr/(r-p)\)，其对应范数就是 \(\|f\|_p^{1-p/r}\)；第三因子同理，首因子的 \(L^r\) 范数则为 \(I(x)^{1/r}\)。于是
\[
 \int_{\mathbb R^d}|f(x-y)|\cdot|g(y)|\,\mathrm dy
 \leq I(x)^{1/r}\cdot\|f\|_p^{1-p/r}\cdot\|g\|_q^{1-q/r}<\infty.
\]
例如 \(p=1,q=r\) 时第三因子省去；\(q=1,p=r\) 时第二因子省去。若 \(r=1\)，则 \(p=q=1\)，两个因子都省去，这一步就是对 \(I(x)\) 的恒等估计。

至此原复值卷积在上述满测集上绝对收敛。积分三角不等式再给出
\[
 |(f*g)(x)|^r
 \leq I(x)\cdot\|f\|_p^{r-p}\cdot\|g\|_q^{r-q}.
\]
在剩余零测集上赋值为零，对 \(x\) 积分并代入已算出的 \(\int I\)，得到
\[
 \|f*g\|_r^r\leq\|f\|_p^r\cdot\|g\|_q^r.
\]
取 \(r\) 次方根即得所求估计。
\end{proof}

这称为\term[young-juanji-budengshi]{Young 卷积不等式}[Young's convolution inequality]，与标量乘积的 Young 不等式区别在于对象和指数关系。常用情形是 \(L^1*L^p\subset L^p\)，以及 \(L^p*L^{p'}\subset L^\infty\)。后一种情形还可利用平移连续性推出连续代表。

\subsection{平移与卷积}
\label{b2:subsec:030303}

第一卷引理 13.4.5 已证明：若 \(1\leq p<\infty\)，则 \(\|\tau_hf-f\|_p\to0\)。该结论首先对连续紧支集函数成立，再由平移等距性和稠密性推广。无穷端点不在其中；区间示性函数的小平移在 \(L^\infty\) 中始终保持误差 \(1\)。

\begin{proposition}\label{b2:prop:03-convolution-continuity}
设 \(d\) 为正整数，\(1\leq p<\infty\)，\(p'\) 为共轭指数，即 \(1/p+1/p'=1\)，其中 \(p=1\) 时 \(p'=\infty\)。若 \(f\in L^p(\mathbb R^d;\mathbb C)\)、\(g\in L^{p'}(\mathbb R^d;\mathbb C)\)，则对任意固定的有限值可测代表 \(f_0,g_0\)，积分
\[
 F(x)=\int_{\mathbb R^d}f_0(x-y)\cdot g_0(y)\,\mathrm dy
 \qquad(x\in\mathbb R^d)
\]
处处绝对收敛。它不依赖代表的选择，本身就是卷积等价类的有界一致连续代表，并且
\[
 \sup_{x\in\mathbb R^d}|F(x)|\leq\|f\|_p\cdot\|g\|_{p'}.
\]
\end{proposition}
\begin{proof}
对每个 \(x\)，反射和平移保持 \(L^p\) 范数，故 Hölder 给出
\[
 \int_{\mathbb R^d}|f_0(x-y)|\cdot|g_0(y)|\,\mathrm dy
 \leq\|f\|_p\cdot\|g\|_{p'}.
\]
这同时证明处处绝对收敛和普通上确界的估计。若换用另一对代表，对每个固定的 \(x\)，被积函数只在 \(g_0\) 的零测例外集与 \(f_0\) 的零测例外集经反射平移所得的集合上改变，故 \(F(x)\) 不变。

沿用本章 \(\tau_hf(u)=f(u-h)\) 的约定，有
\[
 \begin{aligned}
 |F(x+h)-F(x)|
 &\leq\int_{\mathbb R^d}
 |f_0(x+h-y)-f_0(x-y)|\cdot|g_0(y)|\,\mathrm dy\\
 &\leq\|\tau_{-h}f-f\|_p\cdot\|g\|_{p'}.
 \end{aligned}
\]
右侧与 \(x\) 无关，且由第一卷引理 13.4.5 的有限指数平移连续性趋于零。因此
\[
 \sup_{x\in\mathbb R^d}|F(x+h)-F(x)|\longrightarrow0
 \qquad(h\to0),
\]
逐点积分定义的 \(F\) 本身便有界一致连续。它与 Young 不等式给出的几乎处处卷积相同，因而确为所需代表。整个论证包括 \(p=1,p'=\infty\)；此时使用的是 \(L^1\) 因子的平移连续性。
\end{proof}

连续性在这里由两个因子的整体平移行为取得，并非因为原函数本身连续。另一方面，\(L^1*L^1\) 一般只有可积性，不能仅凭卷积符号就称所得函数光滑。若某一因子为光滑紧支集函数，才可将导数移到该因子，并由可积控制证明结果光滑。


\begin{corollary}\label{b2:cor:03-convolution-c0}
设 \(d\) 为正整数。若 \(1<p<\infty\)、\(p'=p/(p-1)\)、\(f\in L^p(\mathbb R^d;\mathbb C)\)、\(g\in L^{p'}(\mathbb R^d;\mathbb C)\)，则卷积的逐点积分代表属于 \(C_0(\mathbb R^d;\mathbb C)\)。
\end{corollary}
\begin{proof}
取连续紧支集函数 \(f_n\to f\) 于 \(L^p\)、\(g_n\to g\) 于 \(L^{p'}\)。可取两列范数有界。由 Hölder 型卷积界，
\[
 \|f*g-f_n*g_n\|_\infty
 \leq\|f-f_n\|_p\cdot\|g\|_{p'}
 +\|f_n\|_p\cdot\|g-g_n\|_{p'}\longrightarrow0.
\]
每个 \(f_n*g_n\) 连续且支撑在两个紧支集的和中，所以属于 \(C_0\)。一致极限仍属于 \(C_0\)。
\end{proof}

两个指数都有限，使两边都能用紧支集函数逼近。端点 \(L^1*L^\infty\) 不保证在无穷远消失：取 \(g=1\)，则卷积恒为 \(\int f\)。这与已经证明的有界一致连续性不矛盾。卷积的“平滑作用”应拆成可积性、连续性、无穷远行为和可微性分别判断，不能把这些性质合成一个没有条件的结论。

卷积不是无条件的连续化操作。取一维非负函数 \(f(x)=x^{-a}\mathbf1_{(0,1)}(x)\)，其中 \(1/2<a<1\)。两因子都可积；当 \(0<x<1\) 时，令 \(y=xt\)，得到
\[
 (f*f)(x)=\int_0^x(x-y)^{-a}y^{-a}\,\mathrm dy
 =x^{1-2a}\int_0^1(1-t)^{-a}t^{-a}\,\mathrm dt.
\]
最后的常数有限且正，因为两个端点的幂指数都大于 \(-1\)。因此卷积在 \(x\downarrow0\) 时无界，却仍属于 \(L^1\)：新奇点的指数 \(1-2a\) 大于 \(-1\)，而一般范数结论已控制全部积分。这个例子同时保留了 \(L^1*L^1\) 的结论并排除了有界性的误读。

\ProseParagraphBreak

若一因子有足够的光滑性，就可以取得真正的光滑化。下面说明求导与估计如何同时发生。

\begin{proposition}\label{b2:prop:03-smooth-convolution}
设 \(d\) 为正整数，\(1\leq p\leq\infty\)，\(p'\) 为共轭指数，\(1/p+1/p'=1\)，其中 \(1/\infty=0\)；\(f\in L^p(\mathbb R^d;\mathbb C)\)，\(\rho\in C_c^\infty(\mathbb R^d;\mathbb C)\)。则 \(\rho*f\) 具有 \(C^\infty(\mathbb R^d;\mathbb C)\) 代表，且每个多重指标 \(\alpha\) 满足
\[
 \partial^\alpha(\rho*f)=(\partial^\alpha\rho)*f,\quad
 \|\partial^\alpha(\rho*f)\|_\infty
 \leq\|\partial^\alpha\rho\|_{p'}\cdot\|f\|_p.
\]
若 \(f\) 几乎处处支撑在紧集 \(K\)，则该代表支撑在 \(K+\operatorname{supp}\rho\)。
\end{proposition}
\begin{proof}
把卷积写成 \(\int\rho(x-y)\cdot f(y)\,\mathrm dy\)。Hölder 不等式使每个点的积分有限，并给出所写范数界。对 \(x_j\) 求差商时，核的差商在 \(L^{p'}\) 中趋于 \(\partial_j\rho\)，包括 \(p'=1,\infty\)：微积分基本定理把差商误差写为
\[
 \int_0^1\bigl(\partial_j\rho(\cdot+the_j)-\partial_j\rho\bigr)\,\mathrm dt.
\]
当 \(|h|\leq1\) 时，这些函数有共同紧支集；\(\partial_j\rho\) 一致连续，故误差在上确界范数以及每个有限指数范数中都趋零。Hölder 再使卷积差商在全部 \(x\) 上一致趋于 \((\partial_j\rho)*f\)。同样的平移估计说明这一导数连续。对每个 \(\partial^\alpha\rho\) 重复这个论证，得到任意阶求导公式。最后选取在 \(K\) 外为零的代表；若 \(x\notin K+\operatorname{supp}\rho\)，被积函数恒为零。由于这个和为紧集，支集结论成立。
\end{proof}

\subsection{Fourier 变换与卷积}
\label{b2:subsec:030304}

\begin{theorem}[卷积定理]\label{b2:thm:03-convolution-transform}
设 \(d\) 为正整数。若 \(f,g\in L^1(\mathbb R^d;\mathbb C)\)，则对每个 \(\xi\in\mathbb R^d\) 有
\[
 \widehat{f*g}(\xi)=\widehat f(\xi)\cdot\widehat g(\xi).
\]
\end{theorem}
\begin{proof}
定义积分中的绝对值重积分等于 \(\|f\|_1\cdot\|g\|_1\)，故可交换次序。令 \(z=x-y\)，指数分解为 \(e^{-2\pi\mathrm{i}z\cdot\xi}e^{-2\pi\mathrm{i}y\cdot\xi}\)，两个积分分别为 \(\widehat f\) 与 \(\widehat g\)。
\end{proof}

这个\term[juanji-dingli]{卷积定理}[convolution theorem]将叠加平移的操作对角化。现在也可补足单位元不存在的论证：若存在 \(e\in L^1(\mathbb R^d)\) 使对所有 \(f\in L^1(\mathbb R^d)\) 都有 \(e*f=f\)，固定频率 \(\xi\) 后取单位质量非负函数的调制，使 \(\widehat f(\xi)=1\)，便有 \(\widehat e(\xi)=1\)。这与 \(\widehat e\in C_0\) 不相容。下一节改用一族核逐渐接近单位作用。

% !TeX root = ../../Book2.tex
\section{逼近恒等}
\label{b2:sec:0304}

圆周上，Fejér 核把质量集中到原点，同时保持积分为 \(1\)。欧氏空间中的同一思想还要控制无穷远尾部。只要这两种控制成立，核可以有符号，也不必紧支集。

\subsection{逼近恒等的条件}
\label{b2:subsec:030401}

\begin{definition}\label{b2:def:03-approximate-identity}
设 \(d\) 为正整数，\((K_\varepsilon)_{\varepsilon>0}\subset L^1(\mathbb R^d;\mathbb C)\)。若
\begin{enumerate}
\item\label{b2:item:03-identity-mass}\(\int K_\varepsilon=1\)；
\item\label{b2:item:03-identity-bound}\(\sup_{\varepsilon>0}\|K_\varepsilon\|_1<\infty\)；
\item\label{b2:item:03-identity-tail}对每个 \(\delta>0\)，\(\int_{|x|\geq\delta}|K_\varepsilon(x)|\,\mathrm{d}x\to0\) 当 \(\varepsilon\downarrow0\)，
\end{enumerate}
则称它为一族\term[bijin-hengdeng]{逼近恒等}[approximate identity]。第一卷定义 13.4.6 的固定核缩放是其中一种情形。
\end{definition}

确实，若 \(K\in L^1\)、\(\int K=1\)，取 \(K_\varepsilon(x)=\varepsilon^{-d}\cdot K(x/\varepsilon)\)，前三个量分别由换元化为 \(1\)、\(\|K\|_1\) 和 \(\int_{|z|\geq\delta/\varepsilon}|K(z)|\,\mathrm{d}z\)。最后一项因可积性趋零。固定核允许正负抵消；一般核族则必须另写统一 \(L^1\) 界，不能由积分等于 \(1\) 推出。

\subsection{逼近恒等的收敛与误差估计}
\label{b2:subsec:030402}

\begin{theorem}\label{b2:thm:03-approximate-identity}
设 \(d\) 为正整数，\((K_\varepsilon)_{\varepsilon>0}\subset L^1(\mathbb R^d;\mathbb C)\) 是逼近恒等：\(\int K_\varepsilon=1\)、\(\sup_{\varepsilon>0}\|K_\varepsilon\|_1<\infty\)，且对每个 \(\delta>0\)，\(\int_{|x|\geq\delta}|K_\varepsilon(x)|\,\mathrm dx\to0\) 当 \(\varepsilon\downarrow0\)。若 \(1\leq p<\infty\)、\(f\in L^p(\mathbb R^d;\mathbb C)\)，则 \(\|K_\varepsilon*f-f\|_p\to0\)。若 \(f\colon\mathbb R^d\to\mathbb C\) 有界且一致连续，则 \(K_\varepsilon*f\to f\) 一致成立。
\end{theorem}
\begin{proof}
记 \(A=\sup_\varepsilon\|K_\varepsilon\|_1\)。利用积分为 \(1\)，将误差写成 \(\int K_\varepsilon(y)(\tau_yf-f)\,\mathrm{d}y\)。对归一化正权 \(|K_\varepsilon|/\|K_\varepsilon\|_1\) 使用 Jensen，再对 \(x\) 积分，得到
\[
 \|K_\varepsilon*f-f\|_p^p
 \leq A^{p-1}\int|K_\varepsilon(y)|\cdot\|\tau_yf-f\|_p^p\,\mathrm{d}y.
\]
平移连续性使 \(|y|<\delta\) 时范数误差统一小于任给的正数；远处使用 \(\|\tau_yf-f\|_p\leq2\|f\|_p\)，再由核的尾部条件控制。先选 \(\delta\)，再令 \(\varepsilon\downarrow0\)，即可完成证明。有界一致连续情形直接对逐点误差作相同分割，近处用一致连续性，远处用 \(2\|f\|_\infty\)。
\end{proof}

有限 \(p\) 与一致连续性的假设不可随意省略。即使取非负紧支集核，跳跃函数也不能在 \(L^\infty\) 中被连续卷积一致逼近。另一方面，一族只满足总积分为 \(1\) 的核可以平移到无穷远，使其与固定函数的卷积远离原函数；质量的位置和绝对质量都要检查。

\ProseParagraphBreak

\textbf{共同光滑逼近。}磨光函数是这种方法的常见实例。取非负 \(\rho\in C_c^\infty(\mathbb R^d)\)、\(\int\rho=1\)，则 \(\rho_\varepsilon*f\to f\) 于有限指数的 \(L^p\)；\cref{b2:prop:03-smooth-convolution}证明了卷积光滑及其求导公式。要把局部光滑化进一步变成光滑紧支集逼近，还需要先做空间截断。下面把两个指数同时逼近的选择写清楚，供下一章延拓使用。

\begin{proposition}\label{b2:prop:03-common-smooth-density}
设 \(d\) 为正整数。若 \(1\leq p,q<\infty\) 且 \(f\in L^p(\mathbb R^d;\mathbb C)\cap L^q(\mathbb R^d;\mathbb C)\)，则存在同一列 \(f_n\in C_c^\infty(\mathbb R^d;\mathbb C)\)，同时在 \(L^p(\mathbb R^d;\mathbb C)\) 与 \(L^q(\mathbb R^d;\mathbb C)\) 中趋于 \(f\)。
\end{proposition}
\begin{proof}
令 \(v_n=f\cdot\mathbf1_{\{|x|\leq n,\,|f(x)|\leq n\}}\)。对 \(r=p,q\)，截掉部分的 \(r\) 次方由 \(|f|^r\) 控制，并几乎处处趋零，故 \(\|v_n-f\|_r\to0\)。每个 \(v_n\) 有界且紧支集。对固定 \(n\)，前一定理使 \(\rho_\varepsilon*v_n\to v_n\) 在两种范数中成立。选 \(0<\varepsilon_n<1/n\)，使两种误差都小于 \(1/n\)，再令 \(f_n=\rho_{\varepsilon_n}*v_n\)。前一节的光滑卷积命题给出 \(f_n\in C_c^\infty\)，支集至多扩大一个 \(O(\varepsilon_n)\) 的距离；三角不等式分别给出两种收敛。两个误差都随尺度趋零，因此每一步都能选取同时满足两项要求的尺度。
\end{proof}


\textbf{一阶误差估计。}上一定理只保证误差趋零，没有给出趋零的速度。若核由同一个可积函数缩放得到，并且它的一阶绝对矩有限，就能把平移误差中的距离因子积分，得到与尺度 \(\varepsilon\) 成正比的上界。

\begin{proposition}\label{b2:prop:03-first-moment-error}
设 \(d\) 为正整数，\(\varepsilon>0\)，\(K\in L^1(\mathbb R^d;\mathbb C)\)，\(K_\varepsilon(x)=\varepsilon^{-d}K(x/\varepsilon)\)、\(\int_{\mathbb R^d}K=1\)，且 \(M_1=\int_{\mathbb R^d}|z|\cdot|K(z)|\,\mathrm{d}z<\infty\)。若 \(f\in C_c^1(\mathbb R^d;\mathbb C)\)、\(1\leq p\leq\infty\)，则
\[
 \|K_\varepsilon*f-f\|_p
 \leq\varepsilon M_1\bigl\||\nabla f|\bigr\|_p.
\]
\end{proposition}
\begin{proof}
微积分基本定理给出
\[
 f(x-\varepsilon z)-f(x)
 =-\varepsilon\int_0^1 z\cdot\nabla f(x-s\varepsilon z)\,\mathrm{d}s.
\]
对绝对值以 \(\varepsilon|z||\nabla f|\) 控制，随后用权重 \(|z||K(z)|\,\mathrm{d}z\,\mathrm{d}s/M_1\) 的 Jensen 不等式并对 \(x\) 积分。平移保持梯度范数，得到误差的 \(p\) 次方至多为 \((\varepsilon M_1)^p\bigl\||\nabla f|\bigr\|_p^p\)。无穷端点直接在梯度上取上确界。若 \(M_1=0\)，则 \(K\) 在原点外几乎处处为零，而原点为零测集，故 \(\int K=0\)，与归一化矛盾；因此 \(M_1>0\)，上述归一化权合法。
\end{proof}

这个估计把定性逼近转成一次阶误差，并说明常数由核的一阶绝对矩和函数梯度控制。若核为偶函数，线性项的积分可能相消；但要把这种相消转成更高阶的范数误差，还需控制二阶余项。后面的 Gauss 核习题会具体计算这一改进。

\subsection{Gauss 核}
\label{b2:subsec:030403}

\begin{theorem}[Gauss 函数的 Fourier 变换]\label{b2:thm:03-gaussian}
设 \(d\) 为正整数，定义实值函数 \(G\colon\mathbb R^d\to\mathbb R\)，\(G(x)=e^{-\pi|x|^2}\)。则 \(G\in L^1(\mathbb R^d;\mathbb R)\)、\(\int_{\mathbb R^d}G=1\)，且 \(\widehat G=G\)。因而对每个 \(\varepsilon>0\)，有
\[
 G_\varepsilon(x)\coloneqq\varepsilon^{-d}e^{-\pi|x|^2/\varepsilon^2},
 \quad \widehat {G_\varepsilon}(\xi)=e^{-\pi\varepsilon^2|\xi|^2}.
\]
\end{theorem}
\begin{proof}
先看一维。若 \(I=\int_{\mathbb R}e^{-\pi x^2}\,\mathrm{d}x\)，则由 Tonelli 和平面极坐标换元，
\[
 I^2=2\pi\int_0^\infty e^{-\pi r^2}r\,\mathrm{d}r=1,
\]
故 \(I=1\)。记 \(g=\widehat G\)。空间矩可积使 \(g\) 可微，而 \(G'=-2\pi xG\)；结合\cref{b2:prop:03-moments}和\cref{b2:prop:03-derivative}得到 \(g'(\xi)=-2\pi\xi g(\xi)\)。所以 \(e^{\pi\xi^2}g(\xi)\) 导数为零，由 \(g(0)=1\) 得 \(g(\xi)=e^{-\pi\xi^2}\)。

高维时 \(G(x)=\prod_j e^{-\pi x_j^2}\)，绝对可积性和 Fubini 将变换分解为一维结果的乘积。最后使用伸缩公式。
\end{proof}

函数 \(G_\varepsilon\) 称为归一化\term[gauss-he]{Gauss 核}[Gaussian kernel]。它有无限支集，但尾部因指数衰减而可控。其 Fourier 变换仍为同类函数，使空间磨光与高频抑制之间出现一个可精确计算的对应：乘子 \(e^{-\pi\varepsilon^2|\xi|^2}\) 在每个有限频率处都严格为正，因而压低高频却不把它直接置零。这与只保留某个有界频率区间、在区间外取零的硬截断不同；下一节反演时将比较两者。
\symidx{\(G_\varepsilon\)：单位质量 Gauss 核}

\subsection{Poisson 核}
\label{b2:subsec:030404}

在 \(\mathbb R^d\) 上另定义
\[
 P_t(x)\coloneqq c_d\dfrac{t}{(t^2+|x|^2)^{(d+1)/2}},
 \quad
 c_d^{-1}=\int_{\mathbb R^d}(1+|z|^2)^{-(d+1)/2}\,\mathrm{d}z.
\]
归一化常数有限且正：单位球内被积函数有界；在半径约 \(2^j\) 的环带上，积分至多为 \(C2^{-j}\)。这些核非负、积分为 \(1\)，并且 \(P_t(x)=t^{-d}P_1(x/t)\)，所以构成逼近恒等。它与上一章的圆周 Poisson 核定义域不同，本节用下标 \(t>0\) 区分欧氏空间的尺度。
\symidx{\(P_t\)：欧氏空间 Poisson 核}

\ProseParagraphBreak

一维时 \(c_1=1/\pi\)，所以 \(P_t(x)=t/(\pi(t^2+x^2))\)。由本章开头双边指数的直接计算，
\[
 \mathcal F(e^{-2\pi t|x|})(\xi)=P_t(\xi).
\]
下一节的反演定理将把这一计算转写为 \(\widehat {P_t}(\xi)=e^{-2\pi t|\xi|}\)。Poisson 核的空间尾部仅按幂衰减，Gauss 核则按指数衰减；两者都逼近恒等，但对远处增长的容许范围不同。

% !TeX root = ../../Book2.tex
\section{Fourier 反演}
\label{b2:sec:0305}

反演的关键不是把指数的符号形式上改回来，而是说明两重振荡积分怎样得到原函数。直接交换两个没有绝对收敛保证的积分没有意义。我们先加一个可积的 Gauss 因子，使换序合法，再让这个因子逐渐趋于 \(1\)。

\subsection{光滑情形}
\label{b2:subsec:030501}

\begin{lemma}\label{b2:lem:03-gaussian-inversion}
设 \(d\) 为正整数，\(f\in L^1(\mathbb R^d;\mathbb C)\)、\(\varepsilon>0\)，并令 \(G_\varepsilon(y)=\varepsilon^{-d}e^{-\pi|y|^2/\varepsilon^2}\)（\(y\in\mathbb R^d\)）。则每个 \(x\in\mathbb R^d\) 满足
\begin{equation}\label{b2:eq:03-gaussian-inversion}
 \int_{\mathbb R^d}\widehat f(\xi)\cdot e^{-\pi\varepsilon^2|\xi|^2}e^{2\pi\mathrm{i}x\cdot\xi}\,\mathrm{d}\xi
 =(f*G_\varepsilon)(x).
\end{equation}
\end{lemma}
\begin{proof}
左侧绝对收敛，因为 \(\widehat f\) 有界而 Gauss 因子可积。代入变换定义，重积分绝对值的积分为 \(\|f\|_1\int e^{-\pi\varepsilon^2|\xi|^2}\,\mathrm{d}\xi<\infty\)，故可用 Fubini。内层正号积分由 Gauss 变换公式及偶性等于 \(G_\varepsilon(x-y)\)，从而得到右侧。
\end{proof}

为何先采用 Gauss 因子，可以与一维的硬截断作比较。对 \(R>0\) 定义
\[
 T_Rf(x)=\int_{-R}^R\widehat f(\xi)\cdot e^{2\pi\mathrm i x\xi}\,\mathrm d\xi.
\]
这个积分对 \(f\in L^1(\mathbb R)\) 有定义。代入变换后，绝对积分至多为 \(2R\|f\|_1\)，故交换次序合法，并得到
\[
 T_Rf(x)=\int_{\mathbb R}f(y)\cdot K_R(x-y)\,\mathrm dy,
 \quad K_R(z)=\frac{\sin(2\pi Rz)}{\pi z},\quad K_R(0)=2R.
\]
这里的卷积逐点绝对收敛，因为 \(|K_R(z)|\leq2R\)，故被积函数的绝对积分至多为 \(2R\|f\|_1\)。但是对每个 \(R>0\)，\(K_R\) 本身不属于 \(L^1(\mathbb R)\)。令 \(u=2\pi Rz\)，则
\[
 \int_{1/(2R)}^\infty|K_R(z)|\,\mathrm dz
 =\frac1\pi\int_\pi^\infty\frac{|\sin u|}{u}\,\mathrm du,
 \quad
 \int_{n\pi}^{(n+1)\pi}\frac{|\sin u|}{u}\,\mathrm du
 \geq\frac{2}{(n+1)\pi}\quad(n\geq1).
\]
最后一个下界使用每个半波的 \(|\sin u|\) 积分等于 \(2\)；把这些下界相加便得到发散的调和级数。因此硬截断核连单个核的绝对可积性都不满足，不能应用\cref{b2:thm:03-approximate-identity}直接推出 \(T_Rf\to f\)。Gauss 因子的逆积分则为非负核，并由上一节的归一化满足 \(\|G_\varepsilon\|_1=\int G_\varepsilon=1\)；其固定半径之外的质量随 \(\varepsilon\downarrow0\) 趋零，故满足逼近恒等的全部条件。两种方法在每个固定 \(R>0\) 或 \(\varepsilon>0\) 下都能合法换序，区别在于让 \(R\uparrow\infty\) 或 \(\varepsilon\downarrow0\) 时，是否已有可用的极限控制。

\ProseParagraphBreak

\Cref{b2:fig:03-hard-cutoff-gaussian-kernels}用一维的实际核比较这两种形状。取 \(R=2\)、\(\varepsilon=1/4\)，两核在原点的值都为 \(4\)，并使用同一组横、纵坐标刻度。硬截断核的正负旁瓣来自锐利的频率边界；Gauss 核始终非负，离开原点后迅速降低。有限画幅不能判断整条实线上的绝对积分是否有限，这一点仍由上面的半波下界和 Gauss 单位质量公式决定。

% 自拟函数图：同峰值4、同横纵尺度；sinc分两侧绘制并显式补入连续延拓值，避免零点商式。
\begin{figure}[htbp]
\centering
\begin{tikzpicture}[x=1.65cm,y=0.70cm,font=\small]
 \begin{scope}
  \draw[gray!60] (-1.5,-1) rectangle (1.5,4.2);
  \draw[gray!55] (-1.5,0)--(1.5,0);
  \foreach \xt in {-1.5,-1,-0.5,0,0.5,1,1.5}{
   \draw (\xt,-1)--(\xt,-0.90);
   \node[below] at (\xt,-1.35) {\(\xt\)};
  }
  \foreach \yt in {-1,0,1,2,3,4}{
   \draw (-1.5,\yt)--(-1.43,\yt);
   \node[left] at (-1.54,\yt) {\(\yt\)};
  }
  \draw[thick,AnalysisBlue] plot[domain=-1.5:-0.001,samples=181]
   (\x,{sin(720*\x)/(pi*\x)});
  \draw[thick,AnalysisBlue] plot[domain=0.001:1.5,samples=181]
   (\x,{sin(720*\x)/(pi*\x)});
  \draw[thick,AnalysisBlue] (-0.001,{sin(-0.72)/(-0.001*pi)})
   --(0,4)--(0.001,{sin(0.72)/(0.001*pi)});
  \node[anchor=south west] at (-1.5,4.4) {\(K_2(x)=\dfrac{\sin(4\pi x)}{\pi x}\)};
  \node[below right] at (1.5,-1.90) {\(x\)};
 \end{scope}
 \begin{scope}[xshift=6.2cm]
  \draw[gray!60] (-1.5,-1) rectangle (1.5,4.2);
  \draw[gray!55] (-1.5,0)--(1.5,0);
  \foreach \xt in {-1.5,-1,-0.5,0,0.5,1,1.5}{
   \draw (\xt,-1)--(\xt,-0.90);
   \node[below] at (\xt,-1.35) {\(\xt\)};
  }
  \foreach \yt in {-1,0,1,2,3,4}{
   \draw (-1.5,\yt)--(-1.43,\yt);
   \node[left] at (-1.54,\yt) {\(\yt\)};
  }
  \draw[thick,AnalysisBlue] plot[domain=-1.5:1.5,samples=201]
   (\x,{4*exp(-16*pi*\x*\x)});
  \node[anchor=south west] at (-1.5,4.4) {\(G_{1/4}(x)=4e^{-16\pi x^2}\)};
  \node[below right] at (1.5,-1.90) {\(x\)};
 \end{scope}
\end{tikzpicture}
\caption[硬频率截断核与 Gauss 核]{硬频率截断核与 Gauss 核。一维取 \(R=2\)、\(\varepsilon=1/4\)，两核在原点均为 \(4\)，均绘于 \(-1.5\leq x\leq1.5\)，纵坐标尺度相同。左图在原点采用连续延拓值，并有正负旁瓣；右图非负。图只显示有限范围内的形状，不代替正文对无穷远尾部的论证。}
\label{b2:fig:03-hard-cutoff-gaussian-kernels}
\end{figure}

\begin{proposition}\label{b2:prop:03-smooth-inversion}
设 \(d\) 为正整数。若 \(f\in C_c^\infty(\mathbb R^d;\mathbb C)\)，则 \(\widehat f\in L^1(\mathbb R^d;\mathbb C)\)，并有在 \(\mathbb R^d\) 上处处成立的绝对收敛反演公式 \(f=\check {\widehat f}\)。
\end{proposition}
\begin{proof}
取整数 \(m>d\)。对每个坐标 \(j\)，\(f\) 及中间导数 \(\partial_j^k f\)（\(1\leq k\leq m\)）都光滑、紧支集，因而可积。逐次应用\cref{b2:prop:03-derivative}，再在非零频率处选取绝对值最大的坐标，得到 \(|\widehat f(\xi)|\leq C|\xi|^{-m}\)。结合 \(|\widehat f|\leq\|f\|_1\)，便有全空间估计 \(|\widehat f(\xi)|\leq C'(1+|\xi|)^{-m}\)，于是变换可积。在\eqref{b2:eq:03-gaussian-inversion}左侧以 \(|\widehat f|\) 控制，得到 \(\check {\widehat f}(x)\)；右侧由有界一致连续函数的逼近恒等结论一致趋于 \(f\)。
\end{proof}

此证明只用到已有的 Gauss 积分、导数变换与逼近恒等，没有预设一般反演。因此它也提供了以后在测试函数空间上建立变换的起点。

\subsection{\texorpdfstring{\(L^1\)}{L¹} 反演}
\label{b2:subsec:030502}

\begin{theorem}[Fourier 反演]\label{b2:thm:03-inversion}
设 \(d\) 为正整数，\(f\in L^1(\mathbb R^d;\mathbb C)\) 且 \(\widehat f\in L^1(\mathbb R^d;\mathbb C)\)。对每个 \(x\in\mathbb R^d\)，绝对收敛积分
\[
 g(x)\coloneqq\int_{\mathbb R^d}\widehat f(\xi)\cdot e^{2\pi\mathrm{i}x\cdot\xi}\,\mathrm{d}\xi
\]
定义了 \(f\) 的唯一连续代表 \(g\in C_0(\mathbb R^d;\mathbb C)\)，即 \(g=f\) 几乎处处。若事先指定的有限值可测代表 \(f_0\colon\mathbb R^d\to\mathbb C\) 在某一点 \(x_0\in\mathbb R^d\) 连续，则 \(f_0(x_0)=g(x_0)\)。
\end{theorem}
\begin{proof}
首先，绝对收敛给出 \(|g(x)|\leq\|\widehat f\|_1\)，而 \(g(x)=\widehat{\widehat f}(-x)\) 使 Riemann--Lebesgue 引理可用于可积函数 \(\widehat f\)，从而 \(g\in C_0(\mathbb R^d)\)。由控制收敛，\eqref{b2:eq:03-gaussian-inversion}左侧一致趋于 \(g=\check {\widehat f}\)，因为误差不超过
\[
 \int|\widehat f(\xi)|\cdot|e^{-\pi\varepsilon^2|\xi|^2}-1|\,\mathrm{d}\xi\longrightarrow0.
\]
右侧由\cref{b2:thm:03-approximate-identity}在 \(L^1\) 中趋于 \(f\)。这与左侧的一致收敛如何合起来识别原函数，可以先在一个有限测度的球 \(B=B(0,N)\) 中比较：
\[
 \int_B|f-g|
 \leq\|f-G_\varepsilon*f\|_1
       +|B|\cdot\|G_\varepsilon*f-g\|_\infty
 \longrightarrow0.
\]
因此 \(f=g\) 在每个这样的球中几乎处处成立。让 \(N\) 遍历正整数，可数个零测例外集的并仍为零测集，便得到全空间的几乎处处等式。

若另有连续代表 \(h=f\) 几乎处处，则连续函数 \(h-g\) 几乎处处为零。如果它在某点不为零，连续性就给出一个球，使其绝对值在该球上有严格正的下界；球的 Lebesgue 测度为正，与几乎处处为零矛盾。这证明了连续代表的唯一性。

最后设 \(f_0\) 是已指定的代表，并在 \(x_0\) 连续。记 \(E=\{\,y\mid f_0(y)\ne g(y)\,\}\)，则 \(E\) 为零测集。每个球 \(B(x_0,1/n)\) 都有正测度，因此可选 \(y_n\in B(x_0,1/n)\setminus E\)。于是 \(y_n\to x_0\) 且 \(f_0(y_n)=g(y_n)\)。分别使用 \(f_0\) 在 \(x_0\) 的连续性和 \(g\) 的连续性，得到 \(f_0(x_0)=g(x_0)\)。
\end{proof}

这称为\term[fourier-fanyan-gongshi]{Fourier 反演公式}[Fourier inversion formula]。条件 \(\widehat f\in L^1\) 保证普通反演积分绝对收敛；它不是所有形式的反演都成立的必要条件。尤其对区间示性函数，此条件失败，但仍可讨论带正则化因子的反演。代表的区别也有实际内容：把零函数只在原点改成 \(1\)，变换仍恒为零，反演积分给出的仍是处处为零的连续代表，不会恢复在原点另行指定的值。

\subsection{点态反演}
\label{b2:subsec:030503}

沿用第一卷定义 14.4.1：设 \(d\) 为正整数，对 \(f\in L^1_{\mathrm{loc}}(\mathbb R^d;\mathbb C)\) 指定一个有限值可测代表 \(f_0\colon\mathbb R^d\to\mathbb C\)。若在 \(x\in\mathbb R^d\) 的小球内满足
\[
 \dfrac1{|B(x,r)|}\int_{B(x,r)}|f_0(y)-f_0(x)|\,\mathrm{d}y\longrightarrow0,
\]
则 \(x\) 为该代表的 Lebesgue 点。第一卷命题 14.4.6 证明了非负径向非增单位质量核在 Lebesgue 点处的逼近；本章以 \(f\in L^1\) 满足该命题对非紧支集核要求的全局可积性。Gauss 核和本章 Poisson 核均满足核的条件。

\begin{corollary}[Gauss 正则化反演]\label{b2:cor:03-pointwise-inversion}
设 \(d\) 为正整数，\(f\in L^1(\mathbb R^d;\mathbb C)\)，并指定一个有限值可测代表 \(f_0\colon\mathbb R^d\to\mathbb C\)。则在 \(f_0\) 的每个 Lebesgue 点 \(x\in\mathbb R^d\) 有
\[
 f_0(x)=\lim_{\varepsilon\downarrow0}
 \int\widehat f(\xi)\cdot e^{-\pi\varepsilon^2|\xi|^2}e^{2\pi\mathrm{i}x\cdot\xi}\,\mathrm{d}\xi.
\]
特别地，该结论几乎处处成立。
\end{corollary}
\begin{proof}
由\eqref{b2:eq:03-gaussian-inversion}，积分就是 \((G_\varepsilon*f_0)(x)\)。可以应用上述第一卷命题；为说明局部平均条件如何控制本章这个无限支集核，这里写出相应的估计。固定 \(f_0\) 的一个 Lebesgue 点 \(x\)，记
\[
 A(r)=\int_{|y|<r}|f_0(x-y)-f_0(x)|\,\mathrm dy,
 \quad g_\varepsilon(r)=\varepsilon^{-d}e^{-\pi r^2/\varepsilon^2}.
\]
给定 \(\eta>0\)，存在 \(\delta>0\)，使 \(A(r)\leq\eta\cdot\omega_dr^d\) 对 \(0<r\leq\delta\) 成立，其中 \(\omega_d\) 是单位球体积。在球内使用
\(g_\varepsilon(|y|)=g_\varepsilon(\delta)+\int_{|y|}^{\delta}(-g_\varepsilon'(r))\,\mathrm dr\)，由非负积分的 Tonelli 换序得
\[
 \begin{aligned}
 &\int_{|y|<\delta}G_\varepsilon(y)\cdot|f_0(x-y)-f_0(x)|\,\mathrm dy\\
 &\quad=g_\varepsilon(\delta)A(\delta)
       +\int_0^\delta(-g_\varepsilon'(r))A(r)\,\mathrm dr\\
 &\quad\leq\eta\left(g_\varepsilon(\delta)\cdot\omega_d\delta^d
       +\int_0^\delta(-g_\varepsilon'(r))\cdot\omega_dr^d\,\mathrm dr\right)\\
 &\quad=\eta\int_{|y|<\delta}G_\varepsilon(y)\,\mathrm dy\leq\eta.
 \end{aligned}
\]
括号中的最后等式把同一个分层积分用于常函数 \(1\)。球外含 \(f_0\) 和含常数 \(f_0(x)\) 的两部分，分别由
\[
 \varepsilon^{-d}e^{-\pi\delta^2/\varepsilon^2}\|f\|_1,
 \qquad |f_0(x)|\int_{|y|\geq\delta}G_\varepsilon(y)\,\mathrm dy
\]
控制；两者在 \(\varepsilon\downarrow0\) 时都趋零。先固定 \(\eta,\delta\)，取尺度极限，再令 \(\eta\downarrow0\)，就得 \((G_\varepsilon*f_0)(x)\to f_0(x)\)。最后用第一卷定理 14.4.2 的几乎处处 Lebesgue 点结论。
\end{proof}

这项结论不要求 \(\widehat f\) 可积，但不能把 Gauss 因子直接删去，也不能未经证明改成任意频率截断形状。一个有两个单侧极限的一维函数，在跳跃点经偶核正则化时趋向中点值：近处分别用左右极限，远处用核在离原点固定距离外的一致衰减估计。点本身的任意赋值不会被积分恢复。

\ProseParagraphBreak

例如对 \(f=\mathbf1_{(0,1)}\)，卷积可以完整算成
\[
 (G_\varepsilon*f)(x)
 =\int_{(x-1)/\varepsilon}^{x/\varepsilon}e^{-\pi z^2}\,\mathrm dz.
\]
当 \(0<x<1\) 时，两端分别趋于负、正无穷，极限为 \(1\)；当 \(x<0\) 或 \(x>1\) 时，积分区间进入同一侧尾部，极限为零。在 \(x=0,1\) 处，偶性使极限为 \(1/2\)。区间端点处的值不论补成什么，都不影响变换；正则化恢复的是左右平均。另一方面，原变换是指数相位乘 sinc 函数，其绝对积分发散，所以这个例子没有满足无正则化反演的假设。

\ProseParagraphBreak

反演还能确定上一节的一维 Poisson 变换。对 \(f(x)=e^{-2\pi t|x|}\)，直接计算已给出 \(\widehat f=P_t\)，且两者均可积。反演与偶性于是给出
\[
 \widehat {P_t}(\xi)=e^{-2\pi t|\xi|}\quad(d=1).
\]
两侧连续，所以最初的几乎处处等式提升为处处等式。

反演建立以后，还能证明与卷积定理互相对应的乘积公式。两个 \(L^1\) 函数的乘积未必可积；下面利用一个因子的可积 Fourier 变换，使这个因子有界，从而保证乘积可积。

\begin{proposition}\label{b2:prop:03-product-transform}
设 \(d\) 为正整数，\(f,g\in L^1(\mathbb R^d;\mathbb C)\)，且 \(\widehat g\in L^1(\mathbb R^d;\mathbb C)\)。选取反演给出的 \(g\) 的有界连续代表，则 \(fg\in L^1(\mathbb R^d;\mathbb C)\)，并有对每个 \(\xi\in\mathbb R^d\) 都成立的公式
\[
 \widehat{fg}(\xi)=\int_{\mathbb R^d}\widehat f(\xi-\eta)\cdot\widehat g(\eta)\,\mathrm d\eta.
\]
右侧的卷积积分绝对收敛。
\end{proposition}
\begin{proof}
反演给出 \(\|g\|_\infty\leq\|\widehat g\|_1\)，所以 \(\|fg\|_1\leq\|f\|_1\cdot\|\widehat g\|_1\)。将 \(g(x)=\int\widehat g(\eta)e^{2\pi\mathrm i x\cdot\eta}\,\mathrm d\eta\) 代入 \(fg\) 的变换。双重积分的绝对积分为 \(\|f\|_1\cdot\|\widehat g\|_1\)，故可用 Fubini 交换次序；内层对 \(x\) 的积分恰为 \(\widehat f(\xi-\eta)\)。反演最初若只几乎处处成立，也不改变乘积的等价类。最后用 \(\|\widehat f\|_\infty\leq\|f\|_1\) 保证所写卷积处处绝对收敛。
\end{proof}

乘积截断会在频率中产生卷积，空间卷积则变成频率乘积。这两个操作都是局部化的方法，但对频率的影响不同：把一个函数限在某一区域，通常会把其频率分散；用一个平滑核平均，则按该核的变换压低频率。后面讨论频率局部化时还会反复使用这两个公式的差别。

\subsection{唯一性与正则性的频率判据}
\label{b2:subsec:030504}

\begin{corollary}[唯一性]\label{b2:cor:03-uniqueness}
设 \(d\) 为正整数。若 \(f,g\in L^1(\mathbb R^d;\mathbb C)\) 且 \(\widehat f=\widehat g\)，则 \(f=g\) 在 \(\mathbb R^d\) 上几乎处处成立。
\end{corollary}
\begin{proof}
差 \(h=f-g\) 的变换恒为零，特别属于 \(L^1\)，故 Fourier 反演使 \(h=0\) 几乎处处。
\end{proof}

零函数在频率侧可积，所以对差函数的零变换使用反演，就能在全部 \(L^1\) 中得到唯一性。例如若 \(f*G_\varepsilon=0\) 对某个固定 \(\varepsilon>0\) 成立，则变换相乘后得到 \(\widehat f(\xi)e^{-\pi\varepsilon^2|\xi|^2}=0\)。Gauss 因子处处非零，因而 \(f=0\)。然而从带误差的卷积数据恢复 \(f\)，却要除以趋零的 Gauss 因子；唯一可恢复与稳定可恢复是不同问题。



\begin{proposition}\label{b2:prop:03-inverse-regularity}
设 \(d\) 为正整数，\(m\) 为非负整数，\(f\in L^1(\mathbb R^d;\mathbb C)\)，且 \(\int_{\mathbb R^d}(1+|\xi|)^m|\widehat f(\xi)|\,\mathrm{d}\xi<\infty\)。则 \(f\) 有一个代表 \(g\in C^m(\mathbb R^d;\mathbb C)\)，且对每个满足 \(|\alpha|\leq m\) 的多重指标 \(\alpha\in\mathbb N_0^d\)，都有 \(\partial^\alpha g\in C_0(\mathbb R^d;\mathbb C)\)。
\end{proposition}
\begin{proof}
零阶权已保证反演适用，取连续代表 \(g=\check {\widehat f}\)。每次对 \(x_j\) 求导，只给频率被积函数增加 \(2\pi\mathrm{i}\xi_j\)。所有相应差商都由假设中的可积矩控制，所以可逐阶求导，得到
\[
 \partial^\alpha g(x)=\int(2\pi\mathrm{i}\xi)^\alpha\widehat f(\xi)e^{2\pi\mathrm{i}x\cdot\xi}\,\mathrm{d}\xi.
\]
其绝对值由频率矩积分控制，连续性与无穷远消失由 Riemann--Lebesgue 引理给出。
\end{proof}

这与空间矩产生频率微分的命题构成对称关系。频率中的可积衰减可以恢复空间光滑性，条件必须作用于带权的绝对积分；仅知道变换逐点趋零，没有提供这些矩。反演因此不仅是一条恢复数值的公式，也是从频率估计提取正则性的工具。


\Cref{b2:tab:03-transform-rules}集中列出本章的运算规则，便于在后续方程计算中同时核对公式与假设。
% 本书自编查阅表；全部公式与条件来自第三章已证明结果。
\begin{table}[H]
\centering\small\renewcommand{\arraystretch}{1.2}\setlength{\tabcolsep}{4pt}
\caption[Fourier 变换的运算规则]{Fourier 变换的运算规则。\(d\) 为正整数，\(x,\xi\in\mathbb R^d\)，统一采用指数 \(e^{-2\pi\mathrm{i}x\cdot\xi}\)；允许复值函数，实值情形包含在内。第三列列出各行的可积性与参数条件。}
\label{b2:tab:03-transform-rules}
\begin{tabularx}{\linewidth}{@{}>{\raggedright\arraybackslash}p{3.4cm}>{\raggedright\arraybackslash}p{6.1cm}>{\raggedright\arraybackslash}X@{}}
\toprule
空间中的操作 & 频率中的表达 & 适用条件 \\
\midrule
\(f(x-h)\) & \(e^{-2\pi\mathrm{i}h\cdot\xi}\cdot\widehat f(\xi)\) & \(f\in L^1(\mathbb R^d)\)，\(h\in\mathbb R^d\) \\
\(e^{2\pi\mathrm{i}\eta\cdot x}\cdot f(x)\) & \(\widehat f(\xi-\eta)\) & \(f\in L^1(\mathbb R^d)\)，\(\eta\in\mathbb R^d\) \\
\(f(Ax)\) & \(|\det A|^{-1}\cdot\widehat f(A^{-\mathsf T}\xi)\) & \(f\in L^1(\mathbb R^d)\)，\(A\in\mathbb R^{d\times d}\)，\(\det A\ne0\) \\
\(\partial_jf\) & \(2\pi\mathrm{i}\xi_j\widehat f(\xi)\) & \(1\leq j\leq d\)，\(f\in L^1(\mathbb R^d)\)，弱导数 \(\partial_jf\in L^1(\mathbb R^d)\) \\
\(x_jf\) & \(\dfrac{\mathrm{i}}{2\pi}\partial_{\xi_j}\widehat f\) & \(1\leq j\leq d\)，\(f,x_jf\in L^1(\mathbb R^d)\) \\
\(f*g\) & \(\widehat f\cdot\widehat g\) & \(f,g\in L^1(\mathbb R^d)\) \\
\(\begin{gathered}G_\varepsilon(x)=\\\varepsilon^{-d}e^{-\pi|x|^2/\varepsilon^2}\end{gathered}\) & \(e^{-\pi\varepsilon^2|\xi|^2}\) & \(\varepsilon>0\)，此核可积且积分为 \(1\) \\
\bottomrule
\end{tabularx}
\end{table}



\begin{chaptersummary}
可积性使 Fourier 变换成为定义在每个频率处的连续函数。平移与调制互换，线性伸缩产生逆转置频率和体积因子；空间矩负责频率微分，空间导数负责频率衰减。卷积在适当指数下存在，其变换是两个变换的乘积。

\ProseParagraphBreak

反演通过逼近恒等取得，而不是交换一个没有收敛依据的双重积分。若变换可积，可以去掉正则化，得到连续代表；若只知道原函数可积，则在 Lebesgue 点保留 Gauss 极限。唯一性仍对整个 \(L^1\) 成立，但高频信息可能在磨光后变得极小，稳定恢复需要额外估计。下一章将以平方能量建立更大的变换空间。
\end{chaptersummary}

\BookExercises

这些题目分别检验相位、卷积的几何形状、逼近条件、误差阶和反演的稳定性。除明确注明的一维题外，定义域均为 \(\mathbb R^d\)。

% 自拟：积分三角不等式的取等条件与调制相联系。
\begin{exercise}\label{b2:ex:03-phase-equality}
设 \(d\) 为正整数，\(0\ne f\in L^1(\mathbb R^d;\mathbb C)\)，\(\xi_0\in\mathbb R^d\)。证明 \(|\widehat f(\xi_0)|=\|f\|_1\) 当且仅当存在 \(c\in\mathbb C\)、\(|c|=1\)，使 \(f(x)=c e^{2\pi\mathrm{i}x\cdot\xi_0}|f(x)|\) 几乎处处。若 \(f\geq0\)，证明对每个 \(\xi_0\ne0\) 都有严格不等式。
\end{exercise}
\begin{solution}
令 \(h(x)=f(x)\cdot e^{-2\pi\mathrm{i}x\cdot\xi_0}\)。若积分模等于绝对积分，选 \(c\) 使 \(\overline c\int h=\int|h|\)。非负函数 \(|h|-\operatorname{Re}(\overline ch)\) 的积分为零，故它几乎处处为零。这迫使 \(h=c|h|\)，反向由代入得到。若 \(f\geq0\) 且 \(\xi_0\ne0\)，取等要求 \(e^{-2\pi\mathrm{i}x\cdot\xi_0}\) 在 \(\{f>0\}\) 上几乎处处固定。其任一水平集包含于可数个平行超平面，Lebesgue 测度为零；而非零非负可积函数的正集具有正测度，矛盾。
\end{solution}

% 自拟：矩形卷积生成三角函数，以反演计算sinc平方积分。
\begin{exercise}\label{b2:ex:03-tent}
在一维令 \(b=\mathbf1_{[-1/2,1/2]}\)。证明 \(b*b(x)=(1-|x|)_+\)，求其 Fourier 变换，并用反演公式计算
\[
 \int_{\mathbb R}\left(\dfrac{\sin(\pi\xi)}{\pi\xi}\right)^2\,\mathrm{d}\xi.
\]
\end{exercise}
\begin{solution}
卷积等于区间 \([-1/2,1/2]\) 与 \([x-1/2,x+1/2]\) 的重叠长度，所以为三角形函数。区间变换和卷积定理给出其变换为 \((\sin(\pi\xi)/(\pi\xi))^2\)，原点补为 \(1\)。该函数有界，且无穷远以 \(|\xi|^{-2}\) 控制，因而可积。在三角形函数的连续点 \(x=0\) 使用反演，所得积分为 \(1\)。
\end{solution}

% 自拟：逼近恒等的绝对质量界不可由总质量和尾部替代。
\begin{optionalexercise}\label{b2:ex:03-signed-kernel}
在一维，令 \(G_\varepsilon(x)=\varepsilon^{-1}e^{-\pi x^2/\varepsilon^2}\)，平移约定为 \(\tau_hu(x)=u(x-h)\)，并设
\[
 K_\varepsilon=G_\varepsilon+\varepsilon^{-1}(\tau_\varepsilon G_\varepsilon-G_\varepsilon),
 \quad0<\varepsilon<1.
\]
证明其积分为 \(1\)，绝对积分尾部在每个固定的原点邻域外趋零，但 \(\|K_\varepsilon\|_1\to\infty\)。再证明，对 \(f\in C_c^2(\mathbb R)\)，有 \(K_\varepsilon*f\to f-f'\) 一致成立。
\end{optionalexercise}
\begin{solution}
平移不改变积分，所以总积分为 \(1\)。作尺度换元，
\[
 \|\tau_\varepsilon G_\varepsilon-G_\varepsilon\|_1
 =\|G(\cdot-1)-G\|_1=c>0,
\]
因此 \(\|K_\varepsilon\|_1\geq c/\varepsilon-1\)。对固定 \(\delta>0\)，当 \(\varepsilon<\delta/2\) 时，两个 Gauss 尾部都由尺度变量中 \(|z|>\delta/(2\varepsilon)\) 的指数尾控制；乘以 \(1+2/\varepsilon\) 后仍趋零。

设 \(h_\varepsilon=G_\varepsilon*f\)，则 \(h_\varepsilon\to f\)、\(h_\varepsilon'\to f'\) 一致成立。又有
\[
 \dfrac{h_\varepsilon(x-\varepsilon)-h_\varepsilon(x)}{\varepsilon}
 =-\int_0^1h_\varepsilon'(x-s\varepsilon)\,\mathrm{d}s\longrightarrow-f'(x)
\]
一致成立，因为 \(f'\) 一致连续。把这一式加到 \(h_\varepsilon\) 上便得到结论。
\end{solution}

% 自拟：利用消失的一阶矩改进逼近误差阶。
% 为题头及首个公式预留空间，避免超链接锚点滞留前页。
\Needspace{6\baselineskip}
\begin{exercise}\label{b2:ex:03-gaussian-error}
设 \(d\) 为正整数，\(M\geq0\)，\(f\in C^2(\mathbb R^d;\mathbb C)\) 有界，且每一点的 Hesse 矩阵的算子范数至多为 \(M\)。对 \(\varepsilon>0\)，记 \(G_\varepsilon(x)=\varepsilon^{-d}e^{-\pi|x|^2/\varepsilon^2}\)。证明
\[
 \|G_\varepsilon*f-f\|_\infty\leq\dfrac{dM}{4\pi}\varepsilon^2.
\]
说明为什么只使用 Lipschitz 条件通常只能得到一次阶误差。
\end{exercise}
\begin{solution}
Taylor 公式给出 \(|f(x-y)-f(x)+\nabla f(x)\cdot y|\leq M|y|^2/2\)。Gauss 核为偶函数，一阶矩为零，因此线性项的积分消失。由一维分部积分，\(\int x^2e^{-\pi x^2}\,\mathrm{d}x=1/(2\pi)\)，所以 \(\int|y|^2G_\varepsilon(y)\,\mathrm{d}y=d\varepsilon^2/(2\pi)\)，得到所求界。只有 Lipschitz 常数 \(L\) 时，直接估计给出 \(L\int|y|G_\varepsilon(y)\,\mathrm{d}y=C_dL\varepsilon\)；缺少可相消的线性近似，不能由这个条件普遍推出二次阶。
\end{solution}

% 自拟：反演与唯一性建立后计算Poisson半群。
\begin{exercise}\label{b2:ex:03-poisson-semigroup}
在一维令 \(P_t(x)=t/[\pi(t^2+x^2)]\)，\(t>0\)。证明对任意 \(t,s>0\) 都有 \(P_t*P_s=P_{t+s}\)，并证明 \(u(t,x)=P_t*f(x)\) 对 \(f\in L^1(\mathbb R;\mathbb C)\) 满足 \(u_{tt}+u_{xx}=0\) 于 \(t>0\)，同时 \(u(t,\cdot)\to f\) 于 \(L^1\)。
\end{exercise}
\begin{solution}
三个核都可积。它们的变换分别为 \(e^{-2\pi t|\xi|}\)、\(e^{-2\pi s|\xi|}\)、\(e^{-2\pi(t+s)|\xi|}\)，由卷积定理与唯一性先得到半群式几乎处处成立。又因 \(P_t\in L^1\)、\(P_s\in L^\infty\)，\cref{b2:prop:03-convolution-continuity}说明左侧的逐点卷积连续；右侧 \(P_{t+s}\) 也连续，所以等式处处成立。对每个 \(t\geq\varepsilon>0\) 的局部范围，\(P_t\) 的二阶 \(t,x\) 导数都由一个有限常数控制，故可在与 \(f\) 的卷积中求导。直接求导有 \((\partial_t^2+\partial_x^2)P_t(x)=0\)，从而得到方程。初值收敛由逼近恒等定理给出，不要求 \(f\) 有逐点边界值。
\end{solution}

% 自拟：卷积幂等元与频率空间的连通性。
\begin{exercise}\label{b2:ex:03-idempotent}
设 \(d\) 为正整数。证明若 \(f\in L^1(\mathbb R^d;\mathbb C)\) 满足 \(f*f=f\)，则 \(f=0\) 几乎处处。与圆周上的常函数比较，说明底空间改变后哪一步不同。
\end{exercise}
\begin{solution}
变换给出 \(\widehat f(\xi)^2=\widehat f(\xi)\)，所以每个值为 \(0\) 或 \(1\)。连续函数从连通空间 \(\mathbb R^d\) 映入这两个点，只能为常数；又因无穷远趋零，它恒为零。最后用唯一性。圆周上的频率空间为离散的 \(\mathbb Z\)，系数可以在某些频率上取 \(1\)、其余取 \(0\)；常函数 \(1\) 就满足周期卷积幂等，而不违反任何连通性结论。
\end{solution}

% 自拟：Gauss磨光的单射性与逆操作的不稳定性。
\begin{optionalexercise}\label{b2:ex:03-unstable-deconvolution}
在一维令 \(G(x)=e^{-\pi x^2}\)，\(f_n(x)=e^{-\pi x^2}e^{2\pi\mathrm{i}nx}\)，\(n=1,2,\ldots\)。计算 \(G*f_n\)，证明 \(\|f_n\|_1=1\) 而 \(\|G*f_n\|_1\to0\)。据此说明，即使 Gauss 卷积是单射，也不存在常数 \(C\) 使 \(\|f\|_1\leq C\|G*f\|_1\) 对所有 \(f\in L^1(\mathbb R;\mathbb C)\) 成立。
\end{optionalexercise}
\begin{solution}
在卷积积分中写 \(y=x/2+z\)，平方项满足 \((x-y)^2+y^2=x^2/2+2z^2\)。用 Gauss 的已知变换，得到
\[
 (G*f_n)(x)=2^{-1/2}e^{-\pi n^2/2}e^{-\pi x^2/2}e^{\pi\mathrm{i}nx}.
\]
因此 \(\|G*f_n\|_1=e^{-\pi n^2/2}\)，而 \(\|f_n\|_1=\int G=1\)。若所述逆估计成立，便有 \(1\leq Ce^{-\pi n^2/2}\) 对全部 \(n\) 成立，矛盾。
\end{solution}

% 自拟：卷积支集估计与几何重叠，避免把包含误认为等号。
\begin{exercise}\label{b2:ex:03-support-convolution}
设 \(d\) 为正整数，\(f,g\in L^1(\mathbb R^d;\mathbb C)\) 分别几乎处处支撑在紧集 \(A,B\subset\mathbb R^d\) 内。证明 \(f*g\) 几乎处处支撑在 \(A+B=\{a+b:a\in A,b\in B\}\) 内。若 \(f,g\) 是非负连续紧支集函数，证明其支集恰为两支集的 Minkowski 和。
\end{exercise}
\begin{solution}
若 \(x\notin A+B\)，则不可能同时有 \(y\in B\) 与 \(x-y\in A\)，所以选取在相应紧集外为零的代表后，卷积积分为零。一般代表给出相同的等价类，故得到第一项。

非负连续情形下，若 \(f(a)>0\)、\(g(b)>0\)，连续性保证在 \(y=b\) 的某个球上同时有 \(g(y)>0\)、\(f(a+b-y)>0\)，因此卷积在 \(a+b\) 严格为正。对支集中的任意 \(a,b\)，用正值点列逼近，得到 \(a+b\) 属于卷积支集。反向已由第一项证明；连续卷积在支集外的开集上若几乎处处为零，便处处为零。非负性在排除抵消时使用，不能省去。
\end{solution}

% !TeX root = ../../Book2.tex
\chapter{\texorpdfstring{\(L^2\)}{L²} 与 \texorpdfstring{\(L^p\)}{Lᵖ} Fourier 理论}
\label{b2:ch:04}
\BookChapterNumber{b2:ch:04}
% Axler2020Measure, Section11C, 印刷375--377页核对L2延拓的定义和证明路线。
% Hausdorff--Young在本章用有限简单函数和已证明三线引理完整推导；不前借第19章。
\begin{chapterabstract}
Fourier 变换不仅处理可积函数，还可以通过能量等式作用在整个平方可积空间。本章证明 Plancherel 延拓的良定义、满射性与反演，再用两端估计和复参数插值建立 Hausdorff--Young 不等式。空间伸缩与分离波包的例子显示，大于二的指数不能按同样方式延拓到共轭指数空间。最后以频率乘子说明酉性怎样简化线性算子，并引出广义变换的需要。
\end{chapterabstract}
\begin{learninggoals}
\begin{itemize}
\item\label{b2:item:04-goal-extension}从共同稠密类构造唯一的平方范数变换，说明定义与逼近列无关。
\item\label{b2:item:04-goal-unitary}用反演证明酉性，区分等距、满射及逐点积分三个不同问题。
\item\label{b2:item:04-goal-interpolation}理解两端估计怎样经参数函数产生中间指数，并核对插值的每个积分和极限。
\item\label{b2:item:04-goal-failure}用尺度及符号平均构造 \(p>2\) 的范数障碍，准确表述其否定范围。
\item\label{b2:item:04-goal-multiplier}在 \(L^2\) 上定义频率乘子，计算范数，说明推广到其他指数还需哪些估计。
\end{itemize}
\end{learninggoals}

测量一个信号时，平方积分常表示它的总能量。如果把信号分成不同频率，能量是否恰好等于频率幅值的平方积分？这个问题对不可积但平方可积的信号同样有意义；然而上一章的绝对收敛积分还不能定义它们的变换。本章从能量等式出发构造变换，并用它求解 \((a-\Delta)u=f\) 这样的方程。之后再问：当输入只有 \(p\) 次可积性时，频率幅值能落在哪个空间？

\ProseParagraphBreak

从可积函数出发得到的反演公式，怎样适用于仅有有限能量的信号？可以先选容易计算的函数逼近信号，再用能量等式控制变换的误差。这个构造将积分公式与平方范数极限联系起来。

\ProseParagraphBreak

\secref{b2:sec:0401}建立能量等式，\secref{b2:sec:0402}再用复参数把可积端点与平方端点连接起来。\secref{b2:sec:0403}说明所得结论的适用边界。Sheldon Axler 的《Measure, Integration \& Real Analysis》\cite[Section 11C, pp. 375--377]{Axler2020Measure}给出了 Plancherel 延拓的另一种讲法；更系统的 \(L^p\) Fourier 理论见 Loukas Grafakos 的《Classical Fourier Analysis》\cite{Grafakos2014Classical}。

% !TeX root = ../../Book2.tex
\section{Plancherel 定理}
\label{b2:sec:0401}

上一章把 Fourier 变换定义为绝对收敛积分，因此要求 \(f\in L^1\)。平方可积函数未必可积，但其能量仍有限。本节先在两个空间的交集上证明能量守恒，再利用完备性延拓变换。延拓所得对象是一个 \(L^2\) 等价类，不能自动解释为每个频率处的普通积分。

\subsection{\texorpdfstring{\(L^1\cap L^2\)}{L¹∩L²} 上的平方范数恒等式}
\label{b2:subsec:040101}

先在同时具有可积变换的函数上计算。

\ProseParagraphBreak

有限区间上的正交展开把平方范数写成系数平方和。在整个空间上，离散的频率指标变为连续变量，所期待的等式是两个积分相等。这里需要先处理一个实际障碍：由 \(f\in L^1\cap L^2\) 只能立即知道 \(\widehat f\) 连续有界，不能立即知道它平方可积。以下第一步先把两侧的积分合法性建立起来，第二步才让正则化消失。

\begin{lemma}\label{b2:lem:04-parseval-core}
设 \(d\) 为正整数，\(f,g\in L^1(\mathbb R^d;\mathbb C)\cap L^2(\mathbb R^d;\mathbb C)\)，且 \(\widehat f\in L^1(\mathbb R^d)\)。则
\[
 \int f(x)\cdot\overline{g(x)}\,\mathrm{d}x
 =\int\widehat f(\xi)\cdot\overline{\widehat g(\xi)}\,\mathrm{d}\xi.
\]
\end{lemma}
\begin{proof}
左侧由 Cauchy--Schwarz 不等式绝对收敛。对 \(f\) 使用上一章的反演定理，使 \(f(x)=\int\widehat f(\xi)e^{2\pi\mathrm{i}x\cdot\xi}\,\mathrm{d}\xi\) 几乎处处。将此式乘以 \(\overline g(x)\) 再积分；重积分绝对值不超过 \(\|\widehat f\|_1\cdot\|g\|_1\)，故可交换次序。内层积分为 \(\overline{\widehat g(\xi)}\)。右侧绝对收敛还可由 \(\widehat f\in L^1\)、\(\|\widehat g\|_\infty\leq\|g\|_1\) 直接确认。
\end{proof}

\begin{theorem}[Plancherel，交空间形式]\label{b2:thm:04-plancherel-core}
设 \(d\) 为正整数。若 \(f\in L^1(\mathbb R^d;\mathbb C)\cap L^2(\mathbb R^d;\mathbb C)\)，则 \(\widehat f\in L^2(\mathbb R^d;\mathbb C)\)，且 \(\|\widehat f\|_2=\|f\|_2\)。
\end{theorem}
\begin{proof}
令 \(f_\varepsilon=G_\varepsilon*f\)，其中 \(G_\varepsilon(x)=\varepsilon^{-d}e^{-\pi|x|^2/\varepsilon^2}\)，积分为 \(1\)。Young 不等式分别给出 \(\|f_\varepsilon\|_1\leq\|f\|_1\)、\(\|f_\varepsilon\|_2\leq\|f\|_2\)；\secref{b2:sec:0304}的近似恒等定理给出 \(f_\varepsilon\to f\) 于 \(L^2\)。卷积定理给出
\[
 \widehat {f_\varepsilon}(\xi)
 =e^{-\pi\varepsilon^2|\xi|^2}\cdot\widehat f(\xi).
\]
右侧是有界函数与可积 Gauss 因子的乘积，所以属于 \(L^1\)，也属于 \(L^2\)。对 \(f_\varepsilon\) 使用前一引理，得到
\[
 \|f_\varepsilon\|_2^2
 =\int e^{-2\pi\varepsilon^2|\xi|^2}|\widehat f(\xi)|^2\,\mathrm{d}\xi.
\]
取任意递减到零的 \(\varepsilon_n\)。范数的反三角不等式使左侧趋于 \(\|f\|_2^2\)；右侧的非负被积函数逐点递增到 \(|\widehat f|^2\)，由单调收敛趋于 \(\int|\widehat f|^2\)，即使预先尚不知道此积分有限也可使用。左侧极限有限，故同时得到平方可积性和等式。由于目标等式不含所选序列，它适用于原来的每个 \(f\)。
\end{proof}

这称为\term[plancherel-dingli]{Plancherel 定理}[Plancherel's theorem]。Gauss 因子先使变换可积，反演与积分换序便可在可积控制下进行，最后再取正则化参数的极限。

\subsection{Plancherel 延拓}
\label{b2:subsec:040102}

\begin{theorem}\label{b2:thm:04-plancherel-extension}
设 \(d\) 为正整数。存在唯一复线性等距映射 \(\mathcal F_2:L^2(\mathbb R^d;\mathbb C)\to L^2(\mathbb R^d;\mathbb C)\)，在 \(L^1(\mathbb R^d;\mathbb C)\cap L^2(\mathbb R^d;\mathbb C)\) 上与可积 Fourier 变换一致。
\end{theorem}
\begin{proof}
对 \(f\in L^2\)，取 \(f_n=f\cdot\mathbf1_{\{|x|\leq n,\ |f|\leq n\}}\)。它有界、有限测度支集，故属于 \(L^1\cap L^2\)，并由控制收敛在 \(L^2\) 中趋于 \(f\)。Plancherel 等式说明 \((\widehat f_n)\) 是 Cauchy 列，定义 \(\mathcal F_2f\) 为其 \(L^2\) 极限。

若使用另一列 \(g_n\in L^1\cap L^2\) 在 \(L^2\) 中逼近 \(f\)，则
\[
 \|\widehat f_n-\widehat g_n\|_2
 =\|f_n-g_n\|_2
 \leq\|f_n-f\|_2+\|g_n-f\|_2\longrightarrow0.
\]
所以定义不依赖选择。特别当 \(f\) 本来属于 \(L^1\cap L^2\) 时，可以选恒定列 \(g_n=f\)，立刻得到新定义与原来的可积变换一致。

对 \(f,h\in L^2\) 分别取逼近列 \(f_n,h_n\)。列 \(af_n+bh_n\) 在交空间内并逼近 \(af+bh\)，因此积分变换的线性性经极限给出延拓的线性性。又有
\[
 \|\mathcal F_2f\|_2
 =\lim_n\|\widehat f_n\|_2
 =\lim_n\|f_n\|_2=\|f\|_2.
\]
任何连续延拓在稠密子集上相等，就在整个空间相等，故唯一。整个构造只使用 \(L^2\) 的完备性，并没有对不同频率逐点选取极限。
\end{proof}
\symidx{\(\mathcal F_2\)：平方可积函数的 Fourier 变换}

当不致混淆时仍记 \(\widehat f=\mathcal F_2f\)。例如一维函数 \((1+|x|)^{-3/4}\) 平方可积却不可积；对它不能用原来的 Lebesgue 积分定义变换，因为振荡因子模为 \(1\)，不改变绝对可积性。它的变换由上述范数极限定义。

\ProseParagraphBreak

若某个具体逼近同时在 \(L^1\) 和 \(L^2\) 中收敛，则两种极限也可直接核对。可积变换给出 \(\|\widehat f_n-\widehat f\|_\infty\leq\|f_n-f\|_1\)；延拓给出 \(\widehat f_n\to\mathcal F_2f\) 于 \(L^2\)。在任意有限球上，一致收敛意味着依测度收敛，而平方范数收敛也意味着依测度收敛，极限唯一。因此二者几乎处处相等。后面处理不同指数的交空间时，仍将使用这个识别步骤。

\subsection{酉性}
\label{b2:subsec:040103}

\begin{theorem}[Plancherel，酉形式]\label{b2:thm:04-unitary}
设 \(d\) 为正整数，\(\mathcal F_2:L^2(\mathbb R^d;\mathbb C)\to L^2(\mathbb R^d;\mathbb C)\) 是积分 \(\widehat f(\xi)=\int f(x)e^{-2\pi\mathrm ix\cdot\xi}\,\mathrm dx\) 从 \(L^1(\mathbb R^d;\mathbb C)\cap L^2(\mathbb R^d;\mathbb C)\) 的等距延拓。映射 \(\mathcal F_2\) 满射并保持内积，因而为\term[you-suanzi]{酉算子}[unitary operator]。若 \(Rf(x)=f(-x)\)，则
\[
 \mathcal F_2^2=R,\quad \mathcal F_2^4=I,
 \quad \mathcal F_2^{-1}=R\mathcal F_2=\mathcal F_2R.
\]
\end{theorem}
\begin{proof}
若 \(f\in C_c^\infty(\mathbb R^d;\mathbb C)\)，取整数 \(m>d\)，\cref{b2:prop:03-higher-derivative-decay}给出 \(\widehat f\) 在 \(|\xi|\geq1\) 上的 \(|\xi|^{-m}\) 衰减；在单位球内它又有界，所以 \(\widehat f\in L^1\)。它属于 \(L^2\) 则由已证的交空间 Plancherel 定理保证。因此可用反演得到 \(\mathcal F_2^2f=Rf\)。光滑紧支集函数在 \(L^2\) 中稠密，且两边都是连续算子，故等式推广到全部 \(L^2\)。于是 \(\mathcal F_2^4=R^2=I\)，特别每个 \(g\) 都是 \(\mathcal F_2(\mathcal F_2^3g)\)，所以满射。线性等距映射保持内积：将 \(f+g\)、\(f-g\)、\(f+\mathrm{i}g\)、\(f-\mathrm{i}g\) 的范数平方展开并作差，可分别恢复内积的实部和虚部。最后由平方等式及反射的交换关系得到逆算子公式；交换关系先由可积函数换元验证，再按稠密性延拓。
\end{proof}

在本书的归一化下，Fourier 变换没有丢失或放大总能量，但一般不保持某个点的值，也不保持函数的支集。酉性属于 Hilbert 空间结构，而不是逐点大小不变。

\subsection{\texorpdfstring{\(L^2\)}{L²} 反演与频率截断}
\label{b2:subsec:040104}

\begin{corollary}\label{b2:cor:04-l2-inversion}
设 \(d\) 为正整数，\(f\in L^2(\mathbb R^d;\mathbb C)\)，记 \(\widehat f=\mathcal F_2f\) 为其 Plancherel 变换。则
\[
 f=\lim_{R\to\infty}\int_{|\xi|\leq R}\widehat f(\xi)\cdot e^{2\pi\mathrm{i}x\cdot\xi}\,\mathrm{d}\xi
\]
在 \(L^2\) 范数中成立。当 \(\varepsilon\downarrow0\) 时，Gauss 正则化也有 \(\mathcal F_2^{-1}(e^{-\pi\varepsilon^2|\xi|^2}\widehat f)\to f\) 于 \(L^2\)。
\end{corollary}
\begin{proof}
记 \(h_R=\mathbf1_{\{|\xi|\leq R\}}\widehat f\)。Cauchy--Schwarz 给出
\[
 \|h_R\|_1\leq |B(0,R)|^{1/2}\|\widehat f\|_2,
\]
所以 \(h_R\in L^1\cap L^2\)。由延拓在交空间上的一致性和逆算子公式，\(\mathcal F_2^{-1}h_R\) 正是所写的正号积分；它有连续代表，但不要求 \(f\) 也预先连续。酉性给出
\[
 \|\mathcal F_2^{-1}h_R-f\|_2^2
 =\|h_R-\widehat f\|_2^2
 =\int_{|\xi|>R}|\widehat f(\xi)|^2\,\mathrm{d}\xi\longrightarrow0.
\]
Gauss 情形的乘子模不超过 \(1\)，并逐点趋于 \(1\)，故频率侧的控制收敛和酉性给出同一结论。
\end{proof}

任何递增穷尽空间的可测频率区域都可替代球；所需条件只是被舍弃部分的平方质量趋零。这个结论没有保证每个点上的截断积分收敛。对于随时间演化的方程，只要频率乘子有界，就可以先在 \(L^2\) 中定义演化，再另行研究点值和正则性。



\subsection{Sobolev 空间的频率刻画}
\label{b2:subsec:040105}

反演恢复的是平方范数意义下的函数。接下来进一步问：怎样从频率一侧判断它是否具有平方可积的弱导数？上一章对可积函数证明了“空间求导变成频率乘法”；这里须把这一对应建立在整个 \(L^2\) 空间上。下文用 \(\widehat f\) 表示由平方范数延拓得到的 Plancherel 变换。

\begin{theorem}[平方可积弱导数的频率刻画]\label{b2:thm:04-weak-derivative}
设 \(d\) 为正整数，\(j\in\{1,\ldots,d\}\)，\(f\in L^2(\mathbb R^d;\mathbb C)\)，并记 \(\widehat f=\mathcal F_2f\)。它具有平方可积的第 \(j\) 个弱导数，当且仅当 \(\xi_j\widehat f(\xi)\in L^2(\mathbb R^d)\)。此时
\[
 \mathcal F_2(\partial_jf)=2\pi\mathrm{i}\xi_j\widehat f,
 \quad\|\partial_jf\|_2=2\pi\|\xi_j\widehat f\|_2.
\]
\end{theorem}
\begin{proof}
先记下两种双线性积分等式，所有变换在此均取 \(L^2\) 意义：
\[
 \int\widehat u\,v=\int u\,\widehat v,
 \quad \int u\,v=\int\widehat u\,\check v,
 \quad \check v=\mathcal F_2^{-1}v.
\]
第一式先对 \(u,v\in L^1\cap L^2\) 由 Fubini 证明，交换重积分的上界为 \(\|u\|_1\|v\|_1\)。用交空间同时逼近两者，Cauchy--Schwarz 和等距性允许两侧取极限。将第一式中的 \(v\) 替换为 \(\check v\)，便得到第二式。这里没有复共轭，与弱导数的双线性测试约定一致；它不同于带有复共轭的内积等式。

若 \(g=\partial_jf\in L^2\)，弱导数测试式原先只用于 \(C_c^\infty\)。对 \(\psi\in C_c^\infty\)，矩公式给出 \(\widehat\psi\in C^\infty\)；又因 \(\psi\) 和 \(\xi_j\psi\) 都在 \(L^1\cap L^2\) 中，交空间 Plancherel 定理使 \(\widehat\psi\)、\(\partial_j\widehat\psi=\widehat{-2\pi\mathrm i\xi_j\psi}\) 都属于 \(L^2\)。取固定光滑截断 \(\chi\)，在单位球上等于 \(1\)、在半径 \(2\) 的球外为零，令 \(\chi_R(x)=\chi(x/R)\)。将 \(\chi_R\widehat\psi\) 代入测试式，得到
\[
 \int g\chi_R\widehat\psi
 =-\int f\chi_R\partial_j\widehat\psi
   -\int f(\partial_j\chi_R)\widehat\psi.
\]
前两项分别由 \(g,f\in L^2\) 和测试函数的平方可积性取极限。最后一项的绝对值至多为
\[
 R^{-1}\|\partial_j\chi\|_\infty\|f\|_2\|\widehat\psi\|_2,
\]
故趋于零。因此也可将 \(\widehat\psi\) 用作测试函数。由于直接对紧支集积分求导有 \(\partial_j\widehat\psi=\widehat{-2\pi\mathrm{i}\xi_j\psi}\)，双线性等式给出
\[
 \int\widehat g\,\psi
 =\int g\,\widehat\psi
 =-\int f\,\partial_j\widehat\psi
 =\int\widehat f\,(2\pi\mathrm{i}\xi_j\psi).
\]
两边的局部可积函数对全部紧支集光滑函数配对相同，故几乎处处相等。这个局部唯一性可由磨光核在 Lebesgue 点的逼近直接得出。于是必要条件成立。

反之，设 \(g=\mathcal F_2^{-1}(2\pi\mathrm{i}\xi_j\widehat f)\)。对 \(\varphi\in C_c^\infty\)，用双线性 Plancherel 等式与正号变换的求导公式，
\[
 \int g\varphi
 =\int2\pi\mathrm{i}\xi_j\widehat f\,\check\varphi
 =-\int\widehat f\,\check{\partial_j\varphi}
 =-\int f\partial_j\varphi.
\]
所以 \(g\) 是弱导数。范数式由酉性得到。
\end{proof}

设 \(m\) 为非负整数。整数阶 Sobolev 空间 \(H^m(\mathbb R^d;\mathbb C)=W^{m,2}(\mathbb R^d;\mathbb C)\) 要求函数及所有 \(|\alpha|\leq m\) 的弱导数都属于 \(L^2\)。对这些导数逐阶使用上一定理，便得到频率乘子 \((2\pi\mathrm i\xi)^\alpha\)。若 \(f\in H^m(\mathbb R^d;\mathbb C)\)，按第一卷定义 21.1.2 中每个多重指标计一次的范数，有
\[
 \|f\|_{H^m}^2
 =\int\left(\sum_{|\alpha|\leq m}(2\pi)^{2|\alpha|}\xi^{2\alpha}\right)|\widehat f(\xi)|^2\,\mathrm{d}\xi.
\]
括号中的多项式与 \((1+4\pi^2|\xi|^2)^m\) 可用只依赖 \(d,m\) 的正数互相控制，但高阶时一般不数值相等：后者的多项式展开带有组合系数。因此用后一权重定义的范数与既有范数等价，却不能静默地把等价写成相等。

\ProseParagraphBreak

这个等价可以直接从多项式系数看出。令 \(t_j=4\pi^2\xi_j^2\)，则第一种权重为 \(P_m(t)=\sum_{|\alpha|\leq m}t^\alpha\)。展开 \((1+t_1+\cdots+t_d)^m\) 后，每个 \(t^\alpha\) 的系数都是一个正整数，且只有有限多个。因此
\[
 P_m(t)\leq(1+t_1+\cdots+t_d)^m\leq C_{d,m}P_m(t)
 \quad(t_j\geq0).
\]
例如 \(d=m=2\) 时，\(P_2=1+t_1+t_2+t_1^2+t_1t_2+t_2^2\)，而另一权重的三个一次或混合项分别带有系数 \(2\)。这正对应了多重指标范数与展开 Laplace 算子的系数差别。

反过来，若 \(f\in L^2(\mathbb R^d;\mathbb C)\) 且 \((1+4\pi^2|\xi|^2)^{m/2}\widehat f\in L^2\)，上述多项式比较使每个 \((2\pi\mathrm i\xi)^\alpha\widehat f\)、\(|\alpha|\leq m\)，都属于 \(L^2\)。当 \(m\geq1\) 时，先用一阶定理构造一阶弱导数，再对已构造的导数重复；每一步需要的频率乘子都已在上述 \(L^2\) 控制之中，故逐阶得到全部高阶弱导数。\(m=0\) 时则只是原来的 \(L^2\) 条件。因此这个加权可积条件不仅是必要的，也是 \(f\in H^m\) 的充分条件。

\subsection{应用：常系数椭圆方程与差商}
\label{b2:subsec:040106}

从这里开始，利用已建立的 Fourier 变换与 Sobolev 空间的对应，通过频率乘子构造常系数椭圆方程的解。先看一个严格正的频率符号：除以它的形式计算，需要同时说明所得函数属于哪个空间、方程如何成立以及解为何唯一。

\begin{proposition}[常系数预解方程]\label{b2:prop:04-resolvent-model}
设 \(d\) 为正整数，\(a>0\)、\(f\in L^2(\mathbb R^d;\mathbb C)\)。方程 \((a-\Delta)u=f\) 有唯一的 \(H^2(\mathbb R^d;\mathbb C)\) 弱解，满足
\[
 \widehat u(\xi)=\frac{\widehat f(\xi)}{a+4\pi^2|\xi|^2},
 \quad \|u\|_2\leq a^{-1}\|f\|_2,
 \quad \|\Delta u\|_2\leq\|f\|_2.
\]
另有 \(\|\nabla u\|_2\leq(2\sqrt a)^{-1}\|f\|_2\)，并有仅依赖 \(a,d\) 的 \(H^2\) 估计。
这里弱解按分布意义理解，即对每个 \(\varphi\in C_c^\infty(\mathbb R^d)\)，均有 \(\int u(a\varphi-\Delta\varphi)=\int f\varphi\)；对 \(H^2\) 解，这也等价于方程在 \(L^2\) 中成立。
\end{proposition}
\begin{proof}
分母至少为 \(a\)，所以右侧属于 \(L^2\)，可由酉逆变换定义 \(u\)。置 \(s=2\pi|\xi|\)。对于零、一、二阶乘子，分别有
\[
 \frac1{a+s^2}\leq\frac1a,
 \quad\frac{s}{a+s^2}\leq\frac1{2\sqrt a},
 \quad\frac{s^2}{a+s^2}\leq1.
\]
中间不等式来自 \(a+s^2\geq2\sqrt a\,s\)。每个 \((2\pi\xi)^\alpha/(a+s^2)\)、\(|\alpha|\leq2\)，都受这三个界之一控制。迭代弱导数刻画得到 \(u\in H^2\)，而平方积分给出所述估计；对梯度先把各分量的平方相加，恰得到分子 \(s^2\)。由 \(\widehat{\Delta u}=-4\pi^2|\xi|^2\widehat u\)，有
\[
 \mathcal F_2((a-\Delta)u)=\widehat f.
\]
酉变换的单射性使等式在 \(L^2\) 中成立，因而也在弱意义下成立。若另一个 \(H^2\) 解存在，对两解之差 \(v\) 取变换，得到 \((a+4\pi^2|\xi|^2)\widehat v=0\)。正分母使 \(\widehat v=0\)，再由单射性得到 \(v=0\)。
\end{proof}

假设 \(a>0\) 使分母在低频也有统一正下界，从而保证对每个 \(L^2\) 右端都能构造解。令 \(a\downarrow0\) 时，\(L^2\) 解算子范数界 \(a^{-1}\) 发散；仅凭 \(f\in L^2\) 不能把分母换成 \(4\pi^2|\xi|^2\) 并断言仍得到平方可积解。这不表示 \(a=0\) 时出现了非零的 \(L^2\) 齐次解：若 \(v\in H^2\) 满足 \(-\Delta v=0\)，则 \(\widehat v\) 在 \(\xi\ne0\) 处为零，而单点 \(\{0\}\) 的测度为零，故 \(v=0\)。这里失去的是对任意右端的有界求解保证。相比之下，二阶导数的界在这个极限中保持一致。区分解本身与其导数的估计，能看出低频与高频承担的不同作用。\chapref{b2:ch:05}将进一步把 \(a=1\) 时的分母解释成基本解的变换。

\ProseParagraphBreak

弱导数的频率刻画还把局部差商与高频尾部连接起来。差商可以直接从函数值计算，因此适合判断一个具体的平方可积函数是否具有有限的梯度能量。

\begin{proposition}[平方范数差商]\label{b2:prop:04-difference-quotients}
设 \(d\) 为正整数，\(f\in L^2(\mathbb R^d;\mathbb C)\)，\(j\in\{1,\ldots,d\}\)，\(e_j\) 为第 \(j\) 个坐标单位向量。对 \(h\in\mathbb R\setminus\{0\}\)，记
\[
 D_{j,h}f(x)\coloneqq\frac{f(x+he_j)-f(x)}h.
\]
若弱导数 \(\partial_jf\in L^2(\mathbb R^d)\)，则 \(\|D_{j,h}f\|_2\leq\|\partial_jf\|_2\)，并且当 \(h\to0\) 时，\(D_{j,h}f\to\partial_jf\) 于 \(L^2\)。反之，若存在 \(h_n\in\mathbb R\setminus\{0\}\)、\(h_n\to0\) 及 \(C\geq0\)，使 \(\sup_n\|D_{j,h_n}f\|_2\leq C\)，则 \(f\) 具有平方可积的第 \(j\) 个弱导数，且 \(\|\partial_jf\|_2\leq C\)。
\end{proposition}
\begin{proof}
平移公式先在 \(L^1\cap L^2\) 上由换元成立，再由平移的等距性和稠密性推广到 \(L^2\)。因此
\[
 \widehat {D_{j,h}f}(\xi)
 =\frac{e^{2\pi\mathrm{i}h\xi_j}-1}{h}\widehat f(\xi).
\]
由 \(e^{\mathrm{i}t}-1=\int_0^t\mathrm{i}e^{\mathrm{i}s}\,\mathrm{d}s\)，乘子的模不超过 \(2\pi|\xi_j|\)，且逐点趋于 \(2\pi\mathrm{i}\xi_j\)。若这个极限乘以 \(\widehat f\) 属于 \(L^2\)，酉性给出差商的范数界。差商与极限之差的平方由 \(4(2\pi\xi_j)^2|\widehat f|^2\) 控制，控制收敛使它在 \(L^2\) 中趋于零，再用弱导数定理识别极限。

反之，若给定序列的差商一致有界，Fatou 引理直接给出
\[
 \int(2\pi\xi_j)^2|\widehat f(\xi)|^2\,\mathrm{d}\xi
 \leq\liminf_n\int\left|\frac{e^{2\pi\mathrm{i}h_n\xi_j}-1}{h_n}\right|^2
       |\widehat f(\xi)|^2\,\mathrm{d}\xi\leq C^2.
\]
所以频率刻画的充分条件成立，给出弱导数和所述界。无需先从差商中提取弱收敛子列。
\end{proof}

在一维，帐篷函数 \(\tau(x)=(1-|x|)_+\) 虽在 \(0\) 处不可微，却属于 \(H^1\)。在 \((-1,0)\)、\((0,1)\) 分别分部积分，端点处由函数连续使边界项相消，得到弱导数在这两段分别为 \(1\)、\(-1\)，其余为零，平方范数为 \(2\)。它也等于 \(\mathbf1_{[-1/2,1/2]}*\mathbf1_{[-1/2,1/2]}\)，所以
\[
 \widehat\tau(\xi)=\left(\frac{\sin(\pi\xi)}{\pi\xi}\right)^2.
\]
于是 \(\xi\widehat\tau\in L^2\)，与弱导数计算一致。相邻斜率的跳跃允许一阶能量有限，但若把帐篷换成矩形 \(v=\mathbf1_{[-1/2,1/2]}\)，便有 \(\xi\widehat v=\sin(\pi\xi)/\pi\notin L^2\)。矩形的函数值跳跃阻止了平方可积的一阶弱导数。这个检验区分了函数的折角与函数本身的不连续。

\ProseParagraphBreak

对于正整数 \(d,m\)、\(f\in H^m(\mathbb R^d;\mathbb C)\) 和 \(R>0\)，高频截断还有定量误差。令 \(S_Rf\coloneqq\mathcal F_2^{-1}(\mathbf1_{\{|\xi|\leq R\}}\widehat f)\)。由 \(|\xi|\geq R\) 在尾部成立，
\[
 \|f-S_Rf\|_2^2
 \leq R^{-2m}\int|\xi|^{2m}|\widehat f(\xi)|^2\,\mathrm{d}\xi
 \leq C_{d,m}R^{-2m}\sum_{|\alpha|=m}\|\partial^\alpha f\|_2^2.
\]
最后一步是多项式 \((\xi_1^2+\cdots+\xi_d^2)^m\) 的展开与弱导数公式。因而每增加一阶平方可积导数，舍弃高频后的范数误差就得到额外一阶尺度衰减。这个估计不把弱导数转换成逐点光滑性；它描述的是总能量在高频中的分布。

% !TeX root = ../../Book2.tex
\section{Hausdorff--Young 不等式}
\label{b2:sec:0402}

变换已经有两个端点估计：可积函数变为有界函数，平方可积函数保持平方范数。中间指数的控制将通过一个连续参数把这两端连起来。调节函数幅值的幂，可以使这个参数在两端分别对应已知的范数；\secref{b2:sec:1704}还将把同一方法用于一般线性算子。

\subsection{已知端点与相容定义}
\label{b2:subsec:04-fourier-endpoints}
% 以下两个旧标签兼容曾分别介绍端点的旧稿；合并后均指本小节。
% 标签后缀不决定打印编号；新引用用语义标签，不把旧040202移给插值小节。
\label{b2:subsec:040201}
\label{b2:subsec:040202}

先把待连接的两个端点和它们的共同定义域列清楚。在本卷 \(e^{-2\pi\mathrm ix\cdot\xi}\) 的约定下，对 \(f\in L^1(\mathbb R^d)\)，积分三角不等式给出 \(\|\widehat f\|_\infty\leq\|f\|_1\)；对 \(f\in L^2(\mathbb R^d)\)，Plancherel 延拓给出 \(\|\mathcal F_2f\|_2=\|f\|_2\)。第一个端点的常数为 \(1\)，非负可积函数在零频率处取等；第二个端点则对每个平方可积函数保持范数。

\ProseParagraphBreak

这里的 \(L^1\) 端点只控制每个频率处的大小，不能给出整个 \(L^1\) 单位球的统一衰减率，见\cref{b2:prop:03-no-uniform-decay}。其局部解释是：把一个非负单位积分的 \(\varphi\in C_c^\infty(\mathbb R^d)\) 调制成 \(f_\eta(x)=e^{2\pi\mathrm ix\cdot\eta}\varphi(x)\)，并不改变 \(L^1\) 范数，却使 \(\widehat f_\eta(\eta)=1\)，而 \(|\eta|\) 可以任意大。这与\cref{b2:thm:03-riemann-lebesgue}对每个固定函数的趋零结论相容：这里允许输入随所测的频率改变。

\ProseParagraphBreak

连接两个端点还需要相容性：\cref{b2:thm:04-plancherel-extension}说明，对每个 \(f\in L^1\cap L^2\)，平方范数变换 \(\mathcal F_2f\) 与绝对积分定义的 \(\widehat f\) 几乎处处相同。因此下面是在改变同一个算子的输入指数。有限测度支集的简单函数同时属于这两个空间；对它们使用复参数改变各个非零取值的幅度，两个边界分别落入 \(L^1\) 与 \(L^2\)，便能在共同定义域内完成插值。所得估计之后再用于一般 \(L^p\) 函数的延拓。

\subsection{Riesz--Thorin 思想}
\label{b2:subsec:04-complex-interpolation}
% 保留旧标识的数学含义，不按当前小节显示序号重新分配。
\label{b2:subsec:040203}

先用三线引理控制复参数的变化。下面出现的整函数都是有限指数和，可以由处处收敛的幂级数直接理解。

\begin{lemma}[三线引理的有界形式]\label{b2:lem:04-three-lines}
设整函数 \(F:\mathbb C\to\mathbb C\) 在闭带 \(0\leq\operatorname{Re}z\leq1\) 上有界，并对每个 \(t\in\mathbb R\) 满足 \(|F(\mathrm{i}t)|\leq1\)、\(|F(1+\mathrm{i}t)|\leq1\)。则 \(|F(\theta)|\leq1\) 对 \(0\leq\theta\leq1\) 成立。
\end{lemma}
\begin{proof}
先回顾幂级数的最大模原理。若中心的最大模为零，函数已在邻域内为零；否则把函数乘以常数相位，使中心值为正实数 \(A\)。在足够小的圆上，幂级数一致绝对收敛，逐项平均合法，所得平均为中心值 \(A\)。圆上实部至多为模，模又至多为 \(A\)，所以连续非负函数 \(A-\operatorname{Re}F\) 的圆周积分为零，它处处为零。此时 \(|F|\leq A\) 还迫使虚部为零，故圆上函数恒为 \(A\)。将圆周值乘以各个负频率后逐项积分，可见所有正阶幂级数系数为零。函数在邻域内常值，并由重叠圆盘延拓到连通域。因此非恒定解析函数在紧矩形上的最大模出现在边界。

给定 \(\varepsilon>0\)，考虑 \(F(z)e^{\varepsilon z^2}\)。在两条竖边上，其模至多为 \(e^\varepsilon\)。若 \(|F|\leq M\)，在高度为 \(\pm R\) 的两条横边上，模至多为 \(Me^{\varepsilon(1-R^2)}\)，对充分大的 \(R\) 也不超过 \(e^\varepsilon\)。在矩形上用刚证明的最大模原理，得 \(|F(\theta)|e^{\varepsilon\theta^2}\leq e^\varepsilon\)。最后令 \(\varepsilon\downarrow0\)。
\end{proof}

这一估计是\term[sanxian-yinli]{三线引理}[three-lines lemma]的特殊形式。它把两条边界上的控制传递到中间，下面让边界对应两种不同的积分指数。

\begin{theorem}[Hausdorff--Young]\label{b2:thm:04-hausdorff-young}
设 \(d\) 为正整数，\(1\leq p\leq2\)，\(q=p'\)，即 \(1/p+1/q=1\)，其中 \(p=1\) 时 \(q=\infty\)。若 \(f\in L^1(\mathbb R^d;\mathbb C)\cap L^p(\mathbb R^d;\mathbb C)\)，则
\[
 \|\widehat f\|_q\leq\|f\|_p.
\]
\end{theorem}
\begin{proof}
端点已证，以下设 \(1<p<2\)，于是 \(2<q<\infty\) 且 \(q'=p\)。先设 \(f\) 为有界有限测度支集的简单函数。零函数直接满足结论；否则除以其 \(L^p\) 范数，先证明 \(\|f\|_p=1\) 的情形。再取任意同类简单函数 \(g\)，满足 \(\|g\|_p=1\)。我们尚不知道 \(\widehat f\in L^q\)，所以先估计配对 \(\int\widehat f g\)，随后在有限球上从这些配对恢复局部 \(L^q\) 范数，最后才扩大到全空间。测试函数的指数是 \(p=q'\)，因此可以对 \(f\) 和 \(g\) 使用相同的变幂构造。

现在说明幂指数如何由目标决定。设幅度的指数取实系数仿射函数 \(a(z)=A+Bz\)，因而在非零取值处，复幂 \(|f|^{a(z)}\) 的模为 \(|f|^{a(\operatorname{Re}z)}\)。要使两条边界恰好使用 \(L^1\) 与 \(L^2\) 端点，并在某个实数 \(0<\theta<1\) 处恢复原函数，必须同时满足
\[
 \begin{aligned}
  \operatorname{Re}z=0:&\quad a(0)=p,
       &&\text{使 }|f_z|=|f|^p,\quad\|f_z\|_1=1;\\
  \operatorname{Re}z=1:&\quad 2a(1)=p,
       &&\text{使 }|f_z|^2=|f|^p,\quad\|f_z\|_2=1;\\
  z=\theta:&\quad a(\theta)=1,
       &&\text{使 }f_\theta=f.
 \end{aligned}
\]
前两式确定 \(A=p\)、\(B=-p/2\)，第三式再确定 \(\theta=2(1-1/p)\)，在 \(1<p<2\) 时确有 \(0<\theta<1\)。这三个条件共同确定了指数 \(a(z)=p(1-z/2)\)；\(g\) 的指数也由相同条件确定。

记非零复数 \(w\) 的相位为 \(\operatorname{phase}(w)=w/|w|\)，零处定义为零。令
\[
 f_z(x)=\operatorname{phase}(f(x))\,|f(x)|^{p(1-z/2)},
 \quad g_z(\xi)=\operatorname{phase}(g(\xi))\,|g(\xi)|^{p(1-z/2)},
\]
零值处仍取零，复幂使用正数 \(|f|\) 和 \(|g|\) 的实对数。对每个 \(z\)，这两个函数仍有界并分别支于原来的有限测度集合，因而下式的绝对值二重积分为 \(\|f_z\|_1\cdot\|g_z\|_1<\infty\)，可按 Fubini 写为 \(\int\widehat {f_z}g_z\)。定义
\[
 F(z)=\int\!\int e^{-2\pi\mathrm{i}x\cdot\xi}f_z(x)\cdot g_z(\xi)\,\mathrm{d}x\,\mathrm{d}\xi
\]
因为函数各只有有限个非零取值，\(F\) 是有限指数和。具体地，将非零取值写成 \(f=\sum_j a_j\mathbf1_{E_j}\)、\(g=\sum_\ell b_\ell\mathbf1_{H_\ell}\)，其中各集合在相应变量中互不相交。令
\[
 c_{j\ell}=\int_{E_j}\int_{H_\ell}e^{-2\pi\mathrm{i}x\cdot\xi}\,\mathrm{d}\xi\,\mathrm{d}x,
 \quad |c_{j\ell}|\leq |E_j|\cdot|H_\ell|.
\]
则 \(F(z)\) 是 \(c_{j\ell}\operatorname{phase}(a_j)\operatorname{phase}(b_\ell)\exp(p(1-z/2)(\log|a_j|+\log|b_\ell|))\) 的有限和。因对数均为实数，每一项沿带的虚方向只改变相位，该项的模仅取决于 \(\operatorname{Re}z\)；对有限项使用三角不等式，便给出 \(F\) 在整个闭带上的统一上界。这同时验证了整性与三线引理要求的全带有界性；仅仅知道两条边上的估计，还不能省略后一条件。

\cref{b2:fig:04-interpolation-strip}把三个指数条件放在同一条带上：虚方向只改变相位，左右边界的幅度分别恢复归一化的 \(L^1\) 与 \(L^2\) 范数。下面逐一验证边界上的配对估计，再用三线引理传到 \(z=\theta\)。

\begin{figure}[htbp]
\centering
\begin{tikzpicture}[x=1cm,y=1cm,font=\small]
 \fill[AnalysisBlue!7] (0,-1.1) rectangle (9,1.1);
 \draw[->,AnalysisGray] (-0.5,0)--(9.7,0) node[right] {\(\operatorname{Re}z\)};
 \draw[<->,thick,AnalysisBlue] (0,-1.2)--(0,1.2);
 \draw[<->,thick,AnalysisBlue] (9,-1.2)--(9,1.2);
 \draw[<->,dashed,AnalysisBlue] (4.5,-1.2)--(4.5,1.2);
 \node[above left] at (0,1.2) {\(\operatorname{Im}z\)};
 \node[below left] at (0,0) {\(0\)};
 \fill[AnalysisBlue] (4.5,0) circle (1.5pt);
 \node[below right] at (4.5,0) {\(\theta\)};
 \node[below right] at (9,0) {\(1\)};
 \node[align=center,anchor=south] at (0,1.65)
   {\(\|f_{\mathrm it}\|_1=\|g_{\mathrm it}\|_1=1\)\\\(L^1\longrightarrow L^\infty\)};
 \node[align=center,anchor=south] at (4.5,1.65)
   {\(f_\theta=f,\quad g_\theta=g\)\\\(\|f\|_p=\|g\|_p=1\)};
 \node[align=center,anchor=south] at (9,1.65)
   {\(\|f_{1+\mathrm it}\|_2=\|g_{1+\mathrm it}\|_2=1\)\\\(L^2\longrightarrow L^2\)};
 \node[below] at (0,-1.35) {\(|F(\mathrm it)|\leq1\)};
 \node[below] at (9,-1.35) {\(|F(1+\mathrm it)|\leq1\)};
 \draw[->,AnalysisBlue] (0.8,-0.65)--(3.35,-0.65);
 \draw[->,AnalysisBlue] (8.2,-0.65)--(5.65,-0.65);
 \node[fill=white,inner sep=3pt] at (4.5,-0.65) {\(|F(\theta)|\leq1\)};
\end{tikzpicture}
\caption[复插值带与端点范数]{复插值带与端点范数。这里只画出无限高带状区域的一段，\(\theta\) 的横向位置仅作示意；两条边界上的估计由三线引理传递到中间。}
\label{b2:fig:04-interpolation-strip}
\end{figure}

当 \(z=\mathrm{i}t\) 时，\(|f_z|=|f|^p\)，故 \(\|f_z\|_1=\int|f|^p=1\)，\(g_z\) 同理。于是
\[
 |F(\mathrm{i}t)|\leq\|\widehat {f_{\mathrm{i}t}}\|_\infty\cdot\|g_{\mathrm{i}t}\|_1\leq1.
\]
当 \(z=1+\mathrm{i}t\) 时，\(|f_z|^2=|f|^p\)，所以 \(\|f_z\|_2=\|g_z\|_2=1\)。利用交空间上的 Plancherel 等式，
\[
 |F(1+\mathrm{i}t)|\leq\|\widehat {f_{1+\mathrm{i}t}}\|_2\cdot\|g_{1+\mathrm{i}t}\|_2=1.
\]
三线引理给出 \(|F(\theta)|\leq1\)；而 \(p(1-\theta/2)=1\)，所以
\[
 \left|\int\widehat f(\xi)g(\xi)\,\mathrm{d}\xi\right|\leq1.
\]
下面从配对估计恢复范数。若直接写全空间 \(L^q\) 的对偶表示，就预先假设了正在证明的 \(\widehat f\in L^q\)。改在有限球上测试可以避开这一点：记 \(h=\widehat f\)，由 \(f\in L^1\) 已知 \(h\) 有界，故对每个 \(R>0\)，它自动属于 \(L^q(B_R)\)，即使尚不知道全空间的积分是否有限。令
\[
 A_R=\left(\int_{B_R}|h|^q\right)^{1/q}.
\]
若 \(A_R=0\) 则已有 \(A_R\leq1\)。否则，为使配对恰好恢复 \(A_R\)，按 Hölder 取等的形式选测试函数：其相位取 \(h\) 相位的共轭以消去振荡，幅度取 \(|h|^{q-1}\)，再除以使 \(L^p\) 范数为 \(1\) 的常数。具体取
\[
 g_R=\frac{\overline{\operatorname{phase}(h)}\,|h|^{q-1}\mathbf1_{B_R}}{A_R^{q-1}}.
\]
因为 \((q-1)p=q\)，直接计算得 \(\|g_R\|_p=1\)、\(\int hg_R=A_R\)。函数 \(g_R\) 有界并支于有限球，可以用同一球上的简单函数在 \(L^p\) 中逼近。此时 \(h\in L^q(B_R)\)，Hölder 不等式保证配对收敛；逼近函数的范数也趋于 \(1\)，归一化后可用已经证明的简单函数估计。因此 \(A_R\leq1\)。最后令 \(R\uparrow\infty\)，单调收敛给出 \(\int|h|^q\leq1\)。去掉最初的归一化，便得简单函数的结论。

最后设 \(f\in L^1\cap L^p\)。先同时截断空间与幅值，再对截断函数量化取值，可得到简单函数 \(s_n\) 同时在 \(L^1\) 与 \(L^p\) 中趋于 \(f\)。第一个收敛使 \(\widehat s_n\) 一致趋于普通变换 \(\widehat f\)，第二个使 \(\|s_n\|_p\to\|f\|_p\)。Fatou 引理遂给出
\[
 \int|\widehat f|^q\leq\liminf_n\int|\widehat s_n|^q
 \leq\lim_n\|s_n\|_p^q=\|f\|_p^q.
\]
因此不等式对每个 \(f\in L^1\cap L^p\) 成立。
\end{proof}

上述证明体现\term[riesz-thorin-chazhi]{Riesz--Thorin 插值}[Riesz--Thorin interpolation]的方法：变动的是函数幅值的幂，使两个边界恰好落在已知空间。结果称为\term[hausdorff-young-budengshi]{Hausdorff--Young 不等式}[Hausdorff--Young inequality]。这里的估计取常数 \(1\)；最佳常数问题还要比较所有非零输入的范数比值。

\subsection{\texorpdfstring{\(1\leq p\leq2\)}{1≤p≤2}}
\label{b2:subsec:04-lp-fourier-extension}
\label{b2:subsec:040204}

\begin{corollary}\label{b2:cor:04-lp-extension}
设 \(d\) 为正整数，\(1<p<2\)，\(p'=p/(p-1)\)，各函数空间均为 \(\mathbb R^d\) 上的复值 Lebesgue 空间。Fourier 变换唯一延拓为连续线性映射 \(\mathcal F_p:L^p(\mathbb R^d;\mathbb C)\to L^{p'}(\mathbb R^d;\mathbb C)\)，且 \(\|\mathcal F_pf\|_{p'}\leq\|f\|_p\)。它在 \(L^1\cap L^p\) 上与可积变换一致，在 \(L^2\cap L^p\) 上与 Plancherel 变换一致。
\end{corollary}
\begin{proof}
取有界有限测度支集的简单函数 \(s_n\to f\) 于 \(L^p\)。前一定理作用于差，给出
\[
 \|\widehat s_n-\widehat s_m\|_{p'}\leq\|s_n-s_m\|_p.
\]
故变换是 \(L^{p'}\) Cauchy 列，以其极限定义 \(\mathcal F_pf\)。两个逼近列的差趋于零，使定义独立于选择；对线性组合选择相应逼近列，再取范数极限，得到线性性及范数界。稠密性保证连续延拓唯一。

若 \(f\in L^p\cap L^r\)、\(r=1\) 或 \(2\)，\cref{b2:prop:03-common-smooth-density}提供同时在两个范数中逼近的 \(C_c^\infty\) 列 \(v_n\)。它们的变换在 \(L^{p'}\) 中趋于 \(\mathcal F_pf\)，在 \(r=1\) 时一致趋于可积变换，在 \(r=2\) 时以平方范数趋于 \(\mathcal F_2f\)。对任意有限球 \(B\)，Chebyshev 不等式使每种有限指数范数收敛都推出依测度收敛，例如
\[
 |\{\xi\in B:|\widehat v_n-\mathcal F_pf|>\delta\}|
 \leq\delta^{-p'}\|\widehat v_n-\mathcal F_pf\|_{p'}^{p'}\longrightarrow0.
\]
若两个极限分别记为 \(H_1,H_2\)，则集合 \(\{|H_1-H_2|>2\delta\}\cap B\) 包含于相应两个误差集合的并，因而测度为零。令 \(\delta\downarrow0\)，再让有限球穷尽空间，得到几乎处处相等。这完成了交空间的一致性核对。
\end{proof}

因此在共同定义域上可以继续使用 \(\widehat f\) 这个记号。若 \(1<p<2\)，变换属于 \(L^{p'}\)，但 \(p'>2\)，不能简单地再对变换使用同一个 Hausdorff--Young 定理来反演；反演的正确表达要看它是否另外满足可积性、平方可积性，或采用后面的分布框架。



\ProseParagraphBreak

一个新定义的必要性，可以从同时不属于两个旧空间的函数看出。固定 \(1<p<2\)，选
\[
 d/2<\alpha<d/p<\beta<d,
\]
并取 \(f(x)=|x|^{-\alpha}\mathbf1_{\{|x|<1\}}+(1+|x|)^{-\beta}\mathbf1_{\{|x|\geq1\}}\)。原点附近由 \(\alpha p<d\) 保证 \(p\) 次可积，无穷远由 \(\beta p>d\) 保证，所以 \(f\in L^p\)。但原点奇性使它不平方可积，无穷远尾部使它不可积。因而它既不能整体代入\chapref{b2:ch:03}的绝对积分，也不能整体代入 Plancherel 定义；Hausdorff--Young 延拓确实增加了对象。

\ProseParagraphBreak

具体的积分检验是 \(\int_0^1 r^{d-1-\alpha p}\,\mathrm{d}r<\infty\) 与 \(\int_1^\infty r^{d-1-\beta p}\,\mathrm{d}r<\infty\)。把指数分别换成 \(2\) 与 \(1\)，第一个积分因 \(2\alpha>d\) 发散，第二个因 \(\beta<d\) 发散。例如一维取 \(p=3/2\)、\(\alpha=3/5\)、\(\beta=4/5\)，就得到可直接检验的具体函数。这里原点与远处是两种不同的障碍，仅靠一次空间截断无法同时移除；逼近时需要空间截断和幅值截断共同作用。

\ProseParagraphBreak

定义这个变换时，只需取在 \(L^p\) 中逼近 \(f\) 的简单函数列。其频率像由不等式组成 \(L^{p'}\) Cauchy 列，不必事先猜测极限的点值公式。若另外将 \(f\) 拆成可积部分与平方可积部分，两个变换的和也应与新定义相容；这种相容性将在下一章的测试配对中统一表达。

% !TeX root = ../../Book2.tex
\section{\texorpdfstring{\(L^p\)}{Lᵖ} Fourier 变换的局限}
\label{b2:sec:0403}

把指数从 \(2\) 往两侧移动，会改变变换所能取值的空间。下面用尺度和相互分离的波包解释为什么 \(p>2\) 不能照搬上一节的结论，并为分布意义下的变换留下明确的问题。

\subsection{\texorpdfstring{\(p>2\)}{p>2} 时的有界延拓障碍}
\label{b2:subsec:040301}

先作尺度检验。假设存在 \(\|\widehat f\|_q\leq C\|f\|_p\) 对全部光滑紧支集函数成立。固定非零 \(\varphi\in C_c^\infty\)，其变换快速下降，故所需的各个范数均有限且非零。对 \(f_a(x)=\varphi(ax)\)、\(a>0\)，伸缩公式给出
\[
 \|f_a\|_p=a^{-d/p}\|\varphi\|_p,
 \quad\|\widehat {f_a}\|_q=a^{-d+d/q}\|\widehat\varphi\|_q,
\]
无穷指数按倒数为零解释。将两式代入假定的估计并消去固定正数，得到
\[
 a^{d(1/p+1/q-1)}\leq C_\varphi\quad(a>0).
\]
若幂为正，让 \(a\to\infty\)；若幂为负，让 \(a\downarrow0\)，均产生矛盾。因此必须有 \(1/p+1/q=1\)。尺度检验只筛去不可能的指数，并不证明剩下的指数可用。

\begin{theorem}\label{b2:thm:04-p-above-two}
设 \(d\) 为正整数，\(2<p\leq\infty\)、\(1/p+1/q=1\)，其中 \(1/\infty=0\)。采用 \(\widehat f(\xi)=\int_{\mathbb R^d}f(x)e^{-2\pi\mathrm{i}x\cdot\xi}\,\mathrm{d}x\) 的归一化，则不存在与 \(f\) 无关的有限常数 \(C\)，使 \(\|\widehat f\|_q\leq C\|f\|_p\) 对所有复值函数 \(f\in C_c^\infty(\mathbb R^d;\mathbb C)\) 成立。
\end{theorem}
\begin{proof}
固定非零 \(\varphi\in C_c^\infty\)，选 \(L\) 足够大，使 \(\varphi(x-jLe_1)\) 的支集互不相交。对符号 \(\epsilon_j\in\{-1,1\}\)，令
\[
 f_N(x)=\sum_{j=1}^N\epsilon_j\varphi(x-jLe_1).
\]
于是 \(\|f_N\|_p=N^{1/p}\|\varphi\|_p\)，\(p=\infty\) 时取 \(N^{1/p}=1\)；其变换为 \(\widehat\varphi(\xi)\sum_{j=1}^N\epsilon_je^{-2\pi\mathrm{i}jL\xi_1}\)。

把全部 \(2^N\) 种符号等权平均。这里 \(\mathbb E\) 只是有限和除以 \(2^N\)，\(\mathbb P\) 只是所占符号组的比例。对任意模为 \(1\) 的复数 \(z_j\)，记 \(Z=\sum_j\epsilon_jz_j\)、\(X=|Z|^2\)。翻转某个符号会使含它奇数次的单项变号，因此这种单项的平均为零。于是
\[
 \mathbb E|Z|^2=\sum_{j,k}z_j\overline{z_k}\,\mathbb E(\epsilon_j\epsilon_k)=N.
\]
展开 \(|Z|^4\) 时，一个四重指标只有每个指标出现偶数次才能保留。它必属于三个成对划分之一：第一位置与第二、第一与第三、或第一与第四配对。每种划分至多有 \(N^2\) 个项，每项系数模为 \(1\)。全相等的指标可能被多次计入，但我们只求上界，故 \(\mathbb EX^2\leq3N^2\)。令 \(E=\{X\geq N/2\}\)，在其补集上 \(X<N/2\)，而在 \(E\) 上使用 Cauchy--Schwarz，得
\[
 N=\mathbb EX\leq N/2+(\mathbb EX^2)^{1/2}\mathbb P(E)^{1/2},
\]
因此 \(\mathbb P(E)\geq1/12\)，进而 \(\mathbb E|Z|^q\geq12^{-1}(N/2)^{q/2}\)。

逐频率取 \(z_j=e^{-2\pi\mathrm{i}jL\xi_1}\)，乘以 \(|\widehat\varphi(\xi)|^q\) 再积分。有限平均可与积分交换，得到
\[
 2^{-N}\sum_{\epsilon\in\{-1,1\}^N}\|\widehat {f_N}\|_q^q
 \geq12^{-1}(N/2)^{q/2}\|\widehat\varphi\|_q^q.
\]
至少一个被平均的非负数不小于平均值，故有一组选定符号使
\[
 \|\widehat {f_N}\|_q\geq c_qN^{1/2}\|\widehat\varphi\|_q.
\]
光滑紧支集函数的变换快速衰减且非零，故右侧范数有限而正。若所求估计成立，则 \(N^{1/2-1/p}\) 应一致有界，与 \(p>2\) 矛盾。
\end{proof}

这个反例甚至使用光滑紧支集函数，故障碍来自范数估计本身，并非某个积分未定义。它否定的是在整个 \(L^p\) 上取值于 \(L^{p'}\) 的有界延拓；并没有否定特殊函数仍可能有很好的变换。

\subsection{普通积分定义的局限}
\label{b2:subsec:040302}

在无限测度空间上，较大的 \(p\) 不意味着 \(L^1\)。例如在 \(\mathbb R^d\) 上，对 \(2<p<\infty\) 取 \(d/p<\alpha\leq d\)，则 \((1+|x|)^{-\alpha}\in L^p\setminus L^1\)；对 \(p=\infty\)，任取 \(0<\alpha\leq d\) 也得到有界但不可积的函数。乘上单位模复指数以后，绝对可积性没有改善，所以定义式不能作为普通 Lebesgue 积分使用。这并不等于某种指定的振荡截断必定发散，而是说普通积分尚未提供定义。

\ProseParagraphBreak

有时可以指定对称截断、Gauss 正则化或其他极限。不同极限的存在性各需证明；一般的 \(L^p\) 等价类也没有预先规定的点值。即使在 \(L^2\) 中，构造的是频率侧的等价类；一句“取每个频率的极限”并不能替代范数完备性给出的定义。

\subsection{\texorpdfstring{\(L^2\)}{L²} Fourier 乘子及其推广障碍}
\label{b2:subsec:040303}

\begin{definition}\label{b2:def:04-multiplier}
设 \(d\) 为正整数。对 \(m\in L^\infty(\mathbb R^d;\mathbb C)\)、\(f\in L^2(\mathbb R^d;\mathbb C)\)，以 \(2\pi\) 归一化的 Plancherel 变换 \(\widehat f=\mathcal F_2f\) 定义
\[
 T_mf\coloneqq\mathcal F_2^{-1}(m\widehat f),
\]
称为符号 \(m\) 的\term[fourier-chengzi-suanzi]{Fourier 乘子算子}[Fourier multiplier operator]。
\end{definition}
\symidx{\(T_m\)：符号为 \(m\) 的 Fourier 乘子算子}

\begin{proposition}\label{b2:prop:04-multiplier-norm}
设 \(d\) 为正整数，\(\mathcal F_2\) 为 \(2\pi\) 归一化的 Plancherel 变换，\(m\in L^\infty(\mathbb R^d;\mathbb C)\)，\(T_mf=\mathcal F_2^{-1}(m\mathcal F_2f)\) 定义于 \(L^2(\mathbb R^d;\mathbb C)\)。有 \(\|T_m\|_{L^2\to L^2}=\|m\|_\infty\)。若 \(K\in L^1(\mathbb R^d;\mathbb C)\)、\(m=\widehat K\) 几乎处处成立，则 \(T_mf=K*f\) 对每个 \(f\in L^2(\mathbb R^d;\mathbb C)\) 成立，等式按 \(L^2\) 等价类理解；卷积公式还定义了每个 \(L^p(\mathbb R^d;\mathbb C)\)、\(1\leq p\leq\infty\) 上的有界算子，范数至多为 \(\|K\|_1\)。
\end{proposition}
\begin{proof}
酉性给出上界。当 \(\|m\|_\infty=0\) 时两边均为零；否则，对任意 \(0<a<\|m\|_\infty\)，集合 \(\{|m|>a\}\) 有正测度。用有限球穷尽空间，可在其中取一个有限正测度子集 \(E\)。由满射性可选 \(f\in L^2\) 使 \(\widehat f=\mathbf1_E\)，于是 \(\|T_mf\|_2>a\|f\|_2\)，令 \(a\uparrow\|m\|_\infty\) 得到等号。卷积表示先对 \(f\in L^1\cap L^2\) 由\cref{b2:thm:03-convolution-transform}成立：Young 不等式保证 \(K*f\) 同时属于 \(L^1\) 与 \(L^2\)，从而可积变换与 Plancherel 变换的一致性可用于两侧。再取 \(f_n\in L^1\cap L^2\) 在 \(L^2\) 中逼近任意 \(f\)；Young 估计使 \(K*f_n\to K*f\)，已证乘子界使 \(T_mf_n\to T_mf\)，故等式推广到 \(L^2\)。其余指数直接由\cref{b2:thm:03-young}给出。
\end{proof}

对 \(p\ne2\)，仅有 \(m\) 有界通常不足以取得同样的乘子结论；符号的平滑程度、奇性与维数会参与估计。\chapref{b2:ch:22}将研究这些性质怎样决定乘子的有界指数范围。


\ProseParagraphBreak

平滑频率截断说明为何符号的导数有用。取 \(m\in C_c^\infty(\mathbb R^d;\mathbb C)\)，令 \(K=\check m\)。\cref{b2:prop:03-higher-derivative-decay}将可积的高阶纯坐标导数转成频率衰减；选取大于维数的阶数，可得 \(K\in L^1\)，所以 \(T_m\) 在每个 \(L^p\) 上都由 Young 不等式控制。对 \(R>0\) 令 \(m_R(\xi)=m(\xi/R)\)，相应核为 \(K_R(x)=R^dK(Rx)\)，且 \(\|K_R\|_1=\|K\|_1\)。于是这一族算子对全部尺度具有同一个范数界。

\begin{proposition}[平滑频率逼近]\label{b2:prop:04-smooth-frequency-approximation}
设 \(d\) 为正整数，\(m\in C_c^\infty(\mathbb R^d;\mathbb C)\)、\(m(0)=1\)，按正号 Fourier 积分定义 \(K=\check m\)，并对每个 \(R>0\) 令 \(K_R(x)=R^dK(Rx)\)。对 \(1\leq p<\infty\)、\(f\in L^p(\mathbb R^d;\mathbb C)\)，当 \(R\to\infty\) 时有 \(K_R*f\to f\) 于 \(L^p\)；若复值函数 \(f\) 在 \(\mathbb R^d\) 上有界且一致连续，则收敛为一致收敛。对 \(f\in L^2(\mathbb R^d;\mathbb C)\)，这个卷积等于 \(T_{m(\cdot/R)}f\)。
\end{proposition}
\begin{proof}
先核对核的可积性。取整数 \(N>d\)，则 \(m\) 及各 \(\partial_j^Nm\) 都属于 \(L^1\)；由\cref{b2:prop:03-higher-derivative-decay}以及 \(\check m(x)=\widehat m(-x)\)，有 \(|K(x)|\leq C|x|^{-N}\) 对 \(|x|\geq1\) 成立，而在单位球内 \(|K|\leq\|m\|_1\)。两部分合起来给出 \(K\in L^1\)。对光滑紧支集的 \(m\) 使用上一章的反演，再作反射换元，得到 \(\widehat K=m\)；这里两边连续，所以等式处处成立。特别 \(\int K=\widehat K(0)=m(0)=1\)。注意 \(K\) 可以有复值且不必非负，所需的是绝对值的统一质量界与尾部集中性。将卷积换元，得到
\[
 K_R*f-f=\int K(y)\bigl(f(\cdot-y/R)-f\bigr)\,\mathrm{d}y.
\]
因而由积分型 Minkowski 不等式，
\[
 \|K_R*f-f\|_p
 \leq\int|K(y)|\,\|f(\cdot-y/R)-f\|_p\,\mathrm{d}y.
\]
对每个固定 \(y\)，有限指数的平移连续性使内层范数趋于零；它又不超过 \(2\|f\|_p\)。控制收敛给出结论。若 \(f\) 有界一致连续，用同一论证把 \(L^p\) 范数换成上确界范数即可。\(L^2\) 上的乘子识别来自前一命题，因为 \(\widehat {K_R}(\xi)=m(\xi/R)\)。
\end{proof}

这个证明也显示为何不能把结论无条件推广到 \(L^\infty\) 范数。在一维取阶跃函数 \(f=\mathbf1_{(0,\infty)}\)，它与光滑核卷积后具有连续代表。若一个连续函数与 \(f\) 的本性上确界距离小于 \(1/2\)，可取统一的 \(c<1/2\)，使它在左侧几乎处处距 \(0\) 至多为 \(c\)，在右侧几乎处处距 \(1\) 至多为 \(c\)。连续性使这些闭圆盘约束在原点也同时成立，但两圆盘不相交。因此 \(\|K_R*f-f\|_\infty\geq1/2\)；这里的收敛限制由函数的跳跃造成，尽管每个算子的 \(L^\infty\) 范数都统一有界。

\ProseParagraphBreak

相反，对 \(R>0\)，一维的硬截断符号 \(m=\mathbf1_{[-R,R]}\) 对应
\[
 K(x)=\dfrac{\sin(2\pi Rx)}{\pi x},
\]
该核不属于 \(L^1\)。对 \(f\in C_c^\infty\)，仍可用有限频率积分和 Fubini 验证 \(T_mf=K*f\)，但不能再用 \(\|K\|_1\) 给出有限的 Young 上界。\(L^2\) 上的范数仍为 \(1\)，因为符号是有界的正交投影。这个例子把两种方法的职责分开：酉性直接控制平方能量，绝对可积核则提供全部指数上的一种充分条件。

\subsection{分布意义的 Fourier 变换的动机}
\label{b2:subsec:040304}

若不能要求变换是一个普通函数，还可以要求它对足够好的测试函数产生确定的数值。对 \(f\in L^p\) 和光滑紧支集的 \(\varphi\)，函数 \(\widehat\varphi\) 快速衰减，因而 \(\int f\widehat\varphi\) 总有意义。若 \(f\) 原来可积，Fubini 又给出
\[
 \int\widehat f(\xi)\varphi(\xi)\,\mathrm{d}\xi
 =\int f(x)\widehat\varphi(x)\,\mathrm{d}x.
\]
右侧因此提供了延拓变换的线索，且对每个 \(1\leq p\leq\infty\) 都适用。下一章将选取在 Fourier 变换下保持不变的 Schwartz 测试空间，使这个配对成为完整定义。

\ProseParagraphBreak

测试函数的下降速度在这里承担了可积性责任。例如取整数 \(N\) 满足 \(Np'>d\)，当 \(p'=\infty\) 时只需取任意正整数。则 Hölder 不等式给出
\[
 \left|\int f\widehat\varphi\right|
 \leq C_{N,p,d}\|f\|_p
 \sup_x(1+|x|)^N|\widehat\varphi(x)|.
\]
其中 \(p'=1\) 也由 \(N>d\) 包含在内。反复分部积分把右侧的加权上确界控制在 \(\varphi\) 的有限多个导数与紧支集大小之内。这样一个可能没有可积频率代表的对象，仍能对每个测试函数给出稳定的线性数值；下一章要精确说明的是这种稳定性应在哪个测试函数拓扑中表述。


\begin{chaptersummary}
Plancherel 等式把变换从绝对收敛积分延拓为能量空间上的酉算子。反演和四次变换等于恒等的关系保证满射性；频率截断的误差则由被舍弃的平方质量精确给出。Hausdorff--Young 把两个端点连接起来，但其共轭指数方向只覆盖 \(1\leq p\leq2\)。

\ProseParagraphBreak

大于二时，空间中相互分离的函数可在频率中产生更大的叠加，破坏所需范数界。应当区分某个好函数的变换存在、变换在整个空间上连续，以及逐点反演这三个层次。\(L^2\) 乘子由符号本性上确界直接控制。一般指数的乘子问题还需另作研究，而全部 \(L^p\) 函数的变换将在温和分布中统一定义。
\end{chaptersummary}

\BookExercises

本组题目以 Plancherel 为计算工具，分别考察尺度、频带、乘子、算子收敛及空间与频率的同时集中。\Cref{b2:ex:04-four-sectors}研究 Fourier 算子的四个特征空间，为下一章 Hermite 函数作为 Fourier 特征函数的讨论提供算子观点。\Cref{b2:ex:04-support-volume}先建立空间投影与频率投影的复合估计，\cref{b2:ex:04-approximate-support}再将它用于有误差的集中；最后的\cref{b2:ex:04-interpolation-sign-moments}直接训练 Hausdorff--Young 延拓、尺度关系与符号平均。

% 自拟：Gauss范数检查非最优Hausdorff--Young常数及尺度不变性。
% 预留题头及首个公式所需空间，避免超链接锚点留在前页而题头移至后页。
\Needspace{6\baselineskip}
\begin{exercise}\label{b2:ex:04-gaussian-ratio}
设 \(d\) 为正整数，\(1<p<2\)、\(q=p/(p-1)\)，并对 \(a>0\) 定义实值函数 \(G_a(x)=e^{-\pi a|x|^2}\)（\(x\in\mathbb R^d\)）。计算 \(\|G_a\|_{L^p(\mathbb R^d)}\)、\(\|\widehat G_a\|_{L^q(\mathbb R^d)}\) 和二者的比值，说明为什么只改变 \(a\) 不能改善这个比值。该计算是否已经证明最佳 Hausdorff--Young 常数？
\end{exercise}
\begin{solution}
尺度公式给出
\[
 \|G_a\|_p=(ap)^{-d/(2p)},\quad
 \widehat G_a(\xi)=a^{-d/2}e^{-\pi|\xi|^2/a},\quad
 \|\widehat G_a\|_q=a^{-d/2}(q/a)^{-d/(2q)}.
\]
因 \(1/p+1/q=1\)，比值为 \(p^{d/(2p)}q^{-d/(2q)}\)，与 \(a\) 无关。这个数是任何普遍常数的下界；只计算一族函数，不能证明所有函数都满足取此常数的上界，因而不能单凭此计算宣称最佳性。
\end{solution}

% 自拟：紧频带产生光滑代表和定量导数界。
\Needspace{6\baselineskip}
\begin{exercise}\label{b2:ex:04-bandlimited}
设 \(d\) 为正整数，\(0<R<\infty\)、\(f\in L^2(\mathbb R^d;\mathbb C)\)，且 Plancherel 变换 \(\widehat f=\mathcal F_2f\) 几乎处处支撑在 \(B(0,R)\) 内。证明 \(f\) 有代表 \(h\in C^\infty(\mathbb R^d;\mathbb C)\)，并对每个多重指标 \(\alpha\in\mathbb N_0^d\) 有
\[
 \|\partial^\alpha h\|_\infty
 \leq(2\pi R)^{|\alpha|}|B(0,R)|^{1/2}\|f\|_2.
\]
\end{exercise}
\begin{solution}
有限频带上的 \(\widehat f\) 可积，乘任意多项式后仍可积。正号积分因此给出一个可任意次逐项求导的连续函数 \(h\)，且其 \(L^2\) 等价类由反演等于 \(f\)。求导后用 Cauchy--Schwarz：
\[
 |\partial^\alpha h(x)|
 \leq(2\pi R)^{|\alpha|}\int_{B(0,R)}|\widehat f|
 \leq(2\pi R)^{|\alpha|}|B(0,R)|^{1/2}\|\widehat f\|_2.
\]
最后用 Plancherel。这里的 \(R\) 有限是保证全部频率矩可积的关键。
\end{solution}

% 自拟：L2反演确实超出绝对可积函数域。
\Needspace{6\baselineskip}
\begin{exercise}\label{b2:ex:04-sinc}
在 \(\mathbb R\) 上使用 Lebesgue 测度，令 \(g(\xi)=\mathbf1_{[-1/2,1/2]}(\xi)\in L^2(\mathbb R;\mathbb C)\)，\(f=\mathcal F_2^{-1}g\)。求 \(f\) 的连续代表，证明 \(f\in L^2(\mathbb R;\mathbb C)\setminus L^1(\mathbb R;\mathbb C)\)，并计算 \(\|f\|_2\)。
\end{exercise}
\begin{solution}
有限频带积分给出 \(f(x)=\sin(\pi x)/(\pi x)\)，原点补为 \(1\)。Plancherel 给出 \(\|f\|_2=\|g\|_2=1\)。在每个 \([n+1/6,n+5/6]\) 上，\(|\sin(\pi x)|\geq1/2\)，因此其绝对积分至少为常数乘 \(\sum_{n\geq1}(n+1)^{-1}\)，故不属于 \(L^1\)。这不妨碍其 \(L^2\) 变换等于区间示性函数。
\end{solution}

% 自拟：平移群强连续却不按算子范数连续。
\Needspace{6\baselineskip}
\begin{exercise}\label{b2:ex:04-translation-operator}
设 \(d\) 为正整数，在 \(L^2(\mathbb R^d;\mathbb C)\) 上令 \((\tau_hf)(x)=f(x-h)\)（\(h\in\mathbb R^d\)），\(I\) 为恒等算子。证明对每个 \(f\in L^2(\mathbb R^d;\mathbb C)\)，\(\|\tau_hf-f\|_2\to0\) 当 \(h\to0\)；但对每个 \(h\ne0\)，\(\|\tau_h-I\|_{L^2\to L^2}=2\)。
\end{exercise}
\begin{solution}
平移的变换公式按稠密性推广到 \(L^2\)。于是
\[
 \|\tau_hf-f\|_2^2
 =\int|e^{-2\pi\mathrm{i}h\cdot\xi}-1|^2|\widehat f(\xi)|^2\,\mathrm{d}\xi\longrightarrow0
\]
由控制收敛成立。乘子范数则等于 \(e^{-2\pi\mathrm{i}h\cdot\xi}-1\) 的本性上确界；\(h\ne0\) 时，可在正测度频率带上使相位任意接近 \(-1\)，故该本性上确界为 \(2\)。前一结论固定函数，后一结论允许随 \(h\) 改变测试函数。
\end{solution}

% 自拟：频率投影和正性、伴随的分类。
\Needspace{6\baselineskip}
\begin{exercise}\label{b2:ex:04-multiplier-classification}
设 \(d\) 为正整数，\(m\in L^\infty(\mathbb R^d;\mathbb C)\)，并在 \(L^2(\mathbb R^d;\mathbb C)\) 上令 \(T_mf=\mathcal F_2^{-1}(m\cdot\mathcal F_2f)\)。求 \(T_m^*\)，并分别刻画 \(T_m\) 自伴、非负、酉和正交投影的充要条件。这里非负指 \(\langle T_mf,f\rangle\) 对每个 \(f\in L^2(\mathbb R^d;\mathbb C)\) 都是非负实数。
\end{exercise}
\begin{solution}
在频率侧计算内积，得 \(T_m^*=T_{\overline m}\)。所以自伴当且仅当 \(m\) 几乎处处为实数。非负当且仅当 \(m\geq0\) 几乎处处，充分性由积分给出。为证必要性，若虚部不为零，可在某个有限正测度集 \(E\) 上使它同号，且其绝对值有正下界；以 \(\widehat f=\mathbf1_E\) 测试，便使 \(\langle T_mf,f\rangle\) 的虚部非零。因此符号先须几乎处处为实数。若实部在正测度集合上为负，同样取有限正测度的负值水平集作测试，就得到负的二次型，矛盾。又因 \(T_m^*T_m=T_{|m|^2}\)，酉性等价于 \(|m|=1\) 几乎处处，此时逆符号为 \(\overline m\)。正交投影等价于自伴且幂等，故 \(m\) 几乎处处只能取 \(0,1\)，即为某个可测集的示性函数。
\end{solution}

% 自拟：傅里叶算子的四个代数谱分量；无须预借Hermite函数。
\Needspace{6\baselineskip}
\begin{optionalexercise}\label{b2:ex:04-four-sectors}
设 \(d\) 为正整数，\(H=L^2(\mathbb R^d;\mathbb C)\)，\(\mathcal F_2\colon H\to H\) 为 Plancherel 变换，\(I\) 为 \(H\) 上的恒等算子。对 \(\lambda\in\{1,-1,\mathrm{i},-\mathrm{i}\}\)，定义
\[
 Q_\lambda=\dfrac14\sum_{j=0}^3\lambda^{-j}\mathcal F_2^j.
\]
其中 \(\mathcal F_2^0=I\)。证明这些算子是两两正交的正交投影，其和为 \(I\)，且其值域等于 \(\{\,f\in H\mid\mathcal F_2f=\lambda f\,\}\)。说明 \(\lambda=\pm1\) 的值域只含几乎处处意义下的偶函数，\(\lambda=\pm\mathrm{i}\) 的值域只含几乎处处意义下的奇函数。
\end{optionalexercise}
\begin{solution}
利用 \(\mathcal F_2^4=I\) 移动求和指标，可得 \(\mathcal F_2Q_\lambda=\lambda Q_\lambda\)。若 \(\mathcal F_2f=\mu f\)，则 \(Q_\lambda f=\frac14\sum_{j=0}^3(\mu/\lambda)^jf\)，当 \(\mu=\lambda\) 时等于 \(f\)，否则为零。并且四个 \(Q_\lambda\) 的和为 \(I\)，因为四次单位根的非零次幂和为零。这已给出幂等性、相互乘积为零和值域描述。酉性使不同单位模特征值的特征向量正交：\(\langle f,g\rangle=\lambda\overline\mu\langle f,g\rangle\)。故这些投影确为正交投影。最后 \(Rf=\mathcal F_2^2f=\lambda^2f\)，得到奇偶性。
\end{solution}

% 自拟：空间投影与频率投影的复合范数估计，作为近似支集题的算子预备。
\Needspace{6\baselineskip}
\begin{exercise}\label{b2:ex:04-support-volume}
设 \(d\) 为正整数，\(A,B\subset\mathbb R^d\) 是 Lebesgue 可测集，且 \(|A|,|B|<\infty\)。在 \(L^2(\mathbb R^d;\mathbb C)\) 上定义
\[
 P_Af=\mathbf1_A\cdot f,
 \quad Q_Bf=\mathcal F_2^{-1}(\mathbf1_B\cdot\mathcal F_2f).
\]
证明 \(P_A,Q_B\) 是正交投影，并证明
\[
 \|Q_BP_A\|_{L^2\to L^2}\leq|A|^{1/2}\cdot|B|^{1/2}.
\]
说明：如果某个非零 \(f\in L^2(\mathbb R^d;\mathbb C)\) 同时满足 \(P_Af=f\)、\(Q_Bf=f\)，那么这一初等估计只能给出必要条件 \(|A|\cdot|B|\geq1\)，并未构造这样的函数。
\end{exercise}
\begin{solution}
乘以示性函数 \(\mathbf1_A\) 的算子自伴且幂等，所以 \(P_A\) 是正交投影。频率侧乘以 \(\mathbf1_B\) 同样是正交投影；与酉算子 \(\mathcal F_2\) 共轭便得到 \(Q_B\)。因此两者的算子范数都不超过 \(1\)。对任意 \(u\in L^2(\mathbb R^d;\mathbb C)\)，Cauchy--Schwarz 给出
\[
 \|P_Au\|_1\leq|A|^{1/2}\cdot\|P_Au\|_2
 \leq|A|^{1/2}\cdot\|u\|_2.
\]
所以 \(P_Au\in L^1\cap L^2\)，其普通 Fourier 变换与 Plancherel 变换几乎处处相等。由酉性及 \(L^1\to L^\infty\) 估计，
\[
 \begin{aligned}
 \|Q_BP_Au\|_2
 &=\|\mathbf1_B\cdot\widehat{P_Au}\|_2\\
 &\leq|B|^{1/2}\cdot\|\widehat{P_Au}\|_\infty\\
 &\leq|B|^{1/2}\cdot\|P_Au\|_1
 \leq|A|^{1/2}\cdot|B|^{1/2}\cdot\|u\|_2.
 \end{aligned}
\]
对全部单位向量取上确界即得算子范数估计；\(|A|=0\) 或 \(|B|=0\) 的情形也包含在内。若另有题目末尾的非零 \(f\)，则 \(Q_BP_Af=f\)，上式成为 \(\|f\|_2\leq|A|^{1/2}\cdot|B|^{1/2}\cdot\|f\|_2\)，除以正数 \(\|f\|_2\) 即得所述条件。
\end{solution}

作为书外补充，Benedicks 的《The support of functions and distributions with a spectral gap》\cite[p. 286]{Benedicks1984SpectralGap}中 Theorem 1 给出更深的结论：对任意正整数 \(d\)，若 \(f\in L^2(\mathbb R^d;\mathbb C)\) 及 \(\mathcal F_2f\) 分别几乎处处支撑于有限测度可测集，则 \(f=0\) 几乎处处。该定理的证明见所引论文。

% 自拟：把空间、频率投影估计用于有误差的同时集中，假设可由实际函数实现。
\Needspace{6\baselineskip}
\begin{optionalexercise}\label{b2:ex:04-approximate-support}
设 \(d\) 为正整数，\(f\in L^2(\mathbb R^d;\mathbb C)\)、\(\|f\|_2=1\)，\(\widehat f=\mathcal F_2f\)，\(A,B\subset\mathbb R^d\) 是有限 Lebesgue 测度的可测集。设 \(\varepsilon,\delta\geq0\)、\(\varepsilon+\delta<1\)，且
\[
 \|f\cdot\mathbf1_{A^c}\|_2\leq\varepsilon,
 \quad\|\widehat f\cdot\mathbf1_{B^c}\|_2\leq\delta.
\]
证明 \(|A|^{1/2}\cdot|B|^{1/2}\geq1-\varepsilon-\delta\)。
\end{optionalexercise}
\begin{solution}
令 \(P_Af=\mathbf1_A\cdot f\)、\(Q_Bf=\mathcal F_2^{-1}(\mathbf1_B\cdot\widehat f)\)。两者均为正交投影，故
\[
 \|f-Q_BP_Af\|_2
 \leq\|f-Q_Bf\|_2+\|Q_B(f-P_Af)\|_2\leq\delta+\varepsilon.
\]
\Cref{b2:ex:04-support-volume}已经对任意输入证明了复合算子范数的上界。对当前单位向量使用它，得到
\[
 \|Q_BP_Af\|_2
 \leq\|Q_BP_A\|_{L^2\to L^2}\cdot\|f\|_2
 \leq|A|^{1/2}\cdot|B|^{1/2}.
\]
由三角不等式，\(1\leq\varepsilon+\delta+\|Q_BP_Af\|_2\)，得到所求。

这些条件确能由非零函数实现。例如取归一化的非零 \(f\in C_c^\infty(\mathbb R^d;\mathbb C)\)，以一个有界集 \(A\) 包含其支集，空间误差就是零。又因 \(\widehat f\in L^2\)，可取足够大的有限半径球 \(B\)，使频率侧尾部的平方积分任意小。这说明近似集中允许非零对象，而无需要求两个精确支集都有限。
\end{solution}

% 自拟：具体非Gauss尾函数的Hausdorff--Young延拓与尺度；加权符号和的精确四矩及矩插值，不复现正文波包反例。
\Needspace{6\baselineskip}
\begin{exercise}\label{b2:ex:04-interpolation-sign-moments}
\begin{enumerate}
\item\label{b2:item:04-hy-nongaussian-scaling} 在一维令 \(h(x)=(1+|x|)^{-1}\)，\(1<p<2\)、\(q=p/(p-1)\)。证明 \(h\in L^p(\mathbb R;\mathbb C)\setminus L^1(\mathbb R;\mathbb C)\)。记 \(H=\mathcal F_ph\in L^q(\mathbb R;\mathbb C)\) 为 Hausdorff--Young 延拓，证明
\[
 \|H\|_q\leq\left(\frac2{p-1}\right)^{1/p}.
\]
对 \(a>0\) 令 \(h_a(x)=h(ax)\)，在 \(L^q\) 意义下证明 \(\mathcal F_ph_a(\xi)=a^{-1}H(\xi/a)\)，并核对输入、输出范数具有同一个尺度因子 \(a^{-1/p}\)。不能把 \(h\) 直接代入绝对收敛 Fourier 积分。
\item\label{b2:item:04-weighted-sign-moments} （选做）设 \(N,d\) 为正整数，\(c_1,\ldots,c_N\in\mathbb C\) 不全为零，\(S=\sum_{j=1}^N|c_j|^2\)。在全部 \(2^N\) 组 \(\epsilon=(\epsilon_1,\ldots,\epsilon_N)\in\{-1,1\}^N\) 上取等权平均 \(\mathbb E\)，令 \(Z=\sum_j\epsilon_jc_j\)。求 \(\mathbb E|Z|^2\)、\(\mathbb E|Z|^4\)，并用 Hölder 不等式证明
\[
 \mathbb E|Z|\geq\sqrt{S/3}.
\]
进一步，给定任意 \(t_1,\ldots,t_N\in\mathbb R^d\) 及非负 \(w\in L^1(\mathbb R^d;\mathbb R)\)、\(\int_{\mathbb R^d}w>0\)，证明存在同一组固定符号 \(\sigma_j\in\{-1,1\}\)，使
\[
 \int_{\mathbb R^d}w(\xi)\cdot
 \left|\sum_{j=1}^N\sigma_jc_je^{-2\pi\mathrm i t_j\cdot\xi}\right|\,\mathrm d\xi
 \geq\sqrt{S/3}\int_{\mathbb R^d}w.
\]
说明为何不能对每个频率分别选一组符号后直接积分，并解释 \(c_j=1\) 时 \(N^{1/2}\) 增长的来源。

最后取 \(2<r\leq\infty\)、\(1/r+1/s=1\)，以及非零 \(\varphi\in C_c^\infty(\mathbb R^d;\mathbb C)\)。选 \(L>0\) 使 \(\varphi(\cdot-jLe_1)\) 的支集互不相交。将上述结论用于 \(c_j=1\)、\(t_j=jLe_1\)、\(w=|\widehat\varphi|^s\)，证明可以选同一组符号使
\[
 \begin{gathered}
 f_N(x)=\sum_{j=1}^N\sigma_j\varphi(x-jLe_1),\\
 \|f_N\|_r=N^{1/r}\|\varphi\|_r,\qquad
 \|\widehat {f_N}\|_s\geq\sqrt{N/3}\,\|\widehat\varphi\|_s.
 \end{gathered}
\]
其中 \(N^{1/\infty}=1\)。提示：在有限测度 \(|\widehat\varphi|^s\,\mathrm d\xi\) 上使用 Hölder 不等式。
\end{enumerate}
\end{exercise}
\begin{solution}
第一问，直接积分给出
\[
 \|h\|_p^p=2\int_0^\infty(1+x)^{-p}\,\mathrm dx=\frac2{p-1}.
\]
将指数换成 \(1\) 则积分发散。因此所用变换必须按\cref{b2:cor:04-lp-extension}的 \(L^p\) 延拓理解，该推论立即给出所求范数界。为证明尺度公式，令 \(h_n=h\cdot\mathbf1_{[-n,n]}\)，则 \(h_n\in L^1\cap L^p\)、\(h_n\to h\) 于 \(L^p\)，并且 \(\widehat{h_n}\to H\) 于 \(L^q\)。对这些可积截断先用伸缩公式：
\[
 \widehat{h_n(a\,\cdot)}(\xi)=a^{-1}\widehat{h_n}(\xi/a).
\]
左侧因输入在 \(L^p\) 中趋于 \(h_a\)，故在 \(L^q\) 中趋于 \(\mathcal F_ph_a\)；右侧在 \(L^q\) 中趋于 \(a^{-1}H(\cdot/a)\)。这里每个固定 \(a>0\) 的伸缩都连续，因为变量代换给出相应范数的精确尺度因子。极限唯一性便证明了公式，并得到
\[
 \|h_a\|_p=a^{-1/p}\left(\frac2{p-1}\right)^{1/p},
 \quad
 \|\mathcal F_ph_a\|_q
 =a^{-1+1/q}\|H\|_q=a^{-1/p}\|H\|_q.
\]

第二问，符号平均使含某个 \(\epsilon_j\) 奇数次的单项全部消失，所以
\[
 \mathbb E|Z|^2=S,
 \quad
 \mathbb E|Z|^4
 =2S^2+\left|\sum_jc_j^2\right|^2-2\sum_j|c_j|^4
 \leq3S^2.
\]
四阶等式可以按三种成对划分核验：两种给出 \(S^2\)，另一种给出 \(|\sum_jc_j^2|^2\)；四个指标全相等的项被计了三次，故减去两次 \(\sum_j|c_j|^4\)。最后用 \(|\sum_jc_j^2|\leq\sum_j|c_j|^2=S\) 得到上界。对有限符号空间使用指数 \(3/2\) 与 \(3\) 的 Hölder 不等式，
\[
 S=\mathbb E\bigl(|Z|^{2/3}\cdot|Z|^{4/3}\bigr)
 \leq(\mathbb E|Z|)^{2/3}(\mathbb E|Z|^4)^{1/3}.
\]
因 \(S>0\)，整理即得 \(\mathbb E|Z|\geq S^{3/2}/(\mathbb E|Z|^4)^{1/2}\geq\sqrt{S/3}\)。这一步把二、四矩的控制转成一阶绝对矩的下界，无须重复正文的阈值事件论证。

对每个固定 \(\xi\)，把系数换成 \(c_je^{-2\pi\mathrm i t_j\cdot\xi}\)，它们的模不变，所以同一个 \(\sqrt{S/3}\) 下界仍适用。乘以 \(w(\xi)\) 后积分，再交换有限平均与积分，得到
\[
 2^{-N}\sum_{\epsilon\in\{-1,1\}^N}
 \int_{\mathbb R^d}w(\xi)\cdot
 \left|\sum_j\epsilon_jc_je^{-2\pi\mathrm i t_j\cdot\xi}\right|\,\mathrm d\xi
 \geq\sqrt{S/3}\int_{\mathbb R^d}w.
\]
每个积分都不超过 \((\sum_j|c_j|)\int w<\infty\)。至少一个被平均的积分不小于平均值，选择那一组符号即得结论；这组符号在整个积分中保持固定。逐点另选符号可能随 \(\xi\) 改变，不能保证来自同一个有限和。等系数时 \(S=N\)，平方根增长正是二矩可加而四矩至多与其平方同阶的结果。

对最后的应用，光滑紧支集函数的变换可积且有界，所以 \(w=|\widehat\varphi|^s\in L^1\)；Fourier 唯一性保证 \(W\coloneqq\int w>0\)。支集分离给出 \(\|f_N\|_r=N^{1/r}\|\varphi\|_r\)，包括 \(r=\infty\)。令 \(Z_\sigma(\xi)=\sum_j\sigma_je^{-2\pi\mathrm i jL\xi_1}\)，所选符号满足 \(\int w|Z_\sigma|\geq\sqrt{N/3}\,W\)。当 \(s>1\) 时，Hölder 不等式给出
\[
 \int w|Z_\sigma|
 \leq\left(\int w|Z_\sigma|^s\right)^{1/s}W^{1-1/s};
\]
\(s=1\) 时同一式就是等式。因此
\[
 \|\widehat {f_N}\|_s
 =\left(\int w|Z_\sigma|^s\right)^{1/s}
 \geq\sqrt{N/3}\,W^{1/s}
 =\sqrt{N/3}\,\|\widehat\varphi\|_s.
\]
这由加权符号平均重新得到正文的增长机制，使用的是矩插值而非阈值事件；与输入的 \(N^{1/r}\) 比较，就看出 \(r>2\) 时范数界为何失效。
\end{solution}

% !TeX root = ../../Book2.tex
\chapter{Schwartz 空间与温和分布}
\label{b2:ch:05}
\BookChapterNumber{b2:ch:05}
% 资料：Hunter2001Applied，作者公开ch11.pdf，实际渲染核读291、294、303--304页。
% Schwartz半范数、双线性配对、归一化转换与正定Gauss公式分别核对。
% 本章自足证明Hermite完备性、S'自同构和各维负Laplace基本解。
\begin{chapterabstract}
Fourier 变换要容纳常函数、点质量和可积函数以外的对象，需要一个在变换下保持不变的测试空间。本章先构造 Schwartz 空间，证明半范数完备性及微分、多项式乘法、卷积和变换的连续性，并在选读专题中以 Hermite 函数展示变换的离散特征结构。随后把变换转移到连续对偶，得到温和分布上的统一理论。最后将常系数微分方程化为频率乘法，构造带正零阶项及各维负 Laplace 算子的基本解。
\end{chapterabstract}
\begin{learninggoals}
\begin{itemize}
\item\label{b2:item:05-goal-schwartz}使用具体半范数检查快速下降、收敛和连续操作，区分 Schwartz 收敛与共同紧支集的测试函数收敛。
\item\label{b2:item:05-goal-hermite}选读进阶目标：在\subsecref{b2:subsec:050204}中由 Gauss 函数构造 Hermite 系，理解其正交性、完备性和 Fourier 特征值。
\item\label{b2:item:05-goal-tempered}用连续泛函定义温和分布，保持双线性配对与既有微分符号一致。
\item\label{b2:item:05-goal-examples}计算 Dirac 测度、多项式、导数与齐次分布的变换，判断普通积分何时不能再用。
\item\label{b2:item:05-goal-fundamental}将常系数算子转为多项式符号，计算不同维数的基本解，并说明原点奇性和边界定向在证明中的作用。
\end{itemize}
\end{learninggoals}

一单位质量集中在原点时，普通函数无法同时表示它的总质量和无限集中程度。另一个看似相反的对象是常函数 \(1\)：它均匀铺满空间，却应只含零频率。两者的 Fourier 变换应彼此对应；要把这句话变成计算规则，需要先找到对二者都可测试、又在变换下保持不变的函数空间。本章构造这个空间，然后用点质量的方程 \(-\Delta E=\delta_0\) 说明广义对象怎样产生可用的解。

\ProseParagraphBreak

本章使用\chapref{b2:ch:03}的变换公式与反演、\chapref{b2:ch:04}的平方能量理论，以及\chapref{b2:ch:01}的分部积分和 Green 公式。第一卷定义 20.2.1 的普通分布作为比较对象：其测试函数紧支集，故不要求控制无穷远；温和分布则须连续作用于全部 Schwartz 函数。二者的范围不同，不能只因名称相似就互相替换。

\ProseParagraphBreak

\secref{b2:sec:0501}和\secref{b2:sec:0502}讨论 Schwartz 函数及其变换；\secref{b2:sec:0503}通过对偶配对定义温和分布上的 Fourier 变换；\secref{b2:sec:0504}用它研究常系数方程。John K. Hunter 与 Bruno Nachtergaele 的《Applied Analysis》\cite[Chapter 11]{Hunter2001Applied}介绍测试函数与温和分布。Lars Hörmander 的《The Analysis of Linear Partial Differential Operators I: Distribution Theory and Fourier Analysis》\cite{Hrmander2003Analysis}详细论述一般分布与常系数方程理论，可在学完本章后按具体问题查阅。

% !TeX root = ../../Book2.tex
\section{一阶线性方程}
\label{b2:sec:0701}

给定某一时刻的浓度分布，若介质只作平移而没有扩散，后来的图形应当保持原来的形状。我们先用方程描述这种平移，再沿直线求出解，由此介绍特征线法。本章把空间变量写在前面，时间变量写在后面，即 \(u=u(x,t)\)。向量场表示速度，函数值既可能表示随粒子携带的量，也可能表示单位体积中的密度；这两种含义将在\secref{b2:sec:0703}区分。

\subsection{输运方程}
\label{b2:subsec:070101}

设 \(b\in\mathbb R^d\) 为固定速度。经过初始位置 \(a\) 的粒子在时刻 \(t\) 位于 \(a+tb\)。若沿这一运动携带的数值不变，则链式法则给出
\[
 \frac{\mathrm d}{\mathrm dt}u(a+tb,t)
 =u_t(a+tb,t)+b\cdot\nabla_xu(a+tb,t)=0.
\]
这产生\term[shuyun-fangcheng]{输运方程}[transport equation]
\begin{equation}\label{b2:eq:07-constant-transport}
 u_t+b\cdot\nabla_xu=0.
\end{equation}
其中 \(u_t\) 衡量固定观察点的变化，\(b\cdot\nabla_xu\) 衡量移动造成的变化。二者可以分别非零而恰好相消。例如一维函数 \(u(x,t)=\sin(x-bt)\) 的每个波峰向右或向左平移，方向取决于 \(b\) 的符号。

\ProseParagraphBreak

这里讨论的经典解属于 \(C^1\)，因此上述链式法则有逐点意义。本卷\chapref{b2:ch:04}的分布语言允许进一步处理不光滑对象，但本章先把几何和唯一性建立在光滑范围内；仅在明确写出测试函数恒等式时才转入弱解。

\subsection{常系数情形}
\label{b2:subsec:070102}

\begin{proposition}\label{b2:prop:07-translation-solution}
设 \(I\) 为包含 \(0\) 的区间，\(u\in C^1(\mathbb R^d\times I)\)。函数 \(u\) 满足\eqref{b2:eq:07-constant-transport}，当且仅当存在 \(g\in C^1(\mathbb R^d)\)，使
\[
 u(x,t)=g(x-tb).
\]
其中 \(g(x)=u(x,0)\)，因而初值唯一确定解。
\end{proposition}
\begin{proof}
固定 \((x,t)\)，取 \(a=x-tb\)。函数 \(s\mapsto u(a+sb,s)\) 的导数为零，从 \(0\) 积分到 \(t\) 得 \(u(x,t)=u(a,0)\)。反过来，对所给公式求导，有 \(u_t=-b\cdot\nabla g(x-tb)\)、\(\nabla_xu=\nabla g(x-tb)\)，故方程成立。时间为负时，积分的方向反转，论证不变。
\end{proof}

平移不会提高正则性：初值在哪里存在折角，后来就在相应平移点保留折角。对于任意 \(1\leq p\leq\infty\)，若初值还属于 \(L^p\)，平移不变性给出 \(\|u(\cdot,t)\|_p=\|g\|_p\)。这一恒等式没有衰减因子，与以后热方程的平滑效应不同。

\ProseParagraphBreak

不把某个变量预先叫作时间时，同样的方法适用于 \(\alpha u_x+\beta u_y=0\)，其中 \((\alpha,\beta)\ne(0,0)\)。取
\[
 r=\alpha x+\beta y,\quad s=\beta x-\alpha y,
\]
则变换的行列式为 \(-\alpha^2-\beta^2\ne0\)，算子化为 \((\alpha^2+\beta^2)\partial_r\)。所以在整个平面上，解恰为 \(F(\beta x-\alpha y)\)。若区域与一条直线的交集不连通，各个连通片上的常数未必由同一个全局函数 \(F\) 描述；局部坐标公式不能自动升级为任意区域上的全局表示。

\subsection{非齐次方程}
\label{b2:subsec:070103}

源项使粒子所携带的值随时间累积，零阶项则造成吸收或增长。先考虑 \(\lambda\in\mathbb R\) 为常数的情形。

\begin{theorem}\label{b2:thm:07-constant-source}
设 \(g\in C^1(\mathbb R^d)\)、\(f\in C^1(\mathbb R^d\times[0,T])\)。问题
\[
 u_t+b\cdot\nabla u+\lambda u=f,\quad u(x,0)=g(x)
\]
有唯一的 \(C^1\) 解，且
\begin{equation}\label{b2:eq:07-constant-source}
 u(x,t)=e^{-\lambda t}\cdot g(x-tb)
 +\int_0^t e^{-\lambda(t-s)}\cdot f\bigl(x-(t-s)b,s\bigr)\,\mathrm ds.
\end{equation}
闭时间区间上的光滑性均理解为能局部延伸到邻域，端点导数也可取相应单侧值。
\end{theorem}
\begin{proof}
沿 \(x=a+tb\) 令 \(v(t)=u(a+tb,t)\)，方程成为
\(v'+\lambda v=f(a+tb,t)\)。乘以 \(e^{\lambda t}\) 并从 \(0\) 积分，得到
\[
 v(t)=e^{-\lambda t}\cdot g(a)
 +\int_0^t e^{-\lambda(t-s)}\cdot f(a+sb,s)\,\mathrm ds.
\]
令 \(a=x-tb\) 即为表示公式，故任何经典解只能取此值。公式右侧在每个紧时空集合上的被积函数及一阶导数都有共同界，\secref{b2:sec:0103}的积分求导规则保证它属于 \(C^1\)。把它沿每条直线拉回，所得函数满足刚才的常微分方程；链式法则再给出原方程。取 \(t=0\) 得到初值。
\end{proof}

例如在一维取 \(b=1\)、\(\lambda=0\)、\(g=0\)、\(f(x,t)=x\)，有
\[
 u(x,t)=\int_0^t(x-t+s)\,\mathrm ds=tx-\frac{t^2}{2}.
\]
这里源项须在粒子过去所在的位置取值；若误把积分写成 \(\int_0^tf(x,s)\,\mathrm ds\)，会得到不满足方程的 \(tx\)。

\subsection{初值问题}
\label{b2:subsec:070104}

初值面 \(t=0\) 被每条轨道恰好穿过一次，因此可以给每条轨道一个独立数值。若改在与轨道相切的曲线上给值，存在性或唯一性就会失效。以 \(u_x+u_y=0\) 为例，沿直线 \(y=x\) 解必须恒定；指定 \(u(s,s)=s\) 没有解，而指定 \(u(s,s)=0\) 则允许所有满足 \(F(0)=0\) 的解 \(u(x,y)=F(x-y)\)。方程与数据之间的几何位置是初值问题的一部分。

\ProseParagraphBreak

常系数公式还直接给出数据稳定性。若 \(u,\widetilde u\) 对应 \(g,f\) 与 \(\widetilde g,\widetilde f\)，并且右侧范数有限，则
\[
 \|u(\cdot,t)-\widetilde u(\cdot,t)\|_\infty
 \leq e^{|\lambda|t}\left(\|g-\widetilde g\|_\infty
 +\int_0^t\|f(\cdot,s)-\widetilde f(\cdot,s)\|_\infty\,\mathrm ds\right).
\]
证明只需在\eqref{b2:eq:07-constant-source}中相减并取绝对值。下一节把直线改成由速度场决定的曲线；为使这种替换合法，先要证明曲线存在、彼此不发生碰撞，并且随初始点可微地变化。

% !TeX root = ../../Book2.tex
\section{特征曲线}
\label{b2:sec:0702}

变系数速度场已经给出粒子在每一点的瞬时运动方向，却没有直接给出整条轨道。先解常微分方程，再用轨道作坐标，是本节的构造顺序。以下 \(U\subseteq\mathbb R^d\) 为开集，\(b\in C^1(U;\mathbb R^d)\)；本节先处理不显含时间的速度场。

\subsection{特征常微分方程}
\label{b2:subsec:070201}

\begin{definition}\label{b2:def:07-characteristic}
若 \(X:I\rightarrow U\) 属于 \(C^1\)，并满足
\[
 X'(t)=b\bigl(X(t)\bigr),
\]
则称它为速度场 \(b\) 的\term[tezheng-quxian]{特征曲线}[characteristic curve]，也称特征线。其常微分方程称为\term[tezheng-changweifen-fangcheng]{特征常微分方程}[characteristic ordinary differential equation]。
\end{definition}

沿这样的曲线，复合函数的导数只涉及 \(u_t+b\cdot\nabla u\)。但在应用链式法则之前，要确认对给定初始位置确实有一条且只有一条曲线。下面给出局部构造，并保留稍后需要的定量控制。

\begin{theorem}[Picard--Lindelöf]\label{b2:thm:07-local-ode}
对每个 \(a_0\in U\)，存在 \(r>0\)、\(h>0\)，使每个 \(a\in B_r(a_0)\) 都有唯一的局部解
\[
 X'(t)=b\bigl(X(t)\bigr),\quad X(0)=a,
 \quad |t|\leq h.
\]
各解取值于同一个紧球。若 \(L\) 是该球上 \(\|Db\|\) 的上界，则
\begin{equation}\label{b2:eq:07-flow-lipschitz}
 |X(t,a)-X(t,\widetilde a)|\leq e^{L|t|}\cdot|a-\widetilde a|.
\end{equation}
局部唯一性指任意两条满足同一初值的解在共同定义区间上相等。
\end{theorem}
\begin{proof}
取闭球 \(\overline B_{2r}(a_0)\subseteq U\)，记其上 \(|b|\)、\(\|Db\|\) 的上界为 \(M,L\)。选 \(h>0\) 满足 \(Mh\leq r\)、\(Lh<1\)；上界为零的条件可略去。对固定 \(a\in B_r(a_0)\)，在连续曲线空间中考虑闭集
\[
 E_a=\{\,Y\in C([-h,h];\mathbb R^d)\mid
 \sup_{|t|\leq h}|Y(t)-a|\leq r\,\}.
\]
一致 Cauchy 曲线列逐点有极限，一致极限仍连续且满足闭球约束，所以 \(E_a\) 关于一致范数完备。定义
\[
 (TY)(t)=a+\int_0^t b\bigl(Y(s)\bigr)\,\mathrm ds.
\]
速度界保证 \(T(E_a)\subseteq E_a\)，中值估计保证
\(\|TY-TZ\|_\infty\leq Lh\cdot\|Y-Z\|_\infty\)。从常曲线 \(Y_0=a\) 迭代，连续两次之差不超过 \((Lh)^n\cdot\|Y_1-Y_0\|_\infty\)。几何级数收敛，故迭代一致收敛到 \(X\in E_a\)。算子连续性给出 \(TX=X\)；积分的被积函数连续，所以 \(X\) 可微并满足方程。两个不动点之间的距离不超过原距离的 \(Lh\) 倍，故相等。

为证明稳定估计，先记录一个积分不等式。若连续非负函数 \(q\) 满足 \(q(t)\leq A+L\int_0^tq(s)\,\mathrm ds\)，令 \(Q(t)=A+L\int_0^tq(s)\,\mathrm ds\)，则 \(Q'\leq LQ\)，从而 \((e^{-Lt}Q)'\leq0\)，得到 \(q(t)\leq Ae^{Lt}\)。此结论也覆盖 \(A=0\)。对两个解的积分方程作差即可用它得到\eqref{b2:eq:07-flow-lipschitz}；负时间用 \(s=-t\) 处理。任意两个局部解在某个更短区间上都落入同一个球，故先在那里相等。在共同区间内，相等的时刻构成闭集，局部唯一性又使它为开集，因而它充满包含初始时刻的整个区间。
\end{proof}

上述结果称为\term[picard-lindelof]{Picard--Lindelöf 定理}[Picard--Lindelöf theorem]，证明中的积分估计是\term[gronwall-budengshi]{Gronwall 不等式}[Gronwall inequality]的一个基本形式。只要求速度连续时，唯一性可能失败。例如 \(X'=2\sqrt{|X|}\)、\(X(0)=0\) 除零解外，对每个 \(a\geq0\) 还允许先在 \([0,a]\) 上停留，之后取 \(X(t)=(t-a)^2\)。本章采用 \(C^1\) 速度，正是为了排除这种轨道分叉。

\subsection{流映射}
\label{b2:subsec:070202}

把从 \(a\) 出发的解记为 \(\Phi_t(a)\)，称 \(a\mapsto\Phi_t(a)\) 为\term[liu-yingshe]{流映射}[flow map]。局部解借唯一性在重叠区间上拼接，得到经过 \(a\) 的最大解区间 \(I_a\)。记号 \(\Phi_t(a)\) 只在 \(t\in I_a\) 时使用；不能因为给了一个映射名称便假定它对所有时间存在。
\symidx{\(\Phi_t\)：自治速度场的流映射}

\begin{proposition}\label{b2:prop:07-flow-composition}
凡各项有定义，流满足
\[
 \Phi_0(a)=a,\quad
 \Phi_{t+s}(a)=\Phi_t\bigl(\Phi_s(a)\bigr),\quad
 \Phi_{-t}\bigl(\Phi_t(a)\bigr)=a.
\]
固定 \(t\) 时，\(\Phi_t\) 在其定义域上单射，其逆为相应的 \(\Phi_{-t}\)。
\end{proposition}
\begin{proof}
曲线 \(r\mapsto\Phi_{s+r}(a)\) 与 \(r\mapsto\Phi_r(\Phi_s(a))\) 在 \(r=0\) 具有相同初值，且都满足自治方程。唯一性给出复合关系。令 \(r=-s\) 得逆关系，进而得到单射性。
\end{proof}

\begin{theorem}\label{b2:thm:07-flow-derivative}
流 \((t,a)\mapsto\Phi_t(a)\) 在其开定义域上属于 \(C^1\)。对初始点的导数 \(A(t,a)=D_a\Phi_t(a)\) 满足
\begin{equation}\label{b2:eq:07-variational-ode}
 A'(t,a)=Db\bigl(\Phi_t(a)\bigr)A(t,a),\quad A(0,a)=I.
\end{equation}
因此固定时间的流是其定义域到像集的 \(C^1\) 微分同胚。
\end{theorem}
\begin{proof}
先在上一构造的紧球和短时间内工作。式\eqref{b2:eq:07-flow-lipschitz}与 \(|\partial_t\Phi_t|\leq M\) 保证流联合连续。矩阵积分方程
\[
 A(t)=I+\int_0^t Db\bigl(\Phi_s(a)\bigr)A(s)\,\mathrm ds
\]
可用同样的迭代构造，且 \(\|A(t)\|\leq e^{L|t|}\)。现取小向量 \(v\)，令 \(\delta_v(t)=\Phi_t(a+v)-\Phi_t(a)\)。中值积分给出
\[
 b\bigl(\Phi_t(a+v)\bigr)-b\bigl(\Phi_t(a)\bigr)
 =Db\bigl(\Phi_t(a)\bigr)\delta_v(t)+r_v(t).
\]
导数 \(Db\) 在紧球上一致连续，而 \(|\delta_v(t)|\leq e^{Lh}|v|\)，所以
\(\sup_t|r_v(t)|/|v|\to0\)。将 \(\delta_v-Av\) 的两个积分方程相减，再用刚证的 Gronwall 估计，得到
\[
 \sup_t\frac{|\Phi_t(a+v)-\Phi_t(a)-A(t,a)v|}{|v|}\longrightarrow0.
\]
这证明的是完整 Fréchet 微分，而非只沿坐标求导。若 \(a\) 改为邻近点，矩阵方程的系数一致收敛；作差和 Gronwall 估计又给出 \(A\) 对 \((t,a)\) 连续。因此流为 \(C^1\)。

对一段已存在的轨道，其图像在任意紧时间区间内为紧集。用上述局部构造覆盖它，并取有限子覆盖，连续依赖使邻近初始点的解也能逐段通过这些邻域。故流的定义域为开集，局部可微结论可以沿轨道继续。最后，\cref{b2:prop:07-flow-composition}给出双方均 \(C^1\) 的逆映射，所以导数矩阵可逆，像集也为开集。
\end{proof}

流的空间可微性比轨道本身可微更强。前者控制相邻粒子之间的距离、角度和体积，后者只说明单个粒子的速度存在。后面 Jacobi 行列式的演化式依赖的是\eqref{b2:eq:07-variational-ode}，不能只从每条轨道 \(C^1\) 便直接写下。

\subsection{沿特征线的表示公式}
\label{b2:subsec:070203}

对自治输运问题 \(u_t+b(x)\cdot\nabla u=0\)，沿轨道有
\[
 \frac{\mathrm d}{\mathrm dt}u\bigl(\Phi_t(a),t\bigr)=0.
\]
因此只要轨道可从时刻 \(t\) 返回 \(0\)，解就由
\begin{equation}\label{b2:eq:07-autonomous-representation}
 u(x,t)=g\bigl(\Phi_{-t}(x)\bigr)
\end{equation}
给出。这里的“返回”既包括时间存在范围，也包括轨道始终位于系数定义域；若在有界区域中轨道先碰到边界，单靠内部初值不足以继续决定数值。

\ProseParagraphBreak

例如平面速度 \(b(x,y)=(-y,x)\) 给出
\[
 \Phi_t(a_1,a_2)=(a_1\cos t-a_2\sin t,\ a_1\sin t+a_2\cos t).
\]
解是初始图形的旋转，\(u(x,y,t)=g(x\cos t+y\sin t,-x\sin t+y\cos t)\)。另一例是 \(b(x)=x\)，其流为 \(\Phi_t(a)=e^ta\)，故 \(u(x,t)=g(e^{-t}x)\)。后者虽然函数值不变，体积却按 \(e^{dt}\) 扩张；把 \(u\) 的积分解释成守恒质量就会出错。

\subsection{存在与唯一性}
\label{b2:subsec:070204}

\begin{proposition}\label{b2:prop:07-global-flow}
设 \(b\in C^1(\mathbb R^d;\mathbb R^d)\)，且存在 \(C\geq0\) 使
\( |b(x)|\leq C(1+|x|)\)。则所有轨道在整个实轴上存在。对 \(g\in C^1(\mathbb R^d)\)，式\eqref{b2:eq:07-autonomous-representation}给出整个时空上的唯一 \(C^1\) 输运解。
\end{proposition}
\begin{proof}
对正时间，积分方程与 Gronwall 估计给出
\(1+|\Phi_t(a)|\leq(1+|a|)e^{Ct}\)。负时间对反向速度 \(-b\) 使用同一估计。若最大解区间有有限右端点 \(T\)，此估计使轨道在 \([0,T)\) 落于某个紧球。该球上速度有界，所以 \(|\Phi_t(a)-\Phi_s(a)|\leq M|t-s|\)，从而 \(t\uparrow T\) 时有极限 \(a_T\)。由 \(a_T\) 出发的局部解可与原解拼接；积分方程保证拼接点导数也满足方程，与最大性矛盾。左端点相同。

全局流的 \(C^1\) 性保证表示公式给出 \(C^1\) 函数。沿特征线该函数等于 \(g(a)\)，故满足方程。任意另一经典解沿同一轨道也恒等于 \(g(a)\)，因此唯一。
\end{proof}

线性增长条件可以有许多替代，但不能不加条件地省去。对于 \(b(x)=x^2\)，从 \(a>0\) 出发的轨道为 \(a/(1-at)\)，在 \(t=1/a\) 逃向无穷。局部唯一性始终成立，障碍在于不存在覆盖所有初始点的共同正时间流。另一方面，若只关心一条有限轨道及其附近，紧集上的局部条件已经足够，不必强求整个空间上的有界导数。

% !TeX root = ../../Book2.tex
\section{变系数输运方程}
\label{b2:sec:0703}

本节把时间依赖、源项和体积变化同时纳入输运。这里 \(d\) 为正整数，\(T>0\)。为避免轨道在有限时间逃离空间，先假定 \(b\in C^1(\mathbb R^d\times[0,T];\mathbb R^d)\)，并有 \(|b(x,t)|\leq C(1+|x|)\)、\(C\geq0\)。局部版本只需把讨论限制在能够返回初始面的轨道邻域中。

\subsection{时空特征线}
\label{b2:subsec:070301}

把从时刻 \(s\) 的位置 \(a\) 出发的轨道记为 \(X(t;s,a)\)，即
\[
 \partial_tX(t;s,a)=b\bigl(X(t;s,a),t\bigr),\quad X(s;s,a)=a.
\]
扩充空间中的向量场 \(B(x,t)=(b(x,t),1)\) 是自治的。对 \(B\) 应用\cref{b2:thm:07-local-ode}与\cref{b2:thm:07-flow-derivative}，再以空间线性增长估计延拓，便得到整个 \([0,T]\) 上的轨道及其 \(C^1\) 依赖。特别地，时间坐标严格按单位速度增加，故轨道与每个等时面恰相交一次。
\symidx{\(X(t;s,a)\)：从时刻 \(s\)、位置 \(a\) 出发的轨道}

\ProseParagraphBreak

唯一性给出的复合关系现在是
\[
 X(t;r,X(r;s,a))=X(t;s,a),\quad
 X(s;t,X(t;s,a))=a.
\]
中间时刻不能省略：一般没有仅由 \(t-s\) 决定的单参数流。固定终点 \((x,t)\) 后，曲线 \(r\mapsto X(r;t,x)\) 描述该粒子的过去，初始位置为 \(X(0;t,x)\)。例如 \(b(x,t)=t x\) 时，
\[
 X(t;s,a)=\exp\left(\frac{t^2-s^2}{2}\right)a;
\]
从 \(s=0\) 走一个单位时间与从 \(s=1\) 走一个单位时间，伸缩比例不同。

\begin{proposition}[速度扰动的轨道估计]\label{b2:prop:07-speed-stability}
设 \(d\) 为正整数，\(T>0\)，\(L,\delta\geq0\)，\(b,\widetilde b\in C^1(\mathbb R^d\times[0,T];\mathbb R^d)\) 满足 \(|b(x,t)|+|\widetilde b(x,t)|\leq C(1+|x|)\)（某个 \(C\geq0\)），从 \((a,s)\)、\(a\in\mathbb R^d\) 出发的轨道分别记为 \(X,\widetilde X\)，\(s,t\in[0,T]\)，并且
\[
 \sup_{x,t}\|D_xb(x,t)\|\leq L,
 \quad\sup_{x,t}|b(x,t)-\widetilde b(x,t)|\leq\delta.
\]
则从同一时刻、同一点出发的两条轨道满足
\[
 |X(t;s,a)-\widetilde X(t;s,a)|
 \leq\delta|t-s|e^{L|t-s|}.
\]
\end{proposition}
\begin{proof}
将两条轨道方程相减，在速度差中加减 \(b(\widetilde X,r)\)。中值估计和速度差界给出，若 \(t\geq s\)，
\[
 |X(t)-\widetilde X(t)|
 \leq L\int_s^t|X(r)-\widetilde X(r)|\,\mathrm{d}r+\delta(t-s).
\]
对一个固定终点 \(t\)，把右侧最后一项在整个 \([s,t]\) 上放大为常数 \(\delta(t-s)\)，Gronwall 给出所述界。反向时间对反向速度应用相同论证，所以 \(t<s\) 也成立。
\end{proof}

若同一初值 \(g\in C^1(\mathbb R^d;\mathbb C)\) 有全空间 Lipschitz 常数 \(M\)，分别以这两个速度场求解 \(u_t+b\cdot\nabla u=0\) 与 \(\widetilde u_t+\widetilde b\cdot\nabla\widetilde u=0\)，便有 \(\|u(\cdot,t)-\widetilde u(\cdot,t)\|_\infty\leq M\delta t e^{Lt}\)、\(0\leq t\leq T\)。这里没有源项或零阶项；对终点固定的两条回溯轨道使用命题，再将它们的初始位置代入 \(g\)。例如常速度增加 \(\varepsilon\) 时，相同图形的位置误差为 \(\varepsilon t\)。若初始图形有跳跃，平移再小也可能在上确界范数中产生大小为一的误差；速度稳定性和选择何种数据范数必须同时考虑。

\subsection{源项}
\label{b2:subsec:070302}

\begin{theorem}\label{b2:thm:07-variable-transport}
设 \(d\) 为正整数，\(T>0\)，\(b\in C^1(\mathbb R^d\times[0,T];\mathbb R^d)\)，且 \(|b(x,t)|\leq C(1+|x|)\) 对某个 \(C\geq0\) 成立。记 \(X(t;s,a)\) 为 \(\partial_tX=b(X,t)\)、\(X(s;s,a)=a\) 的轨道。设 \(c,f\in C^1(\mathbb R^d\times[0,T];\mathbb C)\)、\(g\in C^1(\mathbb R^d;\mathbb C)\)。初值问题
\[
 u_t+b\cdot\nabla u+c\cdot u=f,\quad u(x,0)=g(x)
\]
存在唯一的经典解 \(u\in C^1(\mathbb R^d\times[0,T];\mathbb C)\)。记
\[
 E(t,r;x)=\exp\left(-\int_r^t c(X(q;t,x),q)\,\mathrm dq\right),
\]
则解为
\begin{equation}\label{b2:eq:07-variable-transport}
 u(x,t)=E(t,0;x)\cdot g(X(0;t,x))
 +\int_0^tE(t,r;x)\cdot f(X(r;t,x),r)\,\mathrm dr.
\end{equation}
\end{theorem}
\begin{proof}
固定 \(a\)，令 \(v(t)=u(X(t;0,a),t)\)，则链式法则把原方程变成
\[
 v'(t)+c(X(t;0,a),t)\cdot v(t)=f(X(t;0,a),t),\quad v(0)=g(a).
\]
用积分因子求解，再令 \(a=X(0;t,x)\)，即得\eqref{b2:eq:07-variable-transport}。因每条轨道的标量方程唯一，原方程也至多有一个经典解。

反过来，所给公式在紧时空集合上由 \(C^1\) 函数的复合及有限区间积分组成。轨道的连续性使相应轨道集合紧致，因此被积函数及一阶导数均有局部共同界，积分求导合法。沿 \(X(t;0,a)\) 拉回该公式时，复合关系把每条历史轨道化为 \(X(r;0,a)\)，于是恢复刚才的积分因子公式，得到方程和初值。
\end{proof}

若 \(c\) 为实值且 \(c\geq0\)，因子 \(E\) 介于 \(0\) 与 \(1\) 之间。对于非负实值初值和源项，解始终非负；对于有界复值数据，
\[
 \|u(\cdot,t)\|_\infty\leq\|g\|_\infty
 +\int_0^t\|f(\cdot,r)\|_\infty\,\mathrm dr.
\]
对一般复值 \(c\)，若存在 \(C_0\geq0\) 使 \(\operatorname{Re}c\geq-C_0\)，则 \(|E(t,r;x)|\leq e^{C_0(t-r)}\)，只需在上述上界右端补上 \(e^{C_0t}\)。这类估计还可用于两组数据之差，给出固定速度和零阶系数下的稳定性；两种上界均要求右侧的数据范数积分有限。源项按轨道积分，而零阶项以指数积累，二者对解的作用不同。

\subsection{支集传播}
\label{b2:subsec:070303}

在齐次情形 \(f=0\) 中，指数因子永不为零，因而
\begin{equation}\label{b2:eq:07-support-flow}
 \operatorname{supp}u(\cdot,t)
 =X(t;0,\operatorname{supp}g).
\end{equation}
这是闭支集的等式：流是全空间的同胚，它把非零点集合映为非零点集合，也把闭包映为闭包。

\ProseParagraphBreak

有源项时，各时刻产生的数值继续沿轨道传播，但不同来源之间可能抵消。由表示公式只能断言
\[
 \operatorname{supp}u(\cdot,t)\subseteq
 \overline{X(t;0,\operatorname{supp}g)
 \cup\bigcup_{0\leq r\leq t}X(t;r,\operatorname{supp}f(\cdot,r))}.
\]
为核对这一包含，取右侧闭集之外一点及一个与它不交的小邻域。每个邻近终点的回溯轨道都避开初始支集与全部源项支集，表示公式的每一项因此都为零。闭包不能随意去掉，因为不同时间的支集之并一般不闭。

\begin{proposition}\label{b2:prop:07-finite-speed}
设 \(d\) 为正整数，\(T>0\)，固定速度场 \(b\in C^1(\mathbb R^d\times[0,T];\mathbb R^d)\) 和零阶系数 \(c\in C^1(\mathbb R^d\times[0,T];\mathbb C)\)。设 \(M\geq0\) 且 \(|b(x,t)|\leq M\)，记 \(X(t;s,a)\) 为 \(\partial_tX=b(X,t)\)、\(X(s;s,a)=a\) 的轨道。则对 \(s,t\in[0,T]\)、\(a\in\mathbb R^d\)，有 \(|X(t;s,a)-a|\leq M|t-s|\)。

对 \(j=1,2\)，取 \(g_j\in C^1(\mathbb R^d;\mathbb C)\)、\(f_j\in C^1(\mathbb R^d\times[0,T];\mathbb C)\)，令 \(u_j\in C^1(\mathbb R^d\times[0,T];\mathbb C)\) 为同一组系数下 \((u_j)_t+b\cdot\nabla u_j+cu_j=f_j\)、\(u_j(x,0)=g_j(x)\) 的唯一经典解。对固定的 \((x,t)\in\mathbb R^d\times[0,T]\)，若 \(g_1=g_2\) 在 \(\overline B_{Mt}(x)\) 上成立，且 \(f_1=f_2\) 在回溯圆锥
\[
 \{\,(y,r)\mid0\leq r\leq t,\ |y-x|\leq M(t-r)\,\}
\]
上成立，则 \(u_1(x,t)=u_2(x,t)\)。因而在固定 \(b,c\) 后，该点的解值仅依赖于初值球内的 \(g\) 和回溯圆锥内的 \(f\)。
\end{proposition}
\begin{proof}
积分轨道方程并取绝对值得到速度界。对每个回溯时刻 \(r\)，该界给出 \(|X(r;t,x)-x|\leq M(t-r)\)，特别是初始位置属于 \(\overline B_{Mt}(x)\)。两解使用同一个流，也使用由同一 \(c\) 决定的积分因子 \(E\)。在\eqref{b2:eq:07-variable-transport}中相减，得到
\[
 \begin{aligned}
 u_1(x,t)-u_2(x,t)
 &=E(t,0;x)(g_1-g_2)(X(0;t,x))\\
 &\quad+\int_0^t E(t,r;x)(f_1-f_2)(X(r;t,x),r)\,\mathrm dr=0.
 \end{aligned}
\]
初值项和源项分别因球内和圆锥内的数据相同而为零。
\end{proof}

实际依赖集通常比这个圆锥小得多：标量输运只读取一条回溯曲线。圆锥的作用是给出不必显式求流的统一上界。线性增长的速度仍保证有限时间的轨道存在，但没有一个对所有起点通用的有限传播速度；例如 \(b(x)=x\) 在远离原点处可以运动得任意快。

\subsection{Jacobi 行列式与质量输运初步}
\label{b2:subsec:070304}

记
\[
 J(t;s,a)\coloneqq\det D_aX(t;s,a).
\]
初始时刻 \(J(s;s,a)=1\)。正、逆流的可微性保证导数矩阵一直可逆，所以行列式不能变号。下面的公式进一步说明其大小。
\symidx{\(J(t;s,a)\)：流对初始位置的 Jacobi 行列式}

\begin{proposition}[Liouville 公式]\label{b2:prop:07-jacobian}
设 \(d\) 为正整数，\(T>0\)，\(b\in C^1(\mathbb R^d\times[0,T];\mathbb R^d)\)，且 \(|b(x,t)|\leq C(1+|x|)\) 对某个 \(C\geq0\) 成立。记 \(X(t;s,a)\) 为 \(\partial_tX=b(X,t)\)、\(X(s;s,a)=a\) 的轨道。对 \(s,t\in[0,T]\)、\(a\in\mathbb R^d\)，令 \(J(t;s,a)=\det D_aX(t;s,a)\)。
\begin{equation}\label{b2:eq:07-jacobian}
 J(t;s,a)=\exp\left(\int_s^t\operatorname{div}_xb(X(r;s,a),r)\,\mathrm dr\right)>0.
\end{equation}
\end{proposition}
\begin{proof}
令 \(A=D_aX\)、\(B=D_xb(X,t)\)，则可微依赖给出 \(A'=BA\)。对可逆矩阵，行列式的多重线性展开给出
\(\det(A+H)=\det A\cdot(1+\operatorname{tr}(A^{-1}H)+o(\|H\|))\)。事实上先提出 \(A\)，再展开 \(\det(I+K)\)，所有一次项之和为 \(\operatorname{tr}K\)，其余项至少二次。因此
\[
 J'=J\cdot\operatorname{tr}(A^{-1}A')
 =J\cdot\operatorname{tr}(A^{-1}BA)
 =J\cdot\operatorname{tr}B.
\]
最后一步可由 \(\operatorname{tr}(PQ)=\sum_{i,j}p_{ij}q_{ji}=\operatorname{tr}(QP)\) 验证。解此标量方程，并使用初始值 \(J=1\)，得到结论。
\end{proof}

这就是流的\term[liouville-gongshi]{Liouville 公式}[Liouville formula]。散度描述局部体积变化率；散度为零的流保持 Lebesgue 体积。密度 \(\rho\) 要抵消体积变化，因此守恒质量的方程应写成\term[lianxuxing-fangcheng]{连续性方程}[continuity equation]
\begin{equation}\label{b2:eq:07-continuity}
 \rho_t+\operatorname{div}_x(b\rho)=0,
 \quad\text{即}\quad
 \rho_t+b\cdot\nabla\rho+(\operatorname{div}b)\rho=0.
\end{equation}

\begin{theorem}\label{b2:thm:07-mass-transport}
设 \(d\) 为正整数，\(T>0\)，\(b\in C^2(\mathbb R^d\times[0,T];\mathbb R^d)\)，\(|b(x,t)|\leq C(1+|x|)\) 对某个 \(C\geq0\) 成立；\(X(t;0,a)\) 为此速度场的轨道，\(J(t;0,a)=\det D_aX(t;0,a)\)。设 \(\rho_0\in C^1(\mathbb R^d;\mathbb C)\)。则初值问题 \(\rho_t+\operatorname{div}_x(b\rho)=0\)、\(\rho(\cdot,0)=\rho_0\) 有唯一的经典解 \(\rho\in C^1(\mathbb R^d\times[0,T];\mathbb C)\)，并满足
\begin{equation}\label{b2:eq:07-density-jacobian}
 \rho(X(t;0,a),t)\cdot J(t;0,a)=\rho_0(a).
\end{equation}
若 \(\rho_0\) 为非负实值函数，或 \(\rho_0\in L^1(\mathbb R^d;\mathbb C)\)，则对每个 Borel 集 \(E\subseteq\mathbb R^d\) 和 \(t\in[0,T]\)，
\[
 \int_{X(t;0,E)}\rho(x,t)\,\mathrm dx=\int_E\rho_0(a)\,\mathrm da.
\]
非负情形允许两边为 \(+\infty\)。可积情形还有 \(\|\rho(\cdot,t)\|_1=\|\rho_0\|_1\)。
\end{theorem}
\begin{proof}
因为 \(\operatorname{div}b\in C^1\)，可对\eqref{b2:eq:07-continuity}使用\cref{b2:thm:07-variable-transport}，得到经典解的存在唯一性。沿轨道，\(\rho\) 满足的标量方程与 \(J\) 的方程恰有相反的零阶系数，乘积求导为零，得到\eqref{b2:eq:07-density-jacobian}。对微分同胚 \(a\mapsto X(t;0,a)\) 使用\secref{b2:sec:0102}的换元公式，因为 \(J>0\)，得到积分等式。若初始密度非负，可直接用非负换元；若绝对可积，先对绝对值换元，得到可积性及范数恒等式，再对正、负部或实、虚部分别积分。
\end{proof}

这里加 \(C^2\) 条件是为了由显式公式得到 \(C^1\) 密度；\cref{b2:prop:07-jacobian}本身只需要 \(b\in C^1\)。对于速度 \(b(x)=x\)，输运标量为 \(u(x,t)=g(e^{-t}x)\)，而守恒密度为 \(\rho(x,t)=e^{-dt}\cdot\rho_0(e^{-t}x)\)。两者的支集作同样的扩张，后者的振幅同时降低。运动几何、所携带的值与密度三者至此有了各自准确的表达式。

\ProseParagraphBreak

\Cref{b2:tab:07-scalar-density}把没有源项时的两种未知量并列比较。表中 \(d\) 为正整数，\(T>0\)，同一实速度场 \(b\in C^2(\mathbb R^d\times[0,T];\mathbb R^d)\) 满足 \(|b(x,t)|\leq C(1+|x|)\)、\(C\geq0\)；记 \(X(t,a)=X(t;0,a)\)、\(J(t,a)=\det D_aX(t;0,a)>0\)。标量初值取 \(g\in C^1(\mathbb R^d;\mathbb C)\)，密度初值取非负的 \(\rho_0\in C^1(\mathbb R^d;\mathbb R)\cap L^1(\mathbb R^d;\mathbb R)\)。表中的“移动区域”是同一群初始粒子 \(E\) 的像 \(X(t,E)\)，其中 \(E\subseteq\mathbb R^d\) 为 Borel 集。

\begin{table}[htbp]
\centering
\small
\renewcommand{\arraystretch}{1.3}
\caption[输运标量与守恒密度]{输运标量与守恒密度。两列使用同一个流；密度沿轨道的变化抵消了体积因子，标量沿轨道的数值不变则没有这种补偿。}
\label{b2:tab:07-scalar-density}
\begin{tabularx}{\linewidth}{@{}>{\raggedright\arraybackslash}p{1.5cm}>{\raggedright\arraybackslash}X>{\raggedright\arraybackslash}X@{}}
\toprule
 & 输运标量 \(u\) & 守恒密度 \(\rho\) \\
\midrule
方程 & \(u_t+b\cdot\nabla u=0\) & \(\rho_t+\operatorname{div}_x(b\rho)=0\) \\
沿轨道 & \(u(X(t,a),t)=g(a)\) & \(\rho(X(t,a),t)J(t,a)=\rho_0(a)\) \\
\(b(x)=x\) & \(u(x,t)=g(e^{-t}x)\) & \(\rho(x,t)=e^{-dt}\rho_0(e^{-t}x)\) \\
保持的量 & 每个粒子携带的函数值；总积分或 \(L^p\) 范数即使有限，也不必保持。 & 随流移动区域内的质量：\(\int_{X(t,E)}\rho=\int_E\rho_0\)。总质量及 \(L^1\) 范数不变。 \\
\bottomrule
\end{tabularx}
\end{table}

移动区域的质量不变，并不表示每个固定观察区域内的质量也不变；粒子仍可穿过固定区域的边界。这一区别不要求移动区域具有光滑边界，质量恒等式已经由换元对所有 Borel 集给出。

\ProseParagraphBreak

Jacobi 公式同时决定各个积分范数的演化。以下取正整数 \(d\)、\(T>0\)，实速度场 \(b\in C^1(\mathbb R^d\times[0,T];\mathbb R^d)\) 满足 \(|b(x,t)|\leq C(1+|x|)\)、\(C\geq0\)，且 \(c\in C^1(\mathbb R^d\times[0,T];\mathbb R)\)、\(g\in C^1(\mathbb R^d;\mathbb C)\)。对无源问题 \(u_t+b\cdot\nabla u+cu=0\)、\(u(x,0)=g(x)\) 的经典解，沿轨道有
\[
 u(X(t;0,a),t)=g(a)\exp\left(-\int_0^tc(X(r;0,a),r)\,\mathrm{d}r\right).
\]
对 \(1\leq p<\infty\) 作换元，得到精确式
\[
 \|u(\cdot,t)\|_p^p
 =\int |g(a)|^p\exp\left(\int_0^t
   (\operatorname{div}b-pc)(X(r;0,a),r)\,\mathrm{d}r\right)\,\mathrm{d}a.
\]
若存在实常数 \(D,c_0\)，使整个时空条带上 \(\operatorname{div}b\leq D\)、\(c\geq c_0\)，且 \(g\in L^p(\mathbb R^d;\mathbb C)\)，便有 \(\|u(\cdot,t)\|_p\leq e^{(D/p-c_0)t}\|g\|_p\)。无穷范数不读取体积；若 \(g\in L^\infty\)，直接由表示公式得 \(\|u(\cdot,t)\|_\infty\leq e^{-c_0t}\|g\|_\infty\)。保持同一组 \(b,c\)，在有源情形取 \(f\in C^1(\mathbb R^d\times[0,T];\mathbb C)\)，对每个源时刻用同一换元，再用 Minkowski 不等式，给出
\[
 \|u(\cdot,t)\|_p
 \leq e^{(D/p-c_0)t}\|g\|_p
 +\int_0^t e^{(D/p-c_0)(t-r)}\|f(\cdot,r)\|_p\,\mathrm{d}r,
\]
其中 \(p=\infty\) 按 \(D/p=0\) 解释，并要求右侧有限。

\ProseParagraphBreak

对\cref{b2:thm:07-mass-transport}中 \(b\in C^2\) 所给出的守恒密度，零阶系数是 \(c=\operatorname{div}b\)。它的 \(p\) 次质量因子变为 \((1-p)\operatorname{div}b\)。所以只有 \(p=1\) 对任意体积变化恰好抵消；若 \(1\leq p\leq\infty\)、\(\rho_0\in L^p(\mathbb R^d;\mathbb C)\)，且整个时空条带上 \(|\operatorname{div}b|\leq D\)、\(D\geq0\)，则
\[
 \|\rho(\cdot,t)\|_p\leq e^{(1-1/p)Dt}\|\rho_0\|_p.
\]
对于 \(b(x)=x\)，还可直接算出 \(\|\rho(\cdot,t)\|_p=e^{-d(1-1/p)t}\|\rho_0\|_p\)，而输运标量的范数为 \(e^{dt/p}\|g\|_p\)。同一流对不同物理量产生不同范数变化，不能把所有输运都当成保持每个 \(L^p\) 范数的平移。

% !TeX root = ../../Book2.tex
\OptionalSection{一阶非线性方程}{b2:sec:0704}

在线性输运中，先由已知速度场确定轨道，再沿轨道计算数值。若速度依赖未知函数，这两步便相互耦合。本节只构造局部光滑解，并用一个显式例子说明这种解可能在有限时间失效。后面的波动、热与 Laplace 方程不依赖本节。

\subsection{拟线性方程}
\label{b2:subsec:070401}

考虑已在\secref{b2:sec:0104}定义的拟线性类型
\[
 u_t+b(x,t,u)\cdot\nabla u=f(x,t,u),\quad u(x,0)=g(x).
\]
设 \(b,f\) 在相关变量的开集上为 \(C^1\)，\(g\in C^1\)。把未知的函数值也看成一个坐标，解联立系统
\begin{equation}\label{b2:eq:07-quasilinear-system}
 \begin{aligned}
 X'&=b(X,t,Z),&X(0,a)&=a,\\
 Z'&=f(X,t,Z),&Z(0,a)&=g(a).
 \end{aligned}
\end{equation}
对向量 \((X,t,Z)\) 应用局部流定理，得到其对 \((t,a)\) 的 \(C^1\) 依赖。时刻零有 \(D_aX=I\)，故映射 \((a,t)\mapsto(X(t,a),t)\) 在每个 \((a_0,0)\) 附近的导数可逆。

\begin{proposition}\label{b2:prop:07-quasilinear-local}
在上述假设下，对每个初始点 \(a_0\)，存在它在时空中的邻域，在其中拟线性初值问题有唯一的局部 \(C^1\) 解。
\end{proposition}
\begin{proof}
由\cref{b2:thm:01-inverse}，缩小邻域后可反解 \(a=A(x,t)\)。令 \(u(x,t)=Z(t,A(x,t))\)，则
\(u(X(t,a),t)=Z(t,a)\)。对此式求时间导数并代入\eqref{b2:eq:07-quasilinear-system}，得到偏微分方程。初始时 \(A(x,0)=x\)，所以 \(u(x,0)=g(x)\)。若另有经典解 \(v\)，沿满足 \(X'=b(X,t,v(X,t))\) 的曲线令 \(Z=v(X,t)\)，链式法则表明 \((X,Z)\) 满足同一联立系统和同一初值。常微分方程唯一性使二者与构造的 \((X,Z)\) 相同，再由局部反演得到 \(v=u\)。
\end{proof}

这一结论的邻域可以依赖 \(a_0\)。即使联立系统一直有解，其空间投影 \(a\mapsto X(t,a)\) 也可能先失去可逆性。下面的例子把这种障碍显示为一个分母趋零的计算。

\subsection{特征线交叉}
\label{b2:subsec:070402}

\term[burgers-fangcheng]{无黏 Burgers 方程}[inviscid Burgers equation]为
\[
 u_t+u u_x=0.
\]
沿特征线 \(Z'=0\)，所以 \(Z(t,a)=g(a)\)、\(X(t,a)=a+tg(a)\)。若 \(g\in C^1(\mathbb R)\) 有界，且 \(1+tg'(a)>0\) 对所有 \(a\) 成立，则 \(X(t,\cdot)\) 严格递增，又因 \(|X(t,a)-a|\leq t\|g\|_\infty\) 而满射。它是全实轴的微分同胚，解由 \(u(X(t,a),t)=g(a)\) 给出，并满足
\begin{equation}\label{b2:eq:07-burgers-gradient}
 u_x(X(t,a),t)=\frac{g'(a)}{1+tg'(a)}.
\end{equation}
若 \(g'\geq0\)，上述构造对所有正时间有效；若 \(g'\) 的负最小值 \(m\) 在某点取到，则在 \(T_*=-1/m\) 之前有效，而该点轨道上的梯度在 \(t\uparrow T_*\) 时趋于负无穷。

\ProseParagraphBreak

具体取 \(g(a)=-\tanh a\)，有 \(m=-1\)、\(T_*=1\)。函数值始终处于 \((-1,1)\)，但原点处梯度为 \(-1/(1-t)\)。失效不是振幅无限增长，而是相邻初始点的空间间隔被压缩为零。当 \(t>1\) 时，映射 \(a-t\tanh a\) 在原点附近下降，在两端上升，因而不单射。两条不同初始标签的轨道可以投影到同一个时空点，而所携带的值不同，不可能成为单值经典解。\Cref{b2:fig:07-characteristic-compression}区分了光滑阶段与失效后的轨道投影。
% 自绘：g(a)=-tanh(a) 的特征投影；虚线仅表示光滑失效后的轨道投影。
\begin{figure}[htbp]
\centering
\begin{tikzpicture}[x=2.3cm,y=2.2cm,>=stealth]
 \draw[->] (-1.35,0)--(1.35,0) node[right]{\(x\)};
 \draw[->] (-1.3,-0.05)--(-1.3,1.6) node[above]{\(t\)};
 \draw[gray,densely dotted] (-1.25,1)--(1.25,1);
 \node[left] at (-1.3,1){\(1\)};
 \foreach \a/\v in {-1/-0.761594,-0.75/-0.635149,-0.5/-0.462117,-0.25/-0.244919,0/0,0.25/0.244919,0.5/0.462117,0.75/0.635149,1/0.761594}{
   \draw[AnalysisBlue,thick] (\a,0)--({\a-\v},1);
   \draw[AnalysisBlue,dashed] ({\a-\v},1)--({\a-1.45*\v},1.45);
 }
 \foreach \x in {-1,0,1}{\draw (\x,0.025)--(\x,-0.025) node[below]{\(\x\)};}
 \node[anchor=west,align=left] at (0.68,1.28){\small 轨道投影};
 \node[anchor=west] at (0.68,0.72){\small 光滑阶段};
\end{tikzpicture}
\caption[特征线的压缩与投影交叉]{特征线的压缩与投影交叉。初值为 \(g(a)=-\tanh a\)，轨道投影为 \(x=a-t\tanh a\)。在 \(t=1\) 首次发生梯度爆破；虚线延长只展示投影的交叉，不表示单值经典解的延续。}
\label{b2:fig:07-characteristic-compression}
\end{figure}


\subsection{激波形成的预告}
\label{b2:subsec:070403}

在经典解存在时，Burgers 方程等价于
\(u_t+\partial_x(u^2/2)=0\)。对于有跳跃的函数，后一种写法仍可通过分布导数解释。若一般通量 \(F\in C^1(\mathbb R)\)，一条 \(C^1\) 曲线 \(x=s(t)\) 两侧的函数都是常数 \(u_-,u_+\)，把测试积分在曲线两侧分别分部积分，界面项为
\[
 \int \bigl(F(u_+)-F(u_-)-s'(t)(u_+-u_-)\bigr)
 \varphi(s(t),t)\,\mathrm dt.
\]
它对所有紧支集光滑测试函数为零，当且仅当
\[
 s'(t)(u_+-u_-)=F(u_+)-F(u_-).
\]
这称为\term[rankine-hugoniot]{Rankine--Hugoniot 条件}[Rankine--Hugoniot condition]。对不同的 Burgers 左右状态，速度为 \((u_-+u_+)/2\)；这样的移动跳跃称为\term[jibo]{激波}[shock wave]。

\ProseParagraphBreak

界面恒等式只给出弱方程的相容条件，没有证明某个跳跃是光滑解失效后的唯一延续。物理上允许的压缩跳跃与不允许的展开跳跃，都可能满足同一恒等式。要挑选延续还需附加条件，并建立存在、稳定性与唯一性理论；本节的梯度爆破计算本身不承担这些结论。

\subsection{Hamilton--Jacobi 方程}
\label{b2:subsec:070404}

另一类重要的一阶非线性方程是\term[hamilton-jacobi-fangcheng]{Hamilton--Jacobi 方程}[Hamilton--Jacobi equation]
\[
 u_t+H(x,\nabla u)=0.
\]
这里 \(H\) 为给定函数。先假定 \(H\in C^3\)、\(g\in C^3\)，只在初始点附近构造光滑解。取
\begin{equation}\label{b2:eq:07-hamilton-system}
 X'=H_p(X,P),\quad P'=-H_x(X,P),\quad
 Z'=P\cdot H_p(X,P)-H(X,P),
\end{equation}
初始值为 \(X(0,a)=a\)、\(P(0,a)=\nabla g(a)\)、\(Z(0,a)=g(a)\)。前两式描述位置和候选梯度的演化，第三式计算函数值。

\ProseParagraphBreak

这里需要的二阶参数可微性仍来自流构造：对 \(C^2\) 向量场 \(B\)，在\eqref{b2:eq:07-variational-ode}上再取差商，二阶候选导数满足
\[
 V'=DB(Y)V+D^2B(Y)[A v,A w],\quad V(0)=0.
\]
右侧系数连续，矩阵积分方程给出唯一 \(V\)。以 \(D^2B\) 在紧集上的一致连续性控制差商余项，再用 Gronwall 估计，得到余项一致趋零。这与一阶证明的线性余项论证相同，且明确给出二阶导数的方程；所以此处构造的 \(X,P,Z\) 对 \((t,a)\) 为 \(C^2\)。

\begin{proposition}\label{b2:prop:07-hamilton-local}
在上述假设下，系统\eqref{b2:eq:07-hamilton-system}在充分小的邻域中确定一个 \(C^2\) 解 \(u\)，满足 Hamilton--Jacobi 方程及初值 \(g\)。
\end{proposition}
\begin{proof}
从 \(D_aX(0,a)=I\) 出发，局部反演得到 \(a=A(x,t)\)，令 \(u(x,t)=Z(t,A(x,t))\)。关键是验证 \(P=\nabla u\)，不能把候选梯度直接当成真实梯度。对每个初始坐标 \(a_j\)，记
\(W_j=\partial_{a_j}Z-P\cdot\partial_{a_j}X\)。对系统求导，利用 \(X'=H_p\)、\(P'=-H_x\)，得到
\[
 \begin{aligned}
 \partial_{a_j}Z'
 &=P\cdot\partial_{a_j}(H_p)-H_x\cdot\partial_{a_j}X,\\
 (P\cdot\partial_{a_j}X)'
 &=-H_x\cdot\partial_{a_j}X+P\cdot\partial_{a_j}(H_p).
 \end{aligned}
\]
故 \(W_j'=0\)，初值又为 \(\partial_jg-\partial_jg=0\)。因此 \(D_aZ=P^TD_aX\)，由 \(D_aX\) 可逆得到 \(\nabla_xu(X,t)=P\)。最后沿轨道有
\(u_t+\nabla u\cdot X'=Z'=P\cdot H_p-H\)，约去相同项得 \(u_t+H(x,\nabla u)=0\)。
\end{proof}

取 \(H(p)=|p|^2/2\) 时，\(P=\nabla g(a)\)、\(X=a+t\nabla g(a)\)、\(Z=g(a)+t|\nabla g(a)|^2/2\)。可逆性要求 \(I+tD^2g(a)\) 不退化。这与 Burgers 方程的分母具有同一几何含义，但高维中须检查整个矩阵。局部光滑构造失效之后的解概念，本章不再展开。

\subsection{守恒律的位置}
\label{b2:subsec:070405}

形如 \(u_t+\operatorname{div}_xF(u)=0\) 的方程称为\term[shouhenglv]{守恒律}[conservation law]，其中 \(F:\mathbb R\rightarrow\mathbb R^d\)。光滑时，它等价于拟线性方程 \(u_t+F'(u)\cdot\nabla u=0\)；不光滑时，两种写法的意义不同，后一项中的乘积未必已有定义。守恒形式可直接作用于测试函数，通过积分表达局部通量平衡。

\ProseParagraphBreak

线性连续性方程的速度由外部给定，密度与体积的补偿由 Jacobi 公式控制；非线性守恒律的速度随状态变化，投影交叉会使经典表示失效。两者都使用特征线，却不能把线性流的全局微分同胚结论原样移到非线性问题。本节建立了光滑阶段的局部构造和一个明确的失效机制，为以后阅读守恒律与 Hamilton--Jacobi 理论划出入口。


\begin{chaptersummary}
Schwartz 空间用所有阶导数和所有多项式权重控制函数。Fourier 变换在其中交换微分与乘法，反演使它成为连续自同构；Hermite 选读专题则以正交基把这一变换的四种特征值具体呈现出来。温和分布通过作用于测试函数定义，双线性配对把已有的变换公式连续延伸到 Dirac 测度、多项式以及全部 \(L^p\) 函数。

\ProseParagraphBreak

常系数方程的频率表达是多项式乘法，但零点处不能随意相除。基本解的构造必须兼顾局部可积性、原点质量和无穷远增长。负 Laplace 的幂核、对数核和一维线性核各反映一种维数边界；与紧支集光滑右端卷积可解全空间方程，却不自动决定边值条件。下一章将用同一测试空间把离散格点与连续频率联系起来。
\end{chaptersummary}

\BookExercises

以下题目把半范数、差商和频率符号用于具体对象。前几题强调收敛与符号，后几题处理奇性、频率特征结构和方程。\Cref{b2:ex:05-hermite-oscillator}对应选读的 Hermite 专题；\cref{b2:ex:05-polynomial-comb}比较离散测度的多项式增长与指数增长，最后的\cref{b2:ex:05-cutoff-laplace-kernel}考察截断基本解后出现的壳层余项。

% 自拟：所有导数一致趋零仍不控制多项式权重。
% 预留题头及首个公式所需空间，避免题头换页后超链接锚点仍在前页。
\Needspace{6\baselineskip}
\begin{exercise}\label{b2:ex:05-schwartz-convergence}
设非零 \(\rho\in C_c^\infty(\mathbb R;\mathbb C)\)，对正整数 \(n\) 令 \(\varphi_n(x)=n^{-1}\rho(x/n)\)。证明每个固定阶导数都在 \(\mathbb R\) 上一致趋于零，但 \(\varphi_n\) 不在 \(\mathcal S(\mathbb R;\mathbb C)\) 中趋零。它作为温和分布是否逐测试函数配对趋于零？
\end{exercise}
\begin{solution}
有 \(\|\varphi_n^{(k)}\|_\infty=n^{-k-1}\|\rho^{(k)}\|_\infty\to0\)。选 \(x_0\ne0\) 使 \(\rho(x_0)\ne0\)，则 \(p_1(\varphi_n)\geq n^{-1}(1+n|x_0|)|\rho(x_0)|\)，下极限为正，故不按 Schwartz 半范数趋零。另一方面，对 \(\psi\in\mathcal S\)，\(|\int\varphi_n\psi|\leq n^{-1}\|\rho\|_\infty\cdot\|\psi\|_1\to0\)，所以在 \(\mathcal S'\) 中趋零。
\end{solution}

% 自拟：差分点质量的极限与傅里叶符号同时核对。
\Needspace{6\baselineskip}
\begin{exercise}\label{b2:ex:05-delta-differences}
对正整数 \(n\)，在 \(\mathcal S'(\mathbb R;\mathbb C)\) 中按逐测试函数配对收敛求下列分布列的极限，并计算极限的 Fourier 变换，其中 \(\delta_a\) 表示位于 \(a\in\mathbb R\) 的 Dirac 测度：
\[
 n(\delta_{1/n}-\delta_0),\quad
 n^2(\delta_{1/n}+\delta_{-1/n}-2\delta_0).
\]
\end{exercise}
\begin{solution}
对 \(\varphi\in\mathcal S\) 使用一阶与二阶 Taylor 差商，配对极限分别为 \(\varphi'(0)\) 和 \(\varphi''(0)\)，故极限为 \(-\delta_0'\)、\(\delta_0''\)。它们的变换分别为 \(-2\pi\mathrm{i}\xi\)、\(-4\pi^2\xi^2\)。也可从指数差商核对：\(n(e^{-2\pi\mathrm{i}\xi/n}-1)\to-2\pi\mathrm{i}\xi\)，其模由 \(2\pi|\xi|\) 控制；对称二阶差商的模由 \(4\pi^2\xi^2\) 控制，故两者都可对 Schwartz 测试函数用控制收敛。
\end{solution}

% 自拟：乘以x的核及原点导数项，训练分布除法。
\Needspace{6\baselineskip}
\begin{exercise}\label{b2:ex:05-x-times-distribution}
设 \(T\in\mathcal S'(\mathbb R;\mathbb C)\)。证明 \(xT=0\) 当且仅当存在 \(c\in\mathbb C\) 使 \(T=c\delta_0\)。再求方程 \(xT=\delta_0\) 在同一空间中的全部解。
\end{exercise}
\begin{solution}
取 \(\rho\in C_c^\infty\)，在原点附近为 \(1\)。任意 \(\varphi\in\mathcal S\) 可写成 \(\varphi=\varphi(0)\rho+x\psi\)，其中 \(\psi=(\varphi-\varphi(0)\rho)/x\) 在原点用光滑延拓，远处除以 \(x\) 保持快速下降，所以 \(\psi\in\mathcal S\)。若 \(xT=0\)，便有 \(\langle T,\varphi\rangle=\varphi(0)\langle T,\rho\rangle\)，即 \(T=c\delta_0\)。反向直接成立。又因 \(\langle x\delta_0',\varphi\rangle=-(x\varphi)'(0)=-\varphi(0)\)，一个特解为 \(-\delta_0'\)。全部解因而为 \(-\delta_0'+c\delta_0\)。
\end{solution}

% 自拟：主值分布、奇偶性和一阶分布方程组合得到变换常数。
\Needspace{6\baselineskip}
\begin{optionalexercise}\label{b2:ex:05-principal-value}
证明对每个 \(\varphi\in\mathcal S(\mathbb R;\mathbb C)\)，公式
\[
 \langle T,\varphi\rangle
 =\lim_{\varepsilon\downarrow0}\int_{|x|>\varepsilon}\dfrac{\varphi(x)}x\,\mathrm{d}x
\]
定义 \(T\in\mathcal S'(\mathbb R;\mathbb C)\)，称为 \(\operatorname{pv}(1/x)\)。证明 \(xT=1\)，其中右侧是常函数 \(1\) 给出的分布，并计算 \(\widehat T\)。
\end{optionalexercise}
\begin{solution}
在 \(|x|<1\) 内减去 \(\varphi(0)\)，常数项因对称性抵消，剩余绝对值由 \(\|\varphi'\|_\infty\) 控制；在 \(|x|\geq1\) 内用 \(p_2(\varphi)(1+|x|)^{-2}/|x|\) 控制。因此极限存在且满足有限半范数界。代入 \(x\varphi\) 得 \(xT=1\)，并可换元看出 \(T\) 为奇分布。

由变换公式，\(\frac{\mathrm{i}}{2\pi}\partial_\xi\widehat T=\delta_0\)，故 \(\partial_\xi\widehat T=-2\pi\mathrm{i}\delta_0\)。函数 \(-\pi\mathrm{i}\operatorname{sgn}\xi\) 的分布导数正好为右端，所以两者之差的导数为零。导数为零的温和分布是常数：对积分为零的 \(\psi\in\mathcal S\)，其原函数 \(\int_{-\infty}^x\psi\) 属于 \(\mathcal S\)，故该分布在所有零积分测试函数上为零；再减去一个固定单位积分测试函数的倍数即可。于是 \(\widehat T=-\pi\mathrm{i}\operatorname{sgn}\xi+c\)。Fourier 变换保持奇偶性，而 \(T\) 为奇，所以 \(c=0\)。
\end{solution}

% 自拟：求Hermite低阶式，并把频率特征关系与谐振子联系。
\Needspace{6\baselineskip}
\begin{optionalexercise}\label{b2:ex:05-hermite-oscillator}
在 \(\mathbb R\) 上令 \(h_0(x)=2^{1/4}e^{-\pi x^2}\)、\(a^*=(4\pi)^{-1/2}(-\partial_x+2\pi x)\)，并对 \(n\in\mathbb N_0\) 令 \(h_n=(a^*)^nh_0/\sqrt{n!}\)，即采用\subsecref{b2:subsec:050204}的归一化。求 \(h_1,h_2\) 的显式公式，证明对每个 \(n\in\mathbb N_0\) 有
\[
 \left(-\dfrac{\mathrm{d}^2}{\mathrm{d}x^2}+4\pi^2x^2\right)h_n
 =(4n+2)\pi h_n,
\]
并把 \(x^2h_0\) 展成 Hermite 函数的有限线性组合。
\end{optionalexercise}
\begin{solution}
从上升算子计算得 \(h_1=\sqrt{4\pi}\,xh_0\)、\(h_2=2^{-1/2}(4\pi x^2-1)h_0\)。记下降算子 \(a=(4\pi)^{-1/2}(\partial_x+2\pi x)\)，展开算子有 \(4\pi a^*a=-\mathrm{d}^2/\mathrm{d}x^2+4\pi^2x^2-2\pi\)。由 \(a^*ah_n=nh_n\) 得到特征方程。第二个显式公式给出 \(x^2h_0=(h_0+\sqrt2h_2)/(4\pi)\)。其变换为 \((h_0-\sqrt2h_2)/(4\pi)\)，也可与两次频率求导公式核对。
\end{solution}

% 自拟：有界光滑乘子若导数失控，可以破坏S。
\Needspace{6\baselineskip}
\begin{exercise}\label{b2:ex:05-fast-multiplier}
在 \(\mathbb R\) 上令 \(a(x)=e^{\mathrm{i}e^{x^2}}\)、\(\varphi(x)=e^{-x^2}\)。证明 \(a\colon\mathbb R\to\mathbb C\) 光滑且有界，\(\varphi\in\mathcal S(\mathbb R;\mathbb C)\)，但 \(a\varphi\notin\mathcal S(\mathbb R;\mathbb C)\)。再判断乘以 \(b(x)=e^{\mathrm{i}x^3}\) 是否给出 \(\mathcal S(\mathbb R;\mathbb C)\) 的连续自同构，并说明其逆。
\end{exercise}
\begin{solution}
直接求导得 \((a\varphi)'=a(-2xe^{-x^2}+2\mathrm{i}x)\)，其模随 \(|x|\) 无界，故不满足 Schwartz 条件。对 \(b\)，每个导数都是一个多项式乘 \(b\)，因而至多多项式增长。由正文的乘法命题，乘以 \(b\) 连续保持 \(\mathcal S\)；其逆乘子 \(\overline b\) 也有相同性质，所以还给出连续自同构。
\end{solution}

% 自拟：尺度决定屏蔽算子的基本解与L1质量。
\Needspace{6\baselineskip}
\begin{exercise}\label{b2:ex:05-resolvent-scale}
设 \(d\) 为正整数，\(E\in L^1(\mathbb R^d;\mathbb R)\) 为几乎处处由
\[
 E(x)=\int_0^\infty e^{-t}(4\pi t)^{-d/2}e^{-|x|^2/(4t)}\,\mathrm dt
\]
给出的非负等价类；允许这个积分在原点取无穷值。对 \(a>0\)，在几乎处处意义下令 \(E_a(x)=a^{d/2-1}E(\sqrt a\,x)\)。证明 \(E_a\) 是 \(a-\Delta\) 的可积基本解，求其变换和 \(L^1(\mathbb R^d)\) 范数。若 \(1\leq p\leq\infty\)、\(f\in L^p(\mathbb R^d;\mathbb C)\)，证明 \(\|E_a*f\|_p\leq a^{-1}\|f\|_p\)。
\end{exercise}
\begin{solution}
\Cref{b2:prop:05-bessel-fundamental}已证明 \(E\) 的质量为 \(1\)，且 \(\widehat E(\xi)=(1+4\pi^2|\xi|^2)^{-1}\)。尺度公式给出
\[
 \widehat E_a(\xi)=a^{-1}\widehat E(\xi/\sqrt a)
 =\dfrac1{a+4\pi^2|\xi|^2}.
\]
乘以符号 \(a+4\pi^2|\xi|^2\) 后得到 \(1\)，故为基本解。它非负，换元得 \(\|E_a\|_1=a^{-1}\|E\|_1=a^{-1}\)。最后用 Young 卷积不等式。这个估计允许任意 \(1\leq p\leq\infty\)，但不单凭此范数界宣称解具有任意阶强导数。
\end{solution}

% 自拟：离散测度的温和性及其不可微函数式解释。
\Needspace{6\baselineskip}
\begin{exercise}\label{b2:ex:05-polynomial-comb}
设 \(d\) 为正整数，\(c_k\in\mathbb C\)（\(k\in\mathbb Z^d\)），并存在与 \(k\) 无关的常数 \(C,M\geq0\)，使 \(|c_k|\leq C(1+|k|)^M\)。证明 \(T=\sum_{k\in\mathbb Z^d}c_k\delta_k\) 由绝对收敛配对定义温和分布，且截断 \(T_N=\sum_{|k|_\infty\leq N}c_k\delta_k\) 在 \(\mathcal S'(\mathbb R^d;\mathbb C)\) 中逐测试函数配对趋向 \(T\)，其中 \(N\) 遍历正整数，\(|k|_\infty=\max_j|k_j|\)。再在 \(d=1\) 时证明正测度
\[
 \mu=\sum_{k\in\mathbb Z}e^{|k|}\delta_k
\]
虽给出紧支集测试函数上的普通分布，却不能延拓为 \(\mathcal S'(\mathbb R;\mathbb C)\) 中的元素。可选取非负 \(\rho\in C_c^\infty(\mathbb R;\mathbb R)\)，使 \(\operatorname{supp}\rho\subset(-1/2,1/2)\)、\(\rho(0)=1\)，检查 \(\varphi_n(x)=e^{-n/2}\rho(x-n)\) 的全部 Schwartz 半范数与 \(\int\varphi_n\,\mathrm d\mu\)。
\end{exercise}
\begin{solution}
选整数 \(N>M+d\)。则
\[
 \sum_k|c_k\varphi(k)|
 \leq C p_N(\varphi)\sum_k(1+|k|)^{M-N}<\infty.
\]
最后的格点级数可按整数半径或二进环带比较体积：半径约 \(2^j\) 的点数至多为 \(C_d2^{jd}\)，所以所得几何级数收敛。这同时给出有限半范数界。对每个固定测试函数，绝对收敛级数的截断尾部趋零，正好是 \(\mathcal S'\) 收敛。

指数权测度在每个紧集上只有有限个格点，故局部有限；对有共同紧支集的测试函数，积分由该紧集上的有限质量乘函数的上确界控制，所以它确实定义普通分布。取题中 \(\rho\) 和正整数 \(n\)。对每个固定非负整数 \(m\)，\(\varphi_n\) 的支集在 \((n-1/2,n+1/2)\) 内，其各阶导数都是 \(e^{-n/2}\) 乘相应平移导数。因此
\[
 p_m(\varphi_n)
 \leq C_{m,\rho}\cdot e^{-n/2}(1+n)^m\longrightarrow0.
\]
故 \(\varphi_n\to0\) 于 \(\mathcal S(\mathbb R;\mathbb C)\)。另一方面，支集中只有整数 \(n\)，所以
\[
 \int\varphi_n\,\mathrm d\mu
 =e^n\cdot e^{-n/2}\rho(0)=e^{n/2}\longrightarrow\infty.
\]
若存在延拓到 \(\mathcal S\) 的连续泛函，它在这些紧支集测试函数上的值必须保持上述积分值，却又应由连续性趋于零，矛盾。这不只是说明某个增长判别不能套用，而是直接排除了连续延拓。
\end{solution}

% 自拟：移动Laplace极点并截断空间，求平滑壳层余项及其总质量；不重复正文未截断核的验证题。
\Needspace{6\baselineskip}
\begin{exercise}\label{b2:ex:05-cutoff-laplace-kernel}
设 \(d\in\{2,3\}\)、\(a\in\mathbb R^d\)、\(R>0\)。对 \(x\ne a\)，令
\[
 \Phi_{d,a}(x)=
 \begin{cases}
 -\dfrac1{2\pi}\log|x-a|,&d=2,\\[4pt]
 \dfrac1{4\pi|x-a|},&d=3.
 \end{cases}
\]
在 \(x=a\) 处任意赋一个有限值。取实值 \(\chi\in C_c^\infty(\mathbb R^d;\mathbb R)\)，使 \(\chi=1\) 于 \(B(0,1)\)、\(\chi=0\) 于 \(\{|x|\geq2\}\)，并令 \(\chi_R(x)=\chi((x-a)/R)\)、\(K(x)=\chi_R(x)\cdot\Phi_{d,a}(x)\)。
\begin{enumerate}
\item\label{b2:item:05-shifted-harmonic-flux} 证明 \(\Phi_{d,a}\) 在 \(\mathbb R^d\setminus\{a\}\) 上调和，且对每个 \(\varepsilon>0\)，
\[
 -\int_{\partial B(a,\varepsilon)}\partial_n\Phi_{d,a}\,\mathrm dS=1,
\]
其中 \(n\) 是小球的外单位法向。说明删去小球后，剩余区域在内边界处的外法向为何取相反方向。
\item\label{b2:item:05-cutoff-shell-remainder} 证明 \(K\in L^1(\mathbb R^d;\mathbb R)\) 且紧支集。对 \(x\ne a\) 令
\[
 q(x)=-2\nabla\chi_R(x)\cdot\nabla\Phi_{d,a}(x)
       -(\Delta\chi_R(x))\cdot\Phi_{d,a}(x).
\]
证明 \(q\) 在 \(a\) 附近补成零后属于 \(C_c^\infty(\mathbb R^d;\mathbb R)\)，支集包含于 \(\{\,x\mid R\leq|x-a|\leq2R\,\}\)，并通过删球积分和极限证明
\[
 -\Delta K=\delta_a+q\quad\text{于 }\mathcal S'(\mathbb R^d;\mathbb C).
\]
\item\label{b2:item:05-shell-mass} 证明 \(\int_{\mathbb R^d}q=-1\)。解释为何把基本解截成紧支集后，不能仍只留下一个点质量而没有额外余项。
\end{enumerate}
\end{exercise}
\begin{solution}
记 \(r=|x-a|\)、\(\sigma_1=2\pi\)、\(\sigma_2=4\pi\)。两种维数下径向导数均为
\[
 \Phi_{d,a}'(r)=-\frac1{\sigma_{d-1}r^{d-1}}.
\]
再求一次导数，便有 \(\Phi_{d,a}''+(d-1)\Phi_{d,a}'/r=0\)，所以原点移至 \(a\) 后，核仍在极点外调和。在球的外法向 \(n=e_r\) 下，通量为
\[
 \int_{\partial B(a,\varepsilon)}\partial_n\Phi_{d,a}\,\mathrm dS
 =-\frac1{\sigma_{d-1}\varepsilon^{d-1}}
      \cdot\sigma_{d-1}\varepsilon^{d-1}=-1.
\]
因此带负号的通量等于 \(1\)。对删球后的区域，穿过内边界向区域外走是朝向被删去的球心，故它的外法向为 \(-e_r\)，不是 \(e_r\)。

核在极点附近绝对可积：二维只需 \(\int_0^1r|\log r|\,\mathrm dr<\infty\)，三维只需 \(\int_0^1r\,\mathrm dr<\infty\)。乘上 \(\chi_R\) 后支集包含于 \(\overline{B(a,2R)}\)，故 \(K\in L^1\)。\(\chi_R\) 的所有正阶导数在 \(B(a,R)\) 内及 \(B(a,2R)\) 外都为零，因此 \(q\) 只在远离极点的壳层上可能非零，在那里各因子均光滑。补零后 \(q\) 全局光滑且紧支集。极点外使用普通乘积求导公式，得 \(-\Delta K=q\)；点质量则必须另从边界极限中算出。

取任意 \(\varphi\in C_c^\infty(\mathbb R^d;\mathbb C)\)，选足够大的外球包含 \(K\) 和 \(\varphi\) 的支集，再删去 \(\overline{B(a,\varepsilon)}\)，其中 \(0<\varepsilon<\min\{R,1\}\)。外球边界的项全为零。内球附近 \(K=\Phi_{d,a}\)，并且剩余区域的法向为 \(-e_r\)，Green 公式于是给出
\[
 \begin{aligned}
 &\int_{|x-a|>\varepsilon}K(x)(-\Delta\varphi(x))\,\mathrm dx\\
 &\quad=\int_{\mathbb R^d}q(x)\varphi(x)\,\mathrm dx
       +\int_{|x-a|=\varepsilon}
        \left(-\varphi\,\partial_r\Phi_{d,a}
                     +\Phi_{d,a}\partial_r\varphi\right)\,\mathrm dS.
 \end{aligned}
\]
第一项已可在全空间积分，因为 \(q\) 支撑于 \(|x-a|\geq R\)。边界上含 \(-\partial_r\Phi_{d,a}\) 的项正是 \(\varphi\) 的球面平均，趋于 \(\varphi(a)\)。另一项的模由 \(\|\nabla\varphi\|_\infty\) 乘 \(C\varepsilon|\log\varepsilon|\)（二维）或 \(C\varepsilon\)（三维）控制，故趋零。左侧由 \(K\) 的局部可积性趋于 \(\int K(-\Delta\varphi)\)，得到 \(-\Delta K=\delta_a+q\) 在紧支集测试函数上的等式。\(K\in L^1\)、\(q\in C_c^\infty\) 和 \(\delta_a\) 都定义温和分布，再由 \(C_c^\infty\) 在 \(\mathcal S\) 中稠密，将等式推广到 \(\mathcal S'\)。

最后取 \(\eta\in C_c^\infty(\mathbb R^d;\mathbb R)\)，在 \(\overline{B(a,2R)}\) 的一个邻域内等于 \(1\)。\(\Delta\eta\) 在 \(K\) 的支集上为零，而 \(\eta\) 在 \(q\) 的支集上及 \(a\) 点均为 \(1\)，故
\[
 0=\langle-\Delta K,\eta\rangle
  =\eta(a)+\int_{\mathbb R^d}q(x)\eta(x)\,\mathrm dx
  =1+\int_{\mathbb R^d}q.
\]
于是余项的总质量为 \(-1\)，恰好抵消点质量的总质量。空间截断保留了极点附近的单位通量，却必须在截断变化的壳层中产生这项补偿。
\end{solution}

% !TeX root = ../../Book2.tex
\chapter{Poisson 求和公式与 \texorpdfstring{\(\theta\)}{\AnalysisPdfTheta} 函数预备}
\label{b2:ch:06}
\BookChapterNumber{b2:ch:06}
% Hunter2001Applied Section11.11印刷319页实际图像核读：11.45与Example11.46。
% 本章固定单位周期及2π频率约定，完整展开周期化换序、尺度、格与对偶格。
\begin{chapterabstract}
Poisson 求和公式通过周期化函数的 Fourier 系数，把格点求和问题化为对偶格上的求和。本章从光滑周期化出发证明公式，先讨论缩放求和与数值求积，再通过 Gauss 函数得到尺度倒置与圆周热核的两种表达。带限采样与频率混叠单列为选读专题，PDE 路线可以跳过。实参数 \(\theta\) 级数随后把这种对称连接到格点计数和整数平方和。高维选读部分讨论格基、对偶格与基本区域体积，为第三卷研究 \(\theta\) 函数及解析数论作准备。
\end{chapterabstract}
\begin{learninggoals}
\begin{itemize}
\item\label{b2:item:06-goal-periodic}从周期化函数计算 Fourier 系数，逐项核对求和、积分和微分的交换条件。
\item\label{b2:item:06-goal-scale}写出带平移、调制和尺度的 Poisson 公式，解释格点间距与频率间距的倒数关系，分别控制数值求积的频率误差与空间尾部；选读\subsecref{b2:subsec:060105}进一步研究带限恢复和混叠。
\item\label{b2:item:06-goal-theta}由 Gauss 自对偶性推导 \(\theta\) 变换，使用指数尾在大小参数之间转换估计。
\item\label{b2:item:06-goal-count}区分平滑格点计数与球内硬计数，理解平方和生成函数及 Mellin 积分的作用。
\item\label{b2:item:06-goal-lattice}在选读部分计算对偶格与基本区域体积，保留高维公式中的行列式因子。
\end{itemize}
\end{learninggoals}

怎样用等间距的数值 \(f(nh)\) 估计整条实线上的积分？减小 \(h\) 会改善近似，但它也改变频率中被读出的点。对 Gauss 函数，这个误差能精确改写成另一个快速收敛的级数；对带限信号，离散点值甚至能够完全恢复函数。本章以周期化把这些问题放进同一个计算：空间中的格点求和，对应频率中的倒数格点求和。

\ProseParagraphBreak

本章同时使用\chapref{b2:ch:02}的周期系数唯一性与\chapref{b2:ch:05}的 Schwartz 变换不变性。Gauss 变换沿用\chapref{b2:ch:03}已经固定的常数；平方能量观点则帮助解释热核的频率表达。求和公式的证明同时使用离散展开与连续变换。

\ProseParagraphBreak

若以 PDE 为目标，先读\secref{b2:sec:0601}的周期化、Poisson 求和及\subsecref{b2:subsec:060104}的缩放公式，再读\secref{b2:sec:0602}的 Gauss 函数与热核。对离散测量如何恢复连续信号感兴趣的读者，可继续读\subsecref{b2:subsec:060105}的带限采样和混叠；Hermite 自对偶补例则展示更多变换下保持不变的函数。\secref{b2:sec:0603}的平方和计数与 Mellin 算术联系留给需要连接第三卷的读者。\secref{b2:sec:0604}将整数格推广到一般格，说明格的形状与对偶格怎样进入求和公式。John K. Hunter 与 Bruno Nachtergaele 的《Applied Analysis》\cite[Section 11.11, p. 319]{Hunter2001Applied}以周期化证明求和公式，并给出 theta 级数的例子；Fourier 分析在数论中的应用还见于 Elias M. Stein 与 Rami Shakarchi 的《Fourier Analysis: An Introduction》\cite{Stein2003Fourier}。

% !TeX root = ../../Book2.tex
\section{二阶线性方程}
\label{b2:sec:0801}

两个二阶方程可能只差一个符号，解的行为却截然不同。分类首先抓住最高阶导数的组合，把它看成每一点上的二次型。本节使用 \(z=(z_1,\ldots,z_n)\) 表示全部自变量；只有在讨论演化问题时才拆成 \(z=(t,x)\)。因此 \(n\) 不总等于空间维数。

\subsection{主部}
\label{b2:subsec:080101}

在开集 \(\Omega\subseteq\mathbb R^n\) 上考虑实系数算子
\begin{equation}\label{b2:eq:08-linear-operator}
 Lu=\sum_{i,j=1}^n a^{ij}(z)\partial_i\partial_ju
 +\sum_{i=1}^n b^i(z)\partial_i u+c(z)u.
\end{equation}
系数至少连续，经典解至少属于 \(C^2\)。一阶和零阶项称为低阶项；主部已在\secref{b2:sec:0104}引入，此处专指含二阶导数的和式，记为 \(L_2u\)。由于混合偏导相等，把 \(a^{ij}\) 换成 \((a^{ij}+a^{ji})/2\) 不改变算子。以下始终取实对称矩阵 \(A=(a^{ij})\)。

\ProseParagraphBreak

二维中习惯写为
\[
 L_2u=a u_{xx}+2b u_{xy}+c u_{yy},\quad
 A=\begin{pmatrix}a&b\\b&c\end{pmatrix}.
\]
混合项前面的 \(2\) 来自矩阵求和中 \((1,2)\) 和 \((2,1)\) 两项。若方程写成 \(u_{xx}+6u_{xy}+u_{yy}=0\)，则矩阵的非对角元为 \(3\)，不是 \(6\)。本节中的 \(a,b,c\) 是二维主部的临时系数，不与\eqref{b2:eq:08-linear-operator}中的向量低阶系数混用。

\ProseParagraphBreak

散度形式也有相同的最高阶矩阵。在 \(a^{ij}\in C^1\) 时，
\[
 \sum_{i,j}\partial_i(a^{ij}\partial_ju)
 =\sum_{i,j}a^{ij}\partial_i\partial_ju
 +\sum_{i,j}(\partial_i a^{ij})\partial_ju.
\]
它与非散度形式在经典分类上相同，但分部积分时低阶项不同；在粗糙系数或弱解问题中，两种写法更不能只凭主部相同便随意交换。

\subsection{主符号}
\label{b2:subsec:080102}

\begin{definition}\label{b2:def:08-principal-symbol}
设 \(n\geq1\)，\(\Omega\subset\mathbb R^n\) 为开集，实系数算子 \(L=\sum_{i,j}a^{ij}(z)\partial_{ij}+\sum_i b^i(z)\partial_i+c(z)\) 的系数均连续，且 \(A=(a^{ij})\) 为实对称矩阵场。对 \(z\in\Omega\)，二次齐次多项式
\[
 p(z,\xi)\coloneqq\xi^{\mathsf T}A(z)\xi
 =\sum_{i,j=1}^n a^{ij}(z)\xi_i\xi_j
 \quad(\xi\in\mathbb R^n)
\]
称为 \(L\) 的\term[zhufuhao]{主符号}[principal symbol]。本章不把 Fourier 微分规则中的 \(-4\pi^2\) 因子纳入这个几何记号。
\end{definition}
\symidx{\(p(z,\xi)\)：二阶算子的几何主符号}

后一约定与第一编的 \(2\pi\) 归一化相容：对常系数主部，
\[
 L_2e^{2\pi\mathrm i z\cdot\xi}
 =-4\pi^2p(\xi)\cdot e^{2\pi\mathrm i z\cdot\xi}.
\]
因此 \(p\) 记录不同频率方向的二阶作用，整个 Fourier 乘子则是 \(-4\pi^2p\)。不同文献可能把这个负号纳入符号定义；判断椭圆性和特征集合时，必须先核对其约定。

\ProseParagraphBreak

更一般地，对光滑实相位 \(\phi\) 与大参数 \(\lambda\)，直接两次求导给出
\[
 L\bigl(e^{2\pi\mathrm i\lambda\phi}\bigr)
 =e^{2\pi\mathrm i\lambda\phi}
 \left(-4\pi^2\lambda^2p(z,\nabla\phi)
 +2\pi\mathrm i\lambda\left(\sum_{i,j}a^{ij}\partial_{ij}\phi
 +\sum_i b^i\partial_i\phi\right)+c\right).
\]
在固定紧集上，最高的 \(\lambda^2\) 次项由主符号决定。若该系数为零，通常占主导的二阶作用恰好抵消；特征几何正由这类方向产生。这个计算只是代数恒等式，没有断言任意高频函数就是某个方程的近似解。

\subsection{特征余向量}
\label{b2:subsec:080103}

\begin{definition}\label{b2:def:08-characteristic-covector}
设 \(n\geq1\)，\(\Omega\subset\mathbb R^n\) 为开集，实系数算子 \(L=\sum_{i,j}a^{ij}(z)\partial_{ij}+\sum_i b^i(z)\partial_i+c(z)\) 的系数均连续，且 \(A=(a^{ij})\) 为实对称矩阵场。令 \(p(z,\nu)=\nu^{\mathsf T}A(z)\nu\)。在一点 \(z\in\Omega\)，若非零实余向量 \(\nu\in(\mathbb R^n)^*\) 满足 \(p(z,\nu)=0\)，则称它为该点的\term[tezheng-yuxiangliang]{特征余向量}[characteristic covector]。若实值函数 \(\phi\in C^2(V)\) 定义在开集 \(V\subseteq\Omega\) 上，在超曲面 \(\Sigma=\{z\in V:\phi(z)=0\}\) 上处处有 \(\nabla\phi\ne0\)，且 \(p(z,\mathrm d\phi(z))=0\)，则称 \(\Sigma\) 为\term[tezheng-chaqumian]{特征超曲面}[characteristic hypersurface]。在 \(\Sigma\) 上满足 \(p(z,\mathrm d\phi(z))\ne0\) 的点称为该曲面的非特征点。
\end{definition}

这里的 \(\nu\) 表示余向量，也就是对位移作线性求值的法向数据；借欧氏内积可用列向量表示，但不能把它直接当作曲面的切向。例如，\(\mathrm d\phi(z)\) 对切向量 \(w\) 的作用是 \(\mathrm d\phi(z)[w]=\nabla\phi(z)\cdot w\)，在切平面上等于零。下文计算主符号时，以梯度的列向量表示同一个微分。非零倍数 \(\alpha\nu\) 仍为特征余向量，所以也可以讨论它确定的、忽略非零实倍数的余法向方向；二维特征曲线的切方向则是另一个对象。

\ProseParagraphBreak

换一个定义函数 \(\widetilde\phi=\alpha\phi\)，其中 \(\alpha\) 是曲面附近不为零的 \(C^2\) 实函数，在曲面上有 \(\mathrm d\widetilde\phi=\alpha\,\mathrm d\phi\)。主符号齐次，零值条件不变。一般地，同一光滑超曲面的两个正则定义函数在曲面上的非零微分都有同一个核，即切平面，因而成非零比例；所以特征条件不依赖定义函数的选择。一个曲面也可能只有部分点为特征点，不能从一个点的条件推出整张曲面均为特征。

\ProseParagraphBreak

在二维中，若曲线切向为 \((\mathrm dx,\mathrm dy)\)，则其一个法向与 \((-\mathrm dy,\mathrm dx)\) 平行。特征条件因而写作
\begin{equation}\label{b2:eq:08-characteristic-differential}
 a(\mathrm dy)^2-2b\,\mathrm dx\,\mathrm dy+c(\mathrm dx)^2=0.
\end{equation}
当 \(\mathrm dx\ne0\) 时，它变为斜率方程 \(am^2-2bm+c=0\)。这解释了为什么这里混合项的符号与主符号 \(a\xi^2+2b\xi\eta+c\eta^2\) 相反：一个写切向，一个写法向。

\ProseParagraphBreak

非特征条件的代数意义是可以解出法向二阶导数。取坐标 \(y_1=\phi(z)\)，补足其余坐标后，\(\partial_{y_1}^2v\) 的系数就是 \(p(z,\nabla\phi)\)。该系数非零时，可以把它移到方程一侧并相除。但这一步只整理了方程，没有证明给定 Cauchy 数据的解存在，也没有得到稳定性。例如对 Laplace 方程，任意非零实法向都是非特征的，而给定函数值和法向导数的 Cauchy 问题仍可能非常不稳定。

\subsection{坐标不变性}
\label{b2:subsec:080104}

\begin{proposition}\label{b2:prop:08-coordinate-change}
设 \(n\geq1\)，\(\Omega,\widetilde\Omega\subset\mathbb R^n\) 为开集，实系数算子 \(L=\sum_{i,j}a^{ij}(z)\partial_{ij}+\sum_i b^i(z)\partial_i+c(z)\) 的系数均连续，且 \(A=(a^{ij})\) 为实对称矩阵场。设 \(F:\Omega\to\widetilde\Omega\) 为 \(C^2\) 微分同胚，\(\mathbb K\in\{\mathbb R,\mathbb C\}\)，\(v\in C^2(\widetilde\Omega;\mathbb K)\)，\(y=F(z)\)，\(u(z)=v(F(z))\)，并令 \(B=DF(z)\)。则新坐标中的主部矩阵为
\begin{equation}\label{b2:eq:08-congruence}
 \widetilde A(y)=B A(z)B^{\mathsf T},\quad z=F^{-1}(y),
\end{equation}
一阶系数为
\begin{equation}\label{b2:eq:08-lower-transform}
 \widetilde b^k(y)=\sum_{i,j}a^{ij}(z)\partial_{ij}F^k(z)
 +\sum_i b^i(z)\partial_iF^k(z).
\end{equation}
零阶系数为 \(\widetilde c(y)=c(F^{-1}(y))\)，主符号满足 \(\widetilde p(y,\eta)=p(z,B^{\mathsf T}\eta)\)。因此，对 \(f\in C(\Omega;\mathbb K)\)，方程 \(Lu=f\) 等价于
\[
 \sum_{k,l}\widetilde a^{kl}\partial_{kl}v
 +\sum_k\widetilde b^k\partial_kv+\widetilde c\,v
 =\widetilde f,
 \qquad \widetilde f=f\circ F^{-1}.
\]
\end{proposition}
\begin{proof}
链式法则给出
\[
 \partial_i u=\sum_k\partial_kv\cdot\partial_iF^k,
\quad
 \partial_{ij}u=\sum_{k,l}\partial_{kl}v\cdot\partial_iF^k\cdot\partial_jF^l
 +\sum_k\partial_kv\cdot\partial_{ij}F^k.
\]
代入算子并按 \(v\) 的导数阶数整理，得到两条系数公式。再把 \(\eta\) 代入二次型，便有
\(\eta^{\mathsf T}BAB^{\mathsf T}\eta=(B^{\mathsf T}\eta)^{\mathsf T}A(B^{\mathsf T}\eta)\)。
\end{proof}

实对称矩阵在可逆\term[hetong-bianhuan]{合同变换}[congruence transformation]下保持正、负及零特征值的个数，这三项称为其\term[guanxing]{惯性}[inertia]。为说明这一事实，先由有限维谱定理把矩阵正交对角化。二次型为正定的子空间最大维数恰等于正特征值个数：正特征子空间达到这个维数；任何更大子空间必与非正特征子空间非平凡相交，不可能正定。可逆线性变换保持子空间维数，因而保持该最大值。对负号二次型重复此论证可得负特征值个数不变，最后由总维数得到零特征值个数不变。

\ProseParagraphBreak

所以主部分类不依赖合法坐标，但矩阵的具体特征值和统一常数可能改变。设 \(A(z)\geq\lambda I\)，其中 \(\lambda>0\) 与 \(z\) 无关，则
\[
 \eta^{\mathsf T}\widetilde A(F(z))\eta
 \geq\lambda|B(z)^{\mathsf T}\eta|^2.
\]
若还知道 \(\sup_{z\in\Omega}\|B(z)^{-1}\|\leq M<\infty\)，由
\( |\eta|\leq\|B(z)^{-1}\|\,|B(z)^{\mathsf T}\eta|\)
得到 \(\widetilde A\geq\lambda M^{-2}I\)。这里均使用欧氏算子范数，且转置不改变该范数。逆导数的一致控制是从 \(A\geq\lambda I\) 这一已知下界推出新下界的充分条件；缺少它时，这一个下界本身不足以保证结论，仍须检查 \(A\) 与 \(B\) 的共同变化。

\ProseParagraphBreak

例如在 \(\mathbb R^2\) 上取
\[
 A(x,y)=\begin{pmatrix}1&0\\0&(1+y^2)^2\end{pmatrix},
 \qquad F(x,y)=(x,\arctan y).
\]
这是到 \(\mathbb R\times(-\pi/2,\pi/2)\) 的光滑微分同胚，且
\[
 B=\begin{pmatrix}1&0\\0&(1+y^2)^{-1}\end{pmatrix},
 \qquad \|B^{-1}\|=1+y^2,
 \qquad BAB^{\mathsf T}=I.
\]
逆导数虽然无界，原矩阵在同一方向上的增长却恰好抵消坐标压缩，新矩阵仍有统一正下界。反过来，若保持同一个 \(F\) 而取 \(A=I\)，新矩阵第二个特征值为 \((1+y^2)^{-2}\)，确实没有统一正下界。这两个计算同时说明充分条件的用途和它的适用边界。

\ProseParagraphBreak

还须区分统一正下界与\secref{b2:sec:0802}采用的双侧一致椭圆条件。前一个例子的原矩阵没有统一上界，不能称为该双侧意义下的一致椭圆矩阵。若已知 \(\lambda I\leq A\leq\Lambda I\)，再有 \(\sup\|B\|\leq N\)、\(\sup\|B^{-1}\|\leq M\)，则
\[
 \frac{\lambda}{M^2}|\eta|^2
 \leq\eta^{\mathsf T}\widetilde A\eta
 \leq\Lambda N^2|\eta|^2,
\]
从而双侧条件也被传递。点态坐标不变性与全域定量估计是两个层次。非线性坐标还会产生\eqref{b2:eq:08-lower-transform}中的一阶项，忽略它们会改变原方程。

% !TeX root = ../../Book2.tex
\section{Gauss 函数与自对偶性}
\label{b2:sec:0602}

Poisson 求和把一个函数的格点和转成其变换的格点和。Gauss 函数的变换仍为 Gauss 函数，因此可以显式比较大尺度和小尺度，并把这种比较放回热传播。

\subsection{Gauss 函数的 Fourier 变换}
\label{b2:subsec:060201}

记 \(g_t(x)=e^{-\pi t x^2}\)，其中 \(t>0\)。由\cref{b2:thm:03-gaussian}及伸缩公式，
\[
 \widehat {g_t}(\xi)=t^{-1/2}e^{-\pi\xi^2/t}.
\]
与保持单位质量的 \(G_\varepsilon\) 相比，\(g_t\) 在原点的高度固定为 \(1\)，总积分为 \(t^{-1/2}\)。将两类参数混用会漏掉平方根因子；这里分别用 \(t\) 和 \(\varepsilon\) 保留区别。

\ProseParagraphBreak

平移和调制给出更一般的波包公式。对实数 \(a,b\)，
\[
 \mathcal F\bigl(e^{-\pi t(x-a)^2}e^{2\pi\mathrm{i}bx}\bigr)(\xi)
 =t^{-1/2}e^{-2\pi\mathrm{i}a(\xi-b)}e^{-\pi(\xi-b)^2/t}.
\]
因此空间中心 \(a\) 产生相位，调制频率 \(b\) 产生频率中心，而宽度参数倒置。每一项都来自\chapref{b2:ch:03}已经证明的运算规则。

\subsection{自对偶结构}
\label{b2:subsec:060202}

对本卷 Fourier 变换已定义的函数或分布 \(f\)，若 \(\widehat f=f\) 在同一空间中成立，称 \(f\) 对当前 Fourier 变换\term[zi-duiou]{自对偶}[self-dual]。Gauss 函数 \(g_1\) 满足这一条件。一般 \(g_t\) 的参数变为 \(1/t\)，并附带 \(t^{-1/2}\)；只有 \(t=1\) 保持同一个归一化函数。

\ProseParagraphBreak

\subsecref{b2:subsec:050204}的 Hermite 函数提供了正 Gauss 函数之外的自对偶例子。\cref{b2:thm:05-hermite}给出 \(\widehat h_n=(-\mathrm i)^nh_n\)，故对每个 \(j\in\mathbb N_0\)，\(h_{4j}\) 都自对偶，有限复线性组合仍然自对偶。相反，\(h_{4j+2}\) 变换后变号，奇数指标对应虚特征值。自对偶是一项线性变换关系，不意味着非负性或概率密度性质。

\subsection{热核解释}
\label{b2:subsec:060203}

\chapref{b2:ch:05}已使用 \(H_t(x)=(4\pi t)^{-1/2}e^{-x^2/(4t)}\)，其变换为 \(e^{-4\pi^2t\xi^2}\)。在实线上周期化，得到
\begin{equation}\label{b2:eq:06-periodic-heat}
 K_t(x)\coloneqq\sum_{n\in\mathbb Z}H_t(x+n)
 =\sum_{k\in\mathbb Z}e^{-4\pi^2tk^2}e^{2\pi\mathrm{i}kx}.
\end{equation}
第二个表达由 Poisson 求和得到；对固定 \(t>0\)，两种表达都绝对收敛。函数 \(K_t\) 称为圆周上的\term[rehe]{热核}[heat kernel]。
\symidx{\(K_t\)：圆周热核}

\begin{proposition}\label{b2:prop:06-heat-kernel}
对 \(t,s>0\)，令 \(K_t(x)=\sum_{n\in\mathbb Z}(4\pi t)^{-1/2}e^{-(x+n)^2/(4t)}\)，在 \(\mathbb T=\mathbb R/\mathbb Z\) 上取总质量为 \(1\) 的 Lebesgue 测度，\(*\) 表示圆周卷积。圆周热核 \(K_t\) 非负、积分为 \(1\)，满足 \(K_t*K_s=K_{t+s}\) 和 \(\partial_tK_t=\partial_x^2K_t\)。对 \(f\in L^p(\mathbb T;\mathbb C)\)、\(1\leq p<\infty\)，当 \(t\downarrow0\) 时有 \(K_t*f\to f\) 于 \(L^p\)；对 \(f\in C(\mathbb T;\mathbb C)\) 则一致收敛。
\end{proposition}
\begin{proof}
非负性由实线表达给出，非负项的 Tonelli 换序将其圆周积分化为实线 Gauss 积分，故积分为 \(1\)。两个核均可积，圆周卷积的 Fourier 系数等于两核系数之积，即 \(e^{-4\pi^2(t+s)k^2}\)；由\cref{b2:cor:02-uniqueness}的周期系数唯一性及两侧连续性得到半群关系。对 \(t\geq\varepsilon>0\)，任意阶 \(x,t\) 导数只在系数上增加一个固定次幂的 \(k\)，仍被可和的指数尾统一控制，故可逐项求导，并得到热方程。

为证初值逼近，对 \(0<\delta<1/2\)，圆周上离原点至少为 \(\delta\) 的区域在周期化前对应到实线中 \(|y|\geq\delta\) 的子集。因此
\[
 \int_{\mathrm{d}_{\mathbb T}(x,0)\geq\delta}K_t(x)\,\mathrm{d}x
 \leq\int_{|y|\geq\delta}H_t(y)\,\mathrm{d}y\longrightarrow0.
\]
所以它满足\cref{b2:prop:02-positive-approximation}中的三项条件：非负、单位质量以及远离原点的质量趋零。该命题对任意趋零的正时间列都适用，因而得到 \(t\downarrow0\) 时的全部收敛结论。
\end{proof}

小时间时，实线表达用少量邻近格点就能描述局部尖峰；大时间时，频率表达只需保留少量低频，就能看出函数趋向常数 \(1\)。两个公式描述同一个核，却在相反的尺度下各自简洁。下一节的 \(\theta\) 变换正是这种尺度对换在 \(x=0\) 处的表达。

\ProseParagraphBreak

这种比较可以给出两端都可核算的误差。对 \(t\geq t_0>0\)，由 \(k^2\geq1+3(k-1)\)、\(k\geq1\)，频率表达给出
\[
 \|K_t-1\|_\infty
 \leq2\sum_{k=1}^\infty e^{-4\pi^2tk^2}
 \leq\frac{2e^{-4\pi^2t}}{1-e^{-12\pi^2t}}.
\]
因此向均匀分布趋近的首个指数由 \(k=\pm1\) 决定。对 \(f\in L^2(\mathbb T;\mathbb C)\)，记平均值 \(\overline f=\int_0^1f\)，这里上横线表示平均而非复共轭。\cref{b2:cor:02-parseval}将 \(L^2\) 范数平方写成全部 Fourier 系数模平方之和，给出精确的能量表达
\[
 \|K_t*f-\overline f\|_2^2
 =\sum_{k\ne0}e^{-8\pi^2tk^2}|\widehat f(k)|^2
 \leq e^{-8\pi^2t}\|f-\overline f\|_2^2.
\]
取 \(f(x)=e^{2\pi\mathrm{i}x}\) 时达到等号，所以能量范数中的衰减率 \(4\pi^2\) 不能统一增大。常数模态保留平均，其余模态耗散。

\ProseParagraphBreak

对于小时间，固定 \(0<a<1/2\) 并取 \(|x|\leq a\)。零格点项为 \(H_t(x)\)，其余项满足 \(|x+n|\geq|n|-a\)。对 \(n\geq1\)，
\[
 (n-a)^2\geq(1-a)^2+(3-2a)(n-1),
\]
故
\[
 0\leq K_t(x)-H_t(x)
 \leq\frac{2(4\pi t)^{-1/2}e^{-(1-a)^2/(4t)}}
              {1-e^{-(3-2a)/(4t)}}.
\]
当 \(0<t\leq1\) 时，分母有依赖 \(a\) 的统一正下界。于是邻近原点的圆周传播，与实线 Gauss 传播之间只有指数小的误差。若接近半周期端点，就可能有两个相邻格点同样重要；单独保留零格点的近似范围必须随之改变。

% !TeX root = ../../Book2.tex
\section{坐标变换}
\label{b2:sec:0803}

主符号的惯性告诉我们可以期待哪一种模型，坐标变换则把这一判断落实为计算。本节区分三种任务：常系数主部在整个坐标域中的标准化，变系数主部在一点的标准化，以及利用特征曲线在一个邻域中构造坐标。三者使用的条件不同，结论也不相同。

\subsection{二维标准型}
\label{b2:subsec:080301}

\begin{theorem}\label{b2:thm:08-constant-canonical}
设 \(a,b,c\) 为实常数且不全为零。存在可逆实线性坐标变换，并允许把整个方程乘以一个非零常数，使主部 \(a\partial_x^2+2b\partial_x\partial_y+c\partial_y^2\) 依照 \(D=b^2-ac\) 的符号分别成为
\[
 \partial_r^2+\partial_s^2\quad(D<0),\quad
 \partial_r^2-\partial_s^2\quad(D>0),\quad
 \partial_r^2\quad(D=0).
\]
低阶项须按同一个变换保留。双曲标准型还可用 \(\xi=r+s\)、\(\eta=r-s\) 化为 \(4\partial_\xi\partial_\eta\)。
\end{theorem}
\begin{proof}
先设 \(a\ne0\)，取 \(r=x\)、\(s=y-(b/a)x\)。此变换行列式为一，且
\[
 \partial_x=\partial_r-\frac ba\partial_s,\quad
 \partial_y=\partial_s.
\]
代入并展开，混合项的系数 \(-2b+2b\) 抵消，得到
\[
 a\partial_r^2+\delta\partial_s^2,\quad
 \delta=c-\frac{b^2}{a}=-\frac Da.
\]
若 \(\delta\ne0\)，再令 \(R=r/\sqrt{|a|}\)、\(S=s/\sqrt{|\delta|}\)，主部变为
\(\operatorname{sgn}(a)\partial_R^2+\operatorname{sgn}(\delta)\partial_S^2\)。乘以 \(\operatorname{sgn}(a)\)，第二项符号等于 \(-\operatorname{sgn}(D)\)，得到前两种形式。若 \(\delta=0\)，只缩放 \(r\) 即得第三种。

若 \(a=0,c\ne0\)，交换 \(x,y\) 后归入上一情形。若 \(a=c=0,b\ne0\)，取 \(r=(x+y)/\sqrt2\)、\(s=(x-y)/\sqrt2\)，则 \(2b\partial_x\partial_y=b(\partial_r^2-\partial_s^2)\)，再缩放并调整符号即可。最后，\(\partial_r=\partial_\xi+\partial_\eta\)、\(\partial_s=\partial_\xi-\partial_\eta\)，相减平方得到 \(4\partial_\xi\partial_\eta\)。
\end{proof}

例如主部 \(4u_{xx}+4u_{xy}+2u_{yy}\) 取 \(r=x/2\)、\(s=y-x/2\) 后成为 \(v_{rr}+v_{ss}\)。同样的变换把
\[
 4u_{xx}+4u_{xy}+u_{yy}+u_x=0
\]
化为 \(v_{rr}+v_r/2-v_s/2=0\)，即 \(v_s=2v_{rr}+v_r\)。二阶部分秩为一，而原来的一阶项恰好提供了新的时间导数。若删掉原式的 \(u_x\)，新方程只是 \(v_{rr}=0\)，便没有这一演化含义。

\ProseParagraphBreak

上述论证适用于常系数。若每一点分别对 \(A(x,y)\) 选一个对角化矩阵，并不能立即把这些矩阵拼成坐标变换的导数：一个候选矩阵场必须满足混合偏导的相容条件。特别是把位置依赖的向量场平方，会额外产生系数导数；不能把变系数算子当成常系数多项式直接配方而忽略这些项。

\subsection{特征坐标}
\label{b2:subsec:080302}

先处理一个双曲邻域。假定 \(a,b,c\in C^2\)、\(a\ne0\)、\(D=b^2-ac>0\)。二维特征斜率为
\[
 m_\pm=\frac{b\pm\sqrt D}{a}.
\]
两族曲线分别由 \(y'=m_\pm(x,y)\) 给出。\chapref{b2:ch:07}的局部流与初值可微依赖保证：以横截直线 \(x=x_0\) 上的交点作为标签，可构造局部函数 \(\xi,\eta\)，使
\[
 \xi_x+m_+\xi_y=0,\quad
 \eta_x+m_-\eta_y=0,\quad \xi_y\ne0,\ \eta_y\ne0.
\]
它们分别在两族特征线上保持常值，称为\term[tezheng-zuobiao]{特征坐标}[characteristic coordinates]。系数为 \(C^2\) 时，流也是 \(C^2\)：对其一阶变分方程再求一次初值导数，所得二阶导数满足线性方程，其非齐次项为速度的二阶导数作用于两个一阶变分；局部一致连续性控制差商余项，再用 Gronwall 不等式得到导数存在且连续。因此这里的坐标变换符合二阶链式法则的要求。

\begin{proposition}\label{b2:prop:08-hyperbolic-coordinates}
上述 \((\xi,\eta)\) 在充分小的邻域内构成 \(C^2\) 微分同胚，并把方程化为
\[
 \gamma(\xi,\eta)v_{\xi\eta}
 +\beta_1(\xi,\eta)v_\xi+\beta_2(\xi,\eta)v_\eta
 +\kappa(\xi,\eta)v=\widetilde f,
 \quad \gamma\ne0.
\]
\end{proposition}
\begin{proof}
坐标的行列式为
\[
 \xi_x\eta_y-\xi_y\eta_x
 =(m_--m_+)\xi_y\eta_y\ne0,
\]
因此可用逆函数定理。由特征方程，\(p(\nabla\xi)=p(\nabla\eta)=0\)，新矩阵的两个对角元为零。其混合系数为
\[
 \gamma=2\bigl(a\xi_x\eta_x+b(\xi_x\eta_y+\xi_y\eta_x)+c\xi_y\eta_y\bigr).
\]
合同变换保持矩阵可逆，而 \(\left(\begin{smallmatrix}0&\gamma/2\\\gamma/2&0\end{smallmatrix}\right)\) 可逆恰好要求 \(\gamma\ne0\)。一阶与零阶系数按\cref{b2:prop:08-coordinate-change}保留，即得所述方程。
\end{proof}

例如在 \(x>0\) 上，方程 \(u_{xx}-x^2u_{yy}=0\) 的两族特征为 \(y\pm x^2/2\) 恒定。取 \(\xi=y+x^2/2\)、\(\eta=y-x^2/2\)，其行列式为 \(2x\ne0\)。两次求导得到
\[
 u_{xx}-x^2u_{yy}=-4x^2v_{\xi\eta}+v_\xi-v_\eta.
\]
新方程不只是 \(v_{\xi\eta}=0\)；两个一阶项正是坐标曲率带来的贡献。在 \(x=0\) 上行列式退化，所以不能用同一坐标穿过那里。

\ProseParagraphBreak

若 \(D=0\) 在整个邻域成立，且 \(a\ne0\)，令 \(\beta=b/a\)、\(V=\partial_x+\beta\partial_y\)，则
\[
 a\partial_x^2+2b\partial_x\partial_y+c\partial_y^2
 =aV^2-a(\partial_x\beta+\beta\partial_y\beta)\partial_y.
\]
取 \(\xi=x\)，并沿 \(V\) 的流构造 \(\eta\) 使 \(V\eta=0\)、\(\eta_y\ne0\)，则 \(V=\partial_\xi\)。方程成为
\[
 a\,v_{\xi\xi}+\beta_1v_\xi+\beta_2v_\eta+\kappa v=\widetilde f.
\]
这是邻域中的秩一标准形式。只有 \(\beta_2\ne0\) 时才能解出 \(v_\eta\)，要按 \(\eta\) 正方向解释为向前热型，还须 \(-a/\beta_2>0\)。秩为一的二阶部分本身没有决定这些条件。

\ProseParagraphBreak

椭圆型没有两族实特征线，因此上述沿实轨道找坐标的方法不适用。一般变系数椭圆主部的邻域标准化需要另解坐标方程，本章不把它当作已证工具；本节已经完整建立的是常系数标准型及一点处冻结系数的标准型。某些特殊系数可直接处理，例如在全平面上取 \(r=e^x\)、\(s=y\)，有
\[
 u_{xx}+e^{2x}u_{yy}=r^2(v_{rr}+v_{ss})+r v_r.
\]
这个映射是全平面到半平面的微分同胚。除以 \(r^2\) 后主部化为 Laplace 算子，但保留了 \(r^{-1}v_r\)；所获表达式的定义域也随坐标改变。

\subsection{高维主符号}
\label{b2:subsec:080303}

在 \(n\) 维，实对称矩阵可经正交变换化为 \(\operatorname{diag}(\lambda_1,\ldots,\lambda_n)\)，再沿非零特征方向缩放，使非零对角元都成为 \(1\) 或 \(-1\)。因此常系数主部完全由惯性三元组 \((n_+,n_-,n_0)\) 决定，允许整体乘以负数时还可交换前两项。这个结论包含二维的三种标准型，但高维中还有更多可能。

\ProseParagraphBreak

例如 \(\partial_1^2+\partial_2^2-\partial_3^2-\partial_4^2\) 的惯性为 \((2,2,0)\)。它既不椭圆，也不属于只有一个异号方向的双曲类型，更不是抛物型。又如三变量算子 \(\partial_1^2\) 的惯性为 \((1,0,2)\)，有两个完全缺失的二阶方向。对这些算子只说“判别式非正”没有明确含义；一个行列式也不能恢复全部惯性信息。

\ProseParagraphBreak

变系数时，可以在任意固定点 \(z_0\) 选一个线性坐标使 \(A(z_0)\) 标准化。邻近点的矩阵通常只是接近标准矩阵。若该点主部非退化，连续性使惯性在足够小邻域内保持不变；理由是单位球上可逆矩阵的最小伸缩量有正下界，小扰动不可能使某个特征值穿过零。秩退化情形没有这样的稳定性，任意小的系数扰动都可能改变类型。

\subsection{几何意义}
\label{b2:subsec:080304}

主符号作用在法向上，特征曲线却画在自变量空间里。两者之间的联系可以用一个直接计算辨明。若 \(\nu\ne0\)、\(p(z,\nu)=0\)，则
\[
 \nu\cdot A(z)\nu=0.
\]
所以当 \(A\nu\ne0\) 时，向量 \(A\nu\) 位于法向为 \(\nu\) 的切平面中。对于非退化双曲型，\(A\nu\) 必非零；对于秩退化类型，它可能恰为零。二阶主符号不仅记录符号，也记录哪些方向上二阶约束消失。

\ProseParagraphBreak

在一维波动方程中，特征曲线为 \(x\pm ct\) 恒定，两族曲线把初值上的区间与后来一点联系起来，\chapref{b2:ch:09}将据此推导表示公式。高维波动方程对应实特征锥，\chapref{b2:ch:11}进一步研究其传播区域。热方程的时间方向由一阶项选定，不能把秩退化的零方向画成波动锥；Laplace 方程则没有非零实特征法向。坐标标准化明确了这些差别，具体解的性质仍由后续各章在相应模型中建立。

% !TeX root = ../../Book2.tex
\OptionalSection{高维 Poisson 求和}{b2:sec:0604}

整数格的直角坐标并非本质要求。任意全秩格都可以由可逆线性映射得到，频率侧相应出现逆转置映射。除了点的位置，还必须追踪基本区域的体积。

\subsection{格与对偶格}
\label{b2:subsec:060401}

\begin{definition}\label{b2:def:06-lattice}
设 \(d\geq1\) 为整数。若 \(A\) 是可逆实 \(d\times d\) 矩阵，则集合 \(\Lambda=A\mathbb Z^d\) 称为 \(\mathbb R^d\) 中的\term[ge]{格}[lattice]，本节均指全秩格。集合 \(A[0,1)^d\) 为一个基本区域，其体积
\[
 v(\Lambda)\coloneqq|\det A|
\]
称为格的\term[jiben-quyu-tiji]{基本区域体积}[covolume]。定义\term[duiou-ge]{对偶格}[dual lattice]为
\[
 \Lambda^*\coloneqq\{\,\xi\in\mathbb R^d:\xi\cdot v\in\mathbb Z\quad(\forall v\in\Lambda)\,\}.
\]
\end{definition}
\symidx{\(\Lambda^*\)：对偶格}
\symidx{\(v(\Lambda)\)：格的基本区域体积}

\begin{proposition}\label{b2:prop:06-dual-lattice}
设 \(d\geq1\) 为整数，\(A\) 是可逆实 \(d\times d\) 矩阵，\(\Lambda=A\mathbb Z^d\)，\(\Lambda^*=\{\,\xi\in\mathbb R^d:\xi\cdot v\in\mathbb Z\ (v\in\Lambda)\,\}\)，\(v(\Lambda)=|\det A|\)。则 \(\Lambda^*=A^{-\mathsf T}\mathbb Z^d\)、\((\Lambda^*)^*=\Lambda\) 和 \(v(\Lambda^*)=v(\Lambda)^{-1}\)。基本区域体积不依赖格基的选择。
\end{proposition}
\begin{proof}
条件 \(\xi\cdot An\in\mathbb Z\) 对全部整数向量 \(n\) 成立，等价于 \(A^{\mathsf T}\xi\) 的每个坐标都是整数，所以得到对偶格公式。再次取逆转置和行列式便得其余两个等式。若另一基矩阵为 \(B\)，则 \(U=A^{-1}B\) 及 \(U^{-1}\) 都是整数矩阵，因此其行列式为互为倒数的整数，只能为 \(\pm1\)。故 \(|\det B|=|\det A|\)。
\end{proof}

对角矩阵给出矩形格，边长分别为正数 \(a_j\) 时，对偶边长为 \(a_j^{-1}\)。非正交基也遵循逆转置，而不是把每个基向量分别除以其长度；内积的整数条件才是定义。

\Cref{b2:fig:06-dual-lattice}用一个非正交格显示逆转置的作用。箭头端点的内积满足 \(b_i\cdot b_j^*=\delta_{ij}\)，不能用逐个基向量取倒长度代替这一条件。
% 本书原创：同一欧氏比例绘出非正交格与逆转置对偶格。
\begin{figure}[htbp]
\centering
\begin{tikzpicture}[x=1cm,y=1cm,>=stealth,every node/.style={font=\small},line width=.65pt]
\begin{scope}[xshift=0.0cm]
\fill[AnalysisBlue!10] (0.00000,0.00000)--(1.50000,0.00000)--(2.00000,1.00000)--(0.50000,1.00000)--cycle;
\draw[->,gray] (-1.2,0)--(3.4,0);\draw[->,gray] (0,-1.2)--(0,2.6);
\fill[gray] (-1.00000,1.00000) circle(1.15pt);
\fill[gray] (-0.50000,2.00000) circle(1.15pt);
\fill[gray] (-0.50000,-1.00000) circle(1.15pt);
\fill[gray] (0.00000,0.00000) circle(1.15pt);
\fill[gray] (0.50000,1.00000) circle(1.15pt);
\fill[gray] (1.00000,2.00000) circle(1.15pt);
\fill[gray] (1.00000,-1.00000) circle(1.15pt);
\fill[gray] (1.50000,0.00000) circle(1.15pt);
\fill[gray] (2.00000,1.00000) circle(1.15pt);
\fill[gray] (2.50000,2.00000) circle(1.15pt);
\fill[gray] (2.50000,-1.00000) circle(1.15pt);
\fill[gray] (3.00000,0.00000) circle(1.15pt);
\draw[->,AnalysisBlue,thick] (0,0)--(1.50000,0.00000) node[above right] {\(b_1\)};
\draw[->,AnalysisBlue,thick] (0,0)--(0.50000,1.00000) node[above right] {\(b_2\)};
\node at (1,3.0) {\(\Lambda,\quad v(\Lambda)=3/2\)};\end{scope}
\begin{scope}[xshift=6.6cm]
\fill[AnalysisBlue!10] (0.00000,0.00000)--(0.66667,-0.33333)--(0.66667,0.66667)--(0.00000,1.00000)--cycle;
\draw[->,gray] (-1.2,0)--(3.4,0);\draw[->,gray] (0,-1.2)--(0,2.6);
\fill[gray] (-0.66667,-0.66667) circle(1.15pt);
\fill[gray] (-0.66667,0.33333) circle(1.15pt);
\fill[gray] (-0.66667,1.33333) circle(1.15pt);
\fill[gray] (0.00000,-1.00000) circle(1.15pt);
\fill[gray] (0.00000,0.00000) circle(1.15pt);
\fill[gray] (0.00000,1.00000) circle(1.15pt);
\fill[gray] (0.00000,2.00000) circle(1.15pt);
\fill[gray] (0.66667,-0.33333) circle(1.15pt);
\fill[gray] (0.66667,0.66667) circle(1.15pt);
\fill[gray] (0.66667,1.66667) circle(1.15pt);
\fill[gray] (1.33333,-0.66667) circle(1.15pt);
\fill[gray] (1.33333,0.33333) circle(1.15pt);
\fill[gray] (1.33333,1.33333) circle(1.15pt);
\fill[gray] (2.00000,-1.00000) circle(1.15pt);
\fill[gray] (2.00000,0.00000) circle(1.15pt);
\fill[gray] (2.00000,1.00000) circle(1.15pt);
\fill[gray] (2.00000,2.00000) circle(1.15pt);
\fill[gray] (2.66667,-0.33333) circle(1.15pt);
\fill[gray] (2.66667,0.66667) circle(1.15pt);
\fill[gray] (2.66667,1.66667) circle(1.15pt);
\draw[->,AnalysisBlue,thick] (0,0)--(0.66667,-0.33333) node[above right] {\(b_1^*\)};
\draw[->,AnalysisBlue,thick] (0,0)--(0.00000,1.00000) node[above right] {\(b_2^*\)};
\node at (1,3.0) {\(\Lambda^*,\quad v(\Lambda^*)=2/3\)};\end{scope}
\end{tikzpicture}
\caption[非正交格与对偶格]{非正交格与对偶格。左侧基矩阵为 \(\left(\begin{smallmatrix}3/2&1/2\\0&1\end{smallmatrix}\right)\)，右侧取逆转置；淡蓝色区域为相应基本平行四边形。两图采用同一长度单位，区域面积互为倒数。}
\label{b2:fig:06-dual-lattice}
\end{figure}


\subsection{欧氏空间中的求和公式}
\label{b2:subsec:060402}

\begin{theorem}[全秩格上的 Poisson 求和]\label{b2:thm:06-lattice-poisson}
设 \(d\geq1\) 为整数，\(f\in\mathcal S(\mathbb R^d;\mathbb C)\)，\(\Lambda=A\mathbb Z^d\)，其中 \(A\) 是可逆实 \(d\times d\) 矩阵。令 \(v(\Lambda)=|\det A|\)，\(\Lambda^*=A^{-\mathsf T}\mathbb Z^d\)。则
\[
 \sum_{v\in\Lambda}f(v)
 =\dfrac1{v(\Lambda)}\sum_{\xi\in\Lambda^*}\widehat f(\xi).
\]
两侧都绝对收敛。
\end{theorem}
\begin{proof}
先取 \(\Lambda=\mathbb Z^d\)。逐坐标应用一维 Poisson 公式；固定其他坐标时，一维切片为 Schwartz 函数。为核对每步的共同控制，先看第一个坐标的部分变换
\[
 g(\eta,x')=\int_{\mathbb R}f(x_1,x')e^{-2\pi\mathrm{i}x_1\eta}\,\mathrm{d}x_1.
\]
对 \(\eta\) 求导在积分内产生 \(x_1\) 的幂；乘以 \(\eta\) 的幂可改成 \(x_1\) 方向的导数，再分部积分；对 \(x'\) 的导数与权重则直接保留。这样每个目标加权导数都是某个 \(f\) 的加权导数在 \(x_1\) 上的积分。用足够高的原半范数，可使被积量由 \(C(1+|x_1|)^{-2}\) 控制，同时保留任意指定阶的 \(\eta,x'\) 权重。因此 \(g\) 仍为全部 \(d\) 个变量上的 Schwartz 函数。后续坐标变换重复同一论证。

在每一个中间阶段，进一步取足够高的权重，便可用 \(C\prod_j(1+|y_j|)^{-2}\) 控制混合变量函数；\(y_j\) 可以是待积分的实变量，也可以是已取样的整数指标。连续积分和离散级数均可和，故 Fubini 与绝对收敛重排适用。每次一维求和只把一个坐标的整数采样移到该坐标的频率侧，依次得到 \(\sum_{n\in\mathbb Z^d}f(n)=\sum_{k\in\mathbb Z^d}\widehat f(k)\)。

一般情形令 \(g(x)=f(Ax)\)，则 \(g\in\mathcal S\)，且 \(\widehat g(k)=|\det A|^{-1}\cdot\widehat f(A^{-\mathsf T}k)\)。对 \(g\) 使用整数格公式，即得结论。绝对收敛由可逆映射的长度比较及任意高阶的快速下降保证。
\end{proof}

基本区域体积因子不能省略。例如在一维取 \(h>0\)、\(\Lambda=h\mathbb Z\)，对偶格为 \(h^{-1}\mathbb Z\)，公式恰好恢复 \(\sum_nf(hn)=h^{-1}\sum_k\widehat f(k/h)\)。高维中的行列式就是这个因子的体积版本。

\subsection{格的 \texorpdfstring{\(\theta\)}{\AnalysisPdfTheta} 级数}
\label{b2:subsec:060403}

\begin{definition}\label{b2:def:06-lattice-theta}
设 \(d\geq1\) 为整数，\(\Lambda\subset\mathbb R^d\) 为全秩格。对 \(t>0\)，定义它的\term[ge-theta-jishu]{格 \(\theta\) 级数}[lattice theta series]为
\[
 \Theta_\Lambda(t)\coloneqq\sum_{v\in\Lambda}e^{-\pi t|v|^2}.
\]
\end{definition}
\symidx{\(\Theta_\Lambda\)：全秩格的 \(\theta\) 级数}

\begin{corollary}\label{b2:cor:06-lattice-theta}
设 \(d\geq1\) 为整数，\(\Lambda=A\mathbb Z^d\)，其中 \(A\) 是可逆实 \(d\times d\) 矩阵，\(\Lambda^*=A^{-\mathsf T}\mathbb Z^d\) 为其对偶格，\(v(\Lambda)=|\det A|\)。对全秩格 \(L\subset\mathbb R^d\)，记 \(\Theta_L(u)=\sum_{v\in L}e^{-\pi u|v|^2}\)（\(u>0\)）。这些级数收敛，且对每个 \(t>0\)，
\[
 \Theta_\Lambda(t)=v(\Lambda)^{-1}t^{-d/2}\Theta_{\Lambda^*}(1/t).
\]
\end{corollary}
\begin{proof}
在\cref{b2:thm:06-lattice-poisson}中取 \(f(x)=e^{-\pi t|x|^2}\)。对 \(t>0\)，这个函数属于 \(\mathcal S(\mathbb R^d)\)，其变换为 \(t^{-d/2}e^{-\pi|\xi|^2/t}\)。该定理保证两侧级数绝对收敛，代入即得。
\end{proof}

若 \(\Lambda=\mathbb Z^d\)，则对 \(t>0\) 有 \(\Theta_\Lambda(t)=\vartheta(t)^d\)。一般的变换公式同时涉及倒数参数、对偶格及基本区域体积。满足 \(\Lambda^*=\Lambda\) 的格称为\term[ziduiou-ge]{自对偶格}[self-dual lattice]；此时由体积互反可知 \(v(\Lambda)=1\)，两侧直接成为同一个格的级数。自对偶是做到这一点的充分条件。

\ProseParagraphBreak

更一般地，设 \(d\geq1\) 为整数，\(\Lambda\subset\mathbb R^d\) 为全秩格，\(c>0\)，\(U\) 是实 \(d\times d\) 正交矩阵，即 \(U^{\mathsf T}U=I_d\)，并假设 \(\Lambda^*=cU\Lambda\)。正交变换保持长度，伸缩将平方长度乘以 \(c^2\)，所以对 \(u>0\)，由格点的一一对应得到
\[
 \Theta_{\Lambda^*}(u)
 =\sum_{v\in\Lambda}e^{-\pi u c^2|Uv|^2}
 =\Theta_\Lambda(c^2u).
\]
体积也受这一假设约束：由 \(|\det U|=1\)，
\[
 v(\Lambda)^{-1}=v(\Lambda^*)=c^dv(\Lambda),
 \quad v(\Lambda)=c^{-d/2}.
\]
代入\cref{b2:cor:06-lattice-theta}，对每个 \(t>0\) 得
\[
 \Theta_\Lambda(t)
 =v(\Lambda)^{-1}t^{-d/2}\Theta_\Lambda(c^2/t)
 =c^{d/2}t^{-d/2}\Theta_\Lambda(c^2/t).
\]
当 \(c=1\) 时，只要对偶格与原格正交等距，就有 \(\Theta_{\Lambda^*}(t)=\Theta_\Lambda(t)\) 和 \(v(\Lambda)=1\)，从而得到同一格的倒数参数公式；不要求两个格作为点集相同。

\ProseParagraphBreak

例如二维矩形格
\[
 \Lambda=2\mathbb Z\times\tfrac12\mathbb Z,
 \quad \Lambda^*=\tfrac12\mathbb Z\times2\mathbb Z
\]
并不相同，因为 \((0,1/2)\in\Lambda\) 而 \((0,1/2)\notin\Lambda^*\)。交换坐标的正交矩阵 \(U=\left(\begin{smallmatrix}0&1\\1&0\end{smallmatrix}\right)\) 将 \(\Lambda\) 映到 \(\Lambda^*\)，且 \(v(\Lambda)=2\cdot(1/2)=1\)。因此虽然这个格不自对偶，仍有
\[
 \Theta_\Lambda(t)=t^{-1}\Theta_\Lambda(1/t)\quad(t>0).
\]

对可逆实 \(d\times d\) 基矩阵 \(A\)，若用 Gram 矩阵 \(Q=A^{\mathsf T}A\) 表示长度，则对每个 \(n\in\mathbb Z^d\)，
\[
 |An|^2=n^{\mathsf T}Qn,
 \quad |A^{-\mathsf T}n|^2=n^{\mathsf T}Q^{-1}n.
\]
整数换基作用在系数 \(n\in\mathbb Z^d\) 上，欧式正交变换作用在格向量 \(An\in\mathbb R^d\) 上。两者作用的对象不同；即使矩阵写法相近，也须区分。这为第三卷的二次型 \(\theta\) 级数提供几何来源，但不要求在这里提前引入整数性或模形式的附加结构。

\ProseParagraphBreak

作为非直角的可算例子，取
\[
 A=\begin{pmatrix}\sqrt2&1/\sqrt2\\0&\sqrt{3/2}\end{pmatrix},
 \quad Q=A^{\mathsf T}A=\begin{pmatrix}2&1\\1&2\end{pmatrix}.
\]
格基向量的夹角为 \(\pi/3\)，基本区域体积为 \(\sqrt3\)。对 \(t>0\)，记
\[
 B(t)=\sum_{m,n\in\mathbb Z}e^{-2\pi t(m^2+mn+n^2)}.
\]
这是该格的 \(\theta\) 级数。取整数基变换矩阵
\[
 P=\begin{pmatrix}1&0\\0&-1\end{pmatrix}.
\]
它满足 \(P\mathbb Z^2=\mathbb Z^2\)，并且
\[
 Q^{-1}=\frac13\begin{pmatrix}2&-1\\-1&2\end{pmatrix}
 =\frac13P^{\mathsf T}QP.
\]
因此对 \(u>0\)，把整数系数 \(k=(m,n)^{\mathsf T}\) 换成 \(Pk=(m,-n)^{\mathsf T}\)，得到
\[
 \sum_{k\in\mathbb Z^2}e^{-\pi u k^{\mathsf T}Q^{-1}k}
 =\sum_{k\in\mathbb Z^2}e^{-\pi(u/3)(Pk)^{\mathsf T}Q(Pk)}
 =B(u/3).
\]
Poisson 公式于是给出
\[
 B(t)=\frac1{\sqrt3\,t}
 \sum_{m,n\in\mathbb Z}e^{-2\pi(m^2-mn+n^2)/(3t)}
 =\frac1{\sqrt3\,t}B(1/(3t)).
\]
这里参数中的 \(1/3\) 来自二次型的尺度，而 \(P\) 只负责将整数系数重新编号。要在实际欧式空间中写出对应的正交变换，可取
\[
 U_0=\begin{pmatrix}\sqrt3/2&-1/2\\-1/2&-\sqrt3/2\end{pmatrix}.
\]
直接相乘可核对 \(U_0^{\mathsf T}U_0=I_2\) 及 \(A^{-\mathsf T}=3^{-1/2}U_0AP\)。因 \(P\mathbb Z^2=\mathbb Z^2\)，这就说明 \(\Lambda^*=3^{-1/2}U_0\Lambda\)。相似尺度 \(c=1/\sqrt3\) 给出 \(c^2/t=1/(3t)\)，体积关系给出 \(v(\Lambda)=c^{-1}=\sqrt3\)，与上式的两个因子分别相符。

\ProseParagraphBreak

最短非零向量由 \(m^2+mn+n^2=1\) 给出。写成 \((m+n/2)^2+3n^2/4\)，可先知 \(|n|\leq1\)，再逐个求出六组解 \((\pm1,0),(0,\pm1),(1,-1),(-1,1)\)。为明确其余项的统一误差，记 \(q(m,n)=m^2+mn+n^2\)。其余非零向量满足整数条件 \(q(m,n)\geq2\)，并且 \(q(m,n)\geq(m^2+n^2)/2\)。故对 \(t\geq1\)，
\[
 \sum_{q(m,n)\geq2}e^{-2\pi t q(m,n)}
 \leq e^{-4\pi t}\sum_{q(m,n)\geq2}e^{-2\pi(q(m,n)-2)}.
\]
右端的固定级数由二维 Gauss 级数控制而收敛。因此
\[
 B(t)=1+6e^{-2\pi t}+O(e^{-4\pi t})\quad(t\to\infty).
\]
把这一式用于倒数参数，得到
\[
 B(t)=\frac1{\sqrt3\,t}
 \left(1+6e^{-2\pi/(3t)}+O(e^{-4\pi/(3t)})\right)
 \quad(t\downarrow0).
\]
大参数修正读取最短格向量，小参数修正读取最短对偶格向量；主项读取基本区域体积。这三种几何数据在公式中各有确定位置。


\begin{chaptersummary}
周期化把实线上的函数放到圆周上，其 Fourier 系数就是连续变换在整数频率处的值。Poisson 求和从这个事实与系数唯一性得出；尺度变化产生倒数格与体积因子，不能省去对采样点和代表的正则性要求。

\ProseParagraphBreak

\subsecref{b2:subsec:060104}的数值求积把无穷采样误差与有限窗口的尾部误差分开。\subsecref{b2:subsec:060105}另外利用有限频带上的 \(L^2\) 展开研究精确恢复与混叠；这条选读路线可以略过，不影响周期热核及下一编 PDE 内容的学习。

\ProseParagraphBreak

Gauss 变换把参数取倒数，使 \(\theta\) 级数在大小尺度间交换。相同的对称既解释圆周热核的局部尖峰和低频平衡，也把整数平方和与 Mellin 积分联系起来。第三编将在有限区间上使用特征函数展开，在全空间上使用欧氏变换，并用能量范数和分布方程解释低正则解。
\end{chaptersummary}

\BookExercises

题组围绕一维求和、数值求积、周期化的频率采样与平方和计数展开，选做题使用一般格或参数分析。其中的频率采样通过周期化的 Fourier 系数计算；最后的跨章题把周期化与实线、圆周上的 Gauss 平滑联系起来。所有交换极限与求和的步骤均应说明所用的统一控制。

% 自拟：调制Gauss后得到交错theta，避免从复参数模公式直接摘用。
\Needspace{6\baselineskip}
\begin{exercise}\label{b2:ex:06-alternating-theta}
对 \(t>0\)，证明
\[
 \sum_{n\in\mathbb Z}(-1)^ne^{-\pi tn^2}
 =t^{-1/2}\sum_{k\in\mathbb Z}e^{-\pi(k-1/2)^2/t}>0.
\]
求 \(t\downarrow0\) 时的首项，并给出相对误差的指数阶。
\end{exercise}
\begin{solution}
在带调制的 Poisson 公式\cref{b2:cor:06-scaled-poisson}中取 \(f(x)=e^{-\pi tx^2}\)、\(h=1,a=0,b=1/2\)，得到等式；右侧每项为正，故严格为正。距离 \(1/2\) 最近的整数为 \(0,1\)，贡献 \(2t^{-1/2}e^{-\pi/(4t)}\)。将全部半整数距离成对合并，原和式可写成
\[
 2t^{-1/2}e^{-\pi/(4t)}
 \left(1+\sum_{j=1}^\infty e^{-\pi j(j+1)/t}\right).
\]
因为 \(j(j+1)\geq2+4(j-1)\)、\(j\geq1\)，括号中除 \(1\) 外的尾部介于 \(0\) 与 \(e^{-2\pi/t}/(1-e^{-4\pi/t})\) 之间。因此结果为 \(2t^{-1/2}e^{-\pi/(4t)}(1+O(e^{-2\pi/t}))\)。
\end{solution}

% 自拟：theta的自反尺度给出一个精确的带权求和。
\Needspace{6\baselineskip}
\begin{exercise}\label{b2:ex:06-theta-derivative}
证明
\[
 \sum_{n\in\mathbb Z}n^2e^{-\pi n^2}
 =\dfrac1{4\pi}\sum_{n\in\mathbb Z}e^{-\pi n^2}.
\]
\end{exercise}
\begin{solution}
令 \(\vartheta(t)=\sum_{n\in\mathbb Z}e^{-\pi n^2t}\)、\(t>0\)。在 \(t\) 的紧正区间上，导数级数由可和的 Gauss 尾一致控制，故可逐项求导。对\cref{b2:thm:06-theta-transformation}给出的 \(\vartheta(t)=t^{-1/2}\vartheta(1/t)\) 求导并令 \(t=1\)，得 \(\vartheta'(1)=-\vartheta(1)/2-\vartheta'(1)\)，所以 \(\vartheta'(1)=-\vartheta(1)/4\)。再用 \(\vartheta'(1)=-\pi\sum_nn^2e^{-\pi n^2}\) 即得。
\end{solution}

% 自拟：空间截断与频率误差分开控制，可作为实际求积依据。
\Needspace{6\baselineskip}
\begin{exercise}\label{b2:ex:06-quadrature-error}
设 \(h\in(0,1]\)、\(M\) 为非负整数，令 \(Q_{h,M}=h\sum_{|n|\leq M}e^{-\pi h^2n^2}\)。证明
\[
 |Q_{h,M}-1|
 \leq\dfrac{2e^{-\pi/h^2}}{1-e^{-3\pi/h^2}}
 +\dfrac{2h e^{-\pi h^2(M+1)^2}}{1-e^{-\pi h^2(2M+3)}}.
\]
\end{exercise}
\begin{solution}
无穷求和由 Poisson 公式等于 \(\sum_ke^{-\pi k^2/h^2}\)，减去零频项 \(1\) 后，用 \(k^2\geq1+3(k-1)\)、\(k\geq1\) 得到第一个几何级数界。空间截断误差为 \(2h\sum_{n=M+1}^\infty e^{-\pi h^2n^2}\)。写 \(n=M+1+j\)，有 \(n^2\geq(M+1)^2+(2M+3)j\)，得到第二项。两类误差分别来自离散采样和有限截断，最后用三角不等式合并。
\end{solution}

% 改编：Hunter2001Applied, Theorem11.45, p.319 的C1衰减范围；
% 在单位周期下用于非Schwartz有理核，并推导显式双边求和。
\Needspace{6\baselineskip}
\begin{exercise}\label{b2:ex:06-rational-sum}
对 \(a>0\)，以 \(\coth x=(e^x+e^{-x})/(e^x-e^{-x})\)、\(x>0\) 表示本题所需的双曲余切。证明
\[
 \sum_{n\in\mathbb Z}\dfrac1{n^2+a^2}
 =\dfrac\pi a\coth(\pi a).
\]
所用函数不属于 Schwartz 空间，须核对正文的有限阶衰减判据再计算。
\end{exercise}
\begin{solution}
取 \(f(x)=1/(x^2+a^2)\)，其函数和一阶导数分别为 \(O((1+|x|)^{-2})\)、\(O((1+|x|)^{-3})\)。因此满足\cref{b2:prop:06-finite-decay-poisson}的条件，可以使用求和公式。

\subsecref{b2:subsec:030503}的一维 Poisson 核变换说明，\(P_a(x)=a/(\pi(a^2+x^2))\) 满足 \(\widehat{P_a}(\xi)=e^{-2\pi a|\xi|}\)。由 \(f=(\pi/a)P_a\)，得到 \(\widehat f(k)=\pi a^{-1}e^{-2\pi a|k|}\)。在 \(x=0\) 处使用周期化求和公式，得到
\[
 \sum_nf(n)=\dfrac\pi a\left(1+2\sum_{k=1}^\infty e^{-2\pi ak}\right)
 =\dfrac\pi a\dfrac{1+e^{-2\pi a}}{1-e^{-2\pi a}}
 =\dfrac\pi a\coth(\pi a).
\]
\end{solution}

% 自拟：整数采样变换不恢复未周期化的函数。
\Needspace{6\baselineskip}
\begin{exercise}\label{b2:ex:06-periodization-kernel}
构造非零 \(f\in C_c^\infty(\mathbb R;\mathbb C)\)，使周期化 \((\mathcal Pf)(x)=\sum_{n\in\mathbb Z}f(x+n)\) 恒为零。证明其 Fourier 变换在每个整数处为零，但并非处处为零。这是否与变换唯一性矛盾？
\end{exercise}
\begin{solution}
取非零 \(g\in C_c^\infty((-1/4,1/4);\mathbb C)\)，令 \(f(x)=g(x)-g(x-1)\)。两个支集不交，故 \(f\ne0\)，而周期化时两个整数平移和逐项重排后相同，所以 \(\mathcal Pf=0\)。变换为 \((1-e^{-2\pi\mathrm{i}\xi})\widehat g(\xi)\)，在整数处全为零。若它处处为零，\cref{b2:cor:03-uniqueness}会给出 \(f=0\) 几乎处处，连续性再给出 \(f=0\) 处处，矛盾。因此整数频率的采样只确定周期化函数；唯一性定理要求全部连续频率的数据，二者条件不同。
\end{solution}

% 自拟：从平方和生成函数读取有序计数，辨认符号和顺序。
\Needspace{6\baselineskip}
\begin{exercise}\label{b2:ex:06-squares-count}
设 \(d,e\) 为正整数，\(m\) 为非负整数；对每个正整数 \(q\) 和非负整数 \(\ell\)，记 \(r_q(\ell)=\#\{n\in\mathbb Z^q:|n|^2=\ell\}\)，其中符号和坐标顺序均计数。证明 \(r_{d+e}(m)=\sum_{j=0}^mr_d(j)r_e(m-j)\)，并直接计算 \(r_2(5)\)、\(r_4(1)\)、\(r_4(2)\)。解释为何不应再除以排列数。
\end{exercise}
\begin{solution}
把一个 \((d+e)\) 元整数向量分成前 \(d\) 个和后 \(e\) 个坐标，前一组平方和为某个 \(j\)，后一组为 \(m-j\)。对各组独立计数再求和即得公式。平方和为 \(5\) 的二维向量是 \((\pm1,\pm2)\) 及其坐标交换，共 \(8\) 个。四维平方和为 \(1\) 时选一个坐标放 \(\pm1\)，共 \(8\) 个；平方和为 \(2\) 时选两个坐标各放 \(\pm1\)，共 \(\binom42 2^2=24\) 个。theta 乘积的每个因子对应一个有序坐标，因此这些排列已代表不同向量，不应再除掉。
\end{solution}

% 自拟：非正交格的逆转置与最短向量，调用选读正文。
\Needspace{6\baselineskip}
\begin{optionalexercise}\label{b2:ex:06-parity-lattice}
令 \(\Lambda=\{m(2,0)+n(1,1):m,n\in\mathbb Z\}\subset\mathbb R^2\)，并对 \(t>0\) 记 \(\Theta_\Lambda(t)=\sum_{v\in\Lambda}e^{-\pi t|v|^2}\)。求 \(\Lambda^*\)、两格的基本区域体积和最短非零向量长度，写出 \(\theta\) 变换，并求 \(\Theta_\Lambda(t)\) 在 \(t\to\infty\) 时的首个指数修正。
\end{optionalexercise}
\begin{solution}
基矩阵为 \(A=\left(\begin{smallmatrix}2&1\\0&1\end{smallmatrix}\right)\)，故对偶格由 \((1/2,-1/2)\)、\((0,1)\) 生成；基本区域体积分别为 \(2\)、\(1/2\)。原格是两个整数坐标同奇偶的向量：若坐标都为偶数，非零向量长度至少为 \(2\)；若都为奇数，长度至少为 \(\sqrt2\)。因此最短非零向量恰为 \((\pm1,\pm1)\)。对偶格向量可以写成 \((p/2,q-p/2)\)、\(p,q\in\mathbb Z\)。当 \(p\) 为偶数时两坐标都是整数，非零向量长度至少为 \(1\)；当 \(p\) 为奇数时两坐标都是半整数，长度至少为 \(1/\sqrt2\)，取等向量恰为 \((\pm1/2,\pm1/2)\)。因此
\[
 \Theta_\Lambda(t)=\dfrac1{2t}\Theta_{\Lambda^*}(1/t),
 \quad\Theta_\Lambda(t)=1+4e^{-2\pi t}+O(e^{-4\pi t}).
\]
最后的尾部估计可在 \(t\geq1\) 时提出 \(e^{-4\pi t}\)，其余由格上可和的固定 Gauss 级数控制；最短向量后的平方长度至少为 \(4\)。
\end{solution}

% 自拟：正权平方矩的方差解释theta的对数凸性。
\Needspace{6\baselineskip}
\begin{optionalexercise}\label{b2:ex:06-log-convexity}
对 \(t>0\)，令 \(\vartheta(t)=\sum_{n\in\mathbb Z}e^{-\pi n^2t}\)。证明 \(\vartheta\) 在 \((0,\infty)\) 上严格递减，而 \(\log\vartheta\) 在该区间上严格凸。将后一结论解释成以权重 \(w_n(t)=e^{-\pi n^2t}/\vartheta(t)\) 计算的方差为正。
\end{optionalexercise}
\begin{solution}
逐项求导得到 \(\vartheta'(t)=-\pi\sum_nn^2e^{-\pi n^2t}<0\)。权重正且总和为 \(1\)，于是
\[
 (\log\vartheta)''(t)
 =\pi^2\left(\sum_n n^4w_n(t)-\left(\sum_n n^2w_n(t)\right)^2\right).
\]
括号中为随机取值 \(n^2\) 的方差。值 \(0\) 与 \(1\) 都有正权重，所以它不恒定，方差严格为正。所有一、二阶参数求导都由紧正参数区间上的 Gauss 尾控制。
\end{solution}

% 自拟：综合周期系数、实线卷积、Schwartz 衰减和 Poisson 周期化；不以热方程理论为前置。
\Needspace{6\baselineskip}
\begin{exercise}[周期化与 Gauss 平滑的交换]\label{b2:ex:06-periodization-gaussian-commute}
设 \(f\in\mathcal S(\mathbb R;\mathbb C)\)，\(t>0\)，并令
\[
 H_t(x)=(4\pi t)^{-1/2}e^{-x^2/(4t)},\qquad
 K_t=\mathcal PH_t,\qquad
 Q_tf=\mathcal P(H_t*f).
\]
圆周取 \(\mathbb T=\mathbb R/\mathbb Z\) 及总质量为 \(1\) 的 Lebesgue 测度，周期卷积定义为 \((u*_{\mathbb T}v)(x)=\int_0^1u(x-y)v(y)\,\mathrm dy\)。
\begin{enumerate}
\item\label{b2:item:06-gaussian-commute}证明 \(Q_tf=K_t*_{\mathbb T}(\mathcal Pf)\)，并对展开周期化级数后使用的求和、积分交换给出绝对可积控制。
\item\label{b2:item:06-gaussian-coefficients}求 \(Q_tf\) 的全部 Fourier 系数，并证明 \(\int_0^1Q_tf=\int_{\mathbb R}f\)。
\item\label{b2:item:06-gaussian-approximation}对每个非负整数 \(m\)，证明存在只依赖 \(f,m\) 的常数 \(C_{f,m}\)，使 \(0<t\leq1\) 时
\[
 \|Q_tf-\mathcal Pf\|_{C^m(\mathbb T)}\leq C_{f,m}t,
 \qquad
 \|Q_tf-\mathcal Pf-t(\mathcal Pf)''\|_{C^m(\mathbb T)}
 \leq C_{f,m}t^2.
\]
这里 \(\|v\|_{C^m}=\max_{0\leq j\leq m}\|v^{(j)}\|_\infty\)。
\end{enumerate}
\end{exercise}
\begin{solution}
固定 \(t>0\)。由\cref{b2:prop:05-schwartz-convolution}，\(H_t*f\in\mathcal S\)，故其周期化是光滑函数。又 \(K_t\) 是有界的正周期函数。对每个 \(x\)，展开周期卷积时的绝对值总积分满足
\[
 \begin{aligned}
 &\int_0^1\sum_{m,n\in\mathbb Z}|f(y+m)|H_t(x-y+n)\,\mathrm dy\\
 &\quad=\int_{\mathbb R}|f(z)|\sum_{r\in\mathbb Z}H_t(x-z+r)\,\mathrm dz
 \leq\|f\|_1\|K_t\|_\infty<\infty.
 \end{aligned}
\]
其中令 \(z=y+m\)、\(r=n+m\)，各个区间 \([m,m+1)\) 恰好分割实线；对非负项先用 Tonelli 即可得到这个等式。于是对原来的复值积分也可用 Fubini，得到
\[
 \begin{aligned}
 (K_t*_{\mathbb T}\mathcal Pf)(x)
 &=\sum_{r\in\mathbb Z}\int_{\mathbb R}H_t(x-z+r)f(z)\,\mathrm dz\\
 &=\sum_{r\in\mathbb Z}(H_t*f)(x+r)=Q_tf(x).
 \end{aligned}
\]
这说明先把实线输入平滑再周期化，与先周期化再用周期核平滑，得到同一个函数。

\cref{b2:prop:06-periodic-coefficients}和实线卷积定理给出第 \(k\) 个系数
\[
 \widehat{Q_tf}(k)=\widehat{H_t*f}(k)
 =e^{-4\pi^2tk^2}\widehat f(k),\qquad k\in\mathbb Z.
\]
取 \(k=0\) 即得积分恒等式。由于 \(\widehat f\) 仍是 Schwartz 函数，对任意非负整数 \(N\)，级数 \(\sum_k(1+|k|)^N|\widehat f(k)|\) 收敛。因此 Poisson 公式和逐项求导可写为
\[
 Q_tf(x)=\sum_{k\in\mathbb Z}e^{-4\pi^2tk^2}\widehat f(k)e^{2\pi\mathrm ikx},
 \qquad
 \mathcal Pf(x)=\sum_{k\in\mathbb Z}\widehat f(k)e^{2\pi\mathrm ikx}.
\]
对 \(x\geq0\)，有 \(0\leq1-e^{-x}\leq x\) 及 \(|e^{-x}-1+x|\leq x^2/2\)。将这两个不等式分别用于 \(x=4\pi^2tk^2\)，再对 \(j\leq m\) 阶导数逐项估计，所需上界分别为
\[
 4\pi^2t\sum_k(2\pi|k|)^j k^2|\widehat f(k)|,
 \qquad
 8\pi^4t^2\sum_k(2\pi|k|)^j k^4|\widehat f(k)|.
\]
在 \(j=0\) 时导数因子按 \(1\) 取值。这些固定级数均收敛，且不依赖 \(t\)，从而得到两个 \(C^m\) 估计；第二式中的一次项正是 \(t(\mathcal Pf)''\)。
\end{solution}

